% Ordered Associative Chiral Koszul Duality
% Integration-ready chapter file (stripped from standalone amsart draft).
% Uses only \providecommand for macros that may not be in main.tex preamble.

\providecommand{\Assch}{\mathrm{Ass}^{\mathrm{ch}}}
\providecommand{\Barch}{\overline{B}^{\mathrm{ch}}}
\providecommand{\Cobar}{\Omega^{\mathrm{ch}}}
\providecommand{\coHoch}{\operatorname{coHH}}
\providecommand{\Cotor}{\operatorname{Cotor}}
\providecommand{\Coext}{\operatorname{Coext}}
\providecommand{\RHom}{R\!\operatorname{Hom}}
\providecommand{\Tot}{\operatorname{Tot}}
\providecommand{\KK}{\mathbb{K}}
\providecommand{\Dpbw}{D^{\mathrm{pbw}}}
\providecommand{\Dco}{D^{\mathrm{co}}}
\providecommand{\chotimes}{\mathbin{\otimes^{\mathrm{ch}}}}
\providecommand{\wt}{\widetilde}
\providecommand{\eps}{\varepsilon}
\providecommand{\susp}{s}
\providecommand{\coeq}{\operatorname{coeq}}
\providecommand{\id}{\mathrm{id}}
\providecommand{\op}{\mathrm{op}}
\providecommand{\cop}{\mathrm{cop}}
\providecommand{\HH}{\operatorname{HH}}
\providecommand{\Tor}{\operatorname{Tor}}
\providecommand{\Ext}{\operatorname{Ext}}
\providecommand{\gr}{\operatorname{gr}}
\providecommand{\MC}{\operatorname{MC}}
\providecommand{\Ob}{\operatorname{Ob}}
\providecommand{\eq}{\operatorname{eq}}
\providecommand{\Eone}{\mathsf{E}_1}
\providecommand{\Einf}{\mathsf{E}_{\infty}}
\providecommand{\Vlat}{V}
\providecommand{\barBgeom}{\bar{\mathbf{B}}}

\chapter{Ordered Associative Chiral Koszul Duality}
\label{ch:ordered-associative-chiral-kd}

\section{Introduction}

The ordered bar complex $\Barchord(A)$ of a chiral algebra~$A$
is the cofree coalgebra $T^c(s^{-1}\bar{A})$ with differential extracting
collision residues at consecutive points and coproduct from interval-cutting.
Its linear dual, on the chirally Koszul locus, is the dg-shifted Yangian.

\begin{remark}[$\Eone$-chiral notion used in this chapter]
\label{rem:ord-assoc-notion-fixed}
\index{$\mathbb{E}_1$-chiral algebra!notion fixed in chapter}
The phrase ``$E_1$-chiral algebra'' in this chapter refers to
notion \textup{(A)} of
Warning~\ref{warn:multiple-e1-chiral} (a strict
$\Assch$-algebra, equivalently an algebra over the non-$\Sigma$
chiral associative operad), with strictification of higher
operations $m_k = 0$ for $k\ge 3$. The homotopical refinement is
notion \textup{(B)} (an $A_\infty$-chiral structure in
$\End^{\mathrm{ch}}_A$); the strict-vs-homotopy passage is
controlled by the homotopy transfer of
Hypothesis~\ref{hyp:inputs}. The Etingof--Kazhdan
quantum-vertex notion \textup{(C)} produces the same Koszul-dual
on the Koszul locus through the chiral quantum-group equivalence
\textup{(}cf.\ Theorem~\ref{thm:chiral-qg-equiv}, witnessed for
$\fsl_2$ Yangian at the level of Drinfeld second
presentation\textup{)}; the present chapter does \emph{not}
appeal to notions \textup{(D)}, \textup{(E)}.
\end{remark}

\begin{remark}[Bridge to the Koszul Reflection Theorem]
\label{rem:ord-assoc-koszul-reflection-bridge}
\index{Koszul Reflection Theorem!ordered channel}
The bar--cobar adjunction $\Cobar \dashv \Barch$ on
$E_1$-chiral algebras specialised to the strongly admissible
regime (Hypothesis~\ref{hyp:inputs}) is the chain-level shadow
of the Koszul Reflection Theorem
(Theorem~\ref{thm:koszul-reflection} in
Chapter~\ref{chap:theorem-A-infinity-2}): the unit
$\eta_A \colon A \xrightarrow{\sim} \Cobar\Barch(A)$ on the
ordered side is the strict-input channel of the
$(\infty,2)$-equivalence proved at properad level. The
properad-level theorem upgrades the chain-level
quasi-isomorphism studied in this chapter to an
$(\infty,2)$-adjoint equivalence on the conilpotent-complete
locus; specialising back to the strongly admissible regime via
the Francis--Gaitsgory ambient category recovers
Hypothesis~\ref{hyp:inputs}(I) verbatim.
The $R$-twisted descent
(Lemma~\ref{lem:R-twisted-descent}) is the symmetrisation map
$\Barord(A) \to \BarSig(A)$ that relates the present chapter's
ordered bar to the symmetric bar of Theorem~A.
\end{remark}

Three bar complexes coexist and must not be conflated.
The Francis--Gaitsgory bar $\BarchFG(A)$ uses only the zeroth
product $a_{(0)}b$ and sees the chiral Lie structure.
The symmetric bar $\BarchSig(A)$ uses all OPE products and takes
$\Sigma_n$-coinvariants; this is the bar complex of Theorem~A.
The ordered bar $\Barchord(A)$ uses all OPE products but retains
the linear ordering on configurations, with no $\Sigma_n$ quotient.
The natural map $\Barchord\to \BarchSig$ takes
coinvariants; the associated graded of $\BarchSig$ recovers
$\BarchFG$.

Let $C = \Barch(A)$ denote the ordered bar coalgebra.
The coalgebra-side avatar of the diagonal bimodule~${}_A A_A$
is the \emph{diagonal bicomodule} $C_\Delta$, the coalgebra~$C$
equipped with its coproduct on both sides.
Interval composition on the algebra side is derived tensor over~$A$;
on the coalgebra side it is derived cotensor over~$C$.
The annulus is the self-cotensor:
\[
\HH^{\mathrm{ch}}_\bullet(A)\simeq C_\Delta \square_{C^e}^{\mathbf R} C_\Delta.
\]
The ordered bar complex presents the $\Eone$-chiral Koszul dual:
the associative dual already remembers the open boundary formalism.

\section{The strongly admissible regime and the foundational inputs}
\label{sec:setup}

We work over a field $k$ of characteristic $0$. Let $X$ be a smooth curve over $k$, and let
$\mathcal D(X)$ denote the relevant derived category of $D$-modules on $X$ (or on the Ran space of
$X$) endowed with the chiral tensor product
\[
\chotimes := \boxtimes_{D_X}.
\]
An $E_1$-chiral algebra means an algebra over the non-$\Sigma$ operad $\Assch$.

\begin{definition}[Strongly admissible associative chiral algebra]
An augmented $E_1$-chiral algebra $A$ on $X$ is called \emph{strongly admissible} if the following
hold.
\begin{enumerate}[label=(\alph*)]
\item $A$ is pro-nilpotent and equipped with an exhaustive separated descending PBW-type filtration
$F^\bullet A$ which is complete for the induced topology.

\item Every associated graded piece and every finite filtration quotient is bounded below and perfect
enough that completed tensor products, totalizations, and normalized chains behave exact\-ly on these
pieces.

\item The reduced ordered bar coalgebra
\[
C:=\Barch(A)
\]
is conilpotent, and the corresponding coderived category of bicomodules admits enough coflat
resolutions so that derived cotensor products are defined and associative.

\item Whenever linear duality is invoked, each graded/conformal piece of the object under consideration
is finite-dimensional.
\end{enumerate}
\end{definition}

\begin{remark}
These hypotheses ensure that every totalization, completion, normalization, and passage to coderived categories is legitimate. The arguments below require nothing more.
\end{remark}

\begin{hypothesis}[Repository inputs]
\label{hyp:inputs}
We assume the repository has already established the following.
\begin{enumerate}[label=(I)]
\item \textbf{Associative chiral bar--cobar.}
For every augmented pro-nilpotent strongly admissible $E_1$-chiral algebra $A$ there is a
quasi-isomorphism
\[
\Cobar \Barch(A)\xrightarrow{\sim} A,
\]
and, on the appropriate coalgebra side, the corresponding bar--cobar adjunction is an equivalence.

\item \textbf{One-sided module/comodule duality.}
For every such $A$ there is an equivalence between complete left $A$-modules and conilpotent left
$\Barch(A)$-comodules, with the usual coderived/contraderived refinement in the curved or completed
regime.

\item \textbf{Enveloping algebra formalism.}
The chiral enveloping algebra
\[
A^e := A\chotimes A^{\op}
\]
is defined, and $A$-bimodules are equivalently left $A^e$-modules.
\end{enumerate}
\end{hypothesis}

\begin{definition}[Universal twisting cochain]
Let $A$ be strongly admissible and set $C:=\Barch(A)$. The bar--cobar equivalence produces a
canonical twisting cochain
\[
\tau:C\to A
\]
characterized by the Maurer--Cartan equation
\[
d\tau+\tau\star\tau=0.
\]
We call $\tau$ the \emph{universal twisting cochain}.
\end{definition}

\begin{definition}[Opposite and diagonal objects]
Let $A$ be an associative chiral algebra and let $C$ be a coaugmented dg coalgebra.
\begin{enumerate}[label=(\alph*)]
\item $A^{\op}$ is the same underlying object with multiplication reversed.
\item $C^{\cop}$ is the same underlying object with coproduct flipped.
\item The \emph{coalgebraic enveloping object} is
\[
C^e:=C\widehat\otimes C^{\cop}.
\]
\item The \emph{diagonal bicomodule} $C_\Delta$ is $C$ equipped with left and right coactions both
equal to $\Delta_C$.
\end{enumerate}
\end{definition}

\begin{lemma}[Bicomodules as comodules over the enveloping coalgebra]
\ClaimStatusProvedHere
\label{lem:bicom-e}
Let $C$ be a coaugmented conilpotent dg coalgebra. Then left $C^e$-comodules are canonically the
same thing as $C$-bicomodules.
\end{lemma}

\begin{proof}
Given a left $C^e$-coaction
\[
\rho^e:M\to C\otimes C^{\cop}\otimes M,
\]
compose with the counit on the second factor and on the first factor to obtain
\[
\lambda=(\id_C\otimes \eps\otimes \id_M)\rho^e:M\to C\otimes M,
\]
\[
\rho=(\eps\otimes \id_{C^{\cop}}\otimes \id_M)\rho^e:M\to M\otimes C.
\]
These are a left and a right $C$-coaction, and coassociativity of $\rho^e$ is equivalent to the
usual bicomodule compatibility condition.

Conversely, given commuting left and right coactions $\lambda$ and $\rho$, define
\[
\rho^e
=
(\id_C\otimes \tau_{C,M}\otimes \id_M)\circ(\lambda\otimes \id_M)\circ \rho
:
M\to C\otimes C^{\cop}\otimes M,
\]
where $\tau_{C,M}$ swaps the middle two factors. The bicomodule axioms are precisely the
coassociativity and counitality conditions for $\rho^e$.
\end{proof}

\begin{definition}[Chiral Hochschild and coHochschild theories]
Let $A$ be a strongly admissible associative chiral algebra and $C=\Barch(A)$.
For an $A$-bimodule $M$ and a $C$-bicomodule $N$, define
\[
\HH^{\mathrm{ch}}_\bullet(A,M):=\Tor_\bullet^{A^e}(A,M),
\qquad
\HH_{\mathrm{ch}}^\bullet(A,M):=\RHom_{A^e}(A,M),
\]
\[
\coHoch^{\mathrm{ch}}_\bullet(C,N):=\Cotor_\bullet^{C^e}(C_\Delta,N),
\qquad
\coHoch_{\mathrm{ch}}^\bullet(C,N):=\Coext_{C^e}^\bullet(C_\Delta,N).
\]
When $M=A$ and $N=C_\Delta$ we omit coefficients.
\end{definition}

\section{Ordered bar coalgebras and the chiral shuffle theorem}
\label{sec:shuffle}

In the associative chiral regime the basic bar object is ordered. This is the first point at which
the theory diverges sharply from the commutative shadow: \emph{every shuffle survives as a distinct
summand.}

\begin{definition}[The simplicial ordered bar object]
Let $A$ be an augmented associative chiral algebra. Define the simplicial object
$B^\Delta_\bullet(A)$ by
\[
B^\Delta_n(A):=A^{\chotimes n}\qquad (n\geq 0),
\]
with faces given by augmentation at the two ends and multiplication in adjacent slots, and
degeneracies given by unit insertion. Its normalized chains
\[
N(B^\Delta_\bullet(A))
\]
are canonically identified with the reduced ordered chiral bar complex $\Barch(A)$.
\end{definition}

In suspended notation we write a bar word of length $n$ as
\[
[a_1|\cdots|a_n].
\]

\begin{definition}[EZ-admissible pair]
A pair $(A,B)$ of strongly admissible associative chiral algebras is called \emph{EZ-admissible} if
the normalization functor on bisimplicial objects built from $A$ and $B$ commutes, filtration-wise,
with the completed chiral tensor products and totalizations appearing in the proof of the
Eilenberg--Zilber theorem.
\end{definition}

\begin{theorem}[Ordered chiral shuffle theorem]
\ClaimStatusProvedHere
\label{thm:shuffle}
Let $(A,B)$ be an EZ-admissible pair of strongly admissible associative chiral algebras.
There are natural filtered chain maps
\[
\mathsf{Sh}^{\mathrm{ch}}_{A,B}:
\Barch(A)\widehat\otimes \Barch(B)\longrightarrow \Barch(A\chotimes B),
\]
\[
\mathsf{AW}^{\mathrm{ch}}_{A,B}:
\Barch(A\chotimes B)\longrightarrow \Barch(A)\widehat\otimes \Barch(B),
\]
such that
\[
\mathsf{AW}^{\mathrm{ch}}_{A,B}\circ \mathsf{Sh}^{\mathrm{ch}}_{A,B}=\id,
\qquad
\mathsf{Sh}^{\mathrm{ch}}_{A,B}\circ \mathsf{AW}^{\mathrm{ch}}_{A,B}\simeq \id.
\]
Hence there is a filtered quasi-isomorphism of conilpotent dg coalgebras
\[
\Barch(A)\widehat\otimes \Barch(B)\xrightarrow{\sim}\Barch(A\chotimes B).
\]
\end{theorem}

\begin{proof}
Set
\[
K_{\bullet,\bullet}:=B^\Delta_\bullet(A)\chotimes B^\Delta_\bullet(B).
\]
This is a bisimplicial object in the ambient abelian or stable category of $D$-modules.

Its diagonal identifies canonically with the simplicial bar object of the tensor-product algebra:
\[
\operatorname{diag}(K)_n
=
A^{\chotimes n}\chotimes B^{\chotimes n}
\cong
(A\chotimes B)^{\chotimes n}
=
B^\Delta_n(A\chotimes B).
\]
This identification is compatible with faces and degeneracies because multiplication and augmentation
on $A\chotimes B$ are componentwise.

Now apply the classical Eilenberg--Zilber theorem filtration-wise in each exact graded piece.
EZ-admissibility ensures that normalization commutes with the completed tensor constructions used
here, so one obtains a filtered quasi-isomorphism
\[
N(\operatorname{diag}K)\xrightarrow{\sim} \Tot\bigl(NK_{\bullet,\bullet}\bigr),
\]
with inverse given by Alexander--Whitney. Since
\[
N(B^\Delta_\bullet(A))\cong \Barch(A),
\qquad
N(B^\Delta_\bullet(B))\cong \Barch(B),
\qquad
N(B^\Delta_\bullet(A\chotimes B))\cong \Barch(A\chotimes B),
\]
this yields the claimed quasi-isomorphism.

The formulas are the usual normalized Eilenberg--MacLane formulas, now interpreted in the ordered
chiral setting. Explicitly, for bar words
\[
[a_1|\cdots|a_p]\in \Barch_p(A),\qquad [b_1|\cdots|b_q]\in \Barch_q(B),
\]
set $c_i=a_i\otimes 1$ for $1\leq i\leq p$ and $c_{p+j}=1\otimes b_j$ for $1\leq j\leq q$, and define
\[
\mathsf{Sh}^{\mathrm{ch}}_{p,q}
([a_1|\cdots|a_p]\otimes[b_1|\cdots|b_q])
=
\sum_{\sigma\in \operatorname{Sh}(p,q)}
\varepsilon(\sigma)\,
[c_{\sigma^{-1}(1)}|\cdots|c_{\sigma^{-1}(p+q)}].
\]
Similarly,
\[
\mathsf{AW}^{\mathrm{ch}}([x_1|\cdots|x_n])
=
\sum_{p=0}^n
[\pi_A(x_1)|\cdots|\pi_A(x_p)]\otimes[\pi_B(x_{p+1})|\cdots|\pi_B(x_n)],
\]
where
\[
\pi_A=\id_A\otimes \eps_B,\qquad \pi_B=\eps_A\otimes \id_B.
\]
The chain-map property uses only the simplicial identities.

Finally, the shuffle map is compatible with deconcatenation: cutting a shuffled word is equivalent to
cutting the source words and then shuffling the left halves and the right halves. Therefore
$\mathsf{Sh}^{\mathrm{ch}}_{A,B}$ and $\mathsf{AW}^{\mathrm{ch}}_{A,B}$ are coalgebra maps.
\end{proof}

\begin{remark}
The theorem is genuinely associative. In a commutative or $E_\infty$ setting many of the shuffle
summands are identified after the symmetric-group quotient. Here they remain distinct. The ordered
configuration space remembers their full combinatorics.
\end{remark}

\section{The $R$-matrix as ordered-to-unordered descent}
\label{sec:r-matrix-descent}

Whenever the OPE has poles, the ordered bar complex $\Barchord(A)$
carries strictly more information than the unordered $\Barch(A)$.
The surplus is the $R$-matrix: the monodromy of the flat connection
on ordered configuration space, which serves as the descent datum
for the $\Sigma_n$-covering. The unordered bar complex is recovered
as the $R$-twisted $\Sigma_n$-descent.

\subsection*{The covering-space frame}

\begin{construction}[The covering-space frame]
\label{constr:r-matrix-covering-vol1}
\index{R-matrix!descent datum|textbf}
\index{descent!ordered to unordered bar complex}
The projection
\[
\pi_n\colon
\mathrm{Conf}_n^{\mathrm{ord}}(X)
\;\longrightarrow\;
\mathrm{Conf}_n(X)
\;=\;
\mathrm{Conf}_n^{\mathrm{ord}}(X)/\Sigma_n
\]
is a principal $\Sigma_n$-bundle, \'etale with fibre~$\Sigma_n$.
A local system~$\mathcal{L}$ on the base is equivalent to a
$\Sigma_n$-equivariant local system~$\pi_n^*\mathcal{L}$
on the total space: the descent correspondence
\[
\mathrm{LocSys}\bigl(\mathrm{Conf}_n(X)\bigr)
\;\simeq\;
\mathrm{LocSys}^{\Sigma_n}
\bigl(\mathrm{Conf}_n^{\mathrm{ord}}(X)\bigr).
\]
The bar complex at tensor degree~$n$ defines a local system
$\mathcal{F}_n^{\mathrm{ord}}$ on $\mathrm{Conf}_n^{\mathrm{ord}}(X)$
whose stalk at $(z_1,\dots,z_n)$ is
$(s^{-1}\bar A)^{\otimes n}$, with flat connection induced by the
residue differential of the ordered bar construction. To descend
$\mathcal{F}_n^{\mathrm{ord}}$ to a local system on $\mathrm{Conf}_n(X)$,
one must exhibit a $\Sigma_n$-equivariant structure: isomorphisms
\[
R_\sigma\colon
\mathcal{F}_n^{\mathrm{ord}}\big|_{(z_1,\dots,z_n)}
\;\xrightarrow{\;\sim\;}
\mathcal{F}_n^{\mathrm{ord}}\big|_{(z_{\sigma(1)},\dots,z_{\sigma(n)})}
\]
for each $\sigma\in\Sigma_n$, satisfying the cocycle condition
$R_{\sigma\tau}=R_\tau\circ R_\sigma$ and compatible with the
flat connection.

For a pole-free $E_\infty$-chiral algebra, the equivariant structure
is canonical: $R_\sigma=\varepsilon(\sigma)\cdot\tau_\sigma$, where
$\tau_\sigma$ permutes tensor factors and $\varepsilon(\sigma)$ is
the Koszul sign. The descent recovers the ordinary
$\Sigma_n$-coinvariant bar complex.

For an $E_\infty$-chiral algebra with OPE poles (every interesting vertex
algebra: $\widehat{\mathfrak{g}}_k$, $\mathrm{Vir}_c$, $\mathcal{H}_k$, $W$-algebras),
the equivariant structure must be \emph{constructed} from the OPE data.
The construction is the $R$-matrix.
\end{construction}

\subsection*{The spectral Kohno connection}

\begin{construction}[$R$-matrix as monodromy of the Kohno connection]
\label{constr:r-matrix-monodromy-vol1}
\index{R-matrix!from monodromy|textbf}
\index{monodromy!R-matrix}
\index{Kohno connection}
The bar differential on $\Barchord(A)$ extracts, at each
collision $z_i=z_j$, the collision residue
$r_{ij}(z_i-z_j)\in\operatorname{End}(\bar A\otimes\bar A)\otimes
\mathcal{O}(*\Delta)$. By the bar kernel convention (the $d\log$
kernel absorbs one power), the collision residue $r(z)$ has pole
orders one less than the OPE.

These residues define a flat connection on the trivial bundle
over $\mathrm{Conf}_n^{\mathrm{ord}}(\mathbb{C})$ with fibre
$\bar A^{\otimes n}$:
\begin{equation}\label{eq:kohno-connection}
\nabla
\;=\;
d
\;-\;
\sum_{1\leq i<j\leq n}
r_{ij}(z_i-z_j)\,d(z_i-z_j).
\end{equation}
Here $r_{ij}$ acts in the $(i,j)$ tensor slots of~$\bar A^{\otimes n}$
and $r_{ij}(z_{ij})\,dz_{ij}$ is the connection $1$-form
on $\mathrm{Conf}_n^{\mathrm{ord}}(\mathbb{C})$.
For instance, when $r(z)=k/z$ (Heisenberg), this gives $k\,dz/z$;
when $r(z) = \hbar\Omega/z$ (affine Kac--Moody), it gives
$\hbar\Omega\,dz_{ij}/z_{ij}$, the standard Knizhnik--Zamolodchikov form.

\emph{Flatness.}
The curvature
\[
\nabla^2
\;=\;
\sum_{i<j<k}
\bigl(
[r_{ij},r_{ik}]+[r_{ij},r_{jk}]+[r_{ik},r_{jk}]
\bigr)
\,dz_{ij}\wedge dz_{jk}
\]
vanishes if and only if each triple satisfies the classical
Yang--Baxter equation
\begin{equation}\label{eq:cybe-vol1}
[r_{12}(z_{12}),r_{13}(z_{13})]
+[r_{12}(z_{12}),r_{23}(z_{23})]
+[r_{13}(z_{13}),r_{23}(z_{23})]
\;=\;0,
\end{equation}
where $z_{ij}=z_i-z_j$.
The identity $d^2=0$ on the ordered bar complex
$\Barchord(A)$ entails precisely this CYBE,
as the codimension-two strata of the ordered FM
compactification $\overline{\mathrm{FM}}_n^{\mathrm{ord}}(\mathbb{C})$
impose the three-term relation on triple collisions.
Hence $\nabla$ is flat.

\emph{The $R$-matrix.}
At $n=2$, the ordered configuration space
$\mathrm{Conf}_2^{\mathrm{ord}}(\mathbb{C})
=\{(z_1,z_2)\in\mathbb{C}^2\mid z_1\neq z_2\}$
has fundamental group $\pi_1=\mathbb{Z}$, generated by the
loop $\gamma$ in which $z_1$ circles $z_2$.
The \emph{$R$-matrix} is the monodromy of $\nabla$ around~$\gamma$:
\begin{equation}\label{eq:r-matrix-monodromy-vol1}
R(z)
\;:=\;
\operatorname{Mon}_\gamma(\nabla)
\;=\;
\mathcal{P}\exp\!\Bigl(
\oint_\gamma r(z)\,dz
\Bigr)
\;\in\;
\operatorname{End}(\bar A\otimes\bar A)[\![z^{-1}]\!],
\end{equation}
where $z=z_1-z_2$ is the relative coordinate and $\mathcal{P}$ denotes
path-ordering.

At general~$n$, the pure braid group
$P_n=\pi_1(\mathrm{Conf}_n^{\mathrm{ord}}(\mathbb{C}))$
has standard generators $\gamma_{ij}$ ($i<j$, the loop in which $z_i$
winds around $z_j$ with all other points fixed). The flat connection
$\nabla$ induces a monodromy representation
\[
\rho\colon P_n\to\operatorname{GL}(\bar A^{\otimes n}),
\qquad
\gamma_{ij}\mapsto R_{ij}(z_{ij}).
\]
The braid relation $\gamma_{ij}\gamma_{ik}\gamma_{jk}
=\gamma_{jk}\gamma_{ik}\gamma_{ij}$ in $P_n$ maps to the
Yang--Baxter equation~\eqref{eq:cybe-vol1} (its exponentiated form),
so $\rho$ is a well-defined representation.
\end{construction}

\subsection*{The descent theorem}

\begin{proposition}[$R$-matrix twisted descent]
\ClaimStatusProvedHere
\label{sec:r-matrix-descent-vol1}
\index{descent!R-matrix!ordered to unordered|textbf}
\index{R-matrix!twisted descent}
Let $A$ be a strongly admissible associative chiral algebra with
ordered bar coalgebra $C=\Barchord(A)$ and
$R$-matrix $R(z)$ as in
Construction~\textup{\ref{constr:r-matrix-monodromy-vol1}}.
Then the unordered bar complex is the $R$-twisted $\Sigma_n$-descent:
\begin{equation}\label{eq:descent-identification-vol1}
\Barch(A)_n
\;\simeq\;
\bigl(\Barchord(A)_n\bigr)^{R\text{-}\Sigma_n},
\end{equation}
where $\Sigma_n$ acts on adjacent transpositions by
$\sigma_i\cdot(a_1\otimes\cdots\otimes a_n)
=(\tau_{i,i+1}\circ R_{i,i+1}(z_i{-}z_{i+1}))\,
(a_1\otimes\cdots\otimes a_n)$.
The underlying $B_n$-representation is well-defined if and
only if the $R$-matrix satisfies the Yang--Baxter equation;
it factors through $\Sigma_n$ if and only if~$R$
also satisfies strong unitarity
$R_{12}(z)\,R_{21}(-z)=\id$.
\end{proposition}

\begin{proof}
The proof proceeds in three steps: we establish the
$\Sigma_n$-action, verify its well-definedness, and identify
the fixed locus.

\emph{Step~1: the twisted $\Sigma_n$-action.}
The covering $\pi_n\colon\mathrm{Conf}_n^{\mathrm{ord}}(\mathbb{C})
\to\mathrm{Conf}_n(\mathbb{C})$ presents the base as the orbit space
of the free $\Sigma_n$-action on ordered configurations. An element
$\sigma_i=(i,\,i{+}1)\in\Sigma_n$ lifts to a path in
$\mathrm{Conf}_n^{\mathrm{ord}}(\mathbb{C})$ that continuously
exchanges $z_i$ and $z_{i+1}$. The monodromy of the flat connection
$\nabla$ (eq.~\eqref{eq:kohno-connection}) along this path decomposes
into two factors:
\begin{enumerate}[label=(\roman*)]
\item the permutation $\tau_{i,i+1}$ of tensor slots in
 $\bar A^{\otimes n}$, from the coordinate transposition, and
\item the parallel transport $R_{i,i+1}(z_i-z_{i+1})$ of the fibre,
 from the flat connection on ordered configuration space.
\end{enumerate}
Define the \emph{$R$-twisted action} on generators by
\begin{equation}\label{eq:r-twisted-action}
\rho_R(\sigma_i)
\;:=\;
\tau_{i,i+1}\circ R_{i,i+1}(z_i-z_{i+1}).
\end{equation}

\emph{Step~2: well-definedness via YBE and strong unitarity.}
The formula~\eqref{eq:r-twisted-action} defines a representation
$\rho_R\colon B_n\to\operatorname{GL}(\bar A^{\otimes n})$
of the braid group, provided the braid relation holds (verified below).
To factor $\rho_R$ through $\Sigma_n=B_n/\langle\sigma_i^2\rangle$,
we must also verify the Coxeter relations.

\emph{Involutivity} ($\rho_R(\sigma_i)^2=\id$):
this is the strong unitarity condition $R_{12}(z)\,R_{21}(-z)=\id$.
It holds by an algebraic argument: for $E_\infty$-chiral algebras
(including those with OPE poles), the Casimir
$\Omega=\sum r_n\in(\bar A\otimes\bar A)[\![z^{-1}]\!]$ is symmetric
under exchange of tensor factors ($r_{12}=r_{21}$). Path reversal on
$\mathrm{Conf}_2^{\mathrm{ord}}(\mathbb{C})$ sends $z\mapsto -z$ and
exchanges the two sheets, so the parallel transport along the reversed
path is $R_{21}(-z)$. Since the connection form is symmetric,
path reversal inverts the holonomy: $R_{21}(-z)=R_{12}(z)^{-1}$,
whence $R_{12}(z)\,R_{21}(-z)=\id$.

\emph{Distant commutativity}
($\rho_R(\sigma_i)\rho_R(\sigma_j)
=\rho_R(\sigma_j)\rho_R(\sigma_i)$ for $|i-j|\geq 2$):
immediate because $R_{i,i+1}$ and $R_{j,j+1}$ act in disjoint
tensor slots.

The braid relation
$\rho_R(\sigma_i)\rho_R(\sigma_{i+1})\rho_R(\sigma_i)
=\rho_R(\sigma_{i+1})\rho_R(\sigma_i)\rho_R(\sigma_{i+1})$
is the critical constraint. Expanding:
\begin{align*}
&(\tau_{i,i+1}\circ R_{i,i+1})
(\tau_{i+1,i+2}\circ R_{i+1,i+2})
(\tau_{i,i+1}\circ R_{i,i+1}) \\
={}&
(\tau_{i+1,i+2}\circ R_{i+1,i+2})
(\tau_{i,i+1}\circ R_{i,i+1})
(\tau_{i+1,i+2}\circ R_{i+1,i+2}).
\end{align*}
Conjugating the tensor permutations to the left and using the
standard identities
$\tau_{i,i+1}R_{j,j+1}\tau_{i,i+1}^{-1}=R_{\sigma_i(j),\sigma_i(j+1)}$,
this reduces to
\[
R_{12}(z_{12})\,R_{13}(z_{13})\,R_{23}(z_{23})
\;=\;
R_{23}(z_{23})\,R_{13}(z_{13})\,R_{12}(z_{12}),
\]
the Yang--Baxter equation. This holds because the flatness
$\nabla^2=0$ implies the CYBE
(eq.~\eqref{eq:cybe-vol1}), and the path-ordered exponential
converts the infinitesimal CYBE into the exponentiated YBE
(this is the content of the Kohno--Drinfeld theorem: the
monodromy representation of a flat Knizhnik--Zamolodchikov connection
satisfies the Yang--Baxter equation).

\emph{Step~3: identification of the descent.}
The result is a genuine $\Sigma_n$-action $\rho_R$ on
$\Barchord(A)_n$. By the descent equivalence for
principal bundles, a $\Sigma_n$-equivariant local system on
the total space $\mathrm{Conf}_n^{\mathrm{ord}}(\mathbb{C})$
descends to a local system on
$\mathrm{Conf}_n(\mathbb{C})$. The descended local system
is the equaliser of the $\rho_R$-action, which at the level of
global sections is the $R$-twisted $\Sigma_n$-invariant subspace:
\[
\bigl(\Barchord(A)_n\bigr)^{R\text{-}\Sigma_n}
\;:=\;
\bigl\{
x\in\Barchord(A)_n
\;\big|\;
\rho_R(\sigma)\cdot x=x
\;\text{for all }\sigma\in\Sigma_n
\bigr\}.
\]
We claim this equals $\Barch(A)_n$. By construction, the unordered bar
complex $\Barch(A)_n$ is the complex of sections of the descended
local system $\mathcal{F}_n$ on $\mathrm{Conf}_n(\mathbb{C})$
whose pullback to $\mathrm{Conf}_n^{\mathrm{ord}}(\mathbb{C})$ is
$\mathcal{F}_n^{\mathrm{ord}}$ with the $\rho_R$-equivariant structure.
The descent-ascent adjunction for \'etale $\Sigma_n$-coverings gives
\[
\Gamma\bigl(\mathrm{Conf}_n(\mathbb{C}),\,\mathcal{F}_n\bigr)
\;=\;
\Gamma\bigl(
\mathrm{Conf}_n^{\mathrm{ord}}(\mathbb{C}),\,
\mathcal{F}_n^{\mathrm{ord}}\bigr)^{R\text{-}\Sigma_n}.
\]
The left side is $\Barch(A)_n$ and the right side is
$(\Barchord(A)_n)^{R\text{-}\Sigma_n}$,
completing the proof.
\end{proof}

\begin{corollary}[Pole-free descent is naive]
\ClaimStatusProvedHere
\label{cor:pole-free-descent}
\index{descent!pole-free algebras}
For a pole-free $E_\infty$-chiral algebra \textup{(}a commutative chiral
algebra whose chiral product extends across the diagonal without
singularities\textup{)}, the collision residue vanishes:
$r(z)=0$. Hence $\nabla=d$ is the trivial flat connection,
$R(z)=\id$, and~\eqref{eq:descent-identification-vol1} reduces to
the ordinary $\Sigma_n$-coinvariant:
\[
\Barch(A)_n
\;\simeq\;
\bigl(\Barchord(A)_n\bigr)\big/\Sigma_n.
\]
For $E_\infty$-chiral algebras with OPE poles
\textup{(}all vertex algebras with nontrivial singular OPE:
$\widehat{\mathfrak{g}}_k$, $\mathrm{Vir}_c$, $\mathcal{H}_k$,
lattice algebras, $W$-algebras, all of which are $E_\infty$-chiral,
since every vertex algebra is $E_\infty$\textup{)}, the $R$-matrix
$R(z)\neq\id$ carries nontrivial spectral dependence derived from
the local OPE, and the descent is genuinely $R$-twisted.
\end{corollary}

\begin{proof}
If the OPE has no poles, every collision residue vanishes identically,
so the connection form in~\eqref{eq:kohno-connection} is zero and the
monodromy~\eqref{eq:r-matrix-monodromy-vol1} is the identity.
The $R$-twisted $\Sigma_n$-action~\eqref{eq:r-twisted-action}
reduces to the plain permutation $\tau_{i,i+1}$, and the
$R$-twisted invariants coincide with the ordinary
$\Sigma_n$-coinvariants.
\end{proof}

\begin{remark}[Three bar complex variants and their filtrations]
\label{rem:three-bar-variants}%
\index{bar complex!three variants}%
The chiral bar complex $\Barch(A)$, the ordered bar
complex $\Barchord(A)$, and the
Feigin--Ginzburg bar complex $\barB^{\mathrm{FG}}(A)$
(retaining only simple-pole residues) are related by two
filtrations: a pole-order filtration on $\Barch(A)$
whose associated graded recovers $\barB^{\mathrm{FG}}(A)$,
and the $\Sigma_n$-coinvariant descent
$\Barchord(A) \to \Barch(A)$ of
Proposition~\ref{sec:r-matrix-descent-vol1}.
At higher genus all three fuse via the Feynman transform
of the modular operad ($\mathrm{FCom}$ for the chiral and
Feigin--Ginzburg variants, $F\!\Ass$ for the ordered variant).
\end{remark}

\begin{remark}[The $R$-matrix as Maurer--Cartan element]
\label{rem:r-matrix-mc-vol1}
\index{R-matrix!Maurer--Cartan interpretation}
The classical $r$-matrix
$r(z)\in\operatorname{End}(\bar A\otimes\bar A)
\otimes\mathcal{O}(*\Delta)$
is a Maurer--Cartan element in the degree-$2$ component of the
$E_1$ modular convolution algebra
(cf.\ the $E_1$ shadow obstruction tower in the body of the text).
The CYBE~\eqref{eq:cybe-vol1} is the linearised
Maurer--Cartan equation; the YBE on $R(z)$ is its
exponentiated form. The $R$-matrix is the descent datum because the
convolution dg~Lie algebra governing $\Sigma_n$-equivariant structures
on local systems over $\mathrm{Conf}_n$ is precisely the
Gerstenhaber--Schack complex of bilinear operations, and
Maurer--Cartan elements in it parametrise the nontrivial descents.
Thus the entire apparatus of twisted descent is controlled by a single
MC element in the degree-$2$ sector.
\end{remark}

\begin{remark}[Motivic interpretation of ordered-to-symmetric descent]
\label{rem:motivic-interpretation}
\index{motivic interpretation!ordered bar complex}
\index{descent!motivic content}
\index{mixed Tate motives!configuration space periods}
The averaging map
$\mathrm{av}\colon\Barchord(A)\to\Barch(A)$
is a motivic projection: it kills precisely the transcendental
arithmetic content of the configuration-space integrals.

The ordered configuration space
$\mathrm{Conf}_k^{\mathrm{ord}}(\mathbb{C})$
carries a canonical mixed Hodge structure whose periods are
the iterated integrals of the Knizhnik--Zamolodchikov connection.
The $R$-matrix monodromy of
Proposition~\ref{sec:r-matrix-descent-vol1} is computed by these
periods, and the resulting transcendental data is stratified by
degree:
\begin{itemize}
\item At degree~$2$: the monodromy $R(z)$ of the flat connection
$\nabla = d - r(z)\,dz$ involves $2\pi i$, a period of
the Lefschetz motive~$\mathbb{L}$ (weight~$1$).

\item At degree~$3$: the Drinfeld associator
$\Phi_{\mathrm{KZ}}$
(the monodromy of the KZ connection on
$\mathrm{Conf}_3^{\mathrm{ord}}(\mathbb{C})$)
involves multiple zeta values
$\zeta(2),\,\zeta(3),\,\zeta(2,1),\ldots$\
(periods of mixed Tate motives over~$\mathbb{Z}$).

\item At degree~$k$: the iterated monodromy involves
periods of mixed Tate motives of weight~$\leq k-1$,
by the theorem of Brown~\cite{brown-mzv}.
\end{itemize}

\noindent
The symmetric bar complex $\Barch(A)$ retains none of this.
Its invariants (the curvature~$\kappa$, the complementarity
class~$C$, the shadow obstruction tower~$Q$) are rational numbers, periods
of the trivial motive~$\mathbb{Q}(0)$. The kernel
$\ker(\mathrm{av})$ is the motivic surplus: it carries the full
weight-graded mixed Tate content that the symmetric projection
annihilates.

This is the arithmetic shadow of a structural phenomenon.
Drinfeld's Grothendieck--Teichm\"uller group~$\widehat{GT}$
acts on the tower of KZ associators, equivalently on the
pro-nilpotent completion of the ordered bar complex at each
degree. The Deligne--Ihara conjecture identifies
$\mathrm{Lie}(\widehat{GT})$ with the free Lie algebra on
generators in odd degrees~$\geq 3$, one for each odd
zeta value~$\zeta(2n+1)$. The ordered bar complex of a
chiral algebra is a \emph{representation} of this arithmetic
symmetry group: the $\widehat{GT}$-action permutes the
associators that parametrise equivalent ordered-to-symmetric
descents. Different choices of Drinfeld associator give
different $R$-twisted $\Sigma_n$-actions with isomorphic
but not identical descent data.

The ordered bar complex retains the transcendental arithmetic of
the configuration-space integrals. The symmetric bar complex
projects to the rational shadow. Every motivic weight is
visible in the ordered complex; every motivic weight is killed
by the averaging map.
\end{remark}

\section{Opposite-duality, anti-involutions, and enveloping duality}
\label{sec:opposite}

The second foundational phenomenon is order reversal. It transforms opposite multiplication on the
algebra side into opposite deconcatenation on the bar side.

\begin{theorem}[Opposite-duality for ordered bar coalgebras]
\ClaimStatusProvedHere
\label{thm:opposite}
Let $A$ be a strongly admissible associative chiral algebra. There is a canonical isomorphism of
coaugmented conilpotent dg coalgebras
\[
\mathsf{R}_A:\Barch(A^{\op})\xrightarrow{\sim}\Barch(A)^{\cop}.
\]
If $A$ is finite-type Koszul so that linear duality returns an associative chiral algebra, then
\[
(A^{\op})^!\simeq (A^!)^{\op}.
\]
\end{theorem}

\begin{proof}
On suspended bar words define
\[
\mathsf{R}_A[\susp^{-1} a_1|\cdots|\susp^{-1} a_n]
=
(-1)^{\sum_{i<j}(|a_i|+1)(|a_j|+1)}
[\susp^{-1} a_n|\cdots|\susp^{-1} a_1].
\]
Geometrically, this is induced by the order-reversal involution on the ordered configuration space.

The bar differential is the alternating sum of adjacent collision maps. The $i$-th face for
$A^{\op}$ multiplies the $i$-th and $(i+1)$-st slots in the opposite order. Reversal carries that
face to the $(n-i)$-th face for $A$ with exactly the sign written above. Hence $\mathsf{R}_A$ is a
chain map.

The coproduct on $\Barch(A)$ is deconcatenation. Reversal sends a cut after the $p$-th letter to a
cut before the $(n-p)$-th letter. Therefore
\[
\Delta\circ\mathsf{R}_A
=
\tau\circ(\mathsf{R}_A\otimes\mathsf{R}_A)\circ\Delta,
\]
where $\tau$ flips the tensor factors, which is exactly the statement that the target is the
opposite coalgebra. The claim about linear duals follows because linear duality exchanges opposite
coproduct with opposite product.
\end{proof}

\begin{corollary}[Anti-involutions survive duality]
\ClaimStatusProvedHere
\label{cor:anti}
Any anti-involution $\iota:A\to A^{\op}$ induces an anti-involution
\[
\iota^!:A^!\to (A^!)^{\op}.
\]
If $\iota^2=\id$, then $(\iota^!)^2=\id$.
\end{corollary}

\begin{proof}
Apply Theorem~\ref{thm:opposite} to the anti-involution
$\iota\colon A\to A^{\op}$.
Koszul duality is contravariant, so it gives a morphism
$(A^{\op})^!\to A^!$.
Using the canonical identification
$(A^{\op})^!\cong (A^!)^{\op}$ from the same theorem, this is exactly
the required anti-involution
$\iota^!\colon A^!\to (A^!)^{\op}$.
Functoriality gives
$(\iota^!)^2=(\iota^2)^!$, so $\iota^2=\id$ implies
$(\iota^!)^2=\id$.
\end{proof}

\begin{remark}[Ordered Verdier duality does not exist]
\label{rem:ordered-verdier-nonexistence-opposite}
\index{Verdier duality!ordered bar|textbf}
The Verdier intertwining
$\mathbb{D}_{\Ran}\,\barBch(\cA) \simeq \barBch(\cA^!)$
\textup{(}Theorem~\textup{\ref{thm:bar-cobar-isomorphism-main}}\textup{)}
is a statement about the \emph{symmetric} bar on unordered
$\Ran(X)$. The naive ordered analogue
$\mathbb{D}_{\Ran}\,\barB^{\mathrm{ord}}(\cA)$ is
\emph{undefined}: Verdier duality on $\Ran(X)$ requires
the cosheaf/factorization property, which uses the
\emph{symmetric} monoidal structure on disjoint-open covers.
The ordered bar lives on $\mathrm{Conf}^<_n(X)$ (ordered
configurations), which is \emph{not} a factorization space
in the Beilinson--Drinfeld sense.
The correct $\Eone$~analogue of Verdier intertwining is the
\emph{opposite-coalgebra duality}
of Theorem~\textup{\ref{thm:opposite}}: the reversal involution
$\mathsf{R}_A$ provides
$\barB^{\mathrm{ord}}(\cA)^{\mathrm{op}} \simeq
\barB^{\mathrm{ord}}(\cA^{\mathrm{op}})$,
which is an ordered (non-commutative) replacement for Verdier.
The development of a full ``ribbon $\Ran$'' framework in which
$\mathbb{D}_{\Ran}$ extends to ordered configurations is an open
problem.
\end{remark}

\begin{lemma}[Closure of admissibility under opposite and enveloping constructions]
\ClaimStatusProvedHere
\label{lem:closure}
If $A$ is strongly admissible, then so are $A^{\op}$ and $A^e=A\chotimes A^{\op}$.
\end{lemma}

\begin{proof}
Passing from $A$ to $A^{\op}$ does not change the underlying filtered
vector space, augmentation, or completeness data; it only reverses the
order of multiplication.
Hence every clause in the definition of strong admissibility is checked
by the same filtration and continuity estimates for the opposite
product.
For $A^e=A\chotimes A^{\op}$, the completed tensor product of two
strongly admissible filtered objects is again complete, exhaustive, and
compatible with the bar construction degreewise.
Therefore both $A^{\op}$ and $A^e$ satisfy the same admissibility
package.
\end{proof}

\begin{corollary}[Enveloping duality]
\ClaimStatusProvedHere
\label{cor:enveloping}
Let $A$ be strongly admissible, and assume $(A,A^{\op})$ is EZ-admissible. Set $C=\Barch(A)$.
Then there is a canonical quasi-isomorphism of conilpotent dg coalgebras
\[
\Barch(A^e)\xrightarrow{\sim} C^e=C\widehat\otimes C^{\cop}.
\]
If $A$ is finite-type Koszul, then
\[
(A^e)^!\simeq (A^!)^e.
\]
\end{corollary}

\begin{proof}
By the ordered shuffle theorem,
\[
\Barch(A^e)=\Barch(A\chotimes A^{\op})\xrightarrow{\sim}\Barch(A)\widehat\otimes \Barch(A^{\op}).
\]
Now apply Theorem~\ref{thm:opposite} to the second factor:
\[
\Barch(A^{\op})\xrightarrow{\sim}\Barch(A)^{\cop}.
\]
This proves the coalgebra statement. The algebra statement follows after finite-type linear duality.
\end{proof}

\section{The two-sided twisted bicomodule and the tangent bar construction}
\label{sec:tangent}

An $A$-bimodule has two coalgebra-side incarnations:
differentiation of the bar construction at a square-zero extension,
and twisting a two-sided tensor product by the universal twisting cochain.
They coincide.

\begin{definition}[Two-sided twisted bicomodule]
\label{def:Kbi}
Let $A$ be strongly admissible, let $C=\Barch(A)$, and let $M$ be an $A$-bimodule.
Define the \emph{two-sided twisted bicomodule}
\[
\KK_A^{\mathrm{bi}}(M):=C\widehat\otimes M\widehat\otimes C
\]
with differential
\[
d_{\KK}
=
d_C\otimes 1\otimes 1 + 1\otimes d_M\otimes 1 + 1\otimes 1\otimes d_C + d_L+d_R,
\]
where
\[
d_L(c\otimes m\otimes d)
=
\sum (-1)^{|c_{(1)}|}\, c_{(1)}\otimes \tau(c_{(2)})\cdot m\otimes d,
\]
\[
d_R(c\otimes m\otimes d)
=
\sum (-1)^{|c|+|m|}\, c\otimes m\cdot \tau(d_{(1)})\otimes d_{(2)},
\]
with the usual Koszul signs determined by the grading conventions.
The left and right $C$-coactions are induced by comultiplication on the two outer factors.
\end{definition}

\begin{lemma}
\ClaimStatusProvedHere
\label{lem:Kbi-dg}
$\KK_A^{\mathrm{bi}}(M)$ is a dg $C$-bicomodule.
\end{lemma}

\begin{proof}
The untwisted part of the differential squares to zero. The quadratic terms in $d_L$ and in $d_R$
vanish by the Maurer--Cartan equation
\[
d\tau+\tau\star\tau=0.
\]
The mixed commutator $d_Ld_R+d_Rd_L$ vanishes because the left and right $A$-actions on the
bimodule $M$ commute. Compatibility with the bicomodule structure is immediate from coassociativity
of the coproduct on the outer $C$-factors.
\end{proof}

\begin{definition}[Square-zero extension and one-defect bar complex]
Let $M$ be an $A$-bimodule and form the square-zero extension
\[
A_\eps:=A\oplus \eps M,\qquad \eps^2=0.
\]
Inside the reduced ordered bar complex $\Barch(A_\eps)$, let $T_A^{\mathrm{bar}}(M)$ be the
subcomplex spanned by bar words containing exactly one $\eps M$-slot.
We call this the \emph{one-defect ordered bar complex}.
\end{definition}

\begin{proposition}
\ClaimStatusProvedHere
\label{prop:one-defect}
Modulo $\eps^2$, there is a canonical decomposition
\[
\Barch(A_\eps)\cong \Barch(A)\oplus \eps\, T_A^{\mathrm{bar}}(M).
\]
The one-defect piece $T_A^{\mathrm{bar}}(M)$ is naturally a $C$-bicomodule.
\end{proposition}

\begin{proof}
A reduced bar word in $\Barch(A_\eps)$ contains some number of $\eps M$-slots.
Since $\eps^2=0$, any word with at least two such slots vanishes modulo $\eps^2$.
Therefore only the zero-defect and one-defect pieces survive.

The bar differential merges adjacent slots. Multiplying two ordinary $A$-slots remains in $A$.
Multiplying an $A$-slot with the unique $M$-slot remains in $M$ by the bimodule structure.
There is no nonzero product of two defect slots. Hence the one-defect sector is a subcomplex.
Deconcatenation cuts a one-defect word into a defect-free part and a one-defect part, or vice versa,
so the one-defect sector is naturally a bicomodule over the ordinary bar coalgebra $C=\Barch(A)$.
\end{proof}

\begin{theorem}[Tangent identification]
\ClaimStatusProvedHere
\label{thm:tangent=K}
There is a canonical isomorphism of dg $C$-bicomodules
\[
T_A^{\mathrm{bar}}(M)\xrightarrow{\sim}\KK_A^{\mathrm{bi}}(M).
\]
Equivalently, the one-defect ordered bar complex of a square-zero extension is exactly the two-sided
twisted bicomodule of the defect bimodule.
\end{theorem}

\begin{proof}
As a graded object,
\[
C=\bigoplus_{p\geq 0} (\susp^{-1}\bar A)^{\otimes p},
\]
hence
\[
C\widehat\otimes M\widehat\otimes C
\cong
\bigoplus_{p,q\geq 0}
(\susp^{-1}\bar A)^{\otimes p}\otimes M\otimes(\susp^{-1}\bar A)^{\otimes q}.
\]
This is in tautological bijection with one-defect bar words
\[
[a_1|\cdots|a_p|m|a_{p+1}|\cdots|a_{p+q}]
\]
in which the unique defect sits between a left bar segment of length $p$ and a right bar segment of
length $q$.

Under this identification, the ordinary bar differential on the left segment and the right segment
matches the internal bar differential coming from the two outer copies of $C$, while the two bar
faces adjacent to the defect slot match exactly the two twisting terms $d_L$ and $d_R$.
Therefore the differentials agree. Deconcatenation also agrees on the nose, because cutting a
one-defect word is the same as cutting the left and right bar segments around the unique defect.
\end{proof}

\begin{corollary}[Infinitesimal dual coalgebra]
\ClaimStatusProvedHere
\label{cor:infdual}
Modulo $\eps^2$, the bar coalgebra of the square-zero extension is
\[
\Barch(A_\eps)\cong C\oplus \eps\,\KK_A^{\mathrm{bi}}(M).
\]
If $A$ is finite-type Koszul, then after linear duality
\[
(A_\eps)^!\cong A^!\oplus \eps\, M^!
\qquad (\mathrm{mod}\ \eps^2),
\]
where
\[
M^!:=\Omega_C^{\mathrm{bi}}\bigl(\KK_A^{\mathrm{bi}}(M)\bigr).
\]
\end{corollary}

\begin{proof}
Proposition~\ref{prop:one-defect} identifies the
$\eps$-linear summand of $\Barch(A_\eps)$ with the one-defect bar
complex.
Theorem~\ref{thm:tangent=K} identifies that one-defect complex with
the bicomodule $\KK_A^{\mathrm{bi}}(M)$, so modulo $\eps^2$ the bar
coalgebra splits as
$C\oplus \eps\,\KK_A^{\mathrm{bi}}(M)$.
If $A$ is finite-type Koszul, linear duality and first-order cobar
reconstruction carry the $\eps$-summand to
$M^!=\Omega_C^{\mathrm{bi}}(\KK_A^{\mathrm{bi}}(M))$, giving
$(A_\eps)^!\cong A^!\oplus \eps\,M^! \pmod{\eps^2}$.
\end{proof}

\begin{proposition}[Infinitesimal annular variation]
\ClaimStatusProvedHere
\label{prop:infann}
Modulo $\eps^2$ one has
\[
\HH^{\mathrm{ch}}_\bullet(A_\eps)
\cong
\HH^{\mathrm{ch}}_\bullet(A)\oplus \eps\, \HH^{\mathrm{ch}}_\bullet(A,M).
\]
On the coalgebra side this becomes
\[
\coHoch^{\mathrm{ch}}_\bullet(\Barch(A_\eps))
\cong
\coHoch^{\mathrm{ch}}_\bullet(C)\oplus \eps\,
\coHoch^{\mathrm{ch}}_\bullet(C,\KK_A^{\mathrm{bi}}(M)).
\]
\end{proposition}

\begin{proof}
The chiral Hochschild complex of the square-zero extension decomposes, modulo $\eps^2$, into the sector with
no defect and the sector with exactly one defect; the second computes chiral Hochschild homology with
coefficients in $M$. The coalgebra statement follows from
Corollary~\ref{cor:infdual} and the coHochschild identification proved later in
Theorem~\ref{thm:HH-coHH-homology}.
\end{proof}

\section{PBW-complete two-sided Koszul duality}
\label{sec:bimod-bicomod}

One-sided module/comodule duality extends to the
two-sided theory by passing from $A$ to its enveloping algebra $A^e$ and applying
Corollary~\ref{cor:enveloping}.

\begin{definition}[PBW-complete and coderived bicomodule categories]
Let $A$ be strongly admissible and set $C=\Barch(A)$.
\begin{enumerate}[label=(\alph*)]
\item $\Dpbw_{\mathrm{bi}}(A)$ denotes the derived category of PBW-complete $A$-bimodules obtained
from complete semifree $A^e$-resolutions.

\item $\Dco_{\mathrm{bi}}(C)$ denotes the coderived category of conilpotent $C$-bicomodules.
Equivalently, by Lemma~\ref{lem:bicom-e}, it is the coderived category of left $C^e$-comodules.
\end{enumerate}
\end{definition}

\begin{definition}[Two-sided twisted cobar]
Let $N$ be a $C$-bicomodule. Define
\[
\Omega_C^{\mathrm{bi}}(N):=A\widehat\otimes N\widehat\otimes A
\]
with the differential dual to the one in Definition~\ref{def:Kbi}.
\end{definition}

\begin{theorem}[PBW-complete bimodule/bicomodule equivalence]
\ClaimStatusProvedHere
\label{thm:bimod-bicomod}
Let $A$ be strongly admissible, and assume $(A,A^{\op})$ is EZ-admissible.
Set $C=\Barch(A)$.
Then there is an exact equivalence
\[
\KK_A^{\mathrm{bi}}:\Dpbw_{\mathrm{bi}}(A)\xrightarrow{\sim}\Dco_{\mathrm{bi}}(C)
\]
with inverse $\Omega_C^{\mathrm{bi}}$. In the curved regime the same statement holds after replacing
the ordinary derived category on the algebra side by the corresponding completed
contra\-derived/coderived version.
\end{theorem}

\begin{proof}
By the enveloping algebra formalism, $A$-bimodules are left $A^e$-modules.
By Corollary~\ref{cor:enveloping},
\[
\Barch(A^e)\xrightarrow{\sim} C^e.
\]
Apply the one-sided module/comodule Koszul duality from Hypothesis~\ref{hyp:inputs}(II) to the
associative chiral algebra $A^e$. This gives an equivalence between complete $A^e$-modules and
conilpotent $C^e$-comodules. Since left $C^e$-comodules are the same as $C$-bicomodules by
Lemma~\ref{lem:bicom-e}, we obtain the asserted equivalence.

The functor is precisely the two-sided twisted bar construction
\[
M\longmapsto C\widehat\otimes M\widehat\otimes C,
\]
because that is the one-sided bar functor for the algebra $A^e$ rewritten using the enveloping
identification. The inverse is the corresponding two-sided cobar functor.
\end{proof}

\begin{theorem}[Diagonal correspondence]
\ClaimStatusProvedHere
\label{thm:diagonal}
Under the equivalence of Theorem~\ref{thm:bimod-bicomod}, the diagonal bimodule ${}_A A_A$ is sent
to the diagonal bicomodule $C_\Delta$. Equivalently,
\[
\KK_A^{\mathrm{bi}}(A)\simeq C_\Delta
\]
in $\Dco_{\mathrm{bi}}(C)$.
\end{theorem}

\begin{proof}
By definition,
\[
\KK_A^{\mathrm{bi}}(A)=C\widehat\otimes A\widehat\otimes C
\]
with the two-sided twisting differential.

First contract the left half:
the standard bar--cobar counit resolution yields a quasi-isomorphism
\[
C\widehat\otimes_\tau A\xrightarrow{\sim} \mathbf 1
\]
compatible with the right $C$-coaction. Hence
\[
\KK_A^{\mathrm{bi}}(A)\simeq C
\]
as a right $C$-comodule.

Similarly, contracting the right half using
\[
A\widehat\otimes_\tau C\xrightarrow{\sim}\mathbf 1
\]
gives
\[
\KK_A^{\mathrm{bi}}(A)\simeq C
\]
as a left $C$-comodule.

To identify the opposite coactions, filter $\KK_A^{\mathrm{bi}}(A)$ by total coradical degree in the
two outer copies of $C$. On the associated graded the twisting terms disappear, and the bicomodule
structure is literally the left and right coproducts on the two outer factors. Under either of the
two contractions above, the associated graded collapses to a single copy of $C$, and the surviving
left and right coactions are exactly the two legs of $\Delta_C$. Because the filtration is complete
and bounded below in the strongly admissible regime, this associated-graded identification lifts to
the filtered complex. Therefore the resulting bicomodule is $C_\Delta$.
\end{proof}

\begin{corollary}[The diagonal is the unit for composition]
\ClaimStatusProvedHere
\label{cor:unit}
For every $X\in \Dco_{\mathrm{bi}}(C)$ there are canonical isomorphisms
\[
C_\Delta \square_C^{\mathbf R} X \simeq X \simeq X \square_C^{\mathbf R} C_\Delta.
\]
\end{corollary}

\begin{proof}
Transport the tautologies
\[
A\otimes_A^{\mathbf L} M\simeq M\simeq M\otimes_A^{\mathbf L} A
\]
for $A$-bimodules along Theorem~\ref{thm:bimod-bicomod} and use
Theorem~\ref{thm:diagonal}.
\end{proof}

\begin{corollary}[Tensor--cotensor gluing]
\ClaimStatusProvedHere
\label{cor:tensor-cotensor}
For all $M,N\in \Dpbw_{\mathrm{bi}}(A)$ there is a canonical isomorphism
\[
\KK_A^{\mathrm{bi}}(M\otimes_A^{\mathbf L} N)
\simeq
\KK_A^{\mathrm{bi}}(M)\square_C^{\mathbf R}\KK_A^{\mathrm{bi}}(N)
\]
in $\Dco_{\mathrm{bi}}(C)$.
\end{corollary}

\begin{proof}
Theorem~\ref{thm:bimod-bicomod} identifies the derived category of
PBW-complete bimodules with the coderived category of $C$-bicomodules.
Under this equivalence, the relative tensor product over $A$ is sent to
the relative cotensor product over $C$ because both are computed by the
same two-sided bar resolution transported through
$\KK_A^{\mathrm{bi}}$.
Applying the equivalence to $M\otimes_A^{\mathbf L}N$ therefore yields
the displayed isomorphism
$\KK_A^{\mathrm{bi}}(M)\square_C^{\mathbf R}\KK_A^{\mathrm{bi}}(N)$.
\end{proof}

\begin{remark}
Corollary~\ref{cor:tensor-cotensor} is the coalgebraic gluing law for intervals.
The diagonal bicomodule $C_\Delta$ is therefore the universal boundary condition for interval
composition.
\end{remark}

\section{chiral Hochschild--coHochschild duality}
\label{sec:HH-coHH}

Annular traces and centers on the algebra side are identified with coHochschild theories of the bar
coalgebra.

\begin{theorem}[Associative chiral Hochschild/coHochschild homology]
\ClaimStatusProvedHere
\label{thm:HH-coHH-homology}
Let $A$ be strongly admissible, let $(A,A^{\op})$ be EZ-admissible, and set $C=\Barch(A)$.
Then for every complete $A$-bimodule $M$ there is a canonical isomorphism
\[
\HH^{\mathrm{ch}}_\bullet(A,M)\simeq
\coHoch^{\mathrm{ch}}_\bullet(C,\KK_A^{\mathrm{bi}}(M)).
\]
In particular,
\[
\HH^{\mathrm{ch}}_\bullet(A)\simeq \coHoch^{\mathrm{ch}}_\bullet(C)
\simeq
C_\Delta \square_{C^e}^{\mathbf R} C_\Delta.
\]
\end{theorem}

\begin{proof}
By definition,
\[
\HH^{\mathrm{ch}}_\bullet(A,M)=\Tor_\bullet^{A^e}(A,M).
\]
Apply the equivalence of Theorem~\ref{thm:bimod-bicomod} to the pair $(A,M)$.
The diagonal bimodule $A$ goes to $C_\Delta$ by Theorem~\ref{thm:diagonal}, and $M$ goes to
$\KK_A^{\mathrm{bi}}(M)$ by definition. Derived tensor over $A^e$ is transported to derived cotensor
over $C^e$. Therefore
\[
\Tor_\bullet^{A^e}(A,M)\simeq \Cotor_\bullet^{C^e}(C_\Delta,\KK_A^{\mathrm{bi}}(M)),
\]
which is the stated formula. Taking $M=A$ yields the second assertion.
\end{proof}

\begin{theorem}[Associative chiral Hochschild/coHochschild cohomology]
\ClaimStatusProvedHere
\label{thm:HH-coHH-cohomology}
Under the same hypotheses, for every complete $A$-bimodule $M$ there is a canonical isomorphism
\[
\HH_{\mathrm{ch}}^\bullet(A,M)\simeq
\coHoch_{\mathrm{ch}}^\bullet(C,\KK_A^{\mathrm{bi}}(M)).
\]
In particular,
\[
\HH_{\mathrm{ch}}^\bullet(A)\simeq
\coHoch_{\mathrm{ch}}^\bullet(C)
=
\Coext_{C^e}^\bullet(C_\Delta,C_\Delta).
\]
\end{theorem}

\begin{proof}
Derived Hom is preserved by the equivalence of Theorem~\ref{thm:bimod-bicomod}. Therefore
\[
\RHom_{A^e}(A,M)\simeq \RHom_{C^e}(C_\Delta,\KK_A^{\mathrm{bi}}(M)),
\]
which is the desired coHochschild cohomology complex. The specialization to $M=A$ follows from the
diagonal correspondence.
\end{proof}

\begin{corollary}[The annulus as self-cotrace]
\ClaimStatusProvedHere
\label{cor:annulus}
The annular closed object attached to the universal boundary condition $C_\Delta$ is the coHochschild
object
\[
\mathcal H_C^{\mathrm{cl}}
:=
C_\Delta \square_{C^e}^{\mathbf R} C_\Delta
\simeq
\coHoch^{\mathrm{ch}}_\bullet(C)
\simeq
\HH^{\mathrm{ch}}_\bullet(A).
\]
\end{corollary}

\begin{proof}
Take $M=A$ in Theorem~\ref{thm:HH-coHH-homology}.
Then the diagonal bimodule ${}_AA_A$ is carried to the diagonal
bicomodule $C_\Delta$, and the resulting coHochschild object is
$C_\Delta\square_{C^e}^{\mathbf R}C_\Delta$.
This is exactly the closed annulus object
$\mathcal H_C^{\mathrm{cl}}$, so the three descriptions agree.
\end{proof}

\begin{corollary}[Cap action]
\ClaimStatusProvedHere
\label{cor:cap}
There is a canonical action
\[
\coHoch^{\mathrm{ch}}_\bullet(C)\square_C^{\mathbf R} C_\Delta \longrightarrow C_\Delta
\]
obtained by transporting the chiral Hochschild cap action
\[
\HH^{\mathrm{ch}}_\bullet(A)\otimes A\longrightarrow A
\]
across the equivalence of Theorem~\ref{thm:bimod-bicomod}.
\end{corollary}

\begin{proof}
View the classical Hochschild cap action as a morphism in the derived
category of $A$-bimodules.
Theorem~\ref{thm:bimod-bicomod} transports $A$ to $C_\Delta$ and
transports chiral Hochschild homology to
$\coHoch^{\mathrm{ch}}_\bullet(C)$.
Under the same equivalence, the tensor operation on the algebra side
becomes the cotensor operation on the coalgebra side, so the cap action
becomes the displayed map
$\coHoch^{\mathrm{ch}}_\bullet(C)\square_C^{\mathbf R}C_\Delta\to
C_\Delta$.
\end{proof}

\begin{remark}
The cohomology theorem is the center theorem. The homology theorem is the annulus theorem. Together
they say that the algebraic center and algebraic trace of $A$ can be read directly from the
coalgebraic diagonal $C_\Delta$.
\end{remark}

\section{The diagonal bicomodule as universal boundary condition}
\label{sec:open-trace}

The boundary side carries two operations.
Interval composition is realized by derived cotensor
(Corollary~\ref{cor:tensor-cotensor}). The ordered pair-of-pants
is the transport of the external tensor product of bimodules.
The second operation is inherently ordered: without extra structure,
it is \emph{not} symmetric.

\begin{definition}[Ordered fusion product]
Let $X,Y\in \Dco_{\mathrm{bi}}(C)$. Define
\[
X\star Y
:=
\KK_A^{\mathrm{bi}}\Bigl(
\Omega_C^{\mathrm{bi}}(X)\widehat\otimes \Omega_C^{\mathrm{bi}}(Y)
\Bigr),
\]
where $\widehat\otimes$ is the completed tensor product of bimodules with outer $A$-actions:
the left $A$-action is on the first factor and the right $A$-action is on the second factor.
We call $\star$ the \emph{ordered fusion product}.
\end{definition}

\begin{remark}
The adjective ``ordered'' is essential. Unless further symmetry data are supplied, $\star$ is only
a monoidal product reflecting the ordering of incoming open boundaries. It is the algebraic shadow
of an ordered pair-of-pants, not yet of the fully symmetric open cobordism category.
\end{remark}

\begin{theorem}[Ordered pair-of-pants algebra]
\ClaimStatusProvedHere
\label{thm:pair-of-pants}
The diagonal bicomodule $C_\Delta$ carries a canonical associative algebra structure with respect to
the ordered fusion product:
\[
\mu_P:C_\Delta\star C_\Delta\longrightarrow C_\Delta,
\qquad
\eta_P:\mathbf 1_\star\longrightarrow C_\Delta,
\]
where $\mathbf 1_\star:=\KK_A^{\mathrm{bi}}(k)$ is the image of the trivial bimodule.
\end{theorem}

\begin{proof}
By Theorem~\ref{thm:diagonal}, $C_\Delta\simeq \KK_A^{\mathrm{bi}}(A)$. The multiplication and unit
of the algebra $A$,
\[
\mu_A:A\widehat\otimes A\to A,\qquad \eta_A:k\to A,
\]
transport across $\KK_A^{\mathrm{bi}}$ to maps
\[
\mu_P:C_\Delta\star C_\Delta\to C_\Delta,\qquad \eta_P:\mathbf 1_\star\to C_\Delta.
\]
Associativity and unitality follow immediately from associativity and unitality of $A$.
\end{proof}

\begin{definition}[Ordered open trace formalism]
An \emph{ordered genus-zero open trace formalism} consists of:
\begin{enumerate}[label=(\roman*)]
\item a composition product $\square_C^{\mathbf R}$ with unit object $C_\Delta$,
\item an ordered fusion product $\star$,
\item an associative algebra structure $(C_\Delta,\mu_P,\eta_P)$ for $\star$,
\item a closed object $\mathcal H_C^{\mathrm{cl}}$ obtained by annular closure,
\item an annulus action of $\mathcal H_C^{\mathrm{cl}}$ on $C_\Delta$,
\end{enumerate}
subject to the sewing relations induced from the algebra side.
\end{definition}

\begin{theorem}[Ordered genus-zero open trace formalism]
\ClaimStatusProvedHere
\label{thm:ordered-open}
The data
\[
\bigl(\Dco_{\mathrm{bi}}(C),\square_C^{\mathbf R},C_\Delta,\star,\mu_P,\eta_P,\mathcal H_C^{\mathrm{cl}}\bigr)
\]
constitute an ordered genus-zero open trace formalism. More precisely:
\begin{enumerate}[label=(\alph*)]
\item $C_\Delta$ is the unit for interval composition $\square_C^{\mathbf R}$.
\item $(C_\Delta,\mu_P,\eta_P)$ is an associative algebra object for ordered fusion.
\item Annular closure is
\[
\mathcal H_C^{\mathrm{cl}}\simeq \coHoch^{\mathrm{ch}}_\bullet(C).
\]
\item The annulus acts on the boundary object by the action of
Corollary~\ref{cor:cap}.
\item All ordered genus-zero sewing identities are transported from the corresponding identities for
$A$-bimodules.
\end{enumerate}
\end{theorem}

\begin{proof}
Statement (a) is Corollary~\ref{cor:unit}. Statement (b) is
Theorem~\ref{thm:pair-of-pants}. Statement (c) is
Corollary~\ref{cor:annulus}. Statement (d) is Corollary~\ref{cor:cap}.
Finally, every ordered genus-zero sewing relation on the coalgebra side is the image under the
equivalence $\KK_A^{\mathrm{bi}}$ of the corresponding relation on the algebra side.
\end{proof}

\begin{remark}[What is proved and what is not]
The theorem gives the algebraic ordered open package. Upgrading to a fully symmetric
one-brane open TCFT requires additional structure symmetrizing ordered fusion (a
boundary symmetry datum or a geometric realization by bordered configuration spaces).
This is part of the conjectural frontier below.
\end{remark}

\subsection{Calabi--Yau transport}

The ordered open trace formalism acquires a Frobenius structure under a Calabi--Yau hypothesis on
$A$.

\begin{definition}[Complete chiral Calabi--Yau structure]
A strongly admissible associative chiral algebra $A$ is called \emph{$d$-Calabi--Yau} if:
\begin{enumerate}[label=(\alph*)]
\item $A$ is perfect as an $A^e$-module, and
\item there is an isomorphism in $\Dpbw_{\mathrm{bi}}(A)$
\[
\varphi_A:A\xrightarrow{\sim}\RHom_{A^e}(A,A^e)[d].
\]
\end{enumerate}
Equivalently, $A$ carries a nondegenerate $d$-shifted derived trace pairing.
\end{definition}

\begin{theorem}[Shifted ordered Frobenius structure]
\ClaimStatusProvedHere
\label{thm:CY}
Assume $A$ is $d$-Calabi--Yau. Then the underlying complex of $C_\Delta$ carries a canonical
$d$-shifted ordered Frobenius structure. More precisely:
\begin{enumerate}[label=(\roman*)]
\item the ordered pair-of-pants multiplication $\mu_P$ and unit $\eta_P$ are as above;
\item there is a trace
\[
\eps_P:C_\Delta\to k[-d]
\]
transported from the Calabi--Yau trace on $A$;
\item there is a coproduct
\[
\Delta_P:C_\Delta\to C_\Delta\star C_\Delta[d]
\]
adjoint to $\mu_P$ under $\eps_P$;
\item the ordered Frobenius identities
\[
(\mu_P\star \id)\circ(\id\star \Delta_P)
=
\Delta_P\circ \mu_P
=
(\id\star\mu_P)\circ(\Delta_P\star \id)
\]
hold.
\end{enumerate}
\end{theorem}

\begin{proof}
On the algebra side, the Calabi--Yau structure identifies $A$ with its bimodule dual shifted by
$d$. Therefore multiplication
\[
\mu_A:A\widehat\otimes A\to A
\]
admits a unique adjoint coproduct
\[
\Delta_A:A\to A\widehat\otimes A[d]
\]
with respect to the trace pairing $\eps_A:A\to k[-d]$; explicitly one may define $\Delta_A$ as the
adjoint of $\mu_A$ under $\varphi_A$. The ordered Frobenius identities are formal consequences of
associativity of $\mu_A$ and the nondegeneracy of the pairing.

Transport the entire package $(\mu_A,\eta_A,\eps_A,\Delta_A)$ across the equivalence
$\KK_A^{\mathrm{bi}}$ and use Theorem~\ref{thm:diagonal}. The resulting structure maps are
$\mu_P,\eta_P,\eps_P,\Delta_P$, and all identities are preserved because the equivalence preserves
composition.
\end{proof}

\begin{corollary}[Cardy operator on the coalgebra side]
\ClaimStatusProvedHere
\label{cor:cardy}
Under the hypotheses of Theorem~\ref{thm:CY}, the ordered annulus endomorphism of the universal
boundary object is
\[
\mathsf{Ann}_{C_\Delta}:=\mu_P\circ \Delta_P\in \operatorname{End}(C_\Delta)[d].
\]
Its annular closure factors through the closed object
\[
\mathcal H_C^{\mathrm{cl}}\simeq \coHoch^{\mathrm{ch}}_\bullet(C).
\]
\end{corollary}

\begin{proof}
For a Calabi--Yau algebra, the open annulus operator is the composite
$\mu_A\circ\Delta_A$ on the diagonal bimodule.
Applying $\KK_A^{\mathrm{bi}}$ transports the diagonal bimodule to
$C_\Delta$ and transports multiplication and comultiplication to the
ordered pair-of-pants maps $\mu_P$ and $\Delta_P$.
Hence the ordered annulus endomorphism on the coalgebra side is
$\mu_P\circ\Delta_P$.
Corollary~\ref{cor:annulus} identifies its annular closure with
$\mathcal H_C^{\mathrm{cl}}\simeq\coHoch^{\mathrm{ch}}_\bullet(C)$.
\end{proof}

\section{A master theorem and its conceptual form}
\label{sec:master}

The preceding results assemble into a single statement.

\begin{theorem}[Master theorem]
\label{thm:master}
Let $A$ be a strongly admissible augmented associative chiral algebra on $X$, assume
$(A,A^{\op})$ is EZ-admissible, and set
\[
C:=\Barch(A).
\]
Then:
\begin{enumerate}[label=(\arabic*)]
\item \textbf{Ordered shuffle.}
There is a quasi-isomorphism of conilpotent dg coalgebras
\[
\Barch(A\chotimes B)\simeq \Barch(A)\widehat\otimes \Barch(B)
\]
for every EZ-admissible partner $B$.

\item \textbf{Opposite-duality.}
There is a canonical isomorphism
\[
\Barch(A^{\op})\simeq \Barch(A)^{\cop}.
\]

\item \textbf{Enveloping duality.}
There is a canonical quasi-isomorphism
\[
\Barch(A^e)\simeq C^e.
\]

\item \textbf{Two-sided bimodule/bicomodule duality.}
There is an exact equivalence
\[
\Dpbw_{\mathrm{bi}}(A)\simeq \Dco_{\mathrm{bi}}(C).
\]

\item \textbf{Diagonal correspondence.}
Under this equivalence one has
\[
{}_A A_A \longleftrightarrow C_\Delta.
\]

\item \textbf{Composition.}
The diagonal bicomodule is the unit for derived cotensor:
\[
C_\Delta\square_C^{\mathbf R} X\simeq X\simeq X\square_C^{\mathbf R}C_\Delta.
\]

\item \textbf{chiral Hochschild/coHochschild homology.}
For every complete $A$-bimodule $M$,
\[
\HH^{\mathrm{ch}}_\bullet(A,M)
\simeq
\coHoch^{\mathrm{ch}}_\bullet(C,\KK_A^{\mathrm{bi}}(M)).
\]

\item \textbf{chiral Hochschild/coHochschild cohomology.}
For every complete $A$-bimodule $M$,
\[
\HH_{\mathrm{ch}}^\bullet(A,M)
\simeq
\coHoch_{\mathrm{ch}}^\bullet(C,\KK_A^{\mathrm{bi}}(M)).
\]

\item \textbf{Tangent theory.}
For a square-zero extension $A_\eps=A\oplus \eps M$,
\[
\Barch(A_\eps)\cong C\oplus \eps\,\KK_A^{\mathrm{bi}}(M)\qquad (\mathrm{mod}\ \eps^2).
\]

\item \textbf{Ordered open trace formalism.}
The coalgebraic boundary object $C_\Delta$ carries ordered pair-of-pants, annulus, and composition
operations transported from the algebra side, and if $A$ is $d$-Calabi--Yau then the underlying
complex of $C_\Delta$ becomes a $d$-shifted ordered Frobenius object.
\ClaimStatusProvedHere
\end{enumerate}
\end{theorem}

\begin{proof}
This is a summary of Theorems~\ref{thm:shuffle}, \ref{thm:opposite},
\ref{thm:bimod-bicomod}, \ref{thm:diagonal}, \ref{thm:HH-coHH-homology},
\ref{thm:HH-coHH-cohomology}, \ref{thm:tangent=K}, and \ref{thm:ordered-open}, together with
Corollaries~\ref{cor:enveloping}, \ref{cor:unit}, and \ref{cor:infdual}, and
Theorem~\ref{thm:CY}.
\end{proof}



% ================================================================
% ACT II: THE GEOMETRY
% ================================================================

\section{Configuration space geometry of the ordered bar complex}
\label{sec:config-space-geometry}

The master theorem packages the ordered bar complex as an abstract
dg coalgebra. The product
$\FM_k(\mathbb C)\times\mathrm{Conf}_k^{<}(\mathbb R)$
gives it a concrete body: the bar differential comes from the
holomorphic factor, the coproduct from the topological factor.

\subsection{Ordered configuration spaces}
\label{subsec:ordered-config-spaces-geom}

\begin{definition}[Ordered real configuration space]
\ClaimStatusProvedHere
\label{def:ordered-real-config}
For $k\ge 1$, the \emph{ordered real configuration space} is
\[
\mathrm{Conf}_k^{<}(\mathbb R)
\;:=\;
\bigl\{(t_1,\dots,t_k)\in\mathbb R^k
\;\big|\;
t_1<t_2<\cdots<t_k\bigr\}.
\]
This is a single open convex cell in~$\mathbb R^k$,
hence contractible.
\end{definition}

\begin{proposition}[Topology of ordered real configurations;
\ClaimStatusProvedHere]
\label{prop:ordered-real-config-topology}
\begin{enumerate}[label=\textup{(\roman*)}]
\item $\mathrm{Conf}_k^{<}(\mathbb R)$ is contractible
 (it is a connected component of
 $\mathrm{Conf}_k^{\mathrm{ord}}(\mathbb R)$, which
 consists of $k!$~contractible chambers).
\item The natural projection
 $\pi_k\colon \mathrm{Conf}_k^{\mathrm{ord}}(\mathbb R)
 \to \mathrm{Conf}_k(\mathbb R)/\Sigma_k$
 is a $k!$-sheeted covering. The fibre over a point
 $\{t_1,\dots,t_k\}\in\mathrm{Conf}_k(\mathbb R)/\Sigma_k$
 consists of the $k!$~orderings of the~$t_i$.
\item $\mathrm{Conf}_k^{<}(\mathbb R)$ is a fundamental domain
 for the $\Sigma_k$-action: every unordered configuration
 can be uniquely ordered. This discreteness (each fibre
 is a set, not a path-connected space) is the geometric
 reason that the $E_1$-structure (associativity) carries no
 homotopy-coherence data beyond the sequence of
 associahedra.
\end{enumerate}
\end{proposition}

\begin{proof}
Part~\textup{(i)}: $\mathrm{Conf}_k^{<}(\mathbb R)$ is cut out by the
strict linear inequalities
$t_1<t_2<\cdots<t_k$ inside $\mathbb R^k$, so it is an open convex
cell and hence contractible.
Permuting coordinates gives the other $k!$ chambers of
$\mathrm{Conf}_k^{\mathrm{ord}}(\mathbb R)$, and these chambers are
disjoint because two different orderings cannot hold simultaneously.
Part~\textup{(ii)}: the $\Sigma_k$-action on
$\mathrm{Conf}_k^{\mathrm{ord}}(\mathbb R)$ is free, and quotienting by
that action forgets the ordering while retaining the underlying set of
points, so the fibre over an unordered configuration is exactly its
$k!$ orderings.
Part~\textup{(iii)}: every unordered configuration admits a unique
increasing representative, namely the one obtained by sorting its
points from left to right.
Thus $\mathrm{Conf}_k^{<}(\mathbb R)$ is a fundamental domain, and the
fact that the fibre is discrete explains why the ordered theory carries
only $E_1$-type coherence.
\end{proof}

\begin{definition}[Ordered holomorphic configuration space]
\label{def:ordered-hol-config}
For a smooth curve~$X$ and $k\ge 1$, the \emph{ordered
configuration space} is
\[
\mathrm{Conf}_k^{\mathrm{ord}}(X)
\;:=\;
\bigl\{(z_1,\dots,z_k)\in X^k
\;\big|\;
z_i\neq z_j\;\text{for all}\;i\neq j\bigr\}.
\]
The symmetric group~$\Sigma_k$ acts freely by permuting
coordinates. The unordered configuration space is the
quotient:
$\mathrm{Conf}_k(X)
=\mathrm{Conf}_k^{\mathrm{ord}}(X)/\Sigma_k$.
For $X=\mathbb C$:
$\pi_1(\mathrm{Conf}_k^{\mathrm{ord}}(\mathbb C))=P_k$,
the pure braid group;
$\pi_1(\mathrm{Conf}_k(\mathbb C))=B_k$,
the full braid group.
\end{definition}

\subsection{The product space and its operadic role}
\label{subsec:product-space}

\begin{construction}[The SC$^{\mathrm{ch,top}}$ operation space]
\label{constr:sc-operation-space}
The two-coloured Swiss-cheese operad
$\mathrm{SC}^{\mathrm{ch,top}}$ has operation spaces
\[
\mathrm{SC}^{\mathrm{ch,top}}(k,m)
\;=\;
\FM_k(\mathbb C)\times E_1(m),
\]
where $\FM_k(\mathbb C)$ is the Fulton--MacPherson
compactification of $\mathrm{Conf}_k(\mathbb C)$
(the chiral, holomorphic factor) and $E_1(m)$ is
the little-intervals operad (the coassociative, topological factor).
Concretely, $E_1(m)\simeq \mathrm{Conf}_m^{<}(\mathbb R)$
up to translation and scaling, which is contractible.

For the ordered bar complex of an $E_1$-chiral algebra,
the relevant spaces at tensor degree~$k$ are
\[
\FM_k(\mathbb C)\times \mathrm{Conf}_k^{<}(\mathbb R).
\]
A bar element of degree~$k$ is parametrised by this product:
the coordinates $z_1,\dots,z_k\in\mathbb C$ live on the
holomorphic curve, while the ordering
$t_1<t_2<\cdots<t_k$ records the position in the
topological direction~$\mathbb R_t$. The bar differential
acts in the $\mathbb C$-direction (extracting OPE residues);
the deconcatenation coproduct acts in the $\mathbb R$-direction
(cutting the ordered sequence).
\end{construction}

\begin{remark}[The two colours]
\label{rem:two-colours-geom}
The product $\FM_k(\mathbb C)\times\mathrm{Conf}_k^{<}(\mathbb R)$
is the geometric fingerprint of a 3d holomorphic-topological
field theory on $\mathbb C_z\times\mathbb R_t$:
observables factorise holomorphically in~$z$ and
associatively in~$t$. The bar differential is the
chiral (holomorphic) component. The bar coproduct is the
coassociative (topological) component.
\end{remark}

\begin{remark}[Codimension-one generators for
\texorpdfstring{$\mathrm{SC}^{\mathrm{ch,top}}$}{SC}]
The construction above gives the operation space, not yet a minimal
generators-and-relations presentation, and by itself it does not
equip the ordered bar complex with a full Swiss-cheese algebra
structure. On compactified models the codimension-$1$ generator
families are:
\begin{enumerate}[label=\textup{(\roman*)}]
\item closed collision faces $\partial_S\FM_k(\mathbb C)$ with
 $|S|\ge 2$;
\item open interval-splitting faces in a Stasheff compactification
 of $E_1(m)$, equivalently cuts of the ordered open inputs into
 consecutive blocks;
\item in the mixed space $\FM_k(\mathbb C)\times E_1(m)$, the
 product faces $\partial_S\FM_k(\mathbb C)\times E_1(m)$ and
 $\FM_k(\mathbb C)\times F$, with $F$ an interval-splitting face
 as in \textup{(}ii\textup{)}.
\end{enumerate}
The codimension-$2$ relations are the intersections of these faces:
Arnold--Orlik--Solomon in the closed factor, Stasheff associativity
in the open factor, and mixed commuting squares comparing a closed
collision with an open split.
\end{remark}

\subsection{The ordered FM compactification}
\label{subsec:ordered-fm-compact}

\begin{construction}[Ordered Fulton--MacPherson compactification]
\label{constr:ordered-fm-compact}
The ordered FM compactification
$\overline{\FM}_k^{\mathrm{ord}}(\mathbb C)$ is the
iterated real oriented blowup of
$\mathrm{Conf}_k^{\mathrm{ord}}(\mathbb C)$ along the
\emph{ordered} diagonals
$D_{i,i+1}=\{z_i=z_{i+1}\}$ for $i=1,\dots,k-1$.
(In the unordered FM compactification, one blows up along
all diagonals $D_S=\{z_i=z_j:i,j\in S\}$ for all subsets
$S\subseteq\{1,\dots,k\}$ with $|S|\ge 2$. In the ordered
case, only \emph{consecutive} diagonals appear.)

The codimension-one boundary faces correspond to
\emph{consecutive} collisions $z_i\to z_{i+1}$, not to
arbitrary subsets. A non-consecutive collision such as
$z_1\to z_3$ (with $z_2$ remaining separated) is not a
codimension-one face of the ordered compactification.

\medskip
\noindent\emph{Low-degree examples.}

For $k=2$:
$\overline{\FM}_2^{\mathrm{ord}}(\mathbb C)
\cong\overline{\mathbb C}$ (the Riemann sphere), with a
single boundary divisor $D_{1,2}=\{z_1=z_2\}$.

For $k=3$: two codimension-one faces: $D_{1,2}$ (first pair
collides) and $D_{2,3}$ (last pair collides). The face
$D_{1,3}$ (first and third collide without the middle) is
codimension~$2$: it requires both $D_{1,2}$ and $D_{2,3}$ to
degenerate simultaneously. This is the associahedron~$K_3$:
a line segment with two endpoints.
\end{construction}

\subsection{Boundary strata: planted forests}
\label{subsec:boundary-strata-planted}

\begin{construction}[Boundary stratification by planted forests]
\label{constr:planted-forests}
The boundary strata of $\overline{\FM}_k^{\mathrm{ord}}(\mathbb C)$
are indexed by \emph{planted forests}: ordered sequences of
planar rooted trees.

A \emph{planted forest} $F=(T_1,T_2,\dots,T_r)$ is an
\emph{ordered} sequence of planar rooted trees whose leaves,
read from left to right across all trees, are labelled
$1,2,\dots,k$ in order. The trees $T_1,\dots,T_r$ correspond
to the $r$~clusters of points that remain separated at the
given codimension; within each tree, the internal vertices
record the collision history of points within that cluster.

The stratum indexed by~$F$ is the product
\[
\partial_F\overline{\FM}_k^{\mathrm{ord}}(\mathbb C)
\;\cong\;
\prod_{j=1}^r
\overline{\FM}_{|T_j|}^{\mathrm{ord}}(\mathbb C),
\]
where $|T_j|$ is the number of leaves of~$T_j$.

For the codimension-one faces: each is indexed by a forest
with $k{-}1$~trees, $k{-}2$ of which are single leaves and one
of which is a single binary tree with two leaves.
Concretely, the face $D_{i,i+1}$ corresponds to the forest in
which points $i$ and $i{+}1$ form a binary tree (they collide)
and all other points are isolated.

\emph{Contrast with the unordered case.}
In the unordered FM compactification, boundary strata are
indexed by \emph{trees} (rooted trees with unordered children).
In the ordered case, the strata are indexed by
\emph{planted forests} (ordered sequences of planar trees):
the ordering of trees within the forest records the linear
order on the topological direction.
\end{construction}

\subsection{The bar differential from collision residues}
\label{subsec:bar-diff-collision}

\begin{construction}[Bar differential]
\label{constr:bar-diff-collision}
Let $\cA$ be a strongly admissible $E_1$-chiral algebra with
structure constants $e_I{}_{(n)}\,e_J=\sum_K c^K_{IJ;n}\,e_K$
relative to a homogeneous basis $\{e_I\}$ of the
augmentation ideal~$\bar\cA$. The bar differential on
$\Barchord(\cA)$ decomposes as
\[
d_{\Barchord}
\;=\;
d_{\mathrm{int}}+\sum_{i=1}^{k-1} d_{i,i+1},
\]
where:
\begin{itemize}
\item $d_{\mathrm{int}}$ applies the internal differential
 of~$\cA$ to each tensor factor.
\item $d_{i,i+1}$ extracts the collision residue at the
 codimension-one stratum $D_{i,i+1}=\{z_i=z_{i+1}\}$,
 using the ordered logarithmic propagator
 $\eta_{i,i+1}=d\log(z_i-z_{i+1})$.
\end{itemize}

Explicitly, the component $d_{i,i+1}$ acts on a bar word
$[s^{-1}e_{I_1}|\cdots|s^{-1}e_{I_k}]$ by contracting
the $i$-th and $(i{+}1)$-st slots via all OPE products:
\[
d_{i,i+1}[s^{-1}e_{I_1}|\cdots|s^{-1}e_{I_k}]
\;=\;
\sum_{K,n}\pm\,c^K_{I_i I_{i+1};n}\;
[s^{-1}e_{I_1}|\cdots|s^{-1}e_K|\cdots|s^{-1}e_{I_k}]
\;\otimes\;[\delta_{D_{i,i+1}}^{(n)}],
\]
where the contracted pair is replaced by a single slot and
$[\delta_{D_{i,i+1}}^{(n)}]$ is the $n$-th residue class at the
collision divisor.

The sum runs only over \emph{consecutive} collisions
$i,i{+}1$: this is the ordered constraint. In the unordered
bar complex (Vol.~I, Theorem~A), one sums over all pairs
$\{i,j\}$ and then takes $\Sigma_k$-coinvariants.

The identity $d^2=0$ follows from Stokes' theorem applied
to the boundary strata of
$\overline{\FM}_k^{\mathrm{ord}}(\mathbb C)$: the
codimension-two strata receive two contributions from the
two orderings of consecutive collisions, and the algebraic
content of their cancellation is the associativity axiom
for the chiral product
\textup{(}i.e.~the $E_1$-operad relation on the zeroth-mode
component, supplemented by the higher-mode residue classes
$[\delta_D^{(n)}]$ retained in the tensor decoration\textup{)}.
\end{construction}

\begin{remark}[Ordered vs.\ FG bar: the all-modes clause]
\label{rem:ord-vs-fg-all-modes}
The sum over all OPE modes~$n$ in the bar differential
\textup{(}eq.~above\textup{)} distinguishes the ordered bar
from the Francis--Gaitsgory bar, which retains only the
$n=0$ slot. The higher-mode contributions are recorded on
ordered bar words via the tensor decorations
$[\delta_D^{(n)}]$; equivalently, scalar-valued higher-mode
residues contribute to the curvature~$m_0$ of the curved
bar coalgebra \textup{(}see the Heisenberg case,
Theorem~\ref{thm:heisenberg-ordered-bar}\textup{)}, while
field-valued higher-mode residues feed spectral
$A_\infty$-corrections $m_k$ at degree $k\ge3$
\textup{(}see the Virasoro case,
\S\ref{subsec:virasoro-ordered-bar}\textup{)}.
In every case the ordered bar extracts \emph{all} OPE modes;
the convention chosen here differs from the raw-OPE
presentation on the right-hand side of the classical
Laurent expansion by the single $d\log$~absorption step
.
\end{remark}

\subsection{The deconcatenation coproduct}
\label{subsec:deconcatenation}

\begin{construction}[Deconcatenation coproduct]
\label{constr:deconcatenation}
The ordered bar coalgebra carries a coassociative coproduct
$\Delta$ defined by \emph{deconcatenation}: cutting the
ordered sequence at every possible position. For a bar word
of length~$k$:
\[
\Delta[s^{-1}e_{I_1}|\cdots|s^{-1}e_{I_k}]
\;=\;
\sum_{p=0}^{k}
[s^{-1}e_{I_1}|\cdots|s^{-1}e_{I_p}]
\otimes
[s^{-1}e_{I_{p+1}}|\cdots|s^{-1}e_{I_k}],
\]
where the $p=0$ and $p=k$ terms produce the coaugmentation.

This coproduct is the topological structure: it comes from
the ordered configuration space
$\mathrm{Conf}_k^{<}(\mathbb R)$ by cutting the
ordered real line at a point between $t_p$ and $t_{p+1}$.
The $p$ points to the left form one ordered subconfiguration;
the $k{-}p$ points to the right form another. The coassociativity
of~$\Delta$ is the associativity of interval composition.

\emph{Contrast with the holomorphic differential.}
The bar differential~$d$ extracts residues in the holomorphic
direction (collisions of points on~$\mathbb C$). The coproduct
$\Delta$ cuts in the topological direction (splitting the
ordered sequence on~$\mathbb R$). Together, the pair
$(d,\Delta)$ encodes both components of the $E_1$-chiral
coalgebra:
\begin{center}
\begin{tabular}{lll}
\textbf{Structure} & \textbf{Direction} &
 \textbf{Geometric source} \\
\hline
Bar differential~$d$ & holomorphic ($\mathbb C_z$) &
 OPE residues on $\FM_k(\mathbb C)$ \\
Coproduct~$\Delta$ & topological ($\mathbb R_t$) &
 Deconcatenation of $\mathrm{Conf}_k^{<}(\mathbb R)$
\end{tabular}
\end{center}
\end{construction}

\subsection{The covering space frame}
\label{subsec:covering-space-frame}

\begin{construction}[Ordered-to-unordered covering]
\label{constr:covering-space}
The projection
\[
\pi_k\colon
\mathrm{Conf}_k^{\mathrm{ord}}(\mathbb C)
\;\longrightarrow\;
\mathrm{Conf}_k(\mathbb C)
\;=\;
\mathrm{Conf}_k^{\mathrm{ord}}(\mathbb C)/\Sigma_k
\]
is a principal $\Sigma_k$-bundle, \'etale with fibre~$\Sigma_k$.
The fibre is a $k!$-element set.

In the topological direction, the analogous covering
$\mathrm{Conf}_k^{\mathrm{ord}}(\mathbb R)\to
\mathrm{Conf}_k(\mathbb R)/\Sigma_k$
is also a $k!$-sheeted covering, but here it is trivialised
by the total order on~$\mathbb R$: the component
$\mathrm{Conf}_k^{<}(\mathbb R)$ is a fundamental domain.
This triviality is why the topological direction carries
an $E_1$-structure (associative, no higher coherences from
the covering) while the holomorphic direction carries
monodromy (the $R$-matrix).

The bar complex at tensor degree~$k$ defines a local system
$\mathcal{F}_k^{\mathrm{ord}}$ on
$\mathrm{Conf}_k^{\mathrm{ord}}(\mathbb C)$
whose stalk at $(z_1,\dots,z_k)$ is
$(s^{-1}\bar\cA)^{\otimes k}$, with flat connection induced
by the residue differential $d_{\mathrm{res}}$. The
$R$-matrix is the monodromy of this flat connection.

To descend $\mathcal{F}_k^{\mathrm{ord}}$ from ordered to
unordered configurations requires a $\Sigma_k$-equivariant
structure: isomorphisms
$R_\sigma\colon
\mathcal{F}_k^{\mathrm{ord}}\big|_{(z_1,\dots,z_k)}
\xrightarrow{\sim}
\mathcal{F}_k^{\mathrm{ord}}\big|_{(z_{\sigma(1)},\dots,z_{\sigma(k)})}$
for each $\sigma\in\Sigma_k$, satisfying the cocycle condition
and compatible with the flat connection.

For an $E_\infty$-chiral algebra (which is any vertex algebra,
including those with OPE poles), the equivariant structure
is canonical: $R_\sigma=\epsilon(\sigma)\cdot\tau_\sigma$,
where $\tau_\sigma$ permutes tensor factors and
$\epsilon(\sigma)$ is the Koszul sign. The descent then
recovers the ordinary $\Sigma_k$-coinvariant bar complex
(Vol.~I, Theorem~A).

For a genuinely $E_1$-chiral algebra (where the ordering is
primitive, not derived from a local OPE), the equivariant
structure must be \emph{constructed} from the OPE data.
This construction is the $R$-matrix, and its existence
imposes the Yang--Baxter equation (the content of $d^2=0$
at tensor degree~$3$).
\end{construction}

\begin{remark}[Why ordered spaces]
\label{rem:why-ordered-spaces}
The use of ordered versus unordered configuration spaces
reflects the operadic symmetry of the algebra:
\begin{itemize}
\item An $E_\infty$-algebra (any vertex algebra, any BD chiral
 algebra) has operations that are $\Sigma_n$-equivariant:
 the factorisation structure is defined on the \emph{unordered}
 configuration space $\mathrm{Conf}_k(X)/\Sigma_k$. The
 ordering on the bar complex is a computational convenience
 that can be quotiented out.
\item An $E_1$-algebra (an associative chiral algebra, a
 quantum vertex algebra in the sense of Etingof--Kazhdan)
 has operations indexed by \emph{ordered} compositions:
 its bar complex uses
 $\mathrm{Conf}_k^{\mathrm{ord}}(X)$. The ordering
 records the topological direction: in
 $\mathbb C_z\times\mathbb R_t$, points along
 $\mathbb R_t$ are ordered by time, and the
 $E_1$-structure is the time-ordered product.
\item The Swiss-cheese operad $\mathrm{SC}^{\mathrm{ch,top}}$
 uses \emph{both}: unordered configurations in the
 holomorphic direction (chiral component) and ordered
 configurations in the topological direction (coassociative component).
 The ordered bar complex of a vertex algebra, which is
 $E_\infty$-chiral, uses the ordered configuration space
 for the $E_1$/coassociative component of $\mathrm{SC}^{\mathrm{ch,top}}$;
 this is the topological direction, not a claim that the
 algebra is nonlocal. The holomorphic component remains
 $E_\infty$.
\end{itemize}
The $R$-matrix $R(z)\in\operatorname{End}(V\otimes V)(\!(z)\!)$
is the data needed to pass from ordered to unordered
configurations in the holomorphic direction: it is the
monodromy of the local system on
$\mathrm{Conf}_2^{\mathrm{ord}}(\mathbb C)$ around the
generator of $\pi_1=\mathbb Z$.
\end{remark}



% ================================================================
% ACT III: THE MAIN LINE
% ================================================================

\bigskip
What does the machine produce for the simplest chiral algebra
specifically, the Heisenberg at level~$k$?

\section{The Heisenberg ordered bar complex}
\label{subsec:heisenberg-ordered-bar}

Let $V = \bC\cdot J$ be the one-dimensional vector space spanned by
the Heisenberg current, let $k\in\bC$ be the level, and let $\cH_k$
be the Heisenberg chiral algebra with OPE
\begin{equation}\label{eq:heisenberg-ope-ordered}
J(z)\,J(w) \;\sim\; \frac{k}{(z-w)^2}.
\end{equation}
The Heisenberg algebra is $E_\infty$-chiral: it is a local vertex
algebra with $\Sigma_n$-equivariant factorization structure. Its OPE
has a double pole, but poles do not break $E_\infty$-locality.

The ordered bar complex $\Barch(\cH_k)$ is the cofree conilpotent
coalgebra on the desuspension $\susp^{-1}\overline{\cH}_k$, with
the ordering of tensor factors recording the linear order
$t_1 < \cdots < t_n$ on $\bR$. We compute the complex degree by
degree.

\begin{theorem}[The Heisenberg ordered bar complex]
\label{thm:heisenberg-ordered-bar}
The ordered bar complex $\Barch(\cH_k)$ is a \emph{curved}
cofree conilpotent dg coalgebra
$(T^c(\susp^{-1}V),\, d_{\Barch},\, m_0,\, \Delta)$
carrying the following data:
\begin{enumerate}[label=\textup{(\arabic*)}]
\item \textbf{Curvature.}\enspace
The degree-$0$ component is
\begin{equation}\label{eq:heisenberg-curvature-ordered}
m_0 \;=\; k\cdot\lvert 0\rangle
\;\in\;\Barch_0(\cH_k),
\end{equation}
the constant extracted from the double-pole coefficient
$J_{(1)}J=k$. This is the genus-zero Heisenberg curvature
$\kappa(\cH_k)=k$ of the Vol.~I landscape census.

\item \textbf{Degree~$1$.}\enspace
$\Barch_1(\cH_k) = \bC\cdot[\susp^{-1}J]$, with
$d_{\Barch}[\susp^{-1}J] = 0$
\textup{(}the linear bar differential is binary;
a single element has no adjacent pair to collide\textup{)}.

\item \textbf{Degree~$2$.}\enspace
$\Barch_2(\cH_k) = \bC\cdot[\susp^{-1}J \,|\, \susp^{-1}J]$,
a \emph{single} generator
\textup{(}the ordered complex retains the linear order;
there is no $\Sigma_2$-identification\textup{)}.
The linear bar differential vanishes:
\begin{equation}\label{eq:heisenberg-d2-ordered}
d_{\Barch}[\susp^{-1}J \,|\, \susp^{-1}J]
\;=\; 0,
\end{equation}
because the contraction of the two adjacent $J$-slots into a
\emph{single} bar slot must land in the augmentation ideal
$\bar\cH_k$, whereas the only nonvanishing OPE mode of two
$J$'s is the central scalar $J_{(1)}J=k\in\bC$. The scalar
escapes the augmentation ideal entirely and is collected in
the curvature $m_0$ upstream \textup{(}item~\textup{(1)}\textup{)},
not in the linear differential~$d_1$.

\item \textbf{Degree~$n$.}\enspace
$\Barch_n(\cH_k) = \bC\cdot[\susp^{-1}J \,|\,
\cdots \,|\, \susp^{-1}J]$ \textup{(}$n$~copies\textup{)},
with $d_{\Barch} = 0$ on every element. The argument is
uniform: each adjacent collision either lands in a scalar
\textup{(}contributing to $m_0$\textup{)} or vanishes, so no
contribution survives in the augmentation ideal at positive
degree.

\item \textbf{Coproduct.}\enspace The deconcatenation coproduct
$\Delta$ on $\Barch(\cH_k)$ is
\begin{equation}\label{eq:heisenberg-delta-ordered}
\Delta[\susp^{-1}J \,|\, \cdots \,|\, \susp^{-1}J]
\;=\;
\sum_{i=0}^{n}
[\susp^{-1}J \,|\, \cdots \,|\, \susp^{-1}J]_{i}
\;\otimes\;
[\susp^{-1}J \,|\, \cdots \,|\, \susp^{-1}J]_{n-i},
\end{equation}
where the subscript denotes the number of tensor factors.
This is the tensor coalgebra coproduct, encoding interval-cutting
in $\bR$.

\item \textbf{Curved $E_1$-chiral coalgebra structure.}\enspace
The ordered bar coalgebra is the curved cofree object
$(T^c(\susp^{-1}V),\, d_1=0,\, m_0=k\lvert 0\rangle,\,\Delta)$:
the linear differential vanishes identically, but the
degree-$0$ curvature $m_0=k\lvert 0\rangle$ is non-zero. The
associated R-matrix
$R(z)=e^{k\hbar/z}$
\textup{(}Theorem~\ref{thm:heisenberg-rmatrix}\textup{)} is the
exponentiated form of the curvature seen on the ordered
configuration space.
\ClaimStatusProvedHere
\end{enumerate}
\end{theorem}

\begin{proof}
Heisenberg has the single generator~$J$ and the unique
nonvanishing OPE mode $J_{(1)}J=k$, which is a scalar. The
linear bar differential
$d_1\colon\Barch_n\to\Barch_{n-1}$ contracts an adjacent pair
of slots into a single slot of the augmentation ideal
$\bar\cH_k$. The only output available from $J\cdot J$
contractions is the scalar~$k$, which lies outside
$\bar\cH_k$; hence $d_1=0$ on every bar word. The scalar
contribution is collected separately in the degree-$0$
component, giving $m_0=k\lvert 0\rangle$.

The deconcatenation coproduct is manifestly coassociative
(it is the tensor coalgebra structure), and the curved bar
identity reduces to $d_1\circ d_1=[m_0,-]_\circ$, both sides
of which vanish on the augmentation ideal.
\end{proof}

\begin{remark}[Curvature, not differential]
\label{rem:heisenberg-j1j}
The mode $J_{(1)}J = k$ is the genus-zero curvature of the
ordered Heisenberg bar coalgebra: it lives in the degree-$0$
component~$m_0$, not in the linear differential~$d_1$. In the
\emph{symmetric} bar complex $\barB^{\Sigma}(\cH_k)$
(Vol~I, Theorem~A) the same datum is conventionally recorded by
$d_{\barB}[\susp^{-1}J \,|\, \susp^{-1}J] = k$, viewing the
constant~$k$ as the value of $d$ on the ground field copy
inside the symmetric bar. In the ordered complex the constant
is collected upstream in $m_0$, and the corresponding holonomy
on ordered configuration space is the R-matrix
$R(z)=e^{k\hbar/z}$
(Theorem~\ref{thm:heisenberg-rmatrix}).
The three pictures are equivalent presentations of the same
genus-zero datum $\kappa(\cH_k)=k$.
\end{remark}

\begin{theorem}[Collision residue and $R$-matrix;
\ClaimStatusProvedHere]
\label{thm:heisenberg-rmatrix}
Let $z=z_1-z_2$ denote the relative coordinate on
$\Conf_2^{<}(\bC)$. The collision residue of the Heisenberg
double-pole OPE is
\begin{equation}\label{eq:heisenberg-collision-residue}
r^{\mathrm{coll}}(z)
\;=\;
\frac{k}{z},
\end{equation}
a \emph{simple} pole obtained from the OPE double pole
$k/(z_1-z_2)^2$ by one step of $d\log$ absorption ($d\log$ lowers every OPE pole
order by one).

The flat connection on the trivial bundle over
$\Conf_2^{<}(\bC)$ determined by the bar propagator is
\begin{equation}\label{eq:heisenberg-connection}
\nabla \;=\; d \;-\; r^{\mathrm{coll}}(z)\,dz
\;=\; d \;-\; (k/z)\,dz,
\end{equation}
and its monodromy around the origin is
\begin{equation}\label{eq:heisenberg-monodromy}
\operatorname{Mon}_0(\nabla) \;=\; e^{-2\pi ik}
\end{equation}
(treating $k\in\bC$ as a generic continuous parameter).

The spectral $R$-matrix on modules is
\begin{equation}\label{eq:heisenberg-R-matrix}
R(z) \;=\; e^{k\hbar/z}
\qquad (\text{NOT trivial}).
\end{equation}
Since $\cH_k$ is $\Einf$-chiral, this $R$-matrix is
\emph{derived} from the local OPE via the monodromy of
$\nabla$, not independent input. The Yang--Baxter equation
is satisfied trivially because $R(z)$ is scalar
\textup{(}$\cH_k$ has rank~$1$\textup{)}.
\end{theorem}

\begin{proof}
The OPE mode expansion is
$J(z)J(w)=\sum_{n\ge0}\,J_{(n)}J\cdot(z-w)^{-n-1}$,
with the unique nonvanishing mode $J_{(1)}J=k$ giving the
double-pole coefficient. Pairing the OPE density
$k\,(z-w)^{-2}\,dz$ against the bar kernel $d\log(z-w)$ on
$\FM_2(\bC)$ and using $d\log(z-w)=-dw/(z-w)$ absorbs one
power of $(z-w)^{-1}$, producing the simple-pole residue
kernel $k/(z-w)\,dz$ on the ordered configuration space.
Writing $z=z_1-z_2$ for the relative coordinate,
$r^{\mathrm{coll}}(z)=k/z$ as claimed.

The connection $\nabla = d - (k/z)\,dz$ on the trivial
line bundle over $\bC^\times$ has flat sections
$f(z) = z^{k}$; the associated parallel transport around
$\gamma\colon z\mapsto e^{2\pi it}z$ multiplies flat
sections by $e^{2\pi ik}$, so the monodromy acting on
parallel-transported vectors is $e^{-2\pi ik}$.

The spectral $R$-matrix is obtained by the Kohno
exponentiation of the connection. On Fock modules
$\cF_\mu\otimes\cF_\nu$ the Casimir collapses to the scalar
$\mu\nu\,k$, and the path-ordered exponential reduces to
the ordinary exponential
$R(z)=\exp\bigl(\oint\,r^{\mathrm{coll}}(z)\,dz\bigr)
=e^{k\hbar/z}$ in the spectral variable, where
$\hbar=\mu\nu$. The computation is scalar throughout
because $\cH_k$ is rank one.
\end{proof}

\begin{remark}[Non-triviality]
\label{rem:heisenberg-rmatrix-nontrivial}
The Heisenberg R-matrix $R(z)=e^{k\hbar/z}$ is \emph{not}
the trivial permutation $R=\tau$. For a generic level
$k$ and generic Fock charges, $R(z)$ is a non-identity
scalar exponential series in $1/z$. The collision residue
$k/z$ is derived from the OPE double pole ($d\log$
absorption); attempting to read the R-matrix off the raw OPE
without applying the step gives the incorrect answer
$k/z^2$ and hence a wrong pole structure.
\end{remark}

\begin{theorem}[Open-colour Koszul dual: the abelian Yangian;
\ClaimStatusProvedHere]
\label{thm:heisenberg-yangian}
The ordered \textup{(}coassociative\textup{)} Koszul
dual of $\cH_k$ is the abelian Yangian
\begin{equation}\label{eq:heisenberg-yangian}
\cH_k^{!,\mathrm{ord}}
\;\simeq\;
Y\bigl(\mathfrak{u}(1)\bigr)
\;\simeq\;
\bC[h_0, h_1, h_2, \ldots],
\end{equation}
the polynomial algebra on generators $h_n$ of degree~$n+1$,
with product inherited from the linear dual of the
deconcatenation coproduct. The spectral-parameter
presentation identifies $h_n = \partial_u^n h(u)|_{u=0}$,
where $h(u) = \sum_{n\ge 0} h_n u^{-n-1}$ is the generating
current of the Yangian.
\end{theorem}

\begin{proof}
Since the bar differential vanishes on $\Barch(\cH_k)$,
the Koszul dual is the linear dual of the bar coalgebra:
\[
\cH_k^{!,\mathrm{ord}}
\;=\;
\bigl(\Barch(\cH_k)\bigr)^\vee
\;=\;
\bigl(T^c(\susp^{-1}V)\bigr)^\vee
\;\simeq\;
T(\susp V^\vee)
\;\simeq\;
\bC[h],
\]
where $h = (\susp^{-1}J)^\vee$ is the dual generator.
The coalgebra $T^c(\susp^{-1}V)$ has deconcatenation
coproduct; dualizing produces the concatenation
(polynomial) product on $T(\susp V^\vee) = \bC[h]$.

Since $V$ is one-dimensional and $d_{\Barch} = 0$,
the Koszul dual algebra is freely generated by $h$ with
no relations. Restoring the spectral grading from the
FM configuration-space stratification, the degree-$n$
component corresponds to $\FM_n(\bC)$, giving the
generators $h_0, h_1, h_2, \ldots$ of the abelian
Yangian $Y(\mathfrak{u}(1))$.

Equivalently: the Koszul dual is the endomorphism algebra
of the line operator $b$ in the coassociative line category
$\cC_{\mathrm{line}}$. The line category
$\cC_{\mathrm{line}} \simeq \cH_{-k}\text{-}\mathrm{mod}$
is equivalent to modules over the negative-level Heisenberg,
but the ordered Koszul dual $\cH_k^{!,\mathrm{ord}} = Y(\mathfrak{u}(1))
\simeq \bC[h_0, h_1, \ldots]$ is categorically distinct from
$\cH_{-k}$ as a chiral algebra (not $\cH_{-k}$; they share the
same modular characteristic $\kappa = -k$ but differ as algebras).
\end{proof}

\begin{theorem}[Formality: class~G, shadow depth~$2$;
\ClaimStatusProvedHere]
\label{thm:heisenberg-formality}
The Heisenberg ordered bar complex is formal:
the $\Ainf$-structure on $H^\bullet(\Barch(\cH_k))$
has $m_k = 0$ for all $k\ge 3$. The Heisenberg
belongs to class~G \textup{(}shadow depth $r_{\max} = 2$,
the minimum for a nontrivial chiral algebra\textup{)}.
\end{theorem}

\begin{proof}
Since $d_{\Barch} = 0$, the bar complex \emph{is} its own
cohomology. There is nothing to transfer: the $\Ainf$-structure
on cohomology inherits the trivial differential and the
deconcatenation coproduct, with all higher operations vanishing.
The vanishing of $m_{k\ge 3}$ is equivalent to saying that the
OPE has a single pole order (the double pole),
and the collision residue $r^{\mathrm{coll}} = k/z$ has a
simple pole, which falls in class~G.
\end{proof}

\begin{remark}[Contrast with the chiral \textup{(}symmetric\textup{)}
Koszul dual]
\label{rem:heisenberg-chiral-vs-ordered}
The chiral Koszul dual of $\cH_k$ (Vol~I) is the
commutative chiral algebra
\[
\cH_k^{!,\mathrm{ch}}
\;=\;
\operatorname{Sym}^{\mathrm{ch}}(V^*),
\]
a free commutative vertex algebra on one generator.
The ordered Koszul dual $Y(\mathfrak{u}(1))$ carries
strictly more structure: the spectral parameter~$u$
(arising from the FM stratification in the ordered
complex) is invisible in the symmetric Koszul dual.
The passage from ordered to symmetric is
$R$-matrix-twisted $\Sigma_n$-coinvariants: since
$R(z) = e^{k\hbar/z} \neq \tau$ (the Koszul sign),
the descent $\Barch(\cH_k)\to\barB^{\Sigma}(\cH_k)$
is genuinely twisted, and the spectral parameter
information is lost in the quotient (absorbed into
the bar differential of $\barB^{\Sigma}$, where
$d_{\barB}[\susp^{-1}J\,|\,\susp^{-1}J] = k \neq 0$).

This is the three-tier picture in its simplest instance:
the Heisenberg lives in tier~(ii)
($E_\infty$-chiral with $R(z)\neq\tau$, derived from
the local OPE). Tier~(i) would require $R(z) = \tau$
(no OPE poles); tier~(iii) would require $R(z)$ to be
independent input (genuinely $E_1$).
\end{remark}

\begin{remark}[The Heisenberg as prototype]
\label{rem:heisenberg-prototype}
The Heisenberg ordered bar complex is the simplest instance of a
general pattern. Each subsequent family generalises exactly one
aspect of the structure, while preserving the others.
\begin{enumerate}[nosep]
\item \emph{$\beta\gamma$ system
 $(\S\ref{subsec:betagamma-ordered-bar})$:}
 The underlying vector space grows from $\dim V = 1$ to $\dim V = 2$,
 and the exponential $R$-matrix $e^{k\hbar/z}$ becomes nilpotent:
 $R = 1 + 2\pi i\,\Theta$ with $\Theta^2 = 0$. The bar differential
 remains trivial ($d_{\Barch} = 0$) because the zeroth mode product
 $\beta_{(0)}\gamma = 1$ is a scalar, not a new generator. Formality
 persists: shadow depth is $r_{\max} = 4$ (class~C).

\item \emph{Affine $\widehat{\mathfrak{g}}_k$
 $(\S\ref{subsec:km-ordered-bar})$:}
 The scalar collision residue $k/z$ becomes the matrix-valued
 Casimir kernel $\hbar\,\Omega/z \in \End(\mathfrak{g}\otimes\mathfrak{g})$,
 where $\hbar = 1/(k+h^\vee)$.
 The scalar $R$-matrix $e^{k\hbar/z}$ becomes the Yang $R$-matrix
 $1 + \hbar\,\Omega/z + O(\hbar^2)$. The abelian Yangian
 $Y(\mathfrak{u}(1))$ becomes $Y_\hbar(\mathfrak{g})$. The bar
 differential acquires nontrivial components from $[X,Y]_{(0)} \neq 0$,
 but formality persists: shadow depth is $r_{\max} = 3$ (class~L).

\item \emph{Virasoro $(\S\ref{subsec:virasoro-ordered-bar})$:}
 The double OPE pole becomes quartic: $T(z)T(w) \sim (c/2)/(z-w)^4
 + \cdots$. The simple collision residue $k/z$ becomes the
 triple-pole residue $(c/2)/z^3 + 2T/z$. Formality breaks: $m_k \neq 0$ for all $k\ge 3$, so the
 $\Ainf$-structure is genuinely infinite. Shadow depth is
 $r_{\max} = \infty$ (class~M). The nontrivial Yangian structure
 collapses to a single generator $T$, but the non-formal bar
 differential compensates: the complexity migrates from the algebra's
 rank to the differential's depth.
\end{enumerate}

\noindent
The four classes resolve as follows. The Heisenberg is class~G (shadow
depth~$2$, formal, scalar). The $\beta\gamma$ system is class~C (shadow
depth~$4$, formal, nilpotent matrix). The affine is class~L (shadow
depth~$3$, formal, matrix). The Virasoro is class~M (shadow
depth~$\infty$, non-formal, collapsed rank).
Each step from Heisenberg to Virasoro escalates exactly one structural
parameter while the rest of the ordered bar architecture remains
intact.
\end{remark}


\section{The Drinfeld presentation of $Y_\hbar(\mathfrak{sl}_2)$}
\label{subsec:drinfeld-yangian-sl2}

The ordered bar complex $\Barch(V_k(\mathfrak{sl}_2))$ produces the
Yang $R$-matrix, the RTT presentation, and the Gauss decomposition.
The \emph{Drinfeld presentation} (the generating-current
form that reveals the Yangian as a deformation of the current algebra
$U(\mathfrak{sl}_2[t])$) emerges from the ordered bar data. The affine vertex algebra
$V_k(\mathfrak{sl}_2)$ is $\Einf$-chiral (it is a local vertex algebra
with $\Sigma_n$-equivariant factorisation structure); the Yangian
$Y_\hbar(\mathfrak{sl}_2)$ is the ordered Koszul dual living on the
open boundary $\bR_t$.

\subsubsection*{Generators}

The Drinfeld generators are obtained from the Gauss decomposition of the
transfer matrix $T(u) = F(u)\,K(u)\,E(u)$, where $F(u)$ is
lower-triangular unipotent, $K(u) = \operatorname{diag}(k_1(u),k_2(u))$
is diagonal, and $E(u)$ is upper-triangular unipotent. Define the
Drinfeld currents:
\begin{align}
\label{eq:drinfeld-E-current}
E(u) &= \sum_{r \ge 0} E_r\,u^{-r-1}, \\
\label{eq:drinfeld-F-current}
F(u) &= \sum_{r \ge 0} F_r\,u^{-r-1}, \\
\label{eq:drinfeld-H-current}
H(u) &= 1 + \hbar\sum_{r \ge 0} H_r\,u^{-r-1}
\;=\; \frac{k_1(u)}{k_2(u)}.
\end{align}
The modes $\{E_r,\,F_r,\,H_r\}_{r \ge 0}$ generate $Y_\hbar(\mathfrak{sl}_2)$
as an associative algebra. At $\hbar = 0$, the identification
$E_r \mapsto e \otimes t^r$, $F_r \mapsto f \otimes t^r$,
$H_r \mapsto h \otimes t^r$ recovers $U(\mathfrak{sl}_2[t])$.

\subsubsection*{The Drinfeld relations from the bar complex}

\begin{theorem}[Drinfeld presentation of $Y_\hbar(\mathfrak{sl}_2)$]
\ClaimStatusProvedHere
\label{thm:drinfeld-yangian-sl2}
The Yangian $Y_\hbar(\mathfrak{sl}_2)$, with $\hbar = 1/(k+2)$, is the
associative algebra generated by $\{E_r,\,F_r,\,H_r\}_{r \ge 0}$ subject
to the relations:
\begin{enumerate}[label=\textup{(D\arabic*)},leftmargin=3.5em]
\item\label{D1} $[H_r, H_s] = 0$ for all $r,s \ge 0$.

\item\label{D2}
$[H(u), E(v)] \;=\; \dfrac{2\hbar}{u-v}\bigl(E(u) - E(v)\bigr)$.

\item\label{D3} $[H(u), F(v)] \;=\; -\,\dfrac{2\hbar}{u-v}\bigl(F(u) - F(v)\bigr)$.

\item\label{D4} $[E(u), F(v)] \;=\; \dfrac{1}{u-v}\bigl(H(u) - H(v)\bigr)$.

\item\label{D5} $(u{-}v)\,[E(u), E(v)]
\;=\; \hbar\bigl(E(u)\,E(v) + E(v)\,E(u)\bigr)$
\textup{(}deformed Serre relation\textup{)}.

\item\label{D6} $(u{-}v)\,[F(u), F(v)]
\;=\; -\,\hbar\bigl(F(u)\,F(v) + F(v)\,F(u)\bigr)$
\textup{(}deformed Serre relation\textup{)}.
\end{enumerate}
\end{theorem}

\begin{proof}
Relations~\ref{D1}--\ref{D4} are extracted from the Gauss
decomposition of the RTT equation
$R_{12}(u{-}v)\,T_1(u)\,T_2(v) = T_2(v)\,T_1(u)\,R_{12}(u{-}v)$.
The key step is the factorisation $T(u) = F(u)\,K(u)\,E(u)$: the
diagonal entries produce~\ref{D1}, the off-diagonal--diagonal
cross-terms produce~\ref{D2}--\ref{D3}, and the purely off-diagonal
cross-terms produce~\ref{D4}.

For \ref{D1}: the RTT relations $[A(u),A(v)] = 0 = [D(u),D(v)]$
imply that $k_1(u),k_2(u)$ (and hence $H(u) = k_1(u)/k_2(u)$)
generate a commutative subalgebra.

For \ref{D2}: the RTT relation
$(u{-}v)[A(u),B(v)] = A(v)B(u) - A(u)B(v)$, rewritten via Gauss
decomposition in terms of $H(u)$ and $E(u)$, yields the stated
current-algebra relation. Expanding
$A(u) = k_1(u)(1 + F(u)\,k_2(u)\,k_1(u)^{-1}\,E(u))$ and
$B(u) = k_1(u)\,E(u)$,
the diagonal factor cancels between left and right sides, and the
residual identity in the off-diagonal is precisely~\ref{D2}.
Relation~\ref{D3} follows by the same argument applied to
$(u{-}v)[A(u),C(v)] = C(v)A(u) - C(u)A(v)$.

For \ref{D4}: the RTT relation
$(u{-}v)[B(u),C(v)] = D(v)A(u) - D(u)A(v)$
extracts, via Gauss decomposition, to
$(u{-}v)[E(u),F(v)] = H(u) - H(v)$.

For \ref{D5} and \ref{D6}: the RTT relation $[B(u),B(v)] = 0$
encodes $[E(u),E(v)]$ modulo diagonal factors. The
precise computation: $[B(u),B(v)] = 0$ with $B(u) = k_1(u)\,E(u)$
gives, after extracting the non-commuting diagonal,
$(u{-}v)[E(u),E(v)] = \hbar(E(u)E(v) + E(v)E(u))$.
The sign in~\ref{D6} follows from $[C(u),C(v)] = 0$ by the same
argument with $F$ in place of $E$.

Every relation is a direct algebraic consequence of the RTT
presentation, itself extracted from the degree-$2$ bar differential
and the Yang $R$-matrix $R(u) = u\,I + P$.
\end{proof}

\begin{remark}[Provenance from the bar complex]
\label{rem:drinfeld-bar-provenance}
The chain of identifications is:
\[
\text{OPE of } \widehat{\mathfrak{sl}}_2
\;\xrightarrow{\;\text{bar differential}\;}\;
r(z) = \frac{\hbar\,\Omega}{z}
\;\xrightarrow{\;\text{exponentiation}\;}\;
R(u) = u\,I + P
\;\xrightarrow{\;\text{RTT}\;}\;
\text{(D1)--(D6)}.
\]
The first arrow is (the $d\log$ kernel absorbs one pole from the
affine double-pole OPE); the second is the monodromy of the Kohno
connection on $\mathrm{Conf}_2^{\mathrm{ord}}(\bC)$ specialised to the
fundamental representation; the third is the Gauss decomposition. The
entire Drinfeld presentation is a derived algebraic consequence of the
OPE of $V_k(\mathfrak{sl}_2)$.
\end{remark}

\subsubsection*{Gauss decomposition}

\begin{proposition}[Gauss decomposition]
\ClaimStatusProvedHere
\label{prop:gauss-decomposition-sl2}
The transfer matrix $T(u) = \bigl(\begin{smallmatrix}
A(u) & B(u) \\ C(u) & D(u) \end{smallmatrix}\bigr)$
factors uniquely as
\[
T(u) \;=\;
\underbrace{\begin{pmatrix} 1 & 0 \\ F(u) & 1 \end{pmatrix}}_{F(u)}
\;\cdot\;
\underbrace{\begin{pmatrix} k_1(u) & 0 \\ 0 & k_2(u) \end{pmatrix}}_{K(u)}
\;\cdot\;
\underbrace{\begin{pmatrix} 1 & E(u) \\ 0 & 1 \end{pmatrix}}_{E(u)},
\]
with
$A(u) = k_1(u) + F(u)\,k_2(u)\,E(u)$,
$B(u) = k_1(u)\,E(u)$,
$C(u) = F(u)\,k_2(u)$,
$D(u) = k_2(u)(1 + F(u)\,E(u))$.
The Drinfeld Cartan current is the ratio
$H(u) = k_1(u)/k_2(u)$.
\end{proposition}

\begin{proof}
The upper-triangular unipotent factor is determined by
$E(u) = k_1(u)^{-1}\,B(u)$; the lower-triangular unipotent factor by
$F(u) = C(u)\,k_2(u)^{-1}$; the diagonal entries by
$k_2(u) = D(u) - C(u)\,A(u)^{-1}\,B(u)$ (the Schur complement) and
$k_1(u) = A(u) - B(u)\,D(u)^{-1}\,C(u)$. The quantum determinant
$\operatorname{qdet} T(u) = k_1(u)\,k_2(u{-}\hbar)$ is central in
$Y_\hbar(\mathfrak{sl}_2)$, confirming the factorisation is well-defined
over the Yangian.
\end{proof}

\subsubsection*{Twisted coproduct on Drinfeld generators}

\begin{theorem}[Twisted coproduct]
\ClaimStatusProvedHere
\label{thm:twisted-coproduct-sl2}
The twisted coproduct $\Delta_z \colon Y_\hbar(\mathfrak{sl}_2)
\to Y_\hbar(\mathfrak{sl}_2) \otimes Y_\hbar(\mathfrak{sl}_2)$
acts on the Drinfeld generators as:
\begin{align}
\label{eq:twisted-coprod-H}
\Delta_z(H(u)) &= H(u) \otimes H(u{-}z), \\
\label{eq:twisted-coprod-E}
\Delta_z(E(u)) &= E(u) \otimes 1 + H(u{-}z) \otimes E(u{-}z), \\
\label{eq:twisted-coprod-F}
\Delta_z(F(u)) &= 1 \otimes F(u{-}z) + F(u) \otimes H(u{-}z)^{-1}.
\end{align}
At $z = 0$ this recovers the standard Hopf coproduct of $Y_\hbar(\mathfrak{sl}_2)$.
\end{theorem}

\begin{proof}
The RTT coproduct is $\Delta_z(T(u)) = T(u) \dot\otimes T(u{-}z)$
(matrix multiplication in the auxiliary space, with spectral parameter
shifted by $z$ in the second tensor factor). Applying Gauss
decomposition to both sides:
\[
\Delta_z(T(u))
= F(u)\,K(u)\,E(u) \;\dot\otimes\; F(u{-}z)\,K(u{-}z)\,E(u{-}z).
\]
The diagonal entry gives
$\Delta_z(K(u)) = K(u) \otimes K(u{-}z)$,
whence~\eqref{eq:twisted-coprod-H} follows by taking the ratio
$k_1/k_2$. The off-diagonal entry
$\Delta_z(B(u)) = A(u) \otimes B(u{-}z) + B(u) \otimes D(u{-}z)$,
rewritten via $B = k_1\,E$ and $A = k_1 + F\,k_2\,E$,
yields~\eqref{eq:twisted-coprod-E} after extracting the $k_1$
factor. Equation~\eqref{eq:twisted-coprod-F} follows by the analogous
computation on $C(u) = F(u)\,k_2(u)$.
\end{proof}

\begin{remark}[Factorisation-algebraic origin of $\Delta_z$]
\label{rem:twisted-coproduct-origin}
The spectral parameter $z$ is the relative position along $\bR_t$:
it is the topological interval coordinate of the
$\mathrm{SC}^{\mathrm{ch,top}}$-structure. The twisted coproduct
$\Delta_z$ is the factorisation product of line operators at positions
$t_1$ and $t_2 = t_1 + z$ on the topological interval. The formula
$\Delta_z(T(u)) = T(u) \dot\otimes T(u{-}z)$ is the bar-coalgebra
deconcatenation specialised to degree~$2$ in the ordered bar
complex. At $z = 0$ the two insertion points collide and one recovers
the ordinary (untwisted) Hopf coproduct.
\end{remark}

\subsubsection*{PBW basis and Hilbert series}

\begin{theorem}[PBW theorem for $Y_\hbar(\mathfrak{sl}_2)$]
\ClaimStatusProvedHere
\label{thm:PBW-yangian-sl2}
The ordered monomials
\[
\bigl\{\,
f_{r_1}\cdots f_{r_a}\;\;
h_{s_1}\cdots h_{s_b}\;\;
e_{t_1}\cdots e_{t_c}
\;\bigm|\;
0 \le r_1 \le \cdots \le r_a,\;
0 \le s_1 \le \cdots \le s_b,\;
0 \le t_1 \le \cdots \le t_c
\,\bigr\}
\]
form a basis of $Y_\hbar(\mathfrak{sl}_2)$ over~$\bC$. The Hilbert series
with respect to the grading $\deg(X_r) = r$ is
\[
\operatorname{HS}_{Y_\hbar(\mathfrak{sl}_2)}(q)
\;=\; \prod_{d \ge 0} \frac{1}{(1-q^d)^3}.
\]
\end{theorem}

\begin{proof}
The PBW property follows from the Drinfeld relations~\ref{D1}--\ref{D6}
by a standard diamond lemma argument: the leading terms of
\ref{D2}--\ref{D6} (in the degree-lexicographic order
$f_{r_1}\cdots h_{s_1}\cdots e_{t_1}\cdots$) show that every
out-of-order monomial reduces to a sum of ordered monomials of
equal or lower total degree. The three families of generators
$\{e_r\}_{r \ge 0}$, $\{f_r\}_{r \ge 0}$, $\{h_r\}_{r \ge 0}$
each contribute a factor of $\prod_{d \ge 0}(1-q^d)^{-1}$ to the
Hilbert series, giving the stated product formula. The PBW
basis is compatible with the filtration
$F_{\le n}\,Y_\hbar(\mathfrak{sl}_2) = \langle X_{r_1}\cdots X_{r_k}
\mid r_1 + \cdots + r_k \le n \rangle$, and the associated graded
is $\gr\,Y_\hbar(\mathfrak{sl}_2) \simeq U(\mathfrak{sl}_2[t])$.
\end{proof}

\subsubsection*{Classical limit: comparison with $U(\mathfrak{sl}_2[t])$}

\begin{theorem}[Classical limit]
\ClaimStatusProvedHere
\label{thm:classical-limit-sl2}
At $\hbar = 0$, the Yangian $Y_\hbar(\mathfrak{sl}_2)$ degenerates to the
universal enveloping algebra of the current Lie algebra
$\mathfrak{sl}_2[t] = \mathfrak{sl}_2 \otimes \bC[t]$. Explicitly:
\begin{enumerate}[label=\textup{(\roman*)}]
\item Relations \ref{D2}--\ref{D3} at $\hbar = 0$ become
$[H(u), E(v)] = 0 = [H(u), F(v)]$ after clearing denominators
and passing to modes, yielding
$[h_r, e_s] = 2\,e_{r+s}$, $[h_r, f_s] = -2\,f_{r+s}$
the current-algebra Lie brackets
$[h \otimes t^r,\, e \otimes t^s] = 2\,(e \otimes t^{r+s})$.
\item Relation \ref{D4} at $\hbar = 0$ becomes
$[E_r, F_s] = H_{r+s}$, i.e.\
$[e \otimes t^r,\, f \otimes t^s] = h \otimes t^{r+s}$.
\item Relations \ref{D5}--\ref{D6} at $\hbar = 0$ become
$[E(u),E(v)] = 0 = [F(u),F(v)]$ after clearing $(u{-}v)$
the Serre relation for
$\mathfrak{sl}_2[t]$ (abelian root spaces).
\item The twisted coproduct $\Delta_z$ at $\hbar = 0$
degenerates to the cocommutative coproduct
$\Delta(X_r) = X_r \otimes 1 + 1 \otimes X_r$
of $U(\mathfrak{sl}_2[t])$.
\end{enumerate}
The PBW filtration on $Y_\hbar(\mathfrak{sl}_2)$ has associated graded
$\gr\,Y_\hbar(\mathfrak{sl}_2) \cong U(\mathfrak{sl}_2[t])$,
confirming that the Yangian is a filtered deformation of the current
algebra, with deformation parameter $\hbar = 1/(k+2)$ sourced by the
level of the affine input $V_k(\mathfrak{sl}_2)$.
\end{theorem}

\begin{proof}
Setting $\hbar = 0$ in relations~\ref{D1}--\ref{D6} is immediate.
For (i): \ref{D2} becomes $[H(u),E(v)] = 0$ at $\hbar = 0$;
expanding $H(u) = 1 + \hbar\sum H_r u^{-r-1}$ and taking
the leading $\hbar$-linear term recovers
$[h_r, e_s] = 2\,e_{r+s}$.
For (ii): \ref{D4} is $\hbar$-independent.
For (iii): the right-hand side of \ref{D5} vanishes at $\hbar = 0$,
giving $[E(u),E(v)] = 0$.
For (iv): the shift $u \mapsto u - z$ in the second tensor factor of
$\Delta_z$ is the spectral-parameter manifestation of the coproduct;
at $\hbar = 0$ the modes become primitive.
The associated-graded statement follows from the PBW theorem
(Theorem~\ref{thm:PBW-yangian-sl2}).
\end{proof}

\begin{remark}[The deformation parameter and the level]
\label{rem:hbar-level}
The identification $\hbar = 1/(k+2)$ has a precise bar-complex origin:
the collision residue $r(z) = \hbar\,\Omega/z$ inherits $\hbar$ from
the normalisation of the affine OPE $J^a(z)\,J^b(w) \sim
k\,\kappa^{ab}/(z{-}w)^2 + f^{ab}{}_c\,J^c/(z{-}w)$, where the
$d\log$ kernel absorbs one power and the remaining pole is
rescaled by $1/(k + h^\vee)$ with $h^\vee = 2$ for $\mathfrak{sl}_2$.
At the critical level $k = -h^\vee = -2$, the parameter $\hbar$
diverges and the Yangian construction degenerates: this is the
ordered-bar reflection of the Sugawara obstruction at critical level.
\end{remark}


\section{The double bar and what it recovers}
\label{sec:double-bar}

The ordered Koszul dual of $\widehat{\mathfrak{g}}_k$ is
$Y_\hbar(\mathfrak{g})$. What is the ordered Koszul dual of
$Y_\hbar(\mathfrak{g})$?

If ordered Koszul duality were an involution (as the full
$\mathrm{SC}^{\mathrm{ch,top}}$ duality is
(Theorem~\ref{thm:two-colour-double-kd} below), then
$\Barchord(Y_\hbar(\mathfrak{g}))$ would recover
$\widehat{\mathfrak{g}}_k$. It does not.

The ordered double bar recovers $U(\mathfrak{g}[t])$
\emph{without} the central extension. The level~$k$ is invisible
because it enters the Yangian as $\hbar = 1/(k{+}h^\vee)$, not as an
OPE pole: the spectral $R$-matrix $R(u) = 1 + \hbar\,\Omega/u +
O(\hbar^2)$ has only a simple pole, so the $d\log$ propagator of the
ordered bar extracts the Lie bracket of $\mathfrak{g}[t]$ and nothing
more. The double pole that \emph{creates} the central extension lives
in the holomorphic OPE on $\mathrm{FM}_2(\mathbb{C})$ (the closed
colour), and no iteration of the open bar can reach it.

\begin{theorem}[Central extension is invisible to the ordered
double bar]
\ClaimStatusProvedHere
\label{thm:central-extension-invisible}
\label{thm:double-bar-sl2}
Let $\mathfrak{g}$ be a simple Lie algebra and let
$A = \widehat{\mathfrak{g}}_k$ at non-critical level. The ordered
double Koszul dual satisfies
\[
(A^{!}_{\mathrm{line}})^{!}_{\mathrm{line}}
\;\simeq\;
U(\mathfrak{g}[t]),
\qquad
\textup{not}\;\;
\widehat{\mathfrak{g}}_k.
\]
\end{theorem}

\begin{proof}
The level $k$ of $\widehat{\mathfrak{g}}_k$ enters the ordered Koszul dual
$Y_\hbar(\mathfrak{g})$ as the scalar $\hbar = 1/(k + h^\vee)$.
The holomorphic double pole $k\,\delta^{ab}/(z-w)^2$ (the source
of the central term $k\,n\,\delta_{n+m,0}$ in the mode algebra)
does not survive as a double pole in the spectral variable: the
Yangian $R$-matrix $R(u) = 1 + \hbar\,r(u) + O(\hbar^2)$ has its
singularity at $u = 0$ with a \emph{simple} pole $r(u) = \Omega/u$,
where $\Omega = \sum_a T^a \otimes T^a$ is the Casimir tensor.

The ordered bar differential of $Y_\hbar(\mathfrak{g})$ uses the
propagator $d\log(u_i - u_j)$ on
$\mathrm{Conf}_n^{<}(\mathbb{A}^1_u)$. By the $d\log$ absorption
principle, a simple spectral pole yields the zeroth product
$a_{(0)}b$; a double pole would yield a derivative term
$a_{(1)}b$ carrying the central extension. But the spectral OPE has
no double pole. Degree by degree:

\smallskip\noindent\emph{Degree $1$.}
Yangian generators $E_n, F_n, H_n$ map to current modes $J^a_n$.
At the first bar stage, the cohomology is therefore the
desuspended current Lie algebra:
\begin{equation}\label{eq:bar-degree1-yangian}
  H^*(s^{-1}\bar{Y}_\hbar, d_{\mathrm{bar}})
  \;=\;
  s^{-1}\mathfrak{g}
  \qquad
  \textup{(concentrated in cohomological degree~$1$).}
\end{equation}

\smallskip\noindent\emph{Degree $2$.}
The spectral OPE
$[E(u), F(v)] \sim \hbar/(u-v)\cdot(\ldots)$ has only simple poles, so
\[
d_{\mathrm{bar}}(s^{-1}E_n \otimes s^{-1}F_m)
\;=\;
s^{-1}[E_n, F_m]
\;=\;
s^{-1}H_{n+m},
\]
with no central term. For the $E$--$E$ component, the pole at
$u - v = \hbar$ is shifted off the collision diagonal, giving
$d_{\mathrm{res}}(s^{-1}E \otimes s^{-1}E) = 0$, i.e.\
$[e,e] = 0$ in $\mathfrak{g}[t]$.

\smallskip\noindent\emph{Degree $3$.}
$d^2 = 0$ reduces to the classical Yang--Baxter equation for
$r(u) = \Omega/u$, which holds by the Jacobi identity for
$\mathfrak{g}$.

\smallskip
Hence the bar differential recovers only the Lie bracket of
$\mathfrak{g}[t]$, not the central extension.
\end{proof}

The loss is not a failure of the machinery; it is a precise
diagnosis of what each colour carries.

\begin{theorem}[Two-colour double Koszul duality is involutive]
\ClaimStatusProvedHere
\label{thm:two-colour-double-kd}
Let $A$ be a strongly admissible $\Einf$-chiral algebra on a curve
$X$. The full $\mathrm{SC}^{\mathrm{ch,top}}$ double Koszul duality
satisfies $(A^!)^! \simeq A$, recovering all central extensions and
higher-pole data.
\end{theorem}

\begin{proof}
The full duality operates on both components: the \emph{chiral component}
(symmetric bar $\barB^{\Sigma}(A)$ on $\mathrm{Ran}(X)/\Sigma_n$,
dualised by $D_{\mathrm{Ran}}$) carries the holomorphic OPE including
the double pole; the \emph{coassociative component} (ordered bar
$\Barchord(A)$ on $\mathrm{Conf}_n^{<}(\mathbb{R})$,
dualised by linear duality) carries the spectral data. Verdier/Ran
duality preserves the full pole structure, so the double pole
survives. The bar--cobar involution
$\Omega(\Barch(\Omega(\Barch(A)))) \simeq A$
(Hypothesis~\ref{hyp:inputs}(I), applied twice) then recovers $A$ in
its entirety. For $\widehat{\mathfrak{g}}_k$: the chiral Verdier
dual is $\widehat{\mathfrak{g}}_{-k-2h^\vee}$ (Feigin--Frenkel dual level),
and the second application returns to $\widehat{\mathfrak{g}}_k$ with the
central extension intact.
\end{proof}

\begin{corollary}[Non-redundancy of the two colours]
\ClaimStatusProvedHere
\label{cor:two-colours-non-redundant}
\index{double bar!ordered}
\index{central extension!invisible to ordered bar}
For any $\Einf$-chiral algebra $A$ with non-trivial OPE poles, the
open and closed bar complexes carry genuinely non-redundant
information. Neither colour alone recovers $A$ by iterating Koszul
duality: the ordered double bar loses the central extension
\textup{(Theorem~\ref{thm:central-extension-invisible})}, while the
chiral symmetric bar does not determine the $R$-matrix
monodromy for the ordered-to-unordered descent
\textup{(Proposition~\ref{sec:r-matrix-descent-vol1})}.
\end{corollary}

\begin{proof}
The ordered colour alone is insufficient by
Theorem~\ref{thm:central-extension-invisible}: iterating the ordered
bar-cobar construction forgets the double-pole central extension data.
The symmetric chiral colour alone is also insufficient by
Proposition~\ref{sec:r-matrix-descent-vol1}: it retains the holomorphic
OPE data but does not recover the ordered $R$-matrix monodromy needed
for descent from the ordered to the unordered picture.
These two failures are complementary, so the open and closed colours
carry genuinely different information and both are needed to recover
$A$.
\end{proof}

\begin{remark}[The chiral--coassociative architecture at a glance]
\label{rem:double-bar-architecture}
The double-bar computation for $\widehat{\mathfrak{sl}}_{2,k}$ distils the
architecture:
\[
\begin{array}{c|c|c}
& \text{Chiral (holomorphic)} & \text{Coassociative (topological)} \\
\hline
\text{Pole order} & \text{double: } k/(z-w)^2
 & \text{simple: } \hbar/(u-v) \\
\text{Extracted datum} & \text{central extension}
 & \text{Lie bracket of } \mathfrak{g}[t] \\
\text{Double KD output} & \widehat{\mathfrak{g}}_k \text{ (full)}
 & U(\mathfrak{g}[t]) \text{ (no central ext.)}
\end{array}
\]
The two components of the $E_1$-chiral bar carry genuinely non-redundant
information. The coassociative component sees the Yangian; the chiral component sees
the level. Neither alone recovers the algebra; the central extension
is the datum that forces both components to be present.
\end{remark}


\section{The Virasoro ordered bar complex and the gravitational Yangian}
\label{subsec:virasoro-ordered-bar}

The affine Yangian $Y_\hbar(\mathfrak{sl}_2)$ has three independent
generator families $E(u)$, $F(u)$, $H(u)$. The Virasoro algebra has
one generator~$T$. What happens to the Yangian when the input has a
single generator of spin greater than~$1$?

The answer is forced by two structural features (a triple-pole
collision residue and a single-generator mode collapse) whose
interplay produces an object that is simultaneously trivial as a
Lie algebra and infinitely non-formal as an $\Ainf$ algebra.
The Virasoro OPE is
\begin{equation}\label{eq:virasoro-ope-ordered}
T(z)\,T(w) \;\sim\;
\frac{c/2}{(z-w)^4}
+ \frac{2T(w)}{(z-w)^2}
+ \frac{\partial T(w)}{z-w}.
\end{equation}

\subsubsection*{Beat 1: the triple pole}

The bar complex uses the propagator $d\log(z-w) = dz/(z-w)$, which
absorbs one power of the pole in each OPE term. The quartic
pole becomes a triple pole; the double pole becomes a simple pole; the
simple pole is absorbed entirely and feeds the bar differential
$d_{\Barch}[\susp^{-1}T \,|\, \susp^{-1}T]
= \susp^{-1}(T_{(0)}T) = \susp^{-1}(\partial T)$.

\begin{proposition}[Virasoro collision residue; \ClaimStatusProvedHere]
\label{prop:vir-collision-residue}
The collision residue $r$-matrix of\/ $\mathrm{Vir}_c$ is
\begin{equation}\label{eq:vir-r-collision}
r^{\mathrm{coll}}(z)
\;=\;
\frac{c/2}{z^3}\,\mathbf{1}\otimes\mathbf{1}
\;+\;
\frac{2T}{z}\otimes\mathbf{1}.
\end{equation}
This is a \textbf{triple-pole} $r$-matrix: the leading singularity
$z^{-3}$ is unprecedented in the affine case, where
$r^{\mathrm{coll}}_{\mathrm{KM}}(z) = k\,\Omega_{\mathfrak{g}}/z$ is a
simple pole
\textup{(}Theorem~\textup{\ref{thm:heisenberg-ordered-bar})}.
The two depth components (the central scalar
$(c/2)/z^3$ and the field-valued $2T/z$) play entirely different
roles: the central term controls curvature, while the field-valued
term generates the non-associative cascade.
\end{proposition}

\begin{proof}
The Virasoro OPE has singular terms of orders $4$, $2$, and $1$.
Passing to the collision-residue kernel removes one pole order because
the bar propagator contributes a $d\log$ factor, so these become poles
of orders $3$, $1$, and $0$.
The regular part is absorbed into the ordinary bar differential
$d_{\Barch}$.
The remaining singular pieces are exactly the two terms displayed in
\eqref{eq:vir-r-collision}, with the cubic pole coming from the
central term and the simple pole from the field $T$.
\end{proof}

The collision residue satisfies the classical Yang--Baxter equation.

\begin{proposition}[Virasoro CYBE; \ClaimStatusProvedHere]
\label{prop:vir-CYBE-ordered}
The collision residue~\eqref{eq:vir-r-collision} satisfies
\begin{equation}\label{eq:vir-CYBE-ordered}
[r^{12}(u),\, r^{13}(u+v)]
+ [r^{12}(u),\, r^{23}(v)]
+ [r^{13}(u+v),\, r^{23}(v)]
\;=\; 0.
\end{equation}
\end{proposition}

\begin{proof}
Write
$r(z)=\frac{c/2}{z^3}\,\mathbf 1\otimes\mathbf 1+\frac{2T}{z}
+\frac{2}{z}\,\mathbf 1\otimes T$.
The cubic central term commutes with every tensor factor, so it
contributes nothing to the three commutators in
\eqref{eq:vir-CYBE-ordered}.
For the simple-pole part, the Laplace transform identifies the CYBE
with the Jacobi identity for the Virasoro PVA bracket
$\{T{}_\lambda T\}=\partial T+2T\lambda+\frac{c}{12}\lambda^3$.
The three commutators in the CYBE are precisely the three Jacobi terms,
so their sum vanishes.
\end{proof}

\subsubsection*{Beat 2: the collapse}

Now the single-generator structure exerts its force.

\begin{theorem}[Gravitational Yangian collapse; \ClaimStatusProvedHere]
\label{thm:grav-yangian-collapse}
Let $c' = 26 - c$ denote the Koszul dual central charge. Define the
level-$r$ generators on\/ $\mathrm{Vir}_{c'}$ by
\begin{equation}\label{eq:grav-yangian-generators}
\mathsf{T}_n^{(r)}
\;:=\;
\binom{n+r-1}{r}\,L_{-n-r},
\qquad
\Omega_r
\;=\;
\sum_{n \in \mathbb{Z}}
\mathsf{T}_n^{(r)} \otimes L_n,
\end{equation}
so that $r(z) = \sum_{r \ge 0}\Omega_r\,z^{-r-1}$.
Then:
\begin{enumerate}[label=\textup{(\roman*)}]
\item \emph{Linear dependence.}\enspace
For every $r \ge 1$ and $n \in \mathbb{Z}$,
$\mathsf{T}_n^{(r)} = \binom{n+r-1}{r}\,L_{-n-r}$
is a scalar multiple of a single Virasoro mode. There is no
independent Yangian partner at any level.

\item \emph{Lie-algebraic isomorphism.}\enspace
$Y^{\mathrm{dg}}(\mathrm{Vir}_c) \cong \mathrm{Vir}_{c'}$
as Lie algebras. The entire Yangian Lie structure collapses.

\item \emph{All-level bracket.}\enspace
\begin{equation}\label{eq:grav-yangian-bracket}
[\mathsf{T}_m^{(r)},\,\mathsf{T}_n^{(s)}]
\;=\;
\binom{m+r-1}{r}\binom{n+s-1}{s}\,\bigl[
(n{+}s{-}m{-}r)\,L_{-m-n-r-s}
- \tfrac{c'}{12}((m{+}r)^3 - (m{+}r))\,\delta_{m+n+r+s,0}
\bigr],
\end{equation}
which reduces to a single Virasoro commutator.
\end{enumerate}
\end{theorem}

\begin{proof}
Part~(i) is the entire argument: because $T$ has conformal weight~$2$,
the mode $L_{-n-r}$ depends only on the sum $n+r$. The level
index~$r$ splits this sum but adds no new degree of freedom. In the
affine Yangian, level-$1$ generators $(e_1, f_1, h_1)$ are linearly
independent of $(e_0, f_0, h_0)$ because $\dim\mathfrak{g} \ge 3$
generators of spin~$1$ each have an independent mode family. Here
the single spin-$2$ generator forces all levels into the same
Virasoro mode lattice.
Part~(ii) follows: the generators $\{\mathsf{T}_n^{(r)}\}_{n,r}$
span $\{L_m\}_{m \in \mathbb{Z}}$, which is $\mathrm{Vir}_{c'}$.
Part~(iii) is a direct computation using the Virasoro commutator
$[L_{-m-r}, L_{-n-s}] = (n{+}s{-}m{-}r)\,L_{-m-n-r-s}
+ \frac{c'}{12}((m{+}r)^3 - (m{+}r))\,\delta_{m+n+r+s,0}$.
\end{proof}

The collapse is absolute: the gravitational Yangian has no independent
higher-level data, no non-primitive coproduct, no Drinfeld
partners. As a Lie algebra it is just $\mathrm{Vir}_{26-c}$.
But the story does not end here.

\subsubsection*{Beat 3: the infinite tower}

Despite the Lie-algebraic triviality, the bar complex is infinitely
non-formal. The triple pole forces non-vanishing $\Ainf$ operations
at every degree, through a cascade that cannot terminate.

\begin{theorem}[Virasoro non-formality; \ClaimStatusProvedHere]
\label{thm:vir-non-formality}
For generic central charge~$c$, the $\Ainf$ operation $m_k$ of the
ordered bar complex $\Barch(\mathrm{Vir}_c)$ is nonzero for all
$k \ge 3$. The Virasoro algebra is class~$\mathbf{M}$
\textup{(}infinite shadow depth\textup{)}.
\end{theorem}

\begin{proof}
The degree-$3$ associator is computed directly from the PVA
$\lambda$-bracket. The Stasheff identity at degree~$3$ requires
$m_3 = -A_3$ where
\[
A_3(T,T,T;\,\lambda_{12},\lambda_{23})
\;=\;
-\{T{}_{\lambda_{23}}\{T{}_{\lambda_{12}} T\}\}
\;=\;
\tfrac{c}{12}\,\lambda_{23}^3\,(2\lambda_{12} + \lambda_{23})
\;+\;(\text{$T$-dependent terms}).
\]
The scalar contact term $\frac{c}{12}\lambda_{23}^3(2\lambda_{12}
+ \lambda_{23})$ is nonzero for $c \neq 0$, and the $T$-dependent
terms (from the double-pole coefficient $2T$) are nonzero for
\emph{all}~$c$ including $c = 0$. Hence $m_3 \neq 0$.

The Stasheff cascade is now forced. At degree $k+1$, the identity
$\sum_{i+j = k+1} m_i \circ m_j = 0$ produces a nonzero obstruction
from compositions of lower-degree operations; the unique solution is
$m_{k+1} \neq 0$. The triple pole ensures this cascade never
stabilises: the pole depth~$3$ exceeds the formality threshold
(pole depth~$1$), and each degree draws on the full pole structure
to generate new non-associativity.
\end{proof}

\begin{remark}[What non-formality is not]
\label{rem:non-formality-not-E1}
Three distinct things are being kept apart here.
\begin{enumerate}[label=(\roman*),nosep]
\item \emph{Non-formality ($m_k \neq 0$ for all $k$)} is a property
of the $A_\infty$-chiral structure on $\mathrm{Vir}_c$ transferred to
bar cohomology: it says the transferred higher operations
$m_k$ are not formal. It does \emph{not}
say $\mathrm{Vir}_c$ is non-Koszul. $\mathrm{Vir}_c$ is
chirally Koszul at every non-critical central charge: the bar
cohomology $H^\bullet(\Barch(\mathrm{Vir}_c))$ is concentrated
in bar degree~$1$ by PBW universality and single-generator
freeness. Koszulness (bar concentration) and $A_\infty$
formality (vanishing of transferred higher operations) are
different properties. Virasoro is Koszul but $A_\infty$
non-formal.
\item \emph{$E_1$ versus $E_\infty$} is the question of locality.
$\mathrm{Vir}_c$ is and remains $\Einf$-chiral: local, with
$\Sigma_n$-equivariant factorization. The ordered
bar complex applies the $E_1$-Koszul duality \emph{functor}
to the $\Einf$ input; non-formality of the output does not make
the input non-local. Having OPE poles never breaks $\Einf$.
\item \emph{The $R$-matrix is derived}, not independent input:
$R(z)$ is obtained as the monodromy of the Kohno connection
built from the collision residue $r^{\mathrm{coll}}$, and the
$d\log$ absorption step sets the correct pole
structure.
\end{enumerate}
\end{remark}

\subsubsection*{Curvature}

\begin{proposition}[Gravitational Yangian curvature; \ClaimStatusProvedHere]
\label{prop:grav-yangian-curvature}
The curvature of the gravitational dg-shifted Yangian is
\begin{equation}\label{eq:grav-yangian-kappa}
\kappa\bigl(Y^{\mathrm{dg}}(\mathrm{Vir}_c)\bigr)
\;=\;
\frac{26 - c}{2}.
\end{equation}
Two central charges are distinguished:
\begin{enumerate}[label=\textup{(\roman*)}]
\item \emph{Critical string ($c = 26$):}\enspace
$\kappa = 0$. The gravitational Yangian is flat; $F_g = 0$ for all
$g \ge 1$ \textup{(UNIFORM-WEIGHT)}. This is anomaly cancellation.

\item \emph{Koszul self-dual ($c = 13$):}\enspace
$\kappa = 13/2$. The Koszul involution $c \mapsto 26-c$ fixes
$c = 13$, so $\mathrm{Vir}_{13}^! = \mathrm{Vir}_{13}$.
The genus tower is invariant under $c \mapsto 26 - c$
\textup{(}gravitational S-duality\textup{)} with
$F_g(\mathrm{Vir}_{13}) = \frac{13}{2}\lambda_g^{\mathrm{FP}}
\textup{(}UNIFORM-WEIGHT\textup{)}
\neq 0$.
\end{enumerate}
\end{proposition}

\begin{proof}
Volume~I Theorem~D gives the Virasoro curvature formula
$\kappa(\mathrm{Vir}_c)=c/2$.
By Theorem~\ref{thm:grav-yangian-collapse}, the gravitational
dg-shifted Yangian is identified with the Koszul dual
$\mathrm{Vir}_{26-c}$.
Applying the same curvature formula to that dual central charge gives
\[
\kappa\bigl(Y^{\mathrm{dg}}(\mathrm{Vir}_c)\bigr)
=\kappa(\mathrm{Vir}_{26-c})
=\frac{26-c}{2}.
\]
Substituting $c=26$ and $c=13$ yields the critical and self-dual cases
listed in the statement.
\end{proof}

\subsubsection*{The gauge-gravity dichotomy}

The three beats assemble into a single structural dichotomy.

\begin{corollary}[Gauge-gravity complexity dichotomy; \ClaimStatusProvedHere]
\label{cor:gauge-gravity-dichotomy-ordered}
The two standard HT families occupy opposite corners of the
$(\Ainf,\,\Delta)$-complexity plane:
\[
\begin{array}{c|cc}
& \Ainf\textup{ products } m_k\;(k \ge 3) &
 \textup{coproduct } \Delta_z \\
\hline
\textup{Gauge \textup{(}KM, class~$\mathbf{L}$\textup{)}}
 & m_k = 0\;\textup{(formal)}
 & \textup{nontrivial Yangian } \Delta_z \\[3pt]
\textup{Gravity \textup{(}Vir, class~$\mathbf{M}$\textup{)}}
 & m_k \neq 0\;\forall\, k\;\textup{(non-formal)}
 & \textup{strictly primitive } \Delta_z
\end{array}
\]
The dichotomy is controlled by two inputs:
\begin{enumerate}[label=\textup{(\arabic*)}]
\item \emph{Pole order.}\enspace
Double OPE pole $\Rightarrow$ simple collision residue
$\Rightarrow$ class~$\mathbf{L}$ $\Rightarrow$ formal.
Quartic OPE pole $\Rightarrow$ triple collision residue
$\Rightarrow$ class~$\mathbf{M}$ $\Rightarrow$ non-formal.

\item \emph{Generator count.}\enspace
$\dim\mathfrak{g}$ generators of spin~$1$
$\Rightarrow$ independent Yangian partners
$\Rightarrow$ nontrivial coproduct.
One generator of spin~$2$
$\Rightarrow$ collapse
\textup{(}Theorem~\textup{\ref{thm:grav-yangian-collapse})}
$\Rightarrow$ primitive coproduct.
\end{enumerate}
Drinfeld--Sokolov reduction transports class~$\mathbf{L}$ to
class~$\mathbf{M}$: the Sugawara construction manufactures the quartic
pole from the affine double pole, and the resolvent tree formula
$m_k^{\mathcal{W}} = \sum_{\mathrm{Catalan}}
(m_2^{\mathrm{aff}} \circ h_{\mathrm{DS}})$
generates the infinite $\Ainf$ tower from the single binary product.
\end{corollary}

\begin{proof}
The $\Ainf$ column: class~$\mathbf{L}$ formal by pole-order
classification; class~$\mathbf{M}$ non-formal by
Theorem~\ref{thm:vir-non-formality}. The coproduct column: the
Yangian coproduct is non-primitive for gauge theory
(Theorem~\ref{thm:km-yangian}) and strictly primitive for
gravity by Theorem~\ref{thm:grav-coproduct-primitive}.
\end{proof}

\subsubsection*{Gravitational coproduct primitivity}

The corollary's coproduct column rests on a structural
vanishing that deserves its own statement. The DS-transferred
coproduct on $\mathrm{Vir}_c$ has no non-primitive corrections at
any degree; the mechanism is ghost-number filtration, not a
cancellation.

\begin{theorem}[Gravitational coproduct primitivity; \ClaimStatusProvedHere]
\label{thm:grav-coproduct-primitive}
Let\/ $\Delta_z^{\mathrm{Vir}}$ denote the coproduct on the
Virasoro algebra obtained by Drinfeld--Sokolov reduction from the
affine coproduct on~$\widehat{\mathfrak{sl}}_2$. Then\/
$\Delta_z^{\mathrm{Vir}}$ is strictly primitive at all degrees:
\begin{equation}\label{eq:grav-coproduct-primitive}
(p \otimes p) \circ \Delta_z^{\mathrm{Vir}}(T)
\;=\;
T \otimes 1 + 1 \otimes T,
\end{equation}
where $p$ is the projection to ghost number~$0$.
\end{theorem}

\begin{proof}
The BRST charge $h = c + \phi\, b$ carries ghost degree~$-1$. Every
HPL tree correction to the coproduct at degree~$r \ge 2$ passes through
$h$-decorated internal edges, producing elements of ghost number
$\neq 0$ on at least one tensor factor. Source trees (root on
the left) contribute ghost~$+1$; target trees (root on the right)
contribute ghost~$-1$ on one factor. In both cases,
$p \otimes p$ kills the correction.

The vanishing is not an accident of low degree. The ghost-number
filtration is an exact functor on the category of transferred
operations, and DS reduction preserves the filtration strictly.
The statement holds at all degrees simultaneously.
\end{proof}

\begin{remark}[The $(m,\Delta)$-complexity plane]
\label{rem:m-delta-complexity}
The contrast between product and coproduct complexity is absolute.
The transferred product $m_z^{\mathrm{Vir}}$ is infinitely
complicated: the full $\Ainf$ tower is nonzero at all degrees
(class~$\mathbf{M}$, shadow depth $r_{\max} = \infty$).
The transferred coproduct is trivial. In the
$(m,\Delta)$-complexity plane, gravity sits at $(\infty, 0)$:
infinite product complexity, zero coproduct complexity.
Gauge theory (affine Kac--Moody) sits at
$(3, \text{nonzero})$: finite product complexity (class~$\mathbf{L}$),
nontrivial coproduct. The Sugawara embedding confirms this:
the affine coproduct applied to the Sugawara stress tensor
$T = \kappa^{ab}\normord{J_a J_b}/(2(k + h^\vee))$ produces a
nontrivial cross term
\begin{equation}\label{eq:sugawara-coproduct-cross}
\Delta_z^{\mathrm{KM}}(T) - T \otimes 1 - 1 \otimes T
\;=\;
\frac{\Omega}{(k + h^\vee)\, z},
\qquad
\Omega = \kappa^{ab}\, J_a \otimes J_b.
\end{equation}
Primitivity is a genuine consequence of DS reduction, not a
triviality of Virasoro. The physical content: gravity does not
produce particles. It braids them, curves their propagators, and
dresses them with an infinite tower of higher products, but the
coproduct, which governs particle creation from the vacuum, is
trivial. This is the coalgebraic shadow of the no-hair theorem:
the gravitational sector has no nontrivial coalgebra structure at
any degree.
\end{remark}

\begin{remark}[Summary of the gravitational dg-shifted Yangian]
\label{rem:grav-yangian-summary}
$Y^{\mathrm{dg}}(\mathrm{Vir}_c)$ is $\mathrm{Vir}_{26-c}$ equipped
with:
\textup{(a)}~the infinite $\Ainf$ tower $\{m_k\}_{k \ge 3}$
\textup{(}class~$\mathbf{M}$\textup{)};
\textup{(b)}~primitive coproduct
$\Delta_z(x) = \tau_z(x) \otimes 1 + 1 \otimes x$;
\textup{(c)}~triple-pole collision residue
$r(z) = (c'/2)/z^3 + 2T/z$;
\textup{(d)}~curvature $\kappa = (26-c)/2$.
The Lie-algebraic Yangian structure is trivial (all levels collapse);
the entire nontrivial content resides in the $\Ainf$ tower. The
gravitational Yangian is a \emph{curved $\Ainf$-deformation}
of $U(\mathrm{Vir}_{26-c})$, not a filtered deformation of
$U(\mathrm{Witt}[t])$.
\end{remark}



\section{The universal Kac--Moody Yangian theorem}
\label{sec:km-yangian}

The full nonabelian Kac--Moody family yields the definitive result. For every simple Lie algebra~$\mathfrak{g}$
and non-critical level $k\neq -h^\vee$, the $E_1$-chiral Koszul
dual is the Yangian $Y_\hbar(\mathfrak{g})$. This is the
ordered analogue of Theorem~\ref{thm:universal-kac-moody-koszul},
which gives the symmetric Koszul duality
$\widehat{\mathfrak{g}}_k^!\simeq
\widehat{\mathfrak{g}}_{-k-2h^\vee}$.

All affine Kac--Moody algebras $\widehat{\mathfrak{g}}_k$ are
$E_\infty$-chiral: they are local vertex algebras with
$\Sigma_n$-equivariant factorization structure. The OPE poles
do not break $E_\infty$-locality.

\subsection{The ordered bar complex of
$\widehat{\mathfrak{g}}_k$}
\label{subsec:km-ordered-bar}

Fix a simple Lie algebra~$\mathfrak{g}$ with invariant form
$(\cdot,\cdot)$ normalised so that $(\theta|\theta)=2$
for the highest root~$\theta$, dual Coxeter number~$h^\vee$,
and structure constants $f^{ab}{}_{c}$ in an orthonormal basis
$\{T^a\}_{a=1}^{\dim\mathfrak{g}}$.
Let $\widehat{\mathfrak{g}}_k$ be the affine vertex algebra at
level $k\neq -h^\vee$, with OPE
\begin{equation}\label{eq:km-ope-ordered}
J^a(z)\,J^b(w)
\;\sim\;
\frac{k\,(T^a,T^b)}{(z-w)^2}
\;+\;
\frac{f^{ab}{}_{c}\,J^c(w)}{z-w}.
\end{equation}
Write $\hbar = 1/(k+h^\vee)$ for the deformation parameter
(well-defined away from critical level) and
\[
\Omega
\;=\;
\sum_{a=1}^{\dim\mathfrak{g}} T^a\otimes T^a
\;\in\;
\mathfrak{g}\otimes\mathfrak{g}
\]
for the quadratic Casimir tensor.

\begin{theorem}[Universal Kac--Moody Yangian theorem;
\ClaimStatusProvedHere]
\label{thm:km-yangian}
For every simple Lie algebra~$\mathfrak{g}$ and non-critical
level $k\neq -h^\vee$, the $E_1$-chiral Koszul dual of the
affine vertex algebra $\widehat{\mathfrak{g}}_k$ is the
dg-shifted Yangian:
\begin{equation}\label{eq:km-yangian-duality}
A^{!,\mathrm{ord}}_{\mathrm{line}}
(\widehat{\mathfrak{g}}_k)
\;\simeq\;
Y^{\mathrm{dg}}_\hbar(\mathfrak{g})
\;\simeq\;
Y_\hbar(\mathfrak{g}),
\qquad
\hbar = \frac{1}{k+h^\vee},
\end{equation}
where the second equivalence is the spectral Drinfeld
strictification
\textup{(}Theorem~\textup{\ref{thm:complete-strictification-v1})},
which identifies the dg-shifted Yangian with the classical
Yangian as strict associative algebras.
The duality is:
\begin{enumerate}[label=\textup{(\arabic*)}]
\item functorial in~$\mathfrak{g}$
\textup{(}respects Lie algebra homomorphisms\textup{)};

\item compatible with the symmetric Koszul duality
$\widehat{\mathfrak{g}}_k^!\simeq
\widehat{\mathfrak{g}}_{-k-2h^\vee}$
via $R$-twisted $\Sigma_n$-descent
\textup{(}Proposition~\textup{\ref{sec:r-matrix-descent-vol1})};

\item valid for all nine infinite families and
exceptional types \textup{(}Table~\textup{\ref{tab:km-yangian-data})}.
\end{enumerate}
\end{theorem}

\begin{proof}
The proof has five steps.

\emph{Step~1: Collision residue.}
The OPE~\eqref{eq:km-ope-ordered} has a double pole
$k\,(T^a,T^b)/(z-w)^2$ and a simple pole
$f^{ab}{}_{c}\,J^c(w)/(z-w)$. The bar kernel
$d\log(z-w)$ absorbs one power, so the collision
residue on $\mathrm{Conf}_2^{\mathrm{ord}}(\bC)$ is
\begin{equation}\label{eq:km-collision-residue}
r_{12}(z)
\;=\;
\frac{\hbar\,\Omega_{12}}{z}
\;+\;
(r_0)_{12},
\end{equation}
where $r_0 = \sum_{a,b,c} f^{ab}{}_{c}\,T^a\wedge T^b\otimes T_c$
encodes the Lie bracket and $\hbar\,\Omega/z$ is the spectral
term from the double pole. Here
$(r_0)_{12}$ acts in the $(1,2)$ tensor slots.

\emph{Step~2: Bar differential at degree~$2$.}
The ordered bar differential on
$\Barchord_2(\widehat{\mathfrak{g}}_k)$ extracts
the Lie bracket from the simple-pole collision:
\begin{equation}\label{eq:km-bar-d2}
d_{\Barch}[\susp^{-1}J^a \,|\, \susp^{-1}J^b]
\;=\;
f^{ab}{}_{c}\,[\susp^{-1}J^c].
\end{equation}
This is the Lie bracket $[T^a, T^b] = f^{ab}{}_{c}\,T^c$
of~$\mathfrak{g}$. At degree~$2$, the ordered bar complex
reproduces the Chevalley--Eilenberg complex
of~$\mathfrak{g}$. The double-pole term
$k\,(T^a,T^b)/(z-w)^2$ contributes through the collision
residue $\hbar\,\Omega/z$ to the \emph{$R$-matrix}, not
to the bar differential (this is the ordered/symmetric
distinction: in the symmetric bar complex, the same data
appears in the bar differential as the curvature
$m_0 = \frac{k+h^\vee}{2h^\vee}\kappa$).

\emph{Step~3: $d^2=0$ at degree~$3$ is the CYBE.}
The bar differential satisfies $d^2 = 0$ on the ordered bar
complex. At degree~$3$, this condition applied to
$[\susp^{-1}J^a \,|\, \susp^{-1}J^b \,|\, \susp^{-1}J^c]$
is equivalent to the \emph{classical Yang--Baxter equation}
for $r(z)=\hbar\,\Omega/z$:
\begin{equation}\label{eq:cybe}
[r_{12}(z_1{-}z_2),\, r_{13}(z_1{-}z_3)]
\;+\;
[r_{12}(z_1{-}z_2),\, r_{23}(z_2{-}z_3)]
\;+\;
[r_{13}(z_1{-}z_3),\, r_{23}(z_2{-}z_3)]
\;=\; 0.
\end{equation}
Substituting $r_{ij}(z) = \hbar\,\Omega_{ij}/z$, the three
terms become proportional to $f^{a[b}{}_{e}\,f^{c]e}{}_{d}$,
and the vanishing is exactly the Jacobi identity for
$\mathfrak{g}$. Thus $d^2=0$ holds unconditionally for every
simple~$\mathfrak{g}$.

\emph{Step~4: Complete strictification.}
The spectral Drinfeld class $[\delta]\in H^2(\mathfrak{g},
\mathfrak{g}\otimes\mathfrak{g})$ measures the obstruction
to strictifying the $A_\infty$-Yangian to a strict associative
algebra. For every simple Lie algebra~$\mathfrak{g}$, the
root spaces are one-dimensional
(root-space one-dimensionality),
so the relevant $H^2$ group decomposes into root sectors:
each sector is one-dimensional, and the Jacobi identity forces
the cocycle $[\delta]$ to be a coboundary in every sector.
The class therefore vanishes:
\begin{equation}\label{eq:strictification}
[\delta] = 0
\quad\text{for all simple }\mathfrak{g}.
\end{equation}
This is the strictification mechanism: the Yangian
$Y_\hbar(\mathfrak{g})$ is a \emph{strict} associative
algebra, not merely an $A_\infty$-algebra.

\emph{Step~5: Identification of the Koszul dual.}
The cohomology $H^\bullet(\Barchord
(\widehat{\mathfrak{g}}_k))$ with the dual multiplication
(from dualising the deconcatenation coproduct) is the
RTT presentation of the Yangian $Y_\hbar(\mathfrak{g})$:
the degree-$1$ generators are the currents
$\{T^a_n\}_{a,n}$, the degree-$2$ relations are the RTT
relations $R(u-v)\,T_1(u)\,T_2(v) = T_2(v)\,T_1(u)\,R(u-v)$
with $R$-matrix $R(u) = 1 + \hbar\,\Omega/u$, and the
strictification of Step~4 ensures no higher $A_\infty$
corrections. By the Drinfeld theorem, this is
$Y_\hbar(\mathfrak{g})$.
\end{proof}

\subsection{Data for all simple types}

\begin{table}[ht]
\centering
\caption{Ordered Koszul duality data for the simple
Kac--Moody family. All entries follow from
Theorem~\textup{\ref{thm:km-yangian}}.}
\label{tab:km-yangian-data}
\renewcommand{\arraystretch}{1.3}
\begin{tabular}{@{}llcclcl@{}}
\toprule
Type & $\mathfrak{g}$ & $\dim\mathfrak{g}$ & $h^\vee$
 & Fund.\ repr.\ dim & $\hbar$ & $\kappa(Y_\hbar(\mathfrak{g}))$ \\
\midrule
$A_n$ & $\mathfrak{sl}_{n+1}$ & $n^2+2n$ & $n+1$
 & $n+1$ & $\frac{1}{k+n+1}$
 & $-\frac{(n^2+2n)(k+n+1)}{2(n+1)}$ \\[4pt]
$B_n$ & $\mathfrak{so}_{2n+1}$ & $n(2n+1)$ & $2n-1$
 & $2n+1$ & $\frac{1}{k+2n-1}$
 & $-\frac{n(2n+1)(k+2n-1)}{2(2n-1)}$ \\[4pt]
$C_n$ & $\mathfrak{sp}_{2n}$ & $n(2n+1)$ & $n+1$
 & $2n$ & $\frac{1}{k+n+1}$
 & $-\frac{n(2n+1)(k+n+1)}{2(n+1)}$ \\[4pt]
$D_n$ & $\mathfrak{so}_{2n}$ & $n(2n-1)$ & $2n-2$
 & $2n$ & $\frac{1}{k+2n-2}$
 & $-\frac{n(2n-1)(k+2n-2)}{2(2n-2)}$ \\[4pt]
$E_6$ & $\mathfrak{e}_6$ & $78$ & $12$
 & $27$ & $\frac{1}{k+12}$
 & $-\frac{78(k+12)}{24}$ \\[4pt]
$E_7$ & $\mathfrak{e}_7$ & $133$ & $18$
 & $56$ & $\frac{1}{k+18}$
 & $-\frac{133(k+18)}{36}$ \\[4pt]
$E_8$ & $\mathfrak{e}_8$ & $248$ & $30$
 & $248$ & $\frac{1}{k+30}$
 & $-\frac{248(k+30)}{60}$ \\[4pt]
$F_4$ & $\mathfrak{f}_4$ & $52$ & $9$
 & $26$ & $\frac{1}{k+9}$
 & $-\frac{52(k+9)}{18}$ \\[4pt]
$G_2$ & $\mathfrak{g}_2$ & $14$ & $4$
 & $7$ & $\frac{1}{k+4}$
 & $-\frac{14(k+4)}{8}$ \\
\bottomrule
\end{tabular}
\end{table}

\noindent
In every case the curvature of the $E_1$-chiral Koszul dual is
\begin{equation}\label{eq:km-yangian-curvature}
\kappa\bigl(Y_\hbar(\mathfrak{g})\bigr)
\;=\;
-\kappa(\widehat{\mathfrak{g}}_k)
\;=\;
-\frac{\dim(\mathfrak{g})\cdot(k+h^\vee)}{2h^\vee},
\end{equation}
the negation of the affine curvature. The KM complementarity
identity $\kappa(\widehat{\mathfrak{g}}_k) + \kappa(Y_\hbar(\mathfrak{g})) = 0$
holds specifically for Kac--Moody and free-field families and
does \emph{not} generalise: for Virasoro one has
$\kappa(\mathrm{Vir}_c) + \kappa(\mathrm{Vir}_{26-c}) = 13$, and for
$\mathcal{W}_3$ one has
$\kappa(\mathcal{W}_3^k) + \kappa(\mathcal{W}_3^{k'}) = 250/3$
(cf.\ Theorem~D of Volume~I).

\begin{remark}[Casimir tensor for each type]
\label{rem:casimir-types}
In the fundamental representation, the Casimir tensor
$\Omega\in\mathfrak{g}\otimes\mathfrak{g}$ acts as:
\begin{enumerate}[label=\textup{(\alph*)}]
\item \emph{Type $A_n$.}\enspace
$\Omega = P - \frac{1}{n+1}\,I\otimes I$,
where $P\in\End(\bC^{n+1}\otimes\bC^{n+1})$ is the
permutation operator.
\item \emph{Types $B_n$, $D_n$.}\enspace
$\Omega = P - K$,
where $K_{ij,kl} = \delta_{ik'}\delta_{jl'}$
is the contraction with the invariant bilinear form
(using the index convention $i' = 2n+2-i$ for $B_n$,
$i' = 2n+1-i$ for $D_n$).
\item \emph{Type $C_n$.}\enspace
$\Omega = P - K$,
where $K_{ij,kl} = J_{ik}J_{jl}$ involves the
symplectic form~$J$.
\item \emph{Exceptional types.}\enspace
$\Omega$ is determined by the structure constants
$f^{ab}{}_{c}$ via
$\Omega = \sum_a T^a\otimes T^a$
in any orthonormal basis; explicit matrix forms
follow from the Chevalley basis.
\end{enumerate}
\end{remark}

\begin{remark}[Relation to Drinfeld--Kohno and quantum groups]
\label{rem:km-drinfeld-kohno}
The Kohno connection~\eqref{eq:kohno-connection} for
$\widehat{\mathfrak{g}}_k$ takes the form
$\nabla = d - \hbar\sum_{i<j}\Omega_{ij}\,dz_{ij}/z_{ij}$,
which is the Knizhnik--Zamolodchikov connection at level~$k$.
Its monodromy representation $\pi_1(\mathrm{Conf}_n(\bC))
\to\GL(\mathfrak{g}^{\otimes n})$ factors through the
Yangian via the Drinfeld--Kohno theorem: the monodromy of
the KZ connection at $\hbar = 1/(k+h^\vee)$ coincides with
the $R$-matrix of $Y_\hbar(\mathfrak{g})$ acting on evaluation
modules. The ordered bar complex provides the chain-level
refinement of this classical result.
\end{remark}

\begin{remark}[The pole-order dichotomy]
\label{rem:km-pole-dichotomy}
The Kac--Moody OPE has a double pole and a simple pole.
The $d\log$ kernel absorbs one power, leaving a collision
residue with a simple pole $\hbar\,\Omega/z$ (the spectral
term) and a constant $r_0$ (the Lie bracket term). This
places all $\widehat{\mathfrak{g}}_k$ in class~L (double-pole
OPE, simple-pole collision residue, finite shadow depth).
By contrast, the Virasoro OPE has a quartic pole, giving
a triple-pole collision residue and class~M (infinite shadow
depth). Under DS reduction
$\widehat{\mathfrak{g}}_k\to\mathcal{W}_k(\mathfrak{g})$,
the pole-order escalation transports class~L to class~M
via Sugawara.
\end{remark}

\begin{remark}[The Kac--Moody frontier]
\label{rem:km-frontier}
Theorem~\ref{thm:km-yangian} covers all simple Lie
algebras. The genuine frontier is Kac--Moody algebras
$\widehat{\mathfrak{g}}_k$ for~$\mathfrak{g}$ a
\emph{Kac--Moody} Lie algebra (not finite-dimensional
simple): for affine or indefinite Kac--Moody algebras,
root multiplicities can exceed one, the root-space
one-dimensionality argument of Step~4 fails, and the
strictification mechanism requires new input.
\end{remark}



\section{Spectral Drinfeld strictification}
\label{sec:spectral-drinfeld-strictification}

The ordered bar complex of a dg-shifted Yangian carries a filtered
$A_\infty$ structure. The question of whether this $A_\infty$ structure
can be straightened to a genuine associative structure with spectral
parameter is controlled by a cohomological obstruction: the
\emph{spectral Drinfeld class}. This
obstruction vanishes for all simple Lie algebras.

\subsection{The obstruction class}

\begin{definition}[Spectral Drinfeld class at filtration $p$]
\label{def:spectral-drinfeld-class-oackd}
Let $A$ be a filtered dg-shifted Yangian carrying a spectral
quasi-factorization datum with quasi-associator
\[
\Phi(w,z)=1+\phi_p(w,z)\pmod{F^{p+1}},
\]
where $\phi_p$ is the first nontrivial term at filtration~$p$.
The spectral cochain complex is defined by
\[
C^2_{\mathrm{spec}}(A)=\{g(u)\in A^{\widehat\otimes 2}\},
\qquad
C^3_{\mathrm{spec}}(A)=\{\phi(w,z)\in A^{\widehat\otimes 3}\},
\]
with linearized spectral differential
\[
(\delta g)(w,z)
=
1\otimes g(z)
+ (\id\otimes \Delta_z)g(w)
- (\Delta_w\otimes \id)g(z)
- g(w)\otimes 1.
\]
The \emph{spectral Drinfeld class at filtration $p$} is
\[
\mathfrak D_{\mathrm{spec}}(A,\Omega)\big|_{F^p}
:=[\phi_p]
\in H^3_{\mathrm{spec}}(\gr^p A).
\]
\end{definition}

\begin{remark}[Interpretation]
The spectral Drinfeld class measures the failure of the $A_\infty$
bar-complex structure to strictify to a genuine associative product
with spectral parameter. At each filtration $p\ge 2$, the pentagon
relation on the quasi-associator linearizes: $\phi_p$ is a $3$-cocycle,
a spectral twist shifts it by a coboundary, and the class
$[\phi_p]\in H^3_{\mathrm{spec}}(\gr^p A)$ is the intrinsic
obstruction. If $[\phi_p]=0$ for all $p\ge 2$, an iterative
filtered twist produces $\Phi^G=1$, and the quasi-factorization
datum strictifies to an honest spectral factorization $R$-matrix.
\end{remark}

\subsection{Root-space one-dimensionality}

The key structural input is a classical theorem about the root
structure of simple Lie algebras: every root space is one-dimensional.
This forces the spectral Drinfeld class, which is valued in root
spaces, to be a scalar, and a scalar is determined by a single
equation.

\begin{theorem}[Root-space one-dimensionality; \ClaimStatusProvedHere]
\label{thm:root-space-one-dim-v1}
Let $\fg$ be a simple Lie algebra. For any subset
$S=\{\alpha_1,\ldots,\alpha_n\}$ of simple roots such that
$\alpha=\sum_{i=1}^n\alpha_i$ is a root of $\fg$, the multilinear
degree-$n$ space
\[
\operatorname{span}\{\textup{Lie monomials in }E_{\alpha_1},\ldots,
E_{\alpha_n}\textup{ using each exactly once}\}\;\subseteq\;\fg_\alpha
\]
is one-dimensional.
\end{theorem}

\begin{proof}
Every multilinear Lie monomial in $E_{\alpha_1},\ldots,E_{\alpha_n}$
has weight $\alpha_1+\cdots+\alpha_n=\alpha$ (since each bracket
adds the weights of its arguments). Therefore all such monomials
lie in the root space $\fg_\alpha$. In a simple Lie algebra, every
root has multiplicity one: $\dim\fg_\alpha=1$. Hence the
multilinear space is at most one-dimensional.

It remains to verify that at least one multilinear monomial is
nonzero. Since $\alpha$ is a root and the simple roots $\alpha_i$ are
linearly independent, there exists an ordering
$\alpha_{i_1},\ldots,\alpha_{i_n}$ and a sequence of partial sums
$\beta_k=\alpha_{i_1}+\cdots+\alpha_{i_k}$ such that each $\beta_k$
is a root (this follows from the connectedness of the support of
$\alpha$ in the Dynkin diagram and the standard theory of root
strings). The iterated bracket
$[E_{\alpha_{i_1}},[E_{\alpha_{i_2}},\ldots,E_{\alpha_{i_n}}]\cdots]$
is then nonzero (each intermediate bracket lands in a nonzero root
space).
\end{proof}

\subsection{The Jacobi collapse lemma}

When root spaces are one-dimensional, multilinear Lie monomials
in root vectors are necessarily proportional: the target space
has no room for independent contributions. The mechanism by
which the free Lie algebra's higher-dimensional multilinear spaces
collapse to one dimension inside $\fg$ is the Jacobi identity.

The simplest nontrivial case is the $D_4$ star: four simple roots
$\alpha_1,\alpha_2,\alpha_3,\alpha_c$ with $\alpha_c$ adjacent to
each of $\alpha_1,\alpha_2,\alpha_3$, and the $\alpha_i$ mutually
non-adjacent. In the free Lie algebra modulo non-adjacency
commutation, the multilinear degree-$4$ space is two-dimensional.
In $\fg$, the Jacobi identity collapses it to one.

\begin{lemma}[Jacobi collapse for star sectors; \ClaimStatusProvedHere]
\label{lem:jacobi-collapse-v1}
Let $X_1,X_c,X_2,X_3$ be root vectors in a Lie algebra $\fg$ with
$[X_1,X_2]=[X_1,X_3]=[X_2,X_3]=0$ (the $D_4$-star commutation
relations). Define
\[
A = [[X_2,[X_1,X_c]],X_3],
\qquad
B = [[X_3,[X_1,X_c]],X_2].
\]
Then $A=B$.
\end{lemma}

\begin{proof}
Apply the Jacobi identity to $(X_2,X_c,X_3)$:
\[
[X_2,[X_c,X_3]]+[X_c,[X_3,X_2]]+[X_3,[X_2,X_c]]=0.
\]
Since $[X_2,X_3]=0$, the middle term vanishes, giving
\begin{equation}\label{eq:jacobi-star-v1}
[X_2,[X_c,X_3]]=[X_3,[X_c,X_2]].
\end{equation}
Now expand $A$ and $B$ using Jacobi and $[X_1,X_2]=[X_1,X_3]=0$:
\begin{align*}
A &= [[X_2,[X_1,X_c]],X_3]
 = [X_2,[[X_1,X_c],X_3]] - [[X_1,X_c],[X_2,X_3]]\\
 &= [X_2,[[X_1,X_c],X_3]]
 \qquad(\text{since }[X_2,X_3]=0)\\
 &= [X_2,[X_1,[X_c,X_3]]] - [X_2,[X_c,[X_1,X_3]]]\\
 &= [X_2,[X_1,[X_c,X_3]]]
 \qquad(\text{since }[X_1,X_3]=0)\\
 &= [[X_2,X_1],[X_c,X_3]] + [X_1,[X_2,[X_c,X_3]]]\\
 &= [X_1,[X_2,[X_c,X_3]]]
 \qquad(\text{since }[X_1,X_2]=0).
\end{align*}
The identical computation with $2\leftrightarrow 3$ gives
$B=[X_1,[X_3,[X_c,X_2]]]$. By~\eqref{eq:jacobi-star-v1},
$[X_2,[X_c,X_3]]=[X_3,[X_c,X_2]]$, hence $A=B$.
\end{proof}

\begin{remark}[Root multiplicity one as structural inevitability]
The Jacobi collapse lemma is not a clever computation that
happens to kill the obstruction. The point is that the obstruction
\emph{has nowhere to live}.

The spectral Drinfeld class at filtration $n$, restricted to the root
sector $\alpha=\sum\alpha_i$, is a cohomology class valued in
$\fg_\alpha$. Root multiplicity one says
$\dim\fg_\alpha=1$. A cohomology class valued in a one-dimensional
space is a scalar. A scalar is determined by a single equation.
The coproduct rigidity provides that equation.

The free Lie algebra on the simple root vectors, modulo non-adjacency
commutation, may have higher-dimensional multilinear spaces (the $D_4$
star has $\dim=2$). But in $\fg$ the Jacobi identity collapses them
(Lemma~\ref{lem:jacobi-collapse-v1}). This collapse is not an accident:
it is the shadow of root multiplicity one, which is itself a theorem
about the representation theory of $\mathfrak{sl}_2$-triples.
\end{remark}

\subsection{The beta-integral coefficient}

The BCH expansion of $\log(e^{X_1}\cdots e^{X_n})$ produces the
universal coefficient at each filtration. This coefficient is
determined by the Dynkin projection.

\begin{theorem}[Dynkin coefficient via the beta integral; \ClaimStatusProvedHere]
\label{thm:dynkin-beta-integral}
Let $X_1,\ldots,X_n$ satisfy the path commutation relations
$[X_i,X_j]=0$ for $|i-j|\ge 2$. Then the multilinear
$X_1\cdots X_n$-part of the ordered Baker--Campbell--Hausdorff product is
\[
\operatorname{mult}_{X_1\cdots X_n}
\log(e^{X_1}\,e^{X_2}\cdots e^{X_n})
\;=\;
\frac{1}{n}\,[X_1,[X_2,\ldots,[X_{n-1},X_n]\cdots]],
\]
where the coefficient $1/n$ arises from the beta integral
\begin{equation}\label{eq:beta-integral-coeff}
\frac{1}{n}=\int_0^1(1-t)^{n-1}\,dt.
\end{equation}
This is the unique coefficient forced by the BCH formula: the Dynkin
projection $D_n=\frac{1}{n}\ell_n$ applied to the ordered
concatenation $X_1\otimes\cdots\otimes X_n$.
\end{theorem}

\begin{proof}
The multilinear part is $c_n\cdot[X_1,[X_2,\ldots,X_n]\cdots]$ for a
unique scalar $c_n$ (the multilinear Lie space is one-dimensional by
root-space one-dimensionality in $\fg$, or by the path
one-dimensionality theorem in the free Lie algebra quotient). We
determine $c_n$ by extracting the coefficient of the concatenation
monomial $X_1\cdot X_2\cdots X_n$ in the universal enveloping algebra.

\medskip\noindent\textbf{Step~1: Coefficient in the Lie bracket.}
In the universal enveloping algebra, the right-normed bracket
$[X_1,[X_2,\ldots,[X_{n-1},X_n]\cdots]]$ has leading term
$X_1 X_2\cdots X_n$ with coefficient $+1$.

\medskip\noindent\textbf{Step~2: Coefficient in the BCH logarithm.}
The multilinear part of $e^{X_1}\cdots e^{X_n}$ in the free
associative algebra is $X_1\cdots X_n$. The coefficient of
$X_1\cdots X_n$ in $\log(e^{X_1}\cdots e^{X_n})$, computed via
$\log(1+A)=\sum_{k\ge 1}(-1)^{k+1}A^k/k$, counts compositions
of~$n$ into $k$ parts:
\[
\text{coeff of }X_1\cdots X_n\text{ in }\log(P)
=
\sum_{k=1}^{n}\frac{(-1)^{k+1}}{k}\binom{n-1}{k-1}.
\]

\medskip\noindent\textbf{Step~3: The beta integral.}
\begin{align*}
\sum_{k=1}^{n}\frac{(-1)^{k+1}}{k}\binom{n-1}{k-1}
&=
\sum_{j=0}^{n-1}\frac{(-1)^{j}}{j+1}\binom{n-1}{j}
=
\sum_{j=0}^{n-1}\binom{n-1}{j}(-1)^j\int_0^1 t^j\,dt\\
&=
\int_0^1\sum_{j=0}^{n-1}\binom{n-1}{j}(-t)^j\,dt
=
\int_0^1(1-t)^{n-1}\,dt
\;=\;\frac{1}{n}.
\end{align*}
Combining Steps~1--3: $c_n\cdot 1=\frac{1}{n}$, hence
$c_n=\frac{1}{n}$.
\end{proof}

\begin{remark}[The Dynkin projection as coassociative counit]
The coefficient $1/n$ is the Dynkin projection
$D_n\colon T^n(V)\to\mathrm{Lie}^n(V)$,
$D_n=\frac{1}{n}\ell_n$, applied to the ordered concatenation.
The bar-cobar counit $\varepsilon\colon\Omega(\Barch(\cA))\to\cA$
extracts the chiral algebra from its bar complex (a chiral
operation on $\mathrm{FM}_k(\mathbb{C})$). The Dynkin projection
extracts the Lie element from the tensor algebra (a coassociative
operation on $\mathrm{Conf}_n(\mathbb{R})$). Both are extraction
maps from a cofree object to the generating algebra, controlled
by the same combinatorial identity: the logarithm of the ordered
exponential. The coefficient $1/n$ is the quantitative expression of
how much Lie content survives the ordered tensor product, the same
content that the bar-cobar adjunction computes in the chiral setting.
\end{remark}

\subsection{Complete strictification}

\begin{theorem}[Complete strictification for all simple Lie algebras;
\ClaimStatusProvedHere]
\label{thm:complete-strictification-v1}
Let $\fg$ be a simple Lie algebra of any type. The spectral Drinfeld
class vanishes at every filtration:
\[
\mathfrak{D}_{\mathrm{spec}}(A,\Omega)\big|_{F^p}=0
\qquad\text{for all }p\ge 2.
\]
In particular, whenever a residue-bounded complete dg-shifted Yangian
over $\fg$ carries a spectral quasi-factorization datum, that datum
strictifies to an honest spectral factorization realization.
The dg-shifted Yangian $Y^{\mathrm{dg}}_\hbar(\fg)$ carries a strict filtered
algebra structure.
\end{theorem}

\begin{proof}
The proof proceeds by analysing the spectral Drinfeld class sector by
sector. At filtration $p$, each root sector of type
$\alpha=\sum_{i=1}^p\alpha_i$ (where the $\alpha_i$ are simple roots)
contributes independently.

\medskip\noindent\textbf{Case~1: $\alpha$ is not a root.}
If $\alpha=\sum\alpha_i$ is not a root of $\fg$, then the root space
$\fg_\alpha=0$ and every multilinear Lie monomial in
$\{E_{\alpha_i}\}$ vanishes. The sector contributes nothing.

\medskip\noindent\textbf{Case~2: $\alpha$ is a root (path-type
sectors).}
The spectral Drinfeld class restricted to the root sector lives in
$\fg_\alpha$. By Theorem~\ref{thm:root-space-one-dim-v1},
$\dim\fg_\alpha=1$. The multilinear BCH computation
(Theorem~\ref{thm:dynkin-beta-integral}) gives a unique coefficient
$1/p$ for the product side. The coproduct rigidity gives a unique
equation for the coproduct side. Since both sides are scalars in the
same one-dimensional space, equality reduces to a single identity
(the same identity that holds at filtrations $2$ and $3$) and the
obstruction class vanishes.

\medskip\noindent\textbf{Case~3: Non-path sectors (e.g., the $D_4$
star).}
Root-space one-dimensionality
(Theorem~\ref{thm:root-space-one-dim-v1}) forces the multilinear
BCH contribution to be a scalar multiple of the unique generator of
$\fg_\alpha$. The Jacobi collapse lemma
(Lemma~\ref{lem:jacobi-collapse-v1}) ensures that the
two-dimensional free Lie space collapses to one dimension in
$\fg$, so the only surviving contribution is the iterated bracket
$[E_{\alpha_1},[E_{\alpha_2},\ldots]]$, whose coefficient in the
BCH expansion is $1/p$ by the Dynkin formula.

\medskip\noindent\textbf{The inductive strictification.}
Therefore no root sector at any filtration carries an independent
obstruction class. The spectral Drinfeld class vanishes at every
filtration. Starting from any spectral quasi-factorization datum, a
filtered twist kills the first nontrivial quasi-associator term at
each filtration. Iterating filtration by filtration produces a twist
with trivial quasi-associator, hence an honest spectral factorization
realization.
\end{proof}

\begin{remark}[The role of root multiplicity one]
The theorem has the form of a
\emph{root-multiplicity-one rigidity principle}: the spectral
Drinfeld obstruction to strictifying a dg-shifted Yangian vanishes
because every root-sector multilinear space is one-dimensional.

The bar complex $\barB^{\mathrm{ord}}(\cA)$ is an $E_1$ coassociative
coalgebra; the $\SCchtop$ datum emerges at the level of the chiral
derived center, as the pair $(C^\bullet_{\mathrm{ch}}(\cA,\cA),\cA)$
(Theorem~\ref{thm:master}). The complete strictification
theorem says that, once the spectral quasi-factorization package has
been constructed, the passage from quasi to strict is not an additional
datum but an automatic consequence of the root structure. Just as the
Decomposition Theorem (purity of IC complexes) makes the Hecke algebra
arise inevitably from the geometry of $G/B$, root multiplicity one
makes the strictification step arise inevitably from the Lie theory
of~$\fg$.
\end{remark}

\begin{remark}[What fails for Kac--Moody algebras with root multiplicities $>1$]
\label{rem:kac-moody-failure}
For Kac--Moody algebras whose root lattice contains roots with
multiplicity $>1$ (e.g., imaginary roots of indefinite Kac--Moody
algebras), root-space one-dimensionality fails. The Jacobi collapse
mechanism no longer forces multilinear Lie monomials to be
proportional, and the spectral Drinfeld obstruction class can live in
a space of dimension $>1$. In such cases, the coproduct rigidity
provides fewer equations than unknowns, and genuine obstructions
could survive.

The frontier admits a precise stratification: the obstruction vanishes
in all root sectors where root multiplicity one holds (all roots of
simple Lie algebras, real roots of affine Lie algebras), and the
remaining sectors are exactly those with higher multiplicity.

For untwisted affine Lie algebras $\widehat\fg$, the real roots
$\alpha+n\delta$ have multiplicity~$1$ (so the theorem applies
verbatim), while the imaginary roots $n\delta$ have
multiplicity $\ell=\operatorname{rank}(\fg)$. The imaginary sectors
are abelian (elements within a single $\fg_{n\delta}$ commute), and
a Cartan gauge-twist argument should absorb the defect. The mixed
real-plus-imaginary sectors require a separate analysis. Complete
strictification for affine types therefore remains conjectural pending
the detailed imaginary-sector computation.
\end{remark}

\begin{remark}[Scope of the Borcherds lift as strictification device]
\label{rem:borcherds-lift-scope}
The strictification theorem~\ref{thm:complete-strictification-v1}
rests on two logically separable inputs: (i) one-dimensionality of
root spaces, and (ii) Jacobi-driven collapse of multilinear free Lie
spaces onto a single generator
(Lemma~\ref{lem:jacobi-collapse-v1}). Input (ii) is a \emph{structural}
consequence of (i), not merely a coincidence of ranks: whenever
$\dim\fg_\alpha=1$, every multilinear Lie monomial valued in
$\fg_\alpha$ is a scalar multiple of a single iterated bracket, and the
Jacobi identity among the simple roots provides the unique scalar
equation that fixes it. The \emph{mechanism}, Jacobi-star-sector
killing, generalizes to any Lie (super)algebra whose root multiplicity
structure can be encoded so that each root sector becomes
one-dimensional \emph{after an auxiliary grading}. A Borcherds
generalized Kac--Moody superalgebra supplies exactly this grading: a
root of multiplicity $m_\alpha$ on the bosonic Kac--Moody side is
replaced by a single root of multiplicity $1$ on the Borcherds
superalgebra side, with the $m_\alpha$-fold degeneracy absorbed into an
internal $\Lambda$-grading that couples to super-parity
\cite{Borcherds1988,Borcherds1992}. This is the content of the
no-ghost theorem for lattice vertex algebras: physical states carry a
natural $\mathbb{Z}_{\ge 0}$-level grading, and imaginary-root
degeneracies are exactly the dimension of the level-$m_\alpha$ string
Hilbert space, which resolves after quantization by the
$(b,c)$-ghost super-extension.

The scope of this lift is, however, not uniform across Kac--Moody
types. Three regimes must be distinguished:
\begin{enumerate}
\item \emph{Finite and affine types with real roots.} All real roots
$\alpha$ satisfy $\dim\fg_\alpha=1$;
Theorem~\ref{thm:complete-strictification-v1} applies verbatim in the
real-root sectors. The Borcherds lift is trivial (identity) on these
sectors.
\item \emph{Affine types, imaginary roots $n\delta$.} The imaginary
root space is abelian of dimension $\ell=\operatorname{rank}(\fg)$. The
Borcherds lift replaces $\fg_{n\delta}$ by a single generator carrying
a $\mathbb{Z}^\ell$-grading, and the Jacobi-collapse argument extends
verbatim to the super-extension because abelian sectors contribute no
bracket. A Cartan gauge-twist absorbs the diagonal defect; the mixed
real-plus-imaginary star sectors reduce to the real-root argument
after projection.
\item \emph{Hyperbolic and generalized Kac--Moody types.} Imaginary
roots appear with multiplicities $m_\alpha>1$ governed by
Borcherds's denominator identity or the Frenkel--Kac--Segal vertex
construction (e.g.\ the Monster Lie algebra
$\mathfrak{m}$~\cite{Borcherds1992} and the Fake Monster
$\mathfrak{m}_{\mathrm{fake}}$~\cite{Borcherds1990}). In these cases
the super-lift is nontrivial and the Jacobi-collapse identity must
be verified inside the super-bracket after imaginary-root expansion.
\end{enumerate}

\begin{conjecture}[Kac--Moody strictification via the Borcherds lift;
\ClaimStatusConjectured]
\label{conj:kac-moody-strictification-borcherds}
Let $\fg$ be a symmetrizable Kac--Moody algebra (finite, affine,
hyperbolic, or generalized) with root lattice $Q$ and imaginary-root
multiplicity function $m\colon Q\to\Z_{\ge 0}$. Let
$\fg^{\mathrm{super}}$ denote the Borcherds superalgebra obtained by
replacing each imaginary root space $\fg_\alpha$ of multiplicity
$m_\alpha>1$ with a single generator $E_\alpha$ carrying an internal
$\Lambda_\alpha$-grading of rank $m_\alpha$ coupled to super-parity
$|E_\alpha|\equiv m_\alpha\pmod 2$, subject to the Borcherds denominator
identity for the universal enveloping.

Then $\fg^{\mathrm{super}}$ satisfies a super-analogue of root
multiplicity one in the sense that every $\Lambda$-homogeneous root
sector is one-dimensional, and the super-Jacobi identity collapses
every multilinear bracket valued in a $\Lambda$-homogeneous sector to
the unique iterated super-bracket. Consequently the spectral Drinfeld
class of the dg-shifted Yangian $Y^{\mathrm{dg}}_\hbar(\fg^{\mathrm{super}})$
vanishes at every filtration, and the strictification theorem
descends to $Y^{\mathrm{dg}}_\hbar(\fg)$ via the bosonic projection
$\fg^{\mathrm{super}}\twoheadrightarrow\fg$.
\end{conjecture}

Evidence: (i) for untwisted affine $\widehat\fg$, the super-lift is
the trivial extension by the $\mathbb{Z}^\ell$-graded Heisenberg of
the imaginary-root tower, and the Cartan gauge-twist of
Remark~\ref{rem:kac-moody-failure} is exactly the bosonic projection
of a super-Jacobi collapse (verifiable on $\widehat{\mathfrak{sl}_2}$
where the imaginary roots $n\delta$ carry a single $\alpha^\vee$ and
the super-extension is trivial); (ii) for the Fake Monster
$\mathfrak{m}_{\mathrm{fake}}=\operatorname{Lie}(V_\Lambda^{\mathrm{phys}})$
\cite{Borcherds1990}, the lift is supplied by the lattice vertex
algebra $V_\Lambda$ for $\Lambda=\mathrm{II}_{25,1}$ together with the
no-ghost theorem, giving an explicit $\Lambda$-grading on the
imaginary-root spaces; (iii) the Monster Lie algebra $\mathfrak{m}$
admits the Borcherds denominator
$p^{-1}\prod_{(m,n)>0}(1-p^m q^n)^{c(mn)}=j(p)-j(q)$
\cite{Borcherds1992} which exhibits each imaginary-root sector as
$c(mn)$-dimensional and hence as a single super-generator at level
$c(mn)$. The conjecture is open for generic hyperbolic Kac--Moody
algebras where the Borcherds lift has not been made explicit.

\begin{remark}[Tests]
Two verifications suffice to promote the conjecture from heuristic to
an initial theorem in the finite regime: (i) explicit computation on
untwisted $\widehat{\mathfrak{sl}_2}$ at a generic non-critical level,
where imaginary-root multiplicity equals $1$ and the super-lift
collapses to the bosonic theorem, providing a trivial but verifying
instance; (ii) explicit computation on the Monster Lie algebra
$\mathfrak{m}$ at the lowest imaginary-root level (products of the
form $[E_{(1,1)},E_{(1,1)}]$ where $\dim\mathfrak{m}_{(1,1)}=c(1)=196884$),
using the Borcherds superalgebra realization to reduce the
$196884$-fold degeneracy to a single super-generator and verifying
that the super-Jacobi collapse holds on this sector.
\end{remark}
\end{remark}



\section{Worked example: \texorpdfstring{$\mathfrak{sl}_3$}{affine sl(3)} at level \texorpdfstring{$k$}{k}}
\label{sec:sl3-ordered-bar}

Theorem~\ref{thm:km-yangian} applies to every simple Lie algebra.
For $\mathfrak{g}=\mathfrak{sl}_3$ in
explicit matrix-unit coordinates, the rank-$2$ phenomena
invisible in the rank-$1$ ($\mathfrak{sl}_2$) case:
adjacent-root triangle sectors, the coefficient~$\frac12$,
the Serre relations from root-space vanishing, and the $9\times 9$
$R$-matrix.

\subsection{Root system and matrix-unit notation}

The Lie algebra $\mathfrak{sl}_3$ has root system
$\Phi=\{\pm\alpha_1,\pm\alpha_2,\pm(\alpha_1+\alpha_2)\}$
of type~$A_2$, with simple roots $\alpha_1$, $\alpha_2$ and
Cartan matrix
\[
A=\begin{pmatrix}2&-1\\-1&2\end{pmatrix}.
\]
In the $\mathfrak{gl}_3$ matrix-unit basis we set
\[
e_1=E_{12},\quad
e_2=E_{23},\quad
e_{12}=[e_1,e_2]=E_{13},
\]
\[
f_1=E_{21},\quad
f_2=E_{32},\quad
f_{12}=[f_2,f_1]=E_{31},
\]
and $h_1=E_{11}-E_{22}$, $h_2=E_{22}-E_{33}$. The composite root
generators satisfy
\[
[e_1,e_2]=e_{12},
\qquad
[f_2,f_1]=f_{12},
\qquad
[f_1,f_2]=-f_{12}.
\]
The affine vertex algebra $V_k(\mathfrak{sl}_3)$ has OPE
\[
J^a(z)\,J^b(w)\sim\frac{k\,\kappa(a,b)}{(z-w)^2}+\frac{[a,b](w)}{z-w}
\]
where $\kappa$ is the Killing form normalised so that
$\kappa(e_i,f_i)=1$ and $\kappa(e_{12},f_{12})=1$.

\subsection{Degree-\texorpdfstring{$1$}{1} generators: the Drinfeld currents}

The ordered bar complex
${\Barch}^{\mathrm{ord}}_\bullet(V_k(\mathfrak{sl}_3))$
has degree-$1$ generators
$\susp^{-1}e_1$, $\susp^{-1}e_2$, $\susp^{-1}e_{12}$,
$\susp^{-1}f_1$, $\susp^{-1}f_2$, $\susp^{-1}f_{12}$,
$\susp^{-1}h_1$, $\susp^{-1}h_2$.
The corresponding bar cohomology classes organise into
Drinfeld generating functions for $Y_\hbar(\mathfrak{sl}_3)$
with $\hbar = 1/(k+3)$:
\begin{align*}
E_i(u) &= \sum_{r\ge 0} E_i^{(r)}\,u^{-r-1},
\qquad i=1,2,\\
F_i(u) &= \sum_{r\ge 0} F_i^{(r)}\,u^{-r-1},\\
H_i(u) &= 1+\sum_{r\ge 0} H_i^{(r)}\,u^{-r-1},
\end{align*}
where $u=1/z$ is the spectral parameter. The zeroth modes
$E_i^{(0)}$, $F_i^{(0)}$, $H_i^{(0)}$ generate a copy of
$U(\mathfrak{sl}_3)$ inside $Y_\hbar(\mathfrak{sl}_3)$.

\subsection{Degree-\texorpdfstring{$2$}{2} relations: Cartan matrix and Yangian deformation}

The bar differential~\eqref{eq:km-bar-d2} specialises to
\begin{align*}
d_{\Barch}(\susp^{-1}h_i\otimes \susp^{-1}e_j)
&=A_{ij}\,\susp^{-1}e_j,\\
d_{\Barch}(\susp^{-1}e_i\otimes \susp^{-1}f_j)
&=\delta_{ij}\,\susp^{-1}h_i,\\
d_{\Barch}(\susp^{-1}e_1\otimes \susp^{-1}e_2)
&=\susp^{-1}e_{12},\\
d_{\Barch}(\susp^{-1}f_2\otimes \susp^{-1}f_1)
&=\susp^{-1}f_{12}.
\end{align*}
The first two lines reproduce the Cartan relations
$[H_i,E_j]=A_{ij}E_j$ and $[E_i,F_j]=\delta_{ij}H_i$
at leading order in the spectral parameter. The last two express
the composite root as a bar boundary. The Yangian deformation
appears at the next spectral order:
\[
[H_i(u),\,E_j(v)]
=
\frac{A_{ij}}{u-v}\,E_j(v)+\text{(regular)},
\qquad
[E_i(u),\,F_j(v)]
=
\frac{\delta_{ij}}{u-v}\,H_i(v)+\text{(regular)}.
\]

\subsection{Degree-\texorpdfstring{$3$}{3} triangle sectors}

This is the first genuinely rank-$2$ sector: the composite root
$\alpha_1+\alpha_2$ produces degree-$3$ bar elements that
have no $\mathfrak{sl}_2$ analogue.

\begin{definition}[Triangle sectors for $\mathfrak{sl}_3$]
\label{def:sl3-triangle-sectors}
The degree-$3$ ordered bar complex
${\Barch}^{\mathrm{ord}}_3(V_k(\mathfrak{sl}_3))$
contains four \emph{triangle sectors}, indexed by the composite
root $\alpha_1+\alpha_2$ and the two orderings of simple lowering
(respectively raising) operators:
\begin{align*}
T^L_{12}&=E_{13}^{(1)}\,E_{21}^{(2)}\,E_{32}^{(3)},
&
T^L_{21}&=E_{13}^{(1)}\,E_{32}^{(2)}\,E_{21}^{(3)},\\
T^R_{12}&=E_{12}^{(1)}\,E_{23}^{(2)}\,E_{31}^{(3)},
&
T^R_{21}&=E_{23}^{(1)}\,E_{12}^{(2)}\,E_{31}^{(3)}.
\end{align*}
The left triangles have the composite raising operator $e_{12}=E_{13}$ in
the first tensor factor and the two simple lowering operators split
across factors~$2$ and~$3$. The right triangles have the two simple
raising operators in factors~$1$ and~$2$ and the composite lowering
operator $f_{12}=E_{31}$ in factor~$3$.
\end{definition}

\begin{theorem}[Triangle coefficient for $\mathfrak{sl}_3$; \ClaimStatusProvedHere]
\label{thm:sl3-triangle-coefficient}
Let the two-body spectral functions be given by the rational ansatz
$\rho_{\alpha}(u)=c_\alpha/u$, and define the triangle transport
coefficient $\lambda$ by
\[
\Pi_{\langle T^L_{12},\,T^L_{21}\rangle}
\big((\id\otimes\Delta_z)\,X_{\alpha_1+\alpha_2}^{12}(w)\big)
=
\lambda\big(\rho_1(w)\rho_2(w{+}z)\,T^L_{12}
-\rho_2(w)\rho_1(w{+}z)\,T^L_{21}\big).
\]
Then the triangle coefficient is
\[
\boxed{%
\lambda
=
\frac{c_1\,c_2}{2\,c_{12}}
=
\frac{1}{2}
=
\int_0^1(1-t)\,dt,
}
\]
where the last equality identifies the coefficient as the $n=2$ Euler
beta integral $B(1,2)=\int_0^1(1-t)^{2-1}\,dt=1/2$.
In the normalised case $c_1=c_2=c_{12}=1$, the coefficient is~$\frac12$
for both left and right triangles and both orderings:
$\lambda^L_{12}=\lambda^L_{21}=\lambda^R_{12}=\lambda^R_{21}=\frac12$.
\end{theorem}

\begin{proof}
The product side of the filtration-$2$ hexagon defect equation projects
to the triangle sectors by commutator extraction. The only terms that
contribute are
\[
[E_1^{(1)}F_1^{(2)},\,E_2^{(1)}F_2^{(3)}]
=E_{12}^{(1)}\,F_1^{(2)}\,F_2^{(3)}
\]
and
\[
[E_2^{(1)}F_2^{(2)},\,E_1^{(1)}F_1^{(3)}]
=-E_{12}^{(1)}\,F_2^{(2)}\,F_1^{(3)},
\]
where the signs arise from
$[e_1,e_2]=e_{12}$ and $[e_2,e_1]=-e_{12}$.
The logarithm of the ordered product of two-body exponentials gives
\begin{align*}
\Pi_{\langle T^L_{12},\,T^L_{21}\rangle}
\tfrac{1}{2}[S_{12}(w),\,S_{13}(w{+}z)]
&=
\tfrac{1}{2}\rho_1(w)\rho_2(w{+}z)\,T^L_{12}
-\tfrac{1}{2}\rho_2(w)\rho_1(w{+}z)\,T^L_{21}.
\end{align*}
On the coproduct side, $X_{\alpha_1+\alpha_2}$ is the unique composite
root generator that can land in the triangle sector. The twisted
coproduct on $X_{\alpha_1+\alpha_2}^{12}(w)$ projects to the triangle
sector with transport coefficient $\lambda$. Equating the two sides and
substituting the rational ansatz $\rho_\alpha(u)=c_\alpha/u$ gives
\[
\frac{c_{12}}{w}\cdot\frac{\lambda}{w+z}
=
\frac{1}{2}\cdot\frac{c_1}{w}\cdot\frac{c_2}{w+z},
\]
whence $\lambda=c_1 c_2/(2c_{12})$. In the normalised case this
is~$1/2$. The integral representation follows from
$B(1,2)=\int_0^1(1-t)\,dt=1/2$.
\end{proof}

\begin{remark}[The triangle coefficient as BCH coefficient]
\label{rem:triangle-as-bch}
The coefficient $1/2$ is the $n=2$ instance of the universal $n$-gon
coefficient $1/n$ proved in Vol~II (the $n$-gon localisation formula).
At $n=2$ the ordered BCH expansion
$\log(e^X e^Y)=X+Y+\frac12[X,Y]+\cdots$
has $\frac12$ as the coefficient of the first commutator term. This is
not a coincidence: the triangle sector is the first commutator sector
in the ordered bar complex. The integral $\int_0^1(1-t)^{n-1}\,dt=1/n$
provides a uniform derivation for all~$n$, replacing case-by-case
Baker--Campbell--Hausdorff algebra.
\end{remark}

\subsection{Serre relations from root-space vanishing}

\begin{proposition}[Serre relations from root-space vanishing; \ClaimStatusProvedHere]
\label{prop:sl3-serre}
The Serre relations of $Y_\hbar(\mathfrak{sl}_3)$ arise from degree-$3$
bar cohomology by root-space obstruction. Specifically:
\begin{enumerate}[label=(\roman*)]
\item $2\alpha_1+\alpha_2\notin\Phi(\mathfrak{sl}_3)$, so
$(\mathfrak{sl}_3)_{2\alpha_1+\alpha_2}=0$, which forces
$[e_1,[e_1,e_2]]=0$.
\item $\alpha_1+2\alpha_2\notin\Phi(\mathfrak{sl}_3)$, so
$(\mathfrak{sl}_3)_{\alpha_1+2\alpha_2}=0$, which forces
$[e_2,[e_2,e_1]]=0$.
\end{enumerate}
In the ordered bar complex, these correspond to the vanishing of
the degree-$3$ bar differential on the would-be generators
$\susp^{-1}e_1\otimes \susp^{-1}e_1\otimes \susp^{-1}e_2$ and
$\susp^{-1}e_2\otimes \susp^{-1}e_2\otimes \susp^{-1}e_1$:
the target root spaces are empty, so the bar differential has
trivial image and the corresponding bar cohomology classes are zero.
The spectral-parameter deformation promotes these to the Yangian
Serre relations:
\[
\operatorname{Sym}_{u_1,u_2}
\big[E_i(u_1),\,[E_i(u_2),\,E_j(v)]\big]=0,
\qquad
A_{ij}=-1.
\]
\end{proposition}

\begin{proof}
Consider the degree-$3$ bar element
$\susp^{-1}e_1\otimes \susp^{-1}e_1\otimes \susp^{-1}e_2$.
The bar differential evaluates the nested commutator
$[e_1,[e_1,e_2]]=[e_1,e_{12}]$. But $e_1\in(\mathfrak{sl}_3)_{\alpha_1}$
and $e_{12}\in(\mathfrak{sl}_3)_{\alpha_1+\alpha_2}$, so the commutator
would live in $(\mathfrak{sl}_3)_{2\alpha_1+\alpha_2}$. Since
$2\alpha_1+\alpha_2\notin\Phi(\mathfrak{sl}_3)$ (the root system of
$A_2$ contains only roots of height $\le 2$), this root space is zero,
hence $[e_1,e_{12}]=0$. The identical argument with $\alpha_1\leftrightarrow\alpha_2$
gives $[e_2,e_{21}]=[e_2,[e_2,e_1]]=0$.

For the Yangian deformation: the spectral OPE
$E_i(u)E_i(v)$ has a pole shifted to $u-v=\hbar$ (off the collision
diagonal), so the collision residue vanishes and
$d_{\mathrm{res}}(\susp^{-1}E_i\otimes \susp^{-1}E_i)=0$. The degree-$3$
bar cocycle condition $d^2=0$ then forces the symmetrised nested
commutator to vanish, giving the Yangian Serre relation.
\end{proof}

\subsection{The RTT presentation and the \texorpdfstring{$9\times 9$}{9x9} R-matrix}

\begin{theorem}[RTT presentation for $Y_\hbar(\mathfrak{sl}_3)$
from the ordered bar complex; \ClaimStatusProvedHere]
\label{thm:sl3-rtt}
Let $V=\bC^3$ be the defining representation of
$\mathfrak{sl}_3$ and define the transfer matrix
\[
T(u)
=\sum_{a,b=1}^{3}e_{ab}\otimes t^{ab}(u)
\in \operatorname{End}(V)\otimes Y_\hbar(\mathfrak{sl}_3)[\![u^{-1}]\!],
\]
where $t^{ab}(u)$ are the degree-$1$ bar cohomology generating
functions. Then the consistency of the degree-$3$ ordered bar complex
${\Barch}^{\mathrm{ord}}_3$ gives the RTT relation
\begin{equation}\label{eq:sl3-rtt}
R(u{-}v)\,T_1(u)\,T_2(v)
=T_2(v)\,T_1(u)\,R(u{-}v),
\end{equation}
where $R(u)\in\operatorname{End}(V\otimes V)$ is the rational
Yang $R$-matrix
\begin{equation}\label{eq:sl3-R-matrix}
R(u)=1+\frac{\hbar\,\Omega}{u}+O(u^{-2}),
\qquad
\Omega=\sum_{a,b=1}^{3}e_{ab}\otimes e_{ba},
\end{equation}
acting on $V\otimes V\cong\bC^9$. Here
$\Omega$ is the $\mathfrak{sl}_3$-Casimir in the tensor square of
the defining representation
\textup{(}Remark~\textup{\ref{rem:casimir-types}(a)} with $n=2$\textup{)},
and $e_{ab}$ are the standard matrix units of $\mathfrak{gl}_3$.
\end{theorem}

\begin{proof}
This is the $A_2$ specialisation of Theorem~\ref{thm:km-yangian}.
The $R$-matrix arises from the monodromy of the flat connection
$\nabla=d-r(z)\,dz$ on the ordered configuration space
$\Conf_2^<(\bC)$, where $r(z)=\hbar\,\Omega/z$ is the classical
$r$-matrix~\eqref{eq:km-collision-residue} (with the
$d\log$ kernel absorbing one pole order).
The RTT relation~\eqref{eq:sl3-rtt} is equivalent to $d^2=0$ on
${\Barch}^{\mathrm{ord}}_3$: the three-fold bar consistency
intertwines two orderings of the first two tensor factors via the
$R$-matrix and inserts a bar cohomology class in the third slot.
The Yang--Baxter equation
\[
R_{12}(u{-}v)\,R_{13}(u)\,R_{23}(v)
=R_{23}(v)\,R_{13}(u)\,R_{12}(u{-}v)
\]
follows from $d^3=0$ on ${\Barch}^{\mathrm{ord}}_4$ and is verified
directly for $R(u)=1+\hbar\,\Omega/u$ using the $\mathfrak{sl}_3$
Jacobi identity in the tensor cube $V^{\otimes 3}\cong\bC^{27}$.
\end{proof}

\begin{remark}[The $9\times 9$ versus $4\times 4$ $R$-matrix]
\label{rem:9x9-vs-4x4}
For $\mathfrak{sl}_2$ the $R$-matrix acts on
$\bC^2\otimes\bC^2\cong\bC^4$.
For $\mathfrak{sl}_3$ it acts on
$\bC^3\otimes\bC^3\cong\bC^9$.
The decomposition
$\bC^9\cong\mathbf{6}\oplus\mathbf{3}$
(symmetric and exterior squares) is an eigenspace decomposition
for the permutation $P=\sum e_{ab}\otimes e_{ba}$: eigenvalue
$+1$ on $\operatorname{Sym}^2(V)$ and $-1$ on $\bigwedge^2(V)$.
The $R$-matrix $R(u)=1+\hbar\,P/u$ therefore has eigenvalues
$1+\hbar/u$ and $1-\hbar/u$ on these two subspaces.
The four triangle sectors $T^L_{12}$, $T^L_{21}$, $T^R_{12}$, $T^R_{21}$
are the components of the degree-$3$ bar complex that detect the
off-diagonal ($\alpha_1+\alpha_2$) block of this $R$-matrix.
\end{remark}

\subsection{Strictification for \texorpdfstring{$\mathfrak{sl}_3$}{sl(3)}}

\begin{theorem}[Vanishing of the spectral Drinfeld class for $\mathfrak{sl}_3$;
\ClaimStatusProvedHere]
\label{thm:sl3-strictification}
The spectral Drinfeld class $[\delta]\in H^2(\mathfrak{sl}_3,\,
\operatorname{End}(\mathfrak{sl}_3\otimes\mathfrak{sl}_3))$
vanishes. Concretely:
\begin{enumerate}[label=\textup{(\roman*)}]
\item At filtration $1$ \textup{(}Cartan sector\textup{)},
the defect is identically zero because the Cartan matrix $A_{ij}$
has integer entries and the rational ansatz is already exact.
\item At filtration $2$ \textup{(}triangle sector\textup{)},
the obstruction lives in the two-dimensional space spanned by
$T^L_{12}$ and $T^L_{21}$. The triangle coefficient
$\lambda=1/2$ satisfies the rigidity equations of
Theorem~\textup{\ref{thm:sl3-triangle-coefficient}}, so the
defect vanishes.
\item At filtration $\ge 3$, the Serre relation
$[e_1,[e_1,e_2]]=0$ ensures vanishing target.
Root-space one-dimensionality
\textup{(}$\dim(\mathfrak{sl}_3)_\alpha=1$ for all
$\alpha\in\Phi$\textup{)} forces the remaining cochains to be
coboundaries via the Jacobi collapse mechanism
\textup{(Step~$4$ of Theorem~\textup{\ref{thm:km-yangian}})}.
\end{enumerate}
Therefore $[\delta]=0$ at every filtration, and
$Y_\hbar(\mathfrak{sl}_3)$ is a strict associative algebra.
\end{theorem}

\begin{proof}
This is Step~4 of Theorem~\ref{thm:km-yangian} specialised to
$A_2$. The filtration-by-filtration verification above records
what the general root-space one-dimensionality argument looks like
when all root spaces are enumerated: $\Phi$ has exactly $6$ roots,
all of multiplicity one.
\end{proof}

\begin{remark}[Contrast with root multiplicities $>1$]
\label{rem:root-multiplicities}
The vanishing of $[\delta]$ uses $\dim(\mathfrak{sl}_3)_\alpha=1$
in an essential way. For Kac--Moody algebras with root multiplicities
$>1$ (e.g.\ the Monster Lie algebra), the endomorphism space at a
root of multiplicity~$m$ is $\operatorname{End}(\bC^m)$, and
the Jacobi collapse argument fails: the obstruction cochain has
non-scalar entries that need not be coboundaries.
\end{remark}


\section{Higher-rank pattern: the \texorpdfstring{$\mathfrak{sl}_4$}{sl(4)} root quadrilateral}
\label{sec:sl4-quadrilateral}

For $\mathfrak{sl}_4$ the Dynkin diagram is the path
$\alpha_1\text{--}\alpha_2\text{--}\alpha_3$ of type $A_3$, with
positive roots
\[
\Phi^+=\{\alpha_1,\;\alpha_2,\;\alpha_3,\;
\alpha_1{+}\alpha_2,\;\alpha_2{+}\alpha_3,\;
\alpha_1{+}\alpha_2{+}\alpha_3\}.
\]
In matrix-unit notation: $e_1=E_{12}$, $e_2=E_{23}$, $e_3=E_{34}$,
$e_{12}=E_{13}$, $e_{23}=E_{24}$, $e_{123}=E_{14}$.

\begin{definition}[Quadrilateral sectors]
\label{def:sl4-quadrilateral}
The degree-$4$ ordered bar complex exhibits \emph{quadrilateral sectors}
associated to the maximal composite root $\alpha_1+\alpha_2+\alpha_3$.
The left quadrilateral vector is
\[
Q^L_{123}
=E_{14}^{(1)}\,E_{21}^{(2)}\,E_{32}^{(3)}\,E_{43}^{(4)},
\]
with permuted orderings of the simple lowering operators in
factors $2$, $3$, $4$.
\end{definition}

\begin{theorem}[Quadrilateral coefficient for $\mathfrak{sl}_4$; \ClaimStatusProvedHere]
\label{thm:sl4-quadrilateral}
Under the normalised rational ansatz, the quadrilateral transport
coefficient is
\[
\boxed{%
\frac{1}{3}
=
\int_0^1(1-t)^2\,dt
=
B(1,3),
}
\]
the $n=3$ instance of the universal $n$-gon coefficient $1/n$.
\end{theorem}

\begin{proof}
By the same method as Theorem~\ref{thm:sl3-triangle-coefficient},
the product side of the filtration-$3$ hexagon defect gives
\[
\frac{1}{3}\,[X_1,\,[X_2,\,X_3]],
\]
where the coefficient $1/3$ arises from the Dynkin path
BCH formula: $\log(e^{X_1}e^{X_2}e^{X_3})$ restricted to the
multilinear part of the path $\alpha_1\sim\alpha_2\sim\alpha_3$
gives $\frac{1}{3}[X_1,[X_2,X_3]]$ (using $[X_1,X_3]=0$ since
$\alpha_1$ and $\alpha_3$ are non-adjacent). The integral
representation is $\int_0^1(1-t)^2\,dt=\frac{1}{3}$.

On the coproduct side, $X_{\alpha_1+\alpha_2+\alpha_3}$ is the
unique height-$3$ generator whose twisted coproduct projects to the
quadrilateral sector. Equating coefficients gives the claim.

The Serre relations for $\mathfrak{sl}_4$ provide further
constraints: $(\operatorname{ad}e_i)^{1-A_{ij}}(e_j)=0$ with
$A_{ij}=-1$ for adjacent $i,j$ and $A_{13}=0$.
The last condition gives $[e_1,e_3]=0$: non-adjacent simple
roots commute.

The strictification obstruction vanishes for $\mathfrak{sl}_4$ by
root-space one-dimensionality (every root of $A_3$ has
multiplicity~$1$), exactly as for $\mathfrak{sl}_3$.
\end{proof}

\begin{remark}[The beta-integral staircase]
\label{rem:beta-staircase}
The pattern of $n$-gon coefficients
\[
\frac{1}{n}=\int_0^1(1-t)^{n-1}\,dt=B(1,n),
\qquad n=2,3,4,\ldots,
\]
gives a staircase indexed by Dynkin path length:
\[
\begin{array}{c|cccc}
\mathfrak{sl}_N & N=2 & N=3 & N=4 & N=5 \\
\hline
\text{path length }n & 1 & 2 & 3 & 4 \\
\text{coefficient} & 1 & 1/2 & 1/3 & 1/4
\end{array}
\]
At $n=1$ (the $\mathfrak{sl}_2$ case), there is no triangle sector
and the coefficient is~$1$: this is the degree-$2$ bar differential,
which is the Lie bracket itself. The genuine rank-$2$ phenomenon
begins at $n=2$ with the coefficient $1/2$. Each step up the
rank ladder introduces one additional nested commutator and one
additional power of $(1-t)$ in the integrand. The beta integral
provides a single uniform computation that replaces all the
case-by-case Baker--Campbell--Hausdorff algebra.

For general simple $\mathfrak{g}$, the maximal Dynkin path
length is the rank~$r$, and the deepest $n$-gon coefficient is
$1/r$. The strictification obstruction vanishes for all simple
Lie algebras by root-space one-dimensionality (every root in a
finite root system has multiplicity~$1$), confirming
Theorem~\ref{thm:km-yangian}.
\end{remark}


\subsection{Explicit computations and $R$-matrices: types $B_2$, $C_2$, $G_2$}
\label{subsec:BCG-yangians}

Table~\ref{tab:km-yangian-data} records summary data for all types.
We work out the three rank-$2$ families beyond type~$A$, exhibiting
root data, $R$-matrices in the fundamental representation, Serre
relations, and strictification.

\begin{theorem}[Ordered bar complexes and Yangian $R$-matrices for
rank-$2$ non-$A$ types]
\label{thm:B2-ordered-bar}
\label{thm:C2-ordered-bar}
\label{thm:G2-ordered-bar}
\label{thm:r-so5}
\label{thm:r-sp4}
\label{thm:r-g2}
For each of the rank-$2$ types below, the collision residue is
$r(z) = \hbar\,\Omega/z$ with $\hbar = 1/(k{+}h^\vee)$,
root-space one-dimensionality forces complete strictification, and
the RTT relation
$R_{12}(u{-}v)\,T_1(u)\,T_2(v)
= T_2(v)\,T_1(u)\,R_{12}(u{-}v)$
generates the Yangian.

\emph{Type $B_2$ \textup{(}$\mathfrak{so}_5$, $h^\vee=3$,
$V=\mathbb{C}^5$\textup{)}.}
Cartan matrix
$\bigl(\begin{smallmatrix} 2&-1\\-2&2
\end{smallmatrix}\bigr)$;
positive roots $\alpha_1, \alpha_2,
\alpha_1{+}\alpha_2, \alpha_1{+}2\alpha_2$; max height~$3$.
The Casimir image on $V\otimes V$ is $\Omega_V = P - K$,
where $K$ is the trace contraction.
The Zamolodchikov--Fateev $R$-matrix is
\begin{equation}\label{eq:B2-R-matrix}
R(u) = u\cdot\id + P - \frac{Q}{u+\varkappa},
\qquad \varkappa = \tfrac{N-2}{2} = \tfrac{3}{2},
\end{equation}
with eigenvalues $u{+}1$ on
$S^2_0V\cong\mathbf{14}$, $u{-}1$ on
$\Lambda^2V\cong\mathbf{10}$, and
$(u^2{+}\tfrac{5}{2}u{-}\tfrac{7}{2})/(u{+}\tfrac{3}{2})$ on
$\mathbf{1}$.
Serre: $[e_2,[e_2,[e_2,e_1]]]=0$ from
$(\mathfrak{so}_5)_{3\alpha_2+\alpha_1}=0$.

\emph{Type $C_2$ \textup{(}$\mathfrak{sp}_4$, $h^\vee=3$,
$V=\mathbb{C}^4$\textup{)}.}
Cartan matrix
$\bigl(\begin{smallmatrix} 2&-2\\-1&2
\end{smallmatrix}\bigr)$;
positive roots $\alpha_1, \alpha_2,
\alpha_1{+}\alpha_2, 2\alpha_1{+}\alpha_2$; max height~$3$.
The Casimir image uses the symplectic form~$J$:
$\Omega_V = P - K_J$, with
$(K_J)_{ij,kl} = J_{il}\,J_{jk}$.
The symplectic $R$-matrix is
\begin{equation}\label{eq:C2-R-matrix}
R(u) = u\cdot\id + P + \frac{K}{u-\varkappa},
\qquad \varkappa = n{+}1 = 3,
\end{equation}
with eigenvalues $u{+}1$ on $\mathrm{Sym}^2V\cong\mathbf{10}$,
$u{-}1$ on $\Lambda^2_0V\cong\mathbf{5}$, and
$(u^2{-}4u{-}1)/(u{-}3)$ on $\mathbf{1}$.
The sign of the contraction pole is \emph{opposite}
to type~$B$: this is the $R$-matrix--level manifestation
of Langlands duality $B_n \leftrightarrow C_n$.
Serre: $[e_1,[e_1,[e_1,e_2]]]=0$ from
$(\mathfrak{sp}_4)_{3\alpha_1+\alpha_2}=0$.

\emph{Type $G_2$ \textup{(}$\mathfrak{g}_2$, $h^\vee=4$,
$V=\mathbb{C}^7$\textup{)}.}
Cartan matrix
$\bigl(\begin{smallmatrix} 2&-1\\-3&2
\end{smallmatrix}\bigr)$;
positive roots $\alpha_1, \alpha_2, \alpha_1{+}\alpha_2,
2\alpha_1{+}\alpha_2, 3\alpha_1{+}\alpha_2,
3\alpha_1{+}2\alpha_2$; max height~$5$.
The tensor product
$V\otimes V \cong \mathbf{1} \oplus \mathbf{7} \oplus
\mathbf{14} \oplus \mathbf{27}$ decomposes into four
$G_2$-irreducibles with Casimir eigenvalues
$c_1=0$, $c_7=1/4$, $c_{14}=1/2$, $c_{27}=5/12$.
The $R$-matrix involves four equivariant projectors:
\begin{equation}\label{eq:G2-R-matrix}
R(u) = \sum_{\lambda\in\{1,7,14,27\}}
f_\lambda(u)\cdot\Pi_\lambda,
\end{equation}
with $f_{\mathbf{27}}=u{+}1$, $f_{\mathbf{14}}=u{-}1$,
$f_{\mathbf{7}}=(u{+}1)(u{-}2)/(u{+}2)$,
$f_{\mathbf{1}}=(u{+}1)(u{-}4)/(u{+}4)$.
The poles at $u=-2$ and $u=-4$ correspond to the two primitive
$G_2$-invariant tensors beyond $\id$ and~$P$.
The Serre relations follow from root-space vanishing:
$[e_2,[e_2,e_1]]=0$ since $\alpha_1{+}2\alpha_2 \notin
\Phi(G_2)$;
$[e_1,[e_1,[e_1,[e_1,e_2]]]]=0$ since
$4\alpha_1{+}\alpha_2 \notin \Phi(G_2)$.
The Jacobi collapse lemma handles the degree-$5$
star-sector at the longest root $3\alpha_1{+}2\alpha_2$.
\ClaimStatusProvedHere
\end{theorem}

\begin{proof}
All collision residues and root-space one-dimensionality
follow from Theorem~\ref{thm:km-yangian}. The $R$-matrices
are the unique $\mathfrak{g}$-invariant rational YBE solutions
with classical limit $R(u) = u\cdot\id + \hbar\,\Omega_V/u
+ O(u^{-2})$: for types $B$/$C$, the Brauer algebra relations
($P^2=\id$, $Q^2=NQ$ or $K^2=-2nK$) determine
all coefficients; for $G_2$,
$\dim\operatorname{End}_{G_2}(V\otimes V)=4$ and the YBE
plus unitarity $R(u)R_{21}(-u)\propto\id$ fix the eigenvalues
uniquely. Serre relations and strictification
are verified by direct root-space analysis.
\end{proof}

\begin{remark}[Comparison across rank-$2$ types]
\label{rem:rank-2-comparison}
\label{rem:langlands-ordered-bar}
\label{rem:G2-octonion}
\label{rem:b2-c2-isogeny}
\[
\begin{array}{c|ccccc}
\text{Type} & h^\vee & \dim V_{\mathrm{fund}}
& |\Phi^+| & \text{max height} & R\text{-matrix structure} \\
\hline
A_2 & 3 & 3 & 3 & 2 & P \text{ only} \\
B_2 & 3 & 5 & 4 & 3 & P,\, Q \\
C_2 & 3 & 4 & 4 & 3 & P,\, K_J \\
G_2 & 4 & 7 & 6 & 5 & \text{four projectors}
\end{array}
\]
The $R$-matrix complexity increases with
$|\operatorname{Irr}(V_{\mathrm{fund}} \otimes V_{\mathrm{fund}})|$:
type~$A$ has $2$ summands, types $B$/$C$ have $3$, $G_2$
has $4$.
The Langlands duality $B_2^\vee = C_2$ interchanges
$Q \leftrightarrow K_J$ and reverses the contraction-pole sign.
The $G_2$ representation $V = \mathbb{C}^7$ is the imaginary
octonions; the four-projector $R$-matrix encodes octonionic
structure constants.
\end{remark}

% ================================================================
% ACT IV: THE GENUS-1 LIFT
% ================================================================

\bigskip
The preceding sections develop the ordered bar complex at genus zero.
Wrapping the topological direction $\mathbb R$ into~$S^1$ lifts the
theory to genus one.

\section{The annular bar complex and genus-one structure}
\label{sec:annular-bar}

The strip ordered bar complex
$\Barchord(\cA)$
on $\FM_k(\mathbb{C})\times\mathrm{Conf}_k^{<}(\mathbb{R})$ computes
genus-zero data: the topological direction is $\mathbb{R}$, an interval,
and the resulting chain complex classifies the genus-zero
$\Eone$-chiral Koszul dual.
Replacing the strip~$\mathbb{R}$ by the circle
$S^1=\mathbb{R}/\mathbb{Z}$ wraps the topological direction
and produces a genus-one analogue: the \emph{annular bar complex}.
The wrapping converts the two-sided cobar complex into a
coHochschild complex, and the resulting object computes
$\HH^{\mathrm{ch}}_\bullet(\cA)$, connecting the ordered sector to the
modular operad at genus~$1$.

\subsection{Definition of the annular bar complex}

\begin{definition}[Annular bar complex]%
\label{def:annular-bar}%
\index{annular bar complex|textbf}%
\index{bar complex!annular}%
Let $\cA$ be a strongly admissible augmented $\Eone$-chiral
algebra and $C=\Barchord(\cA)$ its ordered
bar coalgebra, with $C_\Delta$ the diagonal bicomodule
(Theorem~\ref{thm:diagonal}).
The \emph{annular bar complex} is the chain complex
\[
B^{\mathrm{ann}}_\bullet(\cA)
\;:=\;
C_\bullet\!\bigl(
 \overline{\FM}{}^{\,\mathrm{ann}}_{r,s};\,
 \omega_{\mathrm{bar}}
\bigr)
\;\cong\;
\bigoplus_{n\ge 0}
 \mathrm{cyc}^n(C,C_\Delta),
\]
where $\mathrm{cyc}^n(C,C_\Delta):=
C_\Delta\widehat\otimes_C C^{\widehat\otimes\,n}$ is the
\emph{cyclic cotensor}:
the completed tensor product that cyclically identifies the
last right $C$-coaction with the first left $C$-coaction,
reflecting the periodicity of the circle.
\end{definition}

The annular configuration space
$\overline{\FM}{}^{\,\mathrm{ann}}_{r,s}$ parametrises
$r$ holomorphic marked points on~$\mathbb{C}$ and $s$ ordered
marked points on~$S^1$, with the Fulton--MacPherson
compactification in the holomorphic direction. It is the
geometric carrier for the genus-one ordered sector: the strip
$\FM_k(\mathbb{C})\times\mathrm{Conf}_k^{<}(\mathbb{R})$ wraps to
$\FM_k(\mathbb{C})\times\mathrm{Conf}_k^{\mathrm{cyc}}(S^1)$,
where $\mathrm{Conf}_k^{\mathrm{cyc}}(S^1)$ denotes cyclically
ordered configurations on the circle. The compactification
$\overline{\FM}{}^{\,\mathrm{ann}}_{r,s}$ is a manifold with
corners whose codimension-one boundary strata encode the four
components of the annular differential below.

\subsection{The annular bar differential}

\begin{theorem}[Annular bar differential;
\ClaimStatusProvedHere]%
\label{thm:annular-bar-differential}%
\index{annular bar complex!differential|textbf}%
The differential on $B^{\mathrm{ann}}_\bullet(\cA)$
decomposes as
\[
d^{\mathrm{ann}}
\;=\;
d_{\mathrm{int}}+d_{\mathrm{res}}
 +d_\Delta+d_{\mathrm{wrap}},
\]
where:
\begin{enumerate}[label=\textup{(\alph*)}]
\item $d_{\mathrm{int}}$: internal differential on each
 tensor factor (the intrinsic differential of~$\cA$).

\item $d_{\mathrm{res}}$: collision residue differential
 arising from OPE residues at interior collisions in the
 holomorphic direction
 (identical to the strip bar differential).

\item $d_\Delta$: comultiplication differential at
 non-wrapping boundary collisions, the coface maps
 $\delta_i$ for $1\le i\le n-1$ that split a tensor factor
 via the ordered coproduct.

\item $d_{\mathrm{wrap}}$: the wrap-around differential
 at the seam of the annulus, given by
\begin{equation}\label{eq:d-wrap}
d_{\mathrm{wrap}}\langle a_1|\cdots|a_n\rangle
\;=\;
\sum_{(a_n)}
(-1)^{\epsilon}\,
\langle a_1|\cdots|a_{n-1}|a_n'|a_n''\rangle
\;+\;
\sum_{(a_1)}
(-1)^{\epsilon'}\,
\langle a_1'|a_1''|a_2|\cdots|a_n\rangle,
\end{equation}
where the sums are over the ordered coproduct
$\Delta(a_n)=\sum a_n'\otimes a_n''$ and
$\Delta(a_1)=\sum a_1'\otimes a_1''$, and the signs
$\epsilon$, $\epsilon'$ are the Koszul signs dictated by the
suspended grading $|s^{-1}a|=|a|-1$.
\end{enumerate}
The identity $(d^{\mathrm{ann}})^2=0$ follows from
$\partial^2=0$ on the manifold with corners
$\overline{\FM}{}^{\,\mathrm{ann}}_{r,s}$: each boundary
face of codimension~$2$ appears as the common boundary
of exactly two codimension-$1$ faces, with opposite orientations.
\end{theorem}

\begin{proof}
The proof parallels that of the strip bar differential
(Construction~\ref{constr:bar-diff-collision}). The
codimension-$1$ boundary strata of
$\overline{\FM}{}^{\,\mathrm{ann}}_{r,s}$ are of four types:
\begin{enumerate}[label=(\roman*)]
\item an internal bubble separates at a single tensor factor
 (producing $d_{\mathrm{int}}$);
\item two or more holomorphic points collide away from the
 seam (producing $d_{\mathrm{res}}$);
\item two consecutive topological points collide away from
 the seam, splitting a tensor factor via the ordered
 coproduct (producing $d_\Delta$);
\item a topological point crosses the seam $0\sim 1$ of $S^1$,
 forcing the first or last tensor factor to split and
 re-enter at the opposite end (producing $d_{\mathrm{wrap}}$).
\end{enumerate}
The identity $(d^{\mathrm{ann}})^2=0$ is the algebraic
manifestation of $\partial^2=0$: every codimension-$2$
face of $\overline{\FM}{}^{\,\mathrm{ann}}_{r,s}$ is a
common boundary of exactly two codimension-$1$ faces, and
the orientations cancel pairwise.
\end{proof}

\begin{remark}[The wrap-around differential as the genus-one novelty]
\label{rem:wrap-around}
The wrap-around differential $d_{\mathrm{wrap}}$ is the
only component without a strip analogue. Algebraically,
it provides the coface maps $\delta_0$ and $\delta_n$ that
close the two-sided cobar complex into the coHochschild
complex: without it the complex computes
$\Cotor^{C^e}(C_\Delta,C_\Delta)$; with it, the
result is the self-cotrace
$C_\Delta\square_{C^e}^{\mathbf{L}}
C_\Delta\simeq
\coHoch^{\mathrm{ch}}_\bullet(C,C_\Delta)$.
\end{remark}

\subsection{Annular chiral Hochschild identification}

\begin{theorem}[Annular bar complex computes chiral Hochschild homology;
\ClaimStatusProvedHere]%
\label{thm:annular-HH}%
\index{annular bar complex!chiral Hochschild homology|textbf}%
\index{chiral Hochschild homology!via annular bar complex}%
Let $\cA$ be a strongly admissible $\Eone$-chiral
algebra with ordered bar coalgebra $C=\Barchord(\cA)$
and diagonal bicomodule $C_\Delta$.
There is a canonical quasi-isomorphism
\[
B^{\mathrm{ann}}_\bullet(\cA)
\;\simeq\;
\coHoch^{\mathrm{ch}}_\bullet(C,C_\Delta)
\;\simeq\;
\HH^{\mathrm{ch}}_\bullet(\cA),
\]
where the second equivalence is
Theorem~\textup{\ref{thm:HH-coHH-homology}}.
\end{theorem}

\begin{proof}
By Definition~\ref{def:annular-bar}, the annular bar complex
is the cyclic cotensor
$\bigoplus_{n\ge 0}\mathrm{cyc}^n(C,C_\Delta)$.
The coHochschild complex of a coalgebra~$C$ with coefficients in a
bicomodule~$M$ is the standard cochain complex
\[
\coHoch^{\mathrm{ch}}_\bullet(C,M)
\;=\;
\bigoplus_{n\ge 0}
M\widehat{\otimes}_C C^{\widehat\otimes\,n}
\]
with differential assembled from the coface maps $\delta_0,\ldots,\delta_n$.
For $M=C_\Delta$, this is precisely the
annular bar complex: the coface maps $\delta_1,\ldots,\delta_{n-1}$
produce $d_\Delta$ (non-wrapping boundary collisions), and the
coface maps $\delta_0$ and $\delta_n$ produce $d_{\mathrm{wrap}}$
(the wrap-around at the seam). The internal differential $d_{\mathrm{int}}$
and collision residue $d_{\mathrm{res}}$ are carried along as the
internal structure of each tensor factor.

The second equivalence is
Theorem~\ref{thm:HH-coHH-homology}, which identifies
$\coHoch^{\mathrm{ch}}_\bullet(C,C_\Delta)\simeq
\HH^{\mathrm{ch}}_\bullet(\cA)$ via the bimodule/bicomodule
equivalence of Theorem~\ref{thm:bimod-bicomod}.
\end{proof}

\subsection{Connection to the modular operad at genus~$1$}

\begin{remark}[The annular bar complex and the modular
$E_1$-chiral structure]
\label{rem:annular-modular}
\index{modular operad!genus-one annular bar}
The annular bar complex provides the ordered sector's
contribution to the genus-$1$ modular $E_1$-chiral operad.
The passage from strip to annulus
($\mathbb{R}\leadsto S^1$) is the topological-direction
analogue of the holomorphic passage from~$\mathbb{C}$ to the
elliptic curve $E_\tau$: both wrap an infinite direction into
a compact one, and both introduce monodromy.

In the modular operad framework:
\begin{enumerate}[label=(\roman*)]
\item The strip bar complex
 $\Barchord(\cA)$ is the genus-$0$ open sector,
 computing the $\Eone$ Koszul dual.
\item The annular bar complex
 $B^{\mathrm{ann}}_\bullet(\cA)$ is the genus-$1$ open
 sector, computing chiral Hochschild homology.
\item The genus-$1$ closed sector carries the curvature
 $d^2=\kappa(\cA)\cdot\omega_1$ from the Hodge bundle
 (Vol~I, Theorem~D).
\item The annular closure of
 Conjecture~\textup{\ref{conj:bordered}(4)} is the geometric
 realisation of the isomorphism
 $B^{\mathrm{ann}}_\bullet\simeq
 \coHoch^{\mathrm{ch}}_\bullet(C)$
 at the configuration-space level.
\end{enumerate}
The open trace formalism (Theorem~\ref{thm:ordered-open})
at genus~$1$ factors through the annular bar complex: the
decategorified trace $\mathrm{tr}\colon K_0(\mathcal{C}_{\mathrm{line}})\to
\HH^{\mathrm{ch}}_0(\cA)$ computes characters of line operators
by wrapping them around the annulus.
\end{remark}

\subsection{Curvature--braiding dichotomy at genus~$1$}

At genus~$0$, the curvature $\kappa(\cA)$ (measuring the failure of
formality, i.e.\ $d^2=\kappa\cdot\omega_g$ in the curved bar complex)
and the spectral $R$-matrix $R(z)$ (measuring the braiding monodromy
on the ordered configuration space) are independent structures:
the curvature is controlled by the chiral differential,
the braiding by the coassociative coproduct.
At genus~$1$, however, they meet: the spectral braiding must be
defined on the elliptic curve $E_\tau$, and the quasi-periodicity
of the elliptic propagator forces the two structures to interact.
The nature of this interaction is controlled entirely by the
lowest-pole coefficient $c_0$ of the $\lambda$-bracket.

\begin{remark}[Convention for this subsection]
\label{rem:lambda-bracket-convention-dichotomy}
Throughout this subsection and through
Definition~\ref{def:decoupled-entangled} we use the
\emph{raw} $\lambda$-bracket expansion
$\{a{}_\lambda b\}=\sum_{n\ge0}c_n(a,b)\,\lambda^n$
\textup{(}so $c_n$ is the $n$-th OPE mode coefficient
$a_{(n)}b$\textup{)} and parameterise the genus-$0$ classical
$r$-matrix by $r(z)=\sum c_n/z^{n+1}$. This differs from the
$d\log$-absorbed collision-residue convention used in
Theorem~\ref{thm:heisenberg-rmatrix}, in the Virasoro
collision-residue proposition below, and in
Theorem~\ref{thm:km-yangian}, where the simple-pole
coefficient of $r^{\mathrm{coll}}$ corresponds to $a_{(1)}b$
rather than $a_{(0)}b$ The two conventions are
related by the shift $n\mapsto n{-}1$ in the pole index.
Within the present subsection we keep the raw convention
because it aligns with the standard genus-$1$ dichotomy
(Heisenberg decoupled because $c_0=a_{(0)}b=0$, affine
entangled because $c_0=f^{ab}{}_cJ^c\ne0$).
\end{remark}

\begin{theorem}[Curvature--braiding dichotomy at genus~$1$;
\ClaimStatusProvedHere]
\label{thm:curvature-braiding-dichotomy}
\index{curvature!braiding dichotomy|textbf}
\index{elliptic spectral dichotomy}
\index{genus 1!curvature--braiding dichotomy}
Let $\cA$ be a logarithmic $\Eone$-chiral algebra and
write the $(-1)$-shifted $\lambda$-bracket on cohomology as
\[
\{a {}_\lambda b\}
\;=\;
c_0(a,b) + c_1(a,b)\,\lambda + c_2(a,b)\,\lambda^2 + \cdots\,,
\]
so that the genus-$0$ classical $r$-matrix has the Laurent
expansion
\[
r(z)
\;=\;
\frac{c_0}{z} + \frac{c_1}{z^2} + \frac{c_2}{z^3} + \cdots
\qquad \text{\textup{(}formal, $z\to\infty$\textup{)}}.
\]
At genus~$1$, the classical $r$-matrix on $E_\tau$ is
obtained by replacing each rational function $1/z^{n+1}$
\textup{(}$n\ge0$\textup{)} by its unique elliptic
counterpart with the same local singularity at $z=0$:
\begin{equation}\label{eq:elliptic-r-matrix}
r^{E_\tau}(z)
\;=\;
c_0\cdot\zeta(z|\tau)
\;+\;
\sum_{n\ge 1}
 \frac{(-1)^{n-1}}{n!}\,
 c_n\cdot\wp^{(n-1)}(z|\tau),
\end{equation}
where $\wp^{(0)}:=\wp$, and $\zeta(z|\tau)$, $\wp(z|\tau)$
are the Weierstrass zeta and $\wp$-functions with periods
$1,\tau$. The first three terms are
$c_0\,\zeta+c_1\,\wp-\tfrac{1}{2}c_2\,\wp'+\cdots$; the
prefactors $(-1)^{n-1}/n!$ come from the Laurent expansion
$\wp^{(m)}(z)=(-1)^m(m+1)!/z^{m+2}+O(1)$ setting $m=n-1$.
Then:
\begin{enumerate}[label=\textup{(\alph*)}]
\item \textup{(Quasi-periodicity.)}
The $B$-cycle monodromy of $r^{E_\tau}$ is
\begin{equation}\label{eq:r-monodromy}
r^{E_\tau}(z+\tau)-r^{E_\tau}(z)
\;=\;
2\eta_\tau\cdot c_0,
\end{equation}
where $\eta_\tau=\zeta(\tau/2|\tau)$ is the quasi-period.
The $A$-cycle monodromy is $2\eta_1\cdot c_0$.

\item \textup{(Decoupling: Cartan type.)}
If $c_0=0$ \textup{(}equivalently, $\{a{}_0 b\}=0$ for all
$a,b\in H^\bullet(\cA,Q)$\textup{)}, then $r^{E_\tau}(z)$
is a meromorphic function on~$E_\tau$
\textup{(}doubly periodic\textup{)}.
The spectral braiding is independent of the $B$-cycle
monodromy that sources the curvature
$\kappa(\cA)\cdot\omega_1$. The genus-$1$ $E_1$-chiral
bar complex splits:
\[
\Barchord(\cA)^{(1)}
\;\simeq\;
(\text{curved chiral differential})
\;\times\;
(\text{flat coassociative coproduct with meromorphic braiding}).
\]

\item \textup{(Entanglement: Yangian type.)}
If $c_0\ne 0$, then $r^{E_\tau}(z)$ is quasi-periodic
with non-trivial $B$-cycle monodromy
$2\eta_\tau\cdot c_0$. The spectral braiding is a
section of a flat connection on the universal cover
of~$E_\tau$ with holonomy~$c_0$, and the genus-$1$
$E_1$-chiral bar complex does \emph{not} split.
The braiding and the curvature share a common geometric
source: the $B$-cycle monodromy of the Weierstrass zeta
function.
\end{enumerate}
\end{theorem}

\begin{proof}
\textbf{Step~1} (Elliptic propagator).
At genus~$0$, the holomorphic propagator is
$K^{\mathbb{C}}(z)=dz/z$ (the $d\log$ kernel). At
genus~$1$, on the elliptic curve~$E_\tau$, the unique meromorphic
one-form with a simple pole at $z=0$ and prescribed quasi-periodicity
is $K^{E_\tau}(z)=\zeta(z|\tau)\,dz$. This has the same local
singularity $\zeta(z|\tau)=1/z+O(z)$ near $z=0$ but acquires
the quasi-periodicity
$\zeta(z+1|\tau)=\zeta(z|\tau)+2\eta_1$ and
$\zeta(z+\tau|\tau)=\zeta(z|\tau)+2\eta_\tau$.

\textbf{Step~2} (Expansion of the elliptic $r$-matrix).
The genus-$0$ $r$-matrix $r(z)=\sum_{n\ge 0}c_n/z^{n+1}$
arises from pairing the propagator $dz/z$ with the
$\lambda$-bracket coefficients. At genus~$1$, each
rational function $1/z^{n+1}$ is replaced by its unique
elliptic counterpart: $1/z\mapsto\zeta(z|\tau)$ and
$1/z^{n+1}\mapsto\frac{(-1)^{n-1}}{n!}\,\wp^{(n-1)}(z|\tau)$
for $n\ge 1$, with $\wp^{(0)}:=\wp$. These are the unique
functions on $\mathbb{C}\setminus(\mathbb{Z}+\mathbb{Z}\tau)$
that match $1/z^{n+1}+O(z)$ near $z=0$ and are doubly
periodic for $n\ge 1$ (or quasi-periodic for $n=0$). The
prefactor comes from the Laurent expansion
$\wp^{(m)}(z)=(-1)^m(m+1)!/z^{m+2}+O(1)$ specialised to
$m=n-1$. This yields formula~\eqref{eq:elliptic-r-matrix}.

\textbf{Step~3} (Dichotomy).
The quasi-periodicity is inherited entirely from $\zeta(z|\tau)$,
since $\wp$ and all its derivatives are doubly periodic.
Therefore the $B$-cycle monodromy of $r^{E_\tau}$ is
$2\eta_\tau\cdot c_0$, proving~(a). When $c_0=0$, the
$\zeta$-term is absent and $r^{E_\tau}$ is meromorphic on
$E_\tau$, hence the spectral braiding defined by the annular bar
differential decouples from the curved chiral differential. When
$c_0\ne 0$, the non-trivial monodromy
forces the braiding to be a section of a flat connection with
holonomy $\exp(2\pi i\cdot c_0)$, and this flat connection
is the same geometric object whose $B$-cycle monodromy
controls the curvature $\kappa(\cA)\cdot\omega_1$.
\end{proof}

\begin{remark}[Classification by pole structure]
\label{rem:annular-pole-classification}
\index{pole-order dichotomy!annular}
The dichotomy partitions all $\Einf$-chiral algebras
(i.e.\ all vertex algebras, which are local and
$\Sigma_n$-equivariant) into two classes at genus~$1$:
\begin{center}
\begin{tabular}{l@{\quad}c@{\quad}c@{\quad}l}
\toprule
\textbf{Class} & $c_0$ & \textbf{Monodromy} & \textbf{Example} \\
\midrule
Cartan (decoupled) & $=0$ & trivial &
 Heisenberg $\mathcal{H}_k$, free fermion \\
Yangian (entangled) & $\ne 0$ & $2\eta_\tau\cdot c_0$ &
 $\widehat{\mathfrak{g}}_k$ ($\mathfrak{g}$ simple),
 Virasoro \\
\bottomrule
\end{tabular}
\end{center}
For the Heisenberg algebra $\mathcal{H}_k$, the OPE
$J(z)J(w)\sim k/(z-w)^2$ has $c_0=\{J{}_{(0)}J\}=0$ (no
simple pole in the $\lambda$-bracket: the lowest nonvanishing
coefficient is $c_1=k$). The genus-$1$ $r$-matrix is
$r^{E_\tau}(z)=k\cdot\wp(z|\tau)$, which is doubly periodic.
For $\widehat{\mathfrak{sl}}_2$ at level~$k$, the OPE
$J^a(z)J^b(w)\sim k\,\kappa^{ab}/(z-w)^2+f^{ab}{}_c\,J^c(w)/(z-w)$
has $c_0=f^{ab}{}_c\,J^c\ne 0$, and the genus-$1$ structure
is genuinely entangled.
\end{remark}



\section{Quantum groups from genus-1 monodromy}
\label{sec:genus1-monodromy}

The preceding sections develop the ordered bar complex and its $R$-matrix
on the genus-zero configuration space $\mathrm{Conf}_n(\mathbb{C})$.
At genus~$1$ the configuration space is replaced by
$\mathrm{Conf}_n(E_\tau)$, where $E_\tau = \mathbb{C}/(\mathbb{Z}
+ \tau\mathbb{Z})$ is the elliptic curve with modulus~$\tau$.
The fundamental group of $\mathrm{Conf}_n(E_\tau)$ is
richer than the pure braid group $P_n$: it contains two
additional generators (the $A$-cycle and $B$-cycle monodromies),
and their interaction with the braid generators produces the quantum
group $U_q(\mathfrak{g})$ from the Yangian $Y_\hbar(\mathfrak{g})$.

The genus-$1$ lift for the affine lineage
$V^k(\mathfrak{g})$ combines the Drinfeld--Kohno theorem with the
one-loop collapse mechanism established in Volume~II, giving a
rigorous passage from the KZ connection to the KZB connection and
thence to the quantum group.

\subsection{The KZ connection on ordered configurations}
\label{sec:kz-connection}

The genus-$0$ starting point is as follows. Let $\mathfrak{g}$ be a simple
Lie algebra and $k \in \mathbb{C} \setminus \{-h^\vee\}$ generic.
Choose an orthonormal basis $\{I^a\}$ for~$\mathfrak{g}$ with respect
to the Killing form normalised so that
$\operatorname{tr}(I^a I^b) = \delta^{ab}$. The Casimir tensor is
$\Omega = \sum_a I^a \otimes I^a \in \mathfrak{g} \otimes \mathfrak{g}$.

\begin{definition}[KZ connection]
\ClaimStatusProvedHere
\label{def:kz-connection}
\index{KZ connection|textbf}
On the trivial bundle over $\mathrm{Conf}_n(\mathbb{C})$ with fibre
$V_1 \otimes \cdots \otimes V_n$ (finite-dimensional
$\mathfrak{g}$-modules), the \emph{Knizhnik--Zamolodchikov connection}
is
\begin{equation}\label{eq:kz}
\nabla_{\mathrm{KZ}}
\;=\;
d
\;-\;
\hbar
\sum_{1 \leq i < j \leq n}
\Omega_{ij}\,d(z_i - z_j)/(z_i - z_j),
\end{equation}
where $\hbar = 1/(k + h^\vee)$ and $\Omega_{ij}$ denotes the
Casimir acting in the $(i,j)$ tensor slots.
\end{definition}

\begin{remark}
The connection~\eqref{eq:kz} is a special case of the Kohno
connection~\eqref{eq:kohno-connection}: for the affine algebra
$V^k(\mathfrak{g})$, the collision residue $r_{ij}(z)$ extracted by
the $d\log$ kernel (cf.\: the kernel absorbs one pole) is
$r_{ij}(z) = \hbar\,\Omega_{ij}/z$, which has the expected
simple pole.
The one-loop collapse mechanism from Volume~II
ensures that the higher $A_\infty$ operations $m_k$ ($k \geq 3$) do not
contribute on evaluation modules, so the full superconnection reduces to
the KZ connection. The flatness $\nabla_{\mathrm{KZ}}^2 = 0$ follows
from $[\Omega_{ij}, \Omega_{ik} + \Omega_{jk}] = 0$, which is the
$\mathfrak{g}$-equivariance of the Casimir.
\end{remark}

\subsection{Genus-0 monodromy: the $R$-matrix}
\label{sec:genus0-monodromy}

At genus~$0$, the fundamental group
$\pi_1(\mathrm{Conf}_n(\mathbb{C}))$ is the pure braid group~$P_n$.
The monodromy of $\nabla_{\mathrm{KZ}}$ around the generator
$\gamma_{ij}$ (in which $z_i$ circles $z_j$) gives the $R$-matrix:
\[
R_{ij}
\;=\;
\operatorname{Mon}_{\gamma_{ij}}(\nabla_{\mathrm{KZ}})
\;=\;
\mathcal{P}\exp\!\Bigl(
\hbar
\oint_{\gamma_{ij}} \Omega_{ij}\,d(z_i - z_j)/(z_i - z_j)
\Bigr)
\;=\;
e^{2\pi i\hbar\,\Omega_{ij}}.
\]
This is the \emph{pure-braid monodromy}: it encodes how parallel
transport twists when one point winds around another on the
genus-$0$ surface.
The Yang--Baxter equation
$R_{12}\,R_{13}\,R_{23} = R_{23}\,R_{13}\,R_{12}$
follows from flatness (eq.~\eqref{eq:cybe-vol1} and its exponentiated form),
as proved in Proposition~\ref{sec:r-matrix-descent-vol1}.

\subsection{The KZB connection on the elliptic curve}
\label{sec:kzb-connection}

At genus~$1$ the base curve is $E_\tau$ and the configuration space
$\mathrm{Conf}_n(E_\tau)$ has fundamental group containing two
distinguished classes of generators beyond the braid generators: the
$A$-cycle (a point traverses the real period~$1$) and the $B$-cycle
(a point traverses the imaginary period~$\tau$). The KZ connection
must be replaced by the \emph{Knizhnik--Zamolodchikov--Bernard} (KZB)
connection, which encodes the genus-$1$ monodromy of conformal blocks.

\begin{definition}[KZB connection]
\ClaimStatusProvedHere
\label{def:kzb-connection}
\index{KZB connection|textbf}
On the trivial bundle over $\mathrm{Conf}_n(E_\tau)$ with fibre
$V_1 \otimes \cdots \otimes V_n$, the \emph{KZB connection} is
\begin{equation}\label{eq:kzb}
\nabla_{\mathrm{KZB}}
\;=\;
d
\;-\;
\hbar
\sum_{1 \leq i < j \leq n}
\Omega_{ij}\,\wp_1(z_i - z_j,\tau)\,dz_{ij}
\;-\;
\frac{\hbar}{2\pi i}
\sum_{i=1}^n
D_i\,d\tau,
\end{equation}
where $\wp_1(z,\tau) = \partial_z \log\theta_1(z,\tau)$ is the
logarithmic derivative of the Jacobi theta function (the elliptic
analogue of $1/z$), $z_{ij} = z_i - z_j$, and $D_i$ is the modular
Hamiltonian
\[
D_i
\;=\;
\sum_{j \neq i}
\bigl(
\wp(z_{ij},\tau)\,\Omega_{ij}
\;+\;
\text{lower order}
\bigr),
\]
with $\wp(z,\tau) = -\partial_z \wp_1(z,\tau)$ the Weierstrass
$\wp$-function.
\end{definition}

\begin{remark}
The KZB connection is a genus-$1$ deformation of the KZ connection. In
the degeneration limit $\tau \to i\infty$ (the nodal rational curve),
$\wp_1(z,\tau) \to \pi\cot(\pi z) \to 1/z + O(z)$, and the KZB
connection degenerates to the KZ connection. The $d\tau$ term,
absent at genus~$0$, records the dependence of conformal blocks on the
modulus~$\tau$ and is the source of the modular anomaly: it measures the
failure of conformal blocks to form a flat bundle over the moduli
space~$\mathcal{M}_{1,n}$.
\end{remark}

\begin{remark}[Bar-complex origin]
\label{rem:kzb-bar}
In the $E_1$-chiral framework, the KZB connection arises from the
ordered bar complex on the elliptic curve $E_\tau$ equipped with its
logarithmic chiral coassociative structure. The collision residue on
$\mathrm{Conf}_n^{\mathrm{ord}}(E_\tau)$ is
$r_{ij}^{E_\tau}(z) = \hbar\,\Omega_{ij}\,\wp_1(z,\tau)$,
which replaces the genus-$0$ collision residue $r_{ij} = \hbar\,\Omega_{ij}/z$
by its elliptic regularisation. The additional $d\tau$ connection
component encodes the curvature $\kappa(A) \cdot \omega_1$ from the
Hodge bundle on $\overline{\mathcal{M}}_{1,n}$: the curved bar
structure $d^2 = \kappa(A) \cdot \omega_1$ at genus~$1$ is precisely
the non-flatness of the naive genus-$0$ connection, corrected by the
Bernard term $D_i\,d\tau$.
\end{remark}

\subsection{$B$-cycle monodromy and the quantum group parameter}
\label{sec:b-cycle-monodromy}

The decisive new phenomenon at genus~$1$ is the $B$-cycle monodromy.

\begin{theorem}[Quantum group from $B$-cycle monodromy; \ClaimStatusProvedHere]
\label{thm:b-cycle-quantum-group}
\index{quantum group!from genus-1 monodromy|textbf}
\index{monodromy!B-cycle|textbf}
Let $\mathfrak{g}$ be a simple Lie algebra and $k \in \mathbb{C}
\setminus \{-h^\vee\}$ generic. Consider the KZB monodromy
representation
\[
\rho_n^{\mathrm{KZB}}
\;\colon\;
\pi_1\bigl(\mathrm{Conf}_n(E_\tau)\bigr)
\;\longrightarrow\;
\operatorname{GL}(V_1 \otimes \cdots \otimes V_n).
\]
\begin{enumerate}[label=\textup{(\roman*)}]
\item \textbf{$A$-cycle monodromy.}
The monodromy around the $A$-cycle generator $\alpha_i$ (in which
$z_i \mapsto z_i + 1$, all other points fixed) gives the $R$-matrix:
\[
\rho_n^{\mathrm{KZB}}(\alpha_i)
\;=\;
R_i(z),
\]
which is the spectral $R$-matrix of
Construction~\textup{\ref{constr:r-matrix-monodromy-vol1}},
now defined via the elliptic collision residue.

\item \textbf{$B$-cycle monodromy.}
The monodromy around the $B$-cycle generator $\beta_i$ (in which
$z_i \mapsto z_i + \tau$, all other points fixed) produces the
quantum group parameter:
\[
\rho_n^{\mathrm{KZB}}(\beta_i)
\;=\;
q^{2H_i}
\;\cdot\;
(\text{braiding corrections}),
\]
where
\begin{equation}\label{eq:q-parameter}
q
\;=\;
e^{\pi i\hbar}
\;=\;
e^{\pi i/(k + h^\vee)}
\end{equation}
and $H_i$ is the Cartan action in the $i$-th tensor slot.

\item \textbf{Quantum group structure.}
The combined $A$-cycle and $B$-cycle monodromies, together with the braid
monodromies, generate a representation of the quantum group
$U_q(\mathfrak{g})$ on $V_1 \otimes \cdots \otimes V_n$ at the parameter
$q = e^{\pi i\hbar}$.
\end{enumerate}
\end{theorem}

\begin{proof}
\textbf{(i)} is immediate from the fact that $A$-cycle monodromy around
$z_i \mapsto z_i + 1$ corresponds to the quasi-periodicity of
$\wp_1(z,\tau)$ along the real period, which introduces the same
braid-type contribution as at genus~$0$.

\medskip
\noindent\textbf{(ii).}
The $B$-cycle monodromy arises from $z_i \mapsto z_i + \tau$. The
quasi-periodicity of the Jacobi theta function gives
$\wp_1(z + \tau, \tau) = \wp_1(z, \tau) - 2\pi i$, so the connection
form acquires a constant shift $-2\pi i \cdot \hbar\,\Omega_{ij}$
from each pair. The parallel transport along the
$B$-cycle path $z_i \mapsto z_i + t\tau$ ($t \in [0,1]$) includes
both the $dz$ contribution (from the connection one-form) and the
$d\tau$ contribution (from the modular Hamiltonian $D_i$). On
evaluation modules, the one-loop collapse ensures that only the
leading Casimir term contributes, and the result is
\[
\rho_n^{\mathrm{KZB}}(\beta_i)
\;=\;
\exp\bigl(2\pi i \hbar \cdot H_i + \text{off-diagonal}\bigr),
\]
where $H_i$ is the Cartan element acting in the $i$-th slot. The
full-cycle exponential $e^{2\pi i \hbar} = q^2$ determines the
standard quantum-group parameter $q = e^{\pi i\hbar}$. The
``braiding corrections'' are conjugations by
$R$-matrices arising from the interaction of the $B$-cycle path with
the other points; they are controlled by the mixed $A$-$B$ relations
in $\pi_1(\mathrm{Conf}_n(E_\tau))$.

\medskip
\noindent\textbf{(iii).}
The Drinfeld--Kohno theorem (genus-$1$ version, see the discussion
below) proves that the monodromy representation of the KZB connection
factors through $U_q(\mathfrak{g})$. The $A$-cycle monodromies give
the $R$-matrix (the braiding in $\mathrm{Rep}_q(\mathfrak{g})$),
the full $B$-cycle monodromies give the Cartan operators $q^{2H}$,
and after choosing the standard half-monodromy convention one recovers
the usual $q^H$ operators. Together with the braid monodromies they
generate the full quantum group action.
\end{proof}

\begin{remark}[Physical origin of $q$]
\label{rem:q-physical}
The parameter $\hbar = 1/(k+h^\vee)$
has a transparent origin in the holomorphic-topological framework.
At genus~$0$, the bar complex on $\mathrm{FM}_n(\mathbb{C}) \times
\mathrm{Conf}_n(\mathbb{R})$ produces the Yangian $Y_\hbar(\mathfrak{g})$,
a deformation over the additive parameter $\hbar = 1/(k + h^\vee)$.
At genus~$1$, the $B$-cycle monodromy exponentiates the additive
parameter~$\hbar$ into the full-cycle factor $e^{2\pi i\hbar}=q^2$.
The standard quantum-group parameter is therefore
$q = e^{\pi i \hbar}$: the passage from the additive formal group
$\widehat{\mathbb{G}}_a$ to the multiplicative formal group
$\widehat{\mathbb{G}}_m$ is the passage from the Yangian to the
quantum group. The elliptic curve $E_\tau$ provides the geometric
mechanism: its $B$-cycle is the Wick-rotated time circle in the
holomorphic-topological theory, and the monodromy around it converts
the infinitesimal deformation parameter into a finite one.
\end{remark}

\begin{remark}[EF comparison]
\label{rem:ef-comparison}
The affine normalization used here matches the standard KZ/KZB
literature after one explicit bridge. In
Remark~\ref{rem:km-collision-residue-rmatrix} we fixed the
\emph{KZ normalization}
\[
r^{\mathrm{KZ}}(z) = \frac{\Omega}{(k+h^\vee)\,z}
= \hbar\,\frac{\Omega}{z},
\qquad
\hbar = \frac{1}{k+h^\vee},
\]
while the landscape census often records the same tensor in the
trace-form convention as $k\,\Omega_{\mathrm{tr}}/z$; the bridge is
$k\,\Omega_{\mathrm{tr}} = \Omega/(k+h^\vee)$. With this choice,
\eqref{eq:kz} is the usual affine KZ connection of
Etingof--Frenkel--Kirillov~\cite[Ch.~1]{EFK98}. Our monodromy formula
uses pure braid loops, hence gives the full-loop operator
$e^{2\pi i\hbar\Omega}$; the Drinfeld--Kohno braid-generator
convention uses the corresponding half-monodromy, so the standard
quantum-group parameter is $q = e^{\pi i\hbar}$.

At genus~$1$, equation~\eqref{eq:kzb} matches Bernard's torus KZB
system~\cite[\S\S3--4]{Bernard88} and the ordinary elliptic KZB
connection of Calaque--Enriquez--Etingof~\cite[Prop.~6.6]{CEE09}:
the $dz$ term carries the coefficient $1/(k+h^\vee)$, while the
$d\tau$ term carries the additional factor $1/(2\pi i)$ forced by the
theta-function heat equation. Drinfeld's original Yangian
$Y(\mathfrak{g})$~\cite{Drinfeld85} is recovered from our ordered bar
construction after passing to degree zero and evaluation modules; the
dg-shifted Yangian is a chain-level enhancement, not a different
rational quantum group.

Felder's elliptic quantum-group picture and the
Etingof--Varchenko theory~\cite{Felder94,EV98} are \emph{dynamical}:
the elliptic $R$-matrix depends on a Cartan variable
$\lambda \in \mathfrak{h}^*$ and satisfies the dynamical, not
ordinary, Yang--Baxter equation. The affine genus-$1$ lane in this
chapter should therefore be compared directly with Bernard/CEE.
Comparison with Felder--Etingof--Varchenko requires adjoining the
dynamical Cartan variable, or equivalently performing a vertex-face
gauge transform. In particular,
Conjecture~\ref{conj:trig-elliptic-ordered} is not an identification
of ordinary Belavin monodromy with the dynamical elliptic quantum
group without this extra datum.
\end{remark}

\begin{observation}[Sesquilinearity forces the additive formal group]
\label{obs:sesquilinearity-additive}
\ClaimStatusProvedHere
\index{formal group!additive, from sesquilinearity}
The spectral parameter $\lambda$ in the $\lambda$-bracket
$\{a_\lambda b\}$ carries the additive formal group law
$F(\lambda,\mu) = \lambda + \mu$. Sesquilinearity
$\{\partial a_\lambda b\} = -\lambda\{a_\lambda b\}$
identifies $\lambda$ as the Fourier dual of translation;
the Fourier dual of the additive group is the additive group.
A multiplicative law $F(\lambda,\mu) = \lambda + \mu + \lambda\mu$
would require
$\{\partial a_\lambda b\} = -(1+\lambda)\{a_\lambda b\} \cdot
\log(1+\lambda)/\lambda$,
which is not polynomial in~$\lambda$.
\end{observation}

\begin{observation}[Endomorphism ring of $\widehat{\mathbb{G}}_a$ and the Wick rotation]
\label{obs:end-additive}
\ClaimStatusProvedHere
\index{formal group!endomorphism ring}
\index{Wick rotation!formal group interpretation}
Over any $\C$-algebra~$R$,
$\End(\widehat{\mathbb{G}}_{a/R}) = R$:
every element $\alpha \in R$ defines an endomorphism
$[\alpha]\colon x \mapsto \alpha x$, and these exhaust all
endomorphisms since $\widehat{\mathbb{G}}_a$ has no nonlinear
endomorphisms in characteristic zero.
In particular, $[i] \in \End(\widehat{\mathbb{G}}_{a/\C})$:
the Wick rotation $\lambda \mapsto i\lambda$
(equivalently $t \mapsto it$ in the topological direction of
$\C_z \times \mathbb{R}_t$) is an automorphism of the
spectral-parameter formal group.
\end{observation}

\begin{remark}[Spectralization contracts the endomorphism ring]
\label{rem:spectralization-contraction}
\ClaimStatusHeuristic
\index{spectralization!endomorphism contraction}
\index{chromatic height!and Wick rotation}
Over $\bF_p$, the additive formal group has
height~$\infty$ ($[p](x) = px = 0$), not height~$1$.
Upon spectralization (passing from chain complexes to spectra,
replacing $\widehat{\mathbb{G}}_a$ by a formal group of finite
height), the endomorphism ring contracts:
\[
\renewcommand{\arraystretch}{1.4}
\begin{array}{@{}lcll@{}}
\widehat{\mathbb{G}}_a & \colon &
 \End = \C & \text{(height $\infty$ over $\bF_p$)} \\
\widehat{\mathbb{G}}_m & \colon &
 \End = \Z_p & \text{(height $1$)} \\
\widehat{E}_\tau & \colon &
 \End = \text{order in } \mathbb{Q}(\sqrt{-d})
 & \text{(height $\leq 2$)}
\end{array}
\]
The ring $\Z_p$ does not contain~$i$ for $p \neq 2$ unless
$p \equiv 1 \pmod{4}$, and even then $i \in \Z_p$ is a
$p$-adic integer, not the complex Wick rotation.
For a generic elliptic curve, $\End(\widehat{E}_\tau) = \Z$,
which excludes~$[i]$ entirely.

The sole exception is the CM curve $E_i$ with $j$-invariant
$j = 1728$: here $\End(\widehat{E}_i) \supseteq \Z[i]$, so $[i]$
is an endomorphism. This is the unique elliptic curve
over~$\C$ (up to isomorphism) whose formal group admits a Wick
rotation: an arithmetic accident, not a structural feature.

The bar-cobar theory of this monograph is a theory of
chain complexes over~$\C$. Every chain complex is an
$H\C$-module, hence rational, hence chromatic height~$0$.
The endomorphism ring $\End(\widehat{\mathbb{G}}_a) = \C$ reflects
this: the theory admits all constant endomorphisms, including
$[i]$. Positive chromatic height enters only after
spectralization, which contracts the endomorphism ring from $\C$
to an arithmetic object ($\Z_p$ or an order in a quadratic
field) and generically excludes the Wick rotation
(cf.\ the chromatic/shadow discussion in Chapter~\ref{chap:concordance}).
\end{remark}

\subsection{The Drinfeld--Kohno theorem}
\label{sec:drinfeld-kohno}

The identification of the KZ/KZB monodromy with the quantum group
action is the content of the Drinfeld--Kohno theorem and its
generalisations.

\begin{theorem}[Drinfeld--Kohno; \ClaimStatusProvedHere{} for the affine lineage]
\label{thm:drinfeld-kohno}
\index{Drinfeld--Kohno theorem|textbf}
\index{monodromy!Drinfeld--Kohno identification}
Let $\mathfrak{g}$ be a simple Lie algebra, $k \in \mathbb{C}
\setminus \{-h^\vee\}$ generic, and $V_1, \dots, V_n$
finite-dimensional $\mathfrak{g}$-modules.
\begin{enumerate}[label=\textup{(\roman*)}]
\item \textbf{Genus $0$.}
The monodromy representation of the KZ
connection~\eqref{eq:kz} on $\mathrm{Conf}_n(\mathbb{C})$
is isomorphic to the braid group representation obtained from the
braided monoidal structure of $\mathrm{Rep}_q(\mathfrak{g})$
via the universal $R$-matrix of $U_q(\mathfrak{g})$:
\begin{equation}\label{eq:dk}
\rho_n^{\mathrm{KZ}}
\;\simeq\;
\rho_n^{U_q(\mathfrak{g})}.
\end{equation}

\item \textbf{Genus $1$.}
The monodromy representation of the KZB
connection~\eqref{eq:kzb} on $\mathrm{Conf}_n(E_\tau)$
factors through $U_q(\mathfrak{g})$ and recovers the
full quantum group action \textup{(}including the Cartan monodromy,
written conventionally as $q^H$ after taking the square root of the
full-cycle operator $q^{2H}$\textup{)} on the tensor product
$V_1 \otimes \cdots \otimes V_n$.
\end{enumerate}
\end{theorem}

\begin{proof}
\textbf{(i)} is the classical Drinfeld--Kohno theorem. The proof
proceeds in three steps.

\emph{Step~1: Drinfeld associator.}
Drinfeld constructs a formal series
$\Phi \in U(\mathfrak{g})^{\widehat\otimes 3}[\![\hbar]\!]$
satisfying the pentagon and hexagon equations, using iterated integrals
of $d\log(z_i - z_j)$ on $\mathrm{Conf}_3(\mathbb{C})$. This
associator converts the infinitesimal braid relations (the CYBE) into
the finite braid relations (the YBE), providing an equivalence of
braided monoidal categories between the representation category of the
Drinfeld--Kohno Lie algebra $\mathfrak{t}_n$ (with its KZ braiding)
and $\mathrm{Rep}_q(\mathfrak{g})$ (with its universal $R$-matrix
braiding).

\emph{Step~2: One-loop collapse.}
For the affine algebra $V^k(\mathfrak{g})$, the one-loop exactness of
the BV-BRST differential (Theorem~\ref*{thm:affine-half-space-bv} of
Volume~II) ensures that the higher $A_\infty$ operations $m_k$
($k \geq 3$) vanish on evaluation modules. The full superconnection
on the bar complex therefore collapses to the KZ connection, and the
bar-complex monodromy $\rho_n^{\mathrm{HT}}$ equals $\rho_n^{\mathrm{KZ}}$
(the one-loop collapse statement of Volume~II).

\emph{Step~3: Identification.}
Combining Steps~1 and~2, the bar-complex monodromy for the affine
lineage equals the quantum-group braid representation:
$\rho_n^{\mathrm{HT}} = \rho_n^{\mathrm{KZ}} \simeq
\rho_n^{U_q(\mathfrak{g})}$.

\medskip
\noindent\textbf{(ii)} follows from the genus-$1$ generalisation of the
Drinfeld--Kohno theorem (Etingof--Schiffmann, Calaque--Enriquez--Etingof):
the KZB monodromy representation on $\mathrm{Conf}_n(E_\tau)$ factors
through $U_q(\mathfrak{g})$, with the full $B$-cycle monodromy providing
the operator $q^{2H}$ and hence, after the standard square-root
convention, the usual $q^H$ operators. The one-loop collapse for the affine lineage ensures
that the bar-complex monodromy on $\mathrm{Conf}_n(E_\tau)$ equals the
KZB monodromy, and the identification follows.
\end{proof}

\subsection{The genus-$1$ deformation: from Yangian to quantum group}
\label{sec:yangian-to-quantum}

The passage from genus~$0$ to genus~$1$ can be summarised as a
deformation of the Yangian $Y_\hbar(\mathfrak{g})$ to the
quantum group $U_q(\mathfrak{g})$.

\begin{theorem}[Yangian--quantum group deformation for the affine lineage]
\ClaimStatusProvedHere
\label{thm:yangian-quantum-group}
\index{Yangian!genus-1 deformation to quantum group|textbf}
\index{quantum group!from Yangian|textbf}
Let $\mathfrak{g}$ be a simple Lie algebra and $k \in \mathbb{C}
\setminus \{-h^\vee\}$ generic.
\begin{enumerate}[label=\textup{(\roman*)}]
\item At genus~$0$, the ordered bar complex of $V^k(\mathfrak{g})$ on
 $\mathrm{FM}_n(\mathbb{C}) \times \mathrm{Conf}_n(\mathbb{R})$
 produces the Yangian $Y_\hbar(\mathfrak{g})$ as the line-operator
 algebra, with $\hbar = 1/(k + h^\vee)$.

\item At genus~$1$, the ordered bar complex on
 $E_\tau \times \mathbb{R}$ produces the quantum group
 $U_q(\mathfrak{g})$ as the line-operator algebra, with
 $q = e^{\pi i \hbar}$.

\item The degeneration $\tau \to i\infty$ \textup{(}pinching the
 $B$-cycle\textup{)} recovers the Yangian from the quantum group:
 $\lim_{\tau \to i\infty} U_q(\mathfrak{g})
 = Y_\hbar(\mathfrak{g})$, in the sense that the KZB connection
 degenerates to the KZ connection and the $B$-cycle monodromy becomes
 trivial.

\item The full $B$-cycle factor $e^{2\pi i \hbar} = q^2$ is the
 exponential of the Yangian parameter~$\hbar$ under the exponential map
 $\widehat{\mathbb{G}}_a \to \widehat{\mathbb{G}}_m$ provided by the
 $B$-cycle of $E_\tau$; the standard quantum-group parameter is its
 square root $q = e^{\pi i \hbar}$.
\end{enumerate}
\end{theorem}

\begin{proof}
Part~(i) is the content of the genus-$0$ programme: the bar differential
on $\mathrm{FM}_n(\mathbb{C})$ encodes the chiral OPE, the $E_1$
coproduct on $\mathrm{Conf}_n(\mathbb{R})$ encodes the associative
composition, and the resulting algebra is the Yangian (see the
holographic dictionary of Volume~II,
Corollary~\ref*{cor:holographic-dictionary}).

Part~(ii) combines Theorem~\ref{thm:b-cycle-quantum-group} with the
Drinfeld--Kohno theorem (Theorem~\ref{thm:drinfeld-kohno}(ii)): the
full monodromy of the KZB connection generates $U_q(\mathfrak{g})$.

Part~(iii) follows from the degeneration
$\wp_1(z, \tau) \to 1/z + O(z)$ as $\tau \to i\infty$: the KZB
connection degenerates to KZ, the $B$-cycle shrinks to a point, and
the full $q^{2H}$ monodromy becomes trivial.

Part~(iv) is a restatement: the full-cycle factor
$e^{2\pi i \hbar} = q^2$ is the exponential of $\hbar$ multiplied by
$2\pi i$, and the $B$-cycle monodromy is the geometric mechanism that
implements this exponentiation.
\end{proof}

\begin{conjecture}[Yangian--quantum group deformation beyond the affine lineage]
\ClaimStatusConjectured
\label{conj:yangian-quantum-group-general}
The Yangian--quantum group deformation of
Theorem~\textup{\ref{thm:yangian-quantum-group}} extends to
non-affine chiral algebras: for any modular Koszul algebra~$\cA$
on~$E_\tau \times \bR$, the genus-$1$ ordered bar complex produces
a quantum-group-type algebra deforming the genus-$0$ Yangian-type
algebra, with the $B$-cycle monodromy implementing the exponential
map $\hbar \mapsto q^2 = e^{2\pi i \hbar}$ and the standard quantum
parameter $q = e^{\pi i \hbar}$.
\end{conjecture}

\subsection{The $\mathfrak{sl}_2$ case and roots of unity}
\label{sec:sl2-roots-unity}

For $\mathfrak{g} = \mathfrak{sl}_2$, the quantum group parameter
specialises to a root of unity at integer level.

\begin{corollary}[$U_q(\mathfrak{sl}_2)$ at roots of unity from affine $\mathfrak{sl}_2$; \ClaimStatusProvedHere]
\label{cor:sl2-root-of-unity}
\index{quantum group!sl2 at roots of unity@$\mathfrak{sl}_2$ at roots of unity|textbf}
For $\mathfrak{g} = \mathfrak{sl}_2$ at positive integer level
$k \in \mathbb{Z}_{>0}$:
\begin{enumerate}[label=\textup{(\roman*)}]
\item The quantum group parameter is
 $q = e^{\pi i/(k+2)}$, and
 $q^{2(k+2)} = 1$.
\item The quantum group $U_q(\mathfrak{sl}_2)$ at this root of
 unity has a finite-dimensional quotient with $k+1$ irreducible
 modules $V_0, V_1, \dots, V_k$ of dimensions $1, 2, \dots, k+1$.
\item These $k+1$ modules are in bijection with the integrable
 highest-weight modules of $\widehat{\mathfrak{sl}}_2$ at level~$k$.
 The truncation of the representation category is the genus-$1$
 avatar of the level-$k$ WZW fusion category.
\end{enumerate}
\end{corollary}

\begin{proof}
Part~(i): $h^\vee = 2$ for $\mathfrak{sl}_2$, so
$\hbar = 1/(k+2)$ and $q = e^{\pi i \hbar} = e^{\pi i/(k+2)}$.
For the root-of-unity property:
$q^{2(k+2)} = e^{2\pi i} = 1$.

Part~(ii) is a classical result: at $q$ a root of unity, the
quantum group $U_q(\mathfrak{sl}_2)$ has a large centre generated
by $E^{k+2}$, $F^{k+2}$, and $K^{2(k+2)}$, and the quotient by
this centre has finitely many simple modules.

Part~(iii) is the Kazhdan--Lusztig equivalence:
$\mathcal{O}_k^{\mathrm{int}}(\widehat{\mathfrak{sl}}_2)
\simeq \mathrm{Rep}_q(\mathfrak{sl}_2)^{\mathrm{fin}}$.
\end{proof}

\subsection{The Jones polynomial from genus-$1$ monodromy}
\label{sec:jones-genus1}

The Jones polynomial emerges as a topological invariant from the
genus-$1$ monodromy applied to the bar complex.

\begin{theorem}[Jones polynomial from genus-$1$ bar-complex monodromy; \ClaimStatusProvedHere]
\label{thm:jones-genus1}
\index{Jones polynomial!from genus-1 monodromy|textbf}
For $\mathfrak{g} = \mathfrak{sl}_2$ at level $k \in \mathbb{Z}_{>0}$,
let $K$ be a framed knot presented as the closure of a braid
$\beta \in B_n$. Then:
\begin{enumerate}[label=\textup{(\roman*)}]
\item The bar complex of $V^k(\mathfrak{sl}_2)$, viewed as a
 logarithmic $\mathrm{SC}^{ch,top}$-algebra, produces the braid
 group representation $\rho_n^{\mathrm{HT}} = \rho_n^{\mathrm{KZ}}$
 on tensor powers of the fundamental representation~$V$.

\item By the Drinfeld--Kohno theorem
 \textup{(}Theorem~\textup{\ref{thm:drinfeld-kohno}}\textup{)},
 this representation factors through
 $U_q(\mathfrak{sl}_2)$ at $q = e^{\pi i/(k+2)}$.

\item The Jones polynomial is recovered by the quantum trace:
 \begin{equation}\label{eq:jones-from-bar}
 J_K(q)
 \;=\;
 \operatorname{tr}_q\bigl(\rho_n^{\mathrm{HT}}(\beta)\bigr)
 \;=\;
 \operatorname{tr}\bigl(q^{2\rho} \cdot \rho_n^{\mathrm{KZ}}(\beta)\bigr),
 \end{equation}
 where $\rho$ is the half-sum of positive roots
 \textup{(}$\rho = 1/2$ for $\mathfrak{sl}_2$\textup{)} and
 $\operatorname{tr}_q$ is the categorical trace in the ribbon
 category $\mathrm{Rep}_q(\mathfrak{sl}_2)$.

\item The integration domain is $\mathrm{FM}_n(\mathbb{C}) \times
 \mathrm{Conf}_n(\mathbb{R})$: the $\mathrm{FM}_n(\mathbb{C})$
 factor carries the bar differential encoding the chiral OPE, and the
 $\mathrm{Conf}_n(\mathbb{R})$ factor carries the $E_1$ coproduct
 implementing the topological interval-cutting. The Jones polynomial
 is a holomorphic-topological invariant: the holomorphic direction
 produces the braiding, and the topological direction produces the
 monoidal structure.
\end{enumerate}
\end{theorem}

\begin{proof}
Parts (i)--(ii) combine the Volume~II one-loop collapse statement
(the bar-complex monodromy equals the KZ monodromy for the affine
lineage) with Theorem~\ref{thm:drinfeld-kohno}(i) (the KZ monodromy
equals the quantum group braid representation).

Part~(iii) is the Reshetikhin--Turaev construction: $\mathrm{Rep}_q(\mathfrak{sl}_2)$
is a ribbon category (the ribbon element is $v = K^{-1}u$ where
$K = q^H$ and $u = \sum S(b_i)a_i$), and the quantum trace
$\operatorname{tr}_q = \operatorname{tr}(q^{2\rho} \,\cdot\,)$
implements the Markov closure converting a braid representation into a
framed link invariant. Since $\rho_n^{\mathrm{HT}} = \rho_n^{\mathrm{KZ}}
\simeq \rho_n^{U_q}$, the quantum trace applied to the bar-complex
monodromy equals the Jones polynomial.

Part~(iv) is a restatement of the $E_1$-chiral factorisation
structure: the bar complex lives on the product of holomorphic and
topological configuration spaces, and the Jones polynomial is the
output of applying the Reshetikhin--Turaev functor to the monodromy of
the connection defined on this product.
\end{proof}

\begin{remark}[The three levels of the construction]
\label{rem:three-levels}
The passage from chiral algebra to knot invariant has three levels,
each proved:
\begin{center}
\renewcommand{\arraystretch}{1.3}
\begin{tabular}{lll}
\textbf{Level} & \textbf{Object} & \textbf{Output} \\
\hline
Genus $0$, ordered bar & KZ connection & Yangian $Y_\hbar(\mathfrak{g})$ \\
Genus $1$, $B$-cycle monodromy & KZB connection & Quantum group $U_q(\mathfrak{g})$ \\
Markov closure & Quantum trace & Jones polynomial $J_K(q)$ \\
\end{tabular}
\end{center}
Each level is proved for the affine lineage by the one-loop collapse
mechanism. Beyond the affine lineage (for $\mathcal{W}$-algebras
obtained by Drinfeld--Sokolov reduction, where higher $A_\infty$
operations contribute), the correspondence between bar-complex
monodromy and quantum group representations remains conjectural.
\end{remark}

\subsection{Scope and frontier}
\label{sec:genus1-frontier}

\begin{remark}[Status]
\label{rem:genus1-status}
The results of this section are proved for the \emph{affine lineage}:
$V^k(\mathfrak{g})$ for simple $\mathfrak{g}$ at generic level~$k$,
acting on finite-dimensional evaluation modules. The proof mechanism
is the one-loop collapse mechanism from Volume~II,
which reduces the full bar-complex superconnection to the KZ/KZB
connection, and the Drinfeld--Kohno theorem (classical at genus~$0$,
Etingof--Schiffmann/Calaque--Enriquez--Etingof at genus~$1$), which
identifies the KZ/KZB monodromy with the quantum group representation.

Beyond the affine lineage, two frontiers are open:
\begin{enumerate}[label=\textup{(\alph*)}]
\item \emph{$\mathcal{W}$-algebras.}
For $\mathcal{W}$-algebras obtained by Drinfeld--Sokolov reduction,
the higher $A_\infty$ operations $m_k$ ($k \geq 3$) do not vanish:
the one-loop collapse mechanism breaks. The full superconnection
contributes to the monodromy, and the genus-$1$ deformation should
produce a \emph{deformed quantum group} or \emph{quantum
$\mathcal{W}$-algebra}. This is the content of
Conjecture~\ref*{conj:rmatrix} in Volume~II.

\item \emph{Higher genus.}
At genus $g \geq 2$, the configuration space
$\mathrm{Conf}_n(\Sigma_g)$ has $2g$ distinguished cycle classes,
and the monodromy acquires $2g$ additional generators. The
curvature $\kappa(A) \cdot \omega_g$ from the Hodge bundle on
$\overline{\mathcal{M}}_{g,n}$ provides genus-by-genus corrections.
The algebraic object controlling the monodromy at genus~$g$ should be
a \emph{higher-genus quantum group} or \emph{mapping-class-group
representation} generalising the braid-group representations at
genus~$0$ and~$1$.
\end{enumerate}
\end{remark}

\subsection{The ordered tridegree decomposition}
\label{subsec:ordered-tridegree}

The symmetric bar complex carries a natural bidegree $(g,n)$
with $g$ the genus and $n$ the degree. The ordered bar complex
refines this to a \emph{tridegree} that records the depth of
consecutive collisions in the Fulton--MacPherson compactification.

\begin{definition}[Ordered tridegree]
\ClaimStatusProvedHere
\label{def:ordered-tridegree}
\index{tridegree!ordered bar complex|textbf}
\index{collision depth!ordered|textbf}
Let $A$ be a strongly admissible associative chiral algebra and
$\Barchord(A)$ its ordered bar complex.
The \emph{ordered tridegree} of a bar element
$\alpha \in \Barchord(A)$ is the triple
$(g,n,d)^{E_1}$ where:
\begin{enumerate}[label=\textup{(\roman*)}]
\item $g \geq 0$ is the genus (the number of handles of the
 underlying curve);
\item $n \geq 1$ is the degree (the number of bar slots);
\item $d \geq 0$ is the \emph{ordered collision depth}: the
 codimension of the planted-forest boundary stratum of
 $\overline{\FM}_n^{\mathrm{ord}}(\mathbb{C})$ at which
 $\alpha$ is supported.
\end{enumerate}
At degree~$n$, the ordered compactification
$\overline{\FM}_n^{\mathrm{ord}}(\mathbb{C})$ has exactly $n{-}1$
codimension-one boundary divisors $D_{i,i+1}$, $1 \leq i \leq n{-}1$,
corresponding to consecutive collisions $z_i \to z_{i+1}$. The
depth~$d$ counts the codimension in this boundary stratification:
$d = 0$ is the interior (no collision), and the maximal depth at
degree~$n$ is $d = n{-}1$ (complete consecutive collapse). In the
unordered FM compactification, degree~$n$ has $\binom{n}{2}$
codimension-one divisors; the restriction to $n{-}1$ consecutive
divisors is the geometric signature of the $E_1$-structure.
\end{definition}

\begin{theorem}[Ordered depth spectrum]
\ClaimStatusProvedHere
\label{thm:ordered-depth-spectrum}
\index{depth spectrum!ordered|textbf}
\index{collision depth!and OPE pole order}
Let $A$ be a strongly admissible associative chiral algebra
whose OPE has pole orders $\{p_1, \ldots, p_s\}$
\textup{(}i.e., $a_{(n)}b \neq 0$ precisely for $n \in \{p_1{-}1, \ldots, p_s{-}1\}$,
with the convention that $a_{(n)}b$ contributes a pole of order $n{+}1$\textup{)}.
Then the depth spectrum at genus~$0$, degree~$2$ is
\begin{equation}\label{eq:depth-spectrum}
\mathrm{Spec}\bigl(d^{E_1}; A\bigr)\big|_{(0,2)}
\;=\;
\bigl\{\,\max(p-1,\,0) : p \text{ is a pole order of the OPE of } A\,\bigr\}.
\end{equation}
The following table records the depth spectrum for the standard
examples:
\begin{center}
\renewcommand{\arraystretch}{1.2}
\begin{tabular}{lll}
\textbf{Algebra} & \textbf{OPE poles} & \textbf{Depth spectrum} \\
\hline
Heisenberg $H_k$ & $\{2\}$ & $\{0,\,1\}$ \\
Affine $V^k(\mathfrak{g})$ & $\{1,\,2\}$ & $\{0,\,1\}$ \\
Virasoro $\mathrm{Vir}_c$ & $\{1,\,2,\,4\}$ & $\{0,\,1,\,3\}$
\end{tabular}
\end{center}
The depth spectrum of the Virasoro algebra exhibits a gap at $d = 2$,
reflecting the absence of a cubic pole in the Virasoro OPE
$T(z)T(w) \sim (c/2)(z{-}w)^{-4} + 2T(w)(z{-}w)^{-2} + \partial T(w)(z{-}w)^{-1}$.
\end{theorem}

\begin{proof}
At degree~$2$, the ordered bar differential acts by extracting the
collision residue at the unique divisor $D_{1,2} = \{z_1 = z_2\}$ via
the ordered logarithmic propagator $\eta = d\log(z_1 - z_2)$.
A pole of order~$p$ in the OPE contributes a term
$c^K_{IJ;p-1}\,(z_1 - z_2)^{-p}$ to the structure constants.
The $d\log$ kernel absorbs one power of the pole
(the collision residue $r(z)$ has pole orders one less than the OPE),
producing a residue of order $p - 1$ in the collision variable.
This residue is supported on the codimension-$(p{-}1)$ boundary
stratum of $\overline{\FM}_2^{\mathrm{ord}}(\mathbb{C})$,
hence contributes at depth $d = p - 1$. The zeroth product
$a_{(0)}b$ (pole of order~$1$) contributes at depth $d = 0$,
living in the interior.

For the Virasoro algebra: the quartic pole gives depth $d = 3$,
the double pole gives depth $d = 1$, and the simple pole gives depth
$d = 0$. The cubic pole is absent, so $d = 2$ does not appear.
\end{proof}

\begin{theorem}[Ordered AOS reduction]
\ClaimStatusProvedHere
\label{thm:ordered-AOS}
\index{AOS relations!ordered reduction|textbf}
\index{Stasheff identity!from Stokes theorem}
\index{associahedron!from ordered FM boundary}
The ordered AOS relations (the identities imposed on the bar
differential $d_{\Barchord}$ by the condition $d^2 = 0$
on $\overline{\FM}_k^{\mathrm{ord}}(\mathbb{C})$) are the Stasheff
associahedron identities. Concretely: at degree~$k$, the boundary of
$\overline{\FM}_k^{\mathrm{ord}}(\mathbb{C})$ has exactly $k{-}1$
codimension-one faces $D_{i,i+1}$ ($1 \leq i \leq k{-}1$),
each corresponding to the consecutive collision $z_i \to z_{i+1}$.
The condition
\begin{equation}\label{eq:ordered-d-squared}
d^2 = 0
\quad\Longleftrightarrow\quad
\sum_{\substack{1 \leq i < j \leq k-1 \\
j = i+1}}
d_{i,i+1} \circ d_{j,j+1}
\;+\;
\sum_{i=1}^{k-1}
d_{i,i+1} \circ d_{\mathrm{int}}
\;=\;0
\end{equation}
is associativity: the $A_\infty$ Stasheff relations
$\sum_{r+s+t=n} m_{r+1+t} \circ (\id^{\otimes r} \otimes m_s
\otimes \id^{\otimes t}) = 0$.

In the unordered FM compactification, $d^2 = 0$ encodes
the Jacobi identity (the Arnold relation among
$\binom{k}{2}$ divisors contributes the antisymmetry). In the
ordered compactification, the Arnold relation is geometrically
redundant: there are no non-consecutive divisors $D_{i,j}$ with
$j > i + 1$, so no three-term antisymmetry relation can form.
\end{theorem}

\begin{proof}
By Stokes' theorem on
$\overline{\FM}_k^{\mathrm{ord}}(\mathbb{C})$, the condition
$d^2 = 0$ is equivalent to the vanishing of the boundary integral
$\int_{\partial\overline{\FM}_k^{\mathrm{ord}}} \eta = 0$
for the bar propagator form~$\eta$. The boundary decomposes as
\[
\partial\overline{\FM}_k^{\mathrm{ord}}(\mathbb{C})
\;=\;
\bigcup_{i=1}^{k-1} D_{i,i+1},
\]
a union of $k{-}1$ consecutive-collision divisors. Each divisor
$D_{i,i+1}$ is isomorphic to
$\overline{\FM}_{k-1}^{\mathrm{ord}}(\mathbb{C})$
(the two colliding points are replaced by their merged image).
The boundary integral along $D_{i,i+1}$ computes the composition
$d_{i,i+1} \circ d$, and the total relation
$\sum_{i=1}^{k-1} d_{i,i+1} \circ d = 0$
is the $k$-th Stasheff identity.

In the unordered setting, $\overline{\FM}_k(\mathbb{C})$ has
$\binom{k}{2}$ codimension-one divisors $D_{i,j}$ for all
$i < j$, and Stokes yields $\binom{k}{2}$ terms:
the Jacobi/Arnold relations. The ordered restriction to
$k{-}1$ consecutive divisors eliminates all non-adjacent
collision terms. The Arnold relation
$\omega_{ij} \wedge \omega_{jk} + \omega_{jk} \wedge \omega_{ki}
+ \omega_{ki} \wedge \omega_{ij} = 0$
cannot be stated when $\omega_{ik}$ does not exist as a
generator (the non-consecutive divisor $D_{i,k}$ is absent),
so it is geometrically vacuous in the ordered compactification.
\end{proof}

\begin{proposition}[Averaging and surplus]
\ClaimStatusProvedHere
\label{prop:averaging-surplus}
\index{averaging map!depth surplus|textbf}
\index{R-matrix!monodromy as averaging kernel}
The averaging map
$\mathrm{av}\colon\Barchord(A) \to \Barch(A)$
intertwines the ordered tridegree with the symmetric bidegree via
\[
\mathrm{av}\colon
\Barchord(A)_{(g,n,d)}
\;\longrightarrow\;
\Barch(A)_{(g,n,d')}
\quad\text{with}\quad
d' \geq d.
\]
The depth can only increase under averaging: symmetrisation may
identify an ordered element supported at depth~$d$ with symmetric
elements supported at deeper strata. At genus~$0$ and fixed
degree~$n$, the kernel
$\ker\bigl(\mathrm{av}|_{(0,n)}\bigr)$ is generated by the $R$-matrix
monodromy data of
Proposition~\textup{\ref{sec:r-matrix-descent-vol1}}.
For $E_\infty$-chiral algebras without OPE poles (the pole-free
subclass), the $R$-matrix is the Koszul sign $R(z) = \tau$,
the averaging map is the naive $\Sigma_n$-quotient, and
$\ker(\mathrm{av}) = 0$.
\end{proposition}

\begin{proof}
The averaging map is the composition
$\Barchord(A)_n
\xrightarrow{\Sigma_n\text{-orbit}} \Barch(A)_n$,
twisted by the $R$-matrix as in
Proposition~\ref{sec:r-matrix-descent-vol1}:
the $\Sigma_n$-action on the ordered bar uses the
$R$-matrix monodromy $R_{i,i+1}(z_i - z_{i+1})$ at each
adjacent transposition. An element $\alpha$ of depth~$d$
is supported on a stratum
$\partial_F\overline{\FM}_n^{\mathrm{ord}}(\mathbb{C})$
of codimension~$d$. Under symmetrisation, the orbit
$\Sigma_n \cdot \alpha$ may map to strata of the unordered
compactification at codimension $d' \geq d$, since the
$\binom{n}{2}$ divisors of the unordered space refine the
$n{-}1$ ordered divisors.

For the kernel: $\mathrm{av}(\alpha) = 0$ if and only if
$\sum_{\sigma \in \Sigma_n} (R\text{-}\sigma) \cdot \alpha = 0$,
which is precisely the condition that $\alpha$ lies in the
$R$-twisted augmentation ideal of the group algebra
$\mathbb{C}[\Sigma_n]$. When $R(z) = \tau$
(pole-free $E_\infty$ case), this reduces to
$\sum_\sigma \mathrm{sgn}(\sigma) \cdot \sigma(\alpha) = 0$,
which vanishes only for the zero element in the symmetric
quotient, giving $\ker(\mathrm{av}) = 0$.
\end{proof}

\begin{remark}[Physical meaning: the soft graviton hierarchy]
\label{rem:tridegree-physical}
\index{soft graviton!and ordered tridegree}
\index{obstruction cascade!depth-3 origin}
The ordered tridegree $(g,n,d)^{E_1}$ organises the soft graviton
hierarchy of the gravity part of Volume~II\@. The correspondence is:
soft order $p$ of a graviton theorem corresponds to degree
$r = p + 2$ in the bar complex. The depth~$d$ records
\emph{which OPE pole is responsible}: a soft theorem at order~$p$
receives contributions from all poles whose depth
$d = \max(p'-1,0)$ satisfies $d \leq r - 1$.

For the Virasoro algebra, the depth-$3$ component
(from the quartic pole, i.e.\ the central charge~$c$) is the
generator of the obstruction cascade. The
Stasheff identity at degree~$4$ propagates the depth-$3$
contribution to $m_4$; at degree~$5$ it forces $m_5$;
and so on in perpetuity. The gap at depth~$2$ in the Virasoro
spectrum (Theorem~\ref{thm:ordered-depth-spectrum}) is the
geometric expression of the absence of a cubic pole,
equivalently the vanishing of the $L_n$-eigenvalue at the wrong
parity: $[L_m, L_n] = (m-n)L_{m+n} + \frac{c}{12}m(m^2-1)\delta_{m+n,0}$
jumps from the linear term ($d = 1$) to the cubic term ($d = 3$)
with no quadratic intermediate.
\end{remark}

\begin{remark}[Motivic interpretation of the tridegree]
\label{rem:tridegree-motivic}
\index{motivic interpretation!tridegree}
\index{mixed Tate motives!and ordered tridegree}
\index{multiple zeta values!and collision depth}
The ordered tridegree refines the motivic content identified in
Remark~\ref{rem:motivic-interpretation}. The averaging map
$\mathrm{av}\colon\Barchord(A) \to \Barch(A)$ is a
motivic projection in the following precise sense: the ordered bar
complex at depth~$d$ computes periods of the mixed Hodge structure
on $\overline{\FM}_n^{\mathrm{ord}}(\mathbb{C})$ that are supported
on strata of codimension~$d$.

Concretely, the three cases are:
\begin{enumerate}[label=\textup{(\roman*)}]
\item \emph{Depth $d = 0$} (interior): the periods are polynomial in
$2\pi i$ and the collision residues. At degree~$2$, these are
the rational multiples of the OPE structure constants.

\item \emph{Depth $d \geq 1$} (boundary strata): the periods are
iterated integrals of the KZ connection along paths in
$\mathrm{Conf}_n^{\mathrm{ord}}(\mathbb{C})$. At degree~$3$ and
depth~$1$, the Drinfeld associator $\Phi_{\mathrm{KZ}}$
contributes multiple zeta values $\zeta(2), \zeta(3), \ldots$,
periods of mixed Tate motives over~$\mathbb{Z}$.

\item \emph{The symmetric image}: the averaging map $\mathrm{av}$
kills all motivic content. The symmetric bar invariants
(curvature~$\kappa$, complementarity class~$C$, shadow
tower~$Q$) are rational numbers, periods of the trivial
motive~$\mathbb{Q}(0)$.
\end{enumerate}

\noindent
The ordered surplus $\ker(\mathrm{av})$ is therefore the entire
mixed Tate content of the configuration-space integrals. The
tridegree decomposition stratifies this content: the depth~$d$
component carries mixed Tate motives of weight~$\leq d$,
and the motivic weight filtration on $\ker(\mathrm{av})$ is
induced by the depth filtration on the ordered bar complex.
This is the precise sense in which the ordered bar complex is
arithmetically richer than the symmetric bar complex: it remembers
the multiple zeta values that the symmetric projection forgets.
\end{remark}

\begin{remark}[Drinfeld centre and the quantum group category]
\label{rem:drinfeld-centre-qg}
The full genus-$1$ ordered chiral homology should
ultimately be understood by working at the categorical rather than algebraic
level. Concretely, the Drinfeld centre
$\cZ(Y_\hbar(\mathfrak{sl}_2)^{\mathrm{ch}}\textrm{-mod})$
should recover the quantum group category
$\mathrm{Rep}_q(\mathfrak{sl}_2)$ at
$q = e^{2\pi i/(k+2)}$. This identification is the
subject of the DK-5 programme; see
\cite[\S6]{Lorgat26I} for the current status in type~$A$.
\end{remark}

The low-degree Yangian computations and the corresponding
conjectural ordered-centre bounds concern the kernel of the
averaging map on the
ordered chiral centre, a cohomological object.
The following proposition establishes the underlying
representation-theoretic input: the dimension of
$\ker(\av_n)$ on the bare tensor power $V^{\otimes n}$,
before passing to bar-complex cohomology.

\begin{proposition}[Kernel of the Reynolds projector:
general simple Lie algebras]
\label{prop:ker-av-schur-weyl}
\ClaimStatusProvedHere
Let $\fg$ be any simple Lie algebra over~$\CC$, let
$V$ be any finite-dimensional representation of~$\fg$ with
$d = \dim V$, and let
$\av_n \colon V^{\otimes n} \to (V^{\otimes n})_{\Sigma_n}$
denote the Reynolds averaging map (projection onto the
$\Sigma_n$-coinvariant quotient).
Then
\begin{equation}\label{eq:ker-av-schur-weyl}
  \dim \ker(\av_n\big|_{V^{\otimes n}})
  \;=\;
  d^n - \binom{n+d-1}{d-1}.
\end{equation}
The formula depends only on $d = \dim V$, not on $\fg$ or
on the $\fg$-module structure of~$V$.

\smallskip
\noindent\textit{Special cases.}
For $\fg = \mathfrak{sl}_2$ with the fundamental
representation $V = \CC^2$:
$\dim\ker(\av_n) = 2^n - (n+1)$.
For any $\fg$ with $d$-dimensional representation~$V$,
at $n = 2$:
$\dim\ker(\av_2) = d^2 - d(d{+}1)/2 = d(d{-}1)/2
= \dim\bigwedge^2 V$.
\end{proposition}

\begin{proof}
The Reynolds projector
$\av_n = \frac{1}{n!}\sum_{\sigma \in \Sigma_n} \sigma$
acts on $V^{\otimes n}$ by permuting tensor factors.
Its image is the $\Sigma_n$-invariant subspace
$\Sym^n(V)$, regardless of any Lie algebra structure
on~$V$. The dimension
$\dim\Sym^n(V) = \binom{n+d-1}{d-1}$ (the number of
degree-$n$ monomials in $d$ variables) depends only on
$d = \dim V$, as the stars-and-bars argument uses no
algebraic structure beyond the vector space.
Since $\dim V^{\otimes n} = d^n$, the kernel has dimension
$d^n - \binom{n+d-1}{d-1}$.

When $\fg = \mathfrak{gl}_r$ with $V = \CC^r$,
Schur--Weyl duality refines this to an isotypic
decomposition. The tensor power decomposes into
$(\Sigma_n, \mathfrak{gl}_r)$-bimodules:
\begin{equation}\label{eq:schur-weyl-decomp}
  V^{\otimes n}
  \;\cong\;
  \bigoplus_{\lambda \,\vdash\, n,\;
    \ell(\lambda) \leq r}
  S_\lambda \otimes V_\lambda(\mathfrak{gl}_r),
\end{equation}
where $\lambda$ runs over partitions of $n$ with at most
$r$ rows, $S_\lambda$ is the corresponding Specht module
of $\Sigma_n$, and $V_\lambda(\mathfrak{gl}_r)$ is the
irreducible $\mathfrak{gl}_r$-representation of highest
weight~$\lambda$.
The projector $\av_n$ selects the trivial isotypic
component $S_{(n)} = \CC$ and annihilates every
$S_\lambda$ with $\lambda \neq (n)$. The image is
\[
  \operatorname{im}(\av_n)
  \;\cong\;
  S_{(n)} \otimes V_{(n)}(\mathfrak{gl}_r)
  \;=\;
  \Sym^n(\CC^r),
\]
and the kernel decomposes by isotypic components:
\[
  \ker(\av_n)
  \;\cong\;
  \bigoplus_{\substack{\lambda \,\vdash\, n \\
    \ell(\lambda) \leq r,\;
    \lambda \neq (n)}}
  S_\lambda \otimes V_\lambda(\mathfrak{gl}_r).
\]
For a general simple~$\fg$ acting on~$V$, the Schur--Weyl
decomposition involves the commutant algebra
$\End_\fg(V^{\otimes n})$ rather than the group algebra of
$\GL(V)$, and the individual summands carry different
$\fg$-representations with different dimensions.
The total kernel dimension, however, depends only on
$\dim\Sym^n(V) = \binom{n+d-1}{d-1}$, which is
determined by $d$ alone.
\end{proof}

\begin{remark}[Explicit values and conjectural ordered-centre bounds]
\label{rem:ker-av-table}
Table~\ref{tab:ker-av-dims} records the kernel dimensions
at low degrees for the fundamental representations of
$\mathfrak{sl}_2$ through $\mathfrak{sl}_5$ (of dimensions
$d = 2, 3, 4, 5$), together with the fundamental
representations of $\mathfrak{so}_5$ ($d = 5$) and
$G_2$ ($d = 7$). Since the formula depends only on
$d = \dim V$, the rows for $\mathfrak{sl}_5$ and
$\mathfrak{so}_5$ are identical.
These are representation-theoretic upper bounds on the
dimensions of the corresponding components of the ordered
chiral centre in the conjectural ordered-centre kernel:
$\ker(\av_n)$ on the ordered chiral centre is a subquotient
of $\ker(\av_n\big|_{V^{\otimes n}})$, obtained by passing
to bar-complex cohomology
(only the cocycles in the non-trivial $\Sigma_n$-isotypics
survive). At degree $2$ for $\mathfrak{sl}_2$: the full
$1$-dimensional kernel $\bigwedge^2 V = \CC$ survives in
the classical Yangian computation.
At degree $n \geq 3$, the kernel on the chiral centre is
in general strictly smaller than the
representation-theoretic bound $d^n - \binom{n+d-1}{d-1}$.

\begin{table}[h]
\centering
\renewcommand{\arraystretch}{1.3}
\begin{tabular}{@{} l c c c c c @{}}
\toprule
  & $n = 2$ & $n = 3$ & $n = 4$ & $n = 5$ & $n = 6$ \\
\midrule
$\mathfrak{sl}_2$ $(d{=}2)$
  & $1$ & $4$ & $11$ & $26$ & $57$ \\
$\mathfrak{sl}_3$ $(d{=}3)$
  & $3$ & $17$ & $66$ & $222$ & $701$ \\
$\mathfrak{sl}_4$ $(d{=}4)$
  & $6$ & $44$ & $221$ & $968$ & $4{,}012$ \\
$\mathfrak{sl}_5$, $\mathfrak{so}_5$ $(d{=}5)$
  & $10$ & $90$ & $555$ & $2{,}999$ & $15{,}415$ \\
$G_2$ \textit{fund.}\ $(d{=}7)$
  & $21$ & $259$ & $2{,}191$ & $16{,}345$ & $116{,}725$ \\
\midrule
\textit{Formula}
  & \multicolumn{5}{c}{$d^n - \binom{n+d-1}{d-1}$} \\
\bottomrule
\end{tabular}
\caption{Dimension of $\ker(\av_n\big|_{V^{\otimes n}})$
for a $d$-dimensional representation $V$ at degrees
$n = 2, \ldots, 6$
\textup{(Proposition~\ref{prop:ker-av-schur-weyl})}.
The formula depends only on $d = \dim V$.}
\label{tab:ker-av-dims}
\end{table}
\end{remark}


% ================================================================
\section{Background: chiral algebras and the $\Eone/\Einf$
hierarchy}
\label{sec:background}
\begin{remark}[The averaging kernel at genus~$1$: quasi-modular anomaly and Hodge curvature]
\label{rem:genus1-ker-av}
\index{averaging map!kernel at genus 1}
\index{quasi-modular forms!and averaging kernel}
\index{Hodge bundle!curvature and ker(av)}
\ClaimStatusConjectured
Proposition~\ref{prop:averaging-surplus} characterises $\ker(\mathrm{av})$
at genus~$0$: it is generated by the $R$-matrix monodromy, carrying
the full mixed Tate content of the configuration-space integrals.
At genus~$1$, a new mechanism enters.

The genus-$1$ bar propagator on $E_\tau$ involves the Weierstrass
$\wp$-function and its derivatives (weight~$1$ after $d\log$ absorption;
cf.\ Theorem~\ref{thm:elliptic-spectral-dichotomy}). The Laurent
expansion of $\wp$ near $z = 0$ generates a fan of Eisenstein
corrections $G_{2k}(\tau)$, of which $G_2$ is distinguished: it is
quasi-modular, not modular. Under $\gamma = \bigl(\begin{smallmatrix}
a & b \\ c & d \end{smallmatrix}\bigr) \in \mathrm{SL}(2,\mathbb{Z})$,
\[
 G_2(\gamma\cdot\tau)
 \;=\;
 (c\tau + d)^2 \, G_2(\tau)
 \;-\;
 2\pi i\, c(c\tau + d).
\]
The anomalous term $-2\pi i\, c(c\tau + d)$ is the shadow of the
Legendre relation $\eta_1 \tau - \eta_\tau = \pi i$ and is
the source of the bar curvature
$d^2 = \kappa(\mathcal{A})\cdot\omega_1$ on
$\overline{\mathcal{M}}_{1,n}$
(Theorem~D, Remark~\ref{rem:eisenstein-fan}).

The passage to the non-holomorphic completion
$\widehat{G}_2(\tau) = G_2(\tau) - \pi / \operatorname{Im}\tau$
restores modularity: $\widehat{G}_2$ transforms as a true weight-$2$
modular form. The difference $G_2 - \widehat{G}_2 = \pi/\operatorname{Im}\tau$
is the anomaly that the symmetric bar complex absorbs into
$\kappa(\mathcal{A})\cdot\omega_1$. In the ordered bar complex,
this anomaly is visible as explicit quasi-modular coefficients in
the elliptic $r$-matrix; the averaging map $\mathrm{av}$ projects
these coefficients to their modular completions, absorbing the
anomaly into the Hodge curvature of the chiral differential.

This suggests the following structural picture for $\ker(\mathrm{av})$
at genus~$1$.

\begin{enumerate}[label=\textup{(\roman*)}]
\item \emph{Rational component} (inherited from genus~$0$):
 the $R$-matrix monodromy on $\mathrm{Conf}_n^{\mathrm{ord}}(E_\tau)$
 contributes elements to $\ker(\mathrm{av})$ by the same mechanism as
 Proposition~\ref{prop:averaging-surplus}. These carry the
 elliptic multiple zeta values of Enriquez.

\item \emph{Quasi-modular component} (new at genus~$1$):
 the $G_2$-anomaly terms in the ordered elliptic $r$-matrix are
 annihilated by $\mathrm{av}$, since the symmetric bar complex sees
 only the completed $\widehat{G}_2$. The resulting kernel elements
 are weighted by $\kappa(\mathcal{A})$: when $\kappa(\mathcal{A}) = 0$
 (uncurved, e.g.\ affine Kac--Moody at critical level $k = -h^\vee$),
 the curvature $d^2 = 0$ and the quasi-modular anomaly contributes
 nothing to the bar differential, so the kernel reduces to the
 genus-$0$ rational component. When
 $\kappa(\mathcal{A}) \neq 0$, the anomaly is nontrivial and scales
 linearly with $\kappa(\mathcal{A})$ .

\item \emph{Interaction}: the two components are not independent.
 The genus-$1$ free energy
 $F_1 = \kappa(\mathcal{A})/24$ measures the leading Hodge class $\lambda_1$ on
 $\overline{\mathcal{M}}_{1,1}$. The $1/24$ is the
 residue of the $\eta$-function ($\eta(\tau) = q^{1/24}\prod_{n \geq 1}(1 - q^n)$),
 which is itself the square root of $1/G_2$ at leading order.
 The rational and quasi-modular components of $\ker(\mathrm{av})$ are
 entangled through this $\eta$-function bridge: the ordered bar complex
 sees the full $q$-expansion of the propagator $E_2^*(\tau)/12$
 (quasi-modular, weight~$2$), while the symmetric bar complex retains
 only the scalar $\kappa(\mathcal{A})/24$ after integration over
 $\overline{\mathcal{M}}_{1,1}$.
\end{enumerate}

\noindent
The genus-$0$ kernel $\ker(\mathrm{av}|_{(0,n)})$ is a representation of
the Grothendieck--Teichm\"uller group $\widehat{GT}$
(Remark~\ref{rem:motivic-interpretation}). At genus~$1$, the
conjectural extension is that $\ker(\mathrm{av}|_{(1,n)})$ carries an
action of the \emph{elliptic} Grothendieck--Teichm\"uller group
$\widehat{GT}^{\mathrm{ell}}$ of Enriquez, with the quasi-modular
anomaly defining the cocycle by which $\widehat{GT}^{\mathrm{ell}}$
extends $\widehat{GT}$. The curvature $\kappa(\mathcal{A})\cdot\omega_1$
enters as the scalar by which this cocycle acts on the chiral differential:
$\widehat{GT}^{\mathrm{ell}}$ permutes the quasi-modular completions
of the elliptic associators, and the failure of the completion to
commute with the $\widehat{GT}$-action is measured by
$\kappa(\mathcal{A})$.
\end{remark}


\section{Curvature--braiding entanglement at genus one}
\label{sec:curvature-braiding-entanglement}

The dichotomy of \S\ref{sec:annular-bar} partitions chiral algebras
according to whether curvature and braiding decouple at genus~$1$.
The sharp form of this dichotomy isolates the mechanism
that forces entanglement and identifies the central open question
to which the entire annular bar theory points.

\begin{theorem}[Elliptic spectral dichotomy, genus-$1$ specialisation;
\ClaimStatusProvedHere]
\label{thm:elliptic-spectral-dichotomy}
\index{elliptic spectral dichotomy|textbf}
\index{genus 1!elliptic spectral dichotomy}
\index{r-matrix@$r$-matrix!elliptic}
Let $\cA$ be a logarithmic $\Eone$-chiral algebra with
$(-1)$-shifted $\lambda$-bracket
$\{a\,{}_\lambda\,b\}=\sum_{n\ge 0}c_n(a,b)\,\lambda^n$
on $H^\bullet(\cA,Q)$.
At genus~$1$, the ordered $r$-matrix on the elliptic curve
$E_\tau=\mathbb{C}/(\mathbb{Z}+\tau\mathbb{Z})$ takes the form
\begin{equation}\label{eq:elliptic-r-dichotomy}
r_1(z;\tau)
\;=\;
c_0\cdot\zeta(z\,|\,\tau)
\;+\;
\sum_{n\ge 1}
 \frac{(-1)^{n-1}}{n!}\,
 c_n\cdot\wp^{(n-1)}(z\,|\,\tau),
\end{equation}
with $\wp^{(0)}:=\wp$, where $\zeta$ and $\wp$ are the
Weierstrass functions with periods $(1,\tau)$. The first two
explicit terms are $c_0\,\zeta+c_1\,\wp$; the next is
$-\tfrac{1}{2}c_2\,\wp'$. The coefficient $c_0=a_{(0)}b$ is
the Coisson bracket \textup{(}zeroth product\textup{)} and
$c_1=a_{(1)}b$ is the conformal pairing. The convention is
the raw $\lambda$-bracket convention of
Remark~\ref{rem:lambda-bracket-convention-dichotomy}
\textup{(}shifted by one from the $d\log$-absorbed
collision-residue convention \textup{)}.

The $B$-cycle monodromy is
\begin{equation}\label{eq:b-cycle-monodromy-sharp}
r_1(z+\tau;\tau) - r_1(z;\tau) \;=\; 2\eta_\tau\cdot c_0,
\end{equation}
with $\eta_\tau=\zeta(\tfrac{\tau}{2}\,|\,\tau)$ the
quasi-period. The $A$-cycle monodromy is $2\eta_1\cdot c_0$,
where $\eta_1=\zeta(\tfrac{1}{2}\,|\,\tau)$. The Legendre
relation $\eta_1\tau-\eta_\tau=\pi i$ constrains the
two monodromies.
\end{theorem}

\begin{proof}
By Theorem~\ref{thm:curvature-braiding-dichotomy}, the
genus-$1$ $r$-matrix is obtained from the genus-$0$ Laurent
expansion $r(z)=\sum c_n/z^{n+1}$ by replacing each rational
function $1/z^{n+1}$ with its unique elliptic counterpart on
$E_\tau$: $1/z\mapsto\zeta(z\,|\,\tau)$ for $n=0$, and
$1/z^{n+1}\mapsto\frac{(-1)^{n-1}}{n!}\,\wp^{(n-1)}(z\,|\,\tau)$
for $n\ge 1$, with $\wp^{(0)}:=\wp$. These are uniquely
determined by matching the local singularity $1/z^{n+1}+O(z)$
near $z=0$ and requiring double periodicity for $n\ge 1$
(or quasi-periodicity for $n=0$).
The monodromy~\eqref{eq:b-cycle-monodromy-sharp} follows from
$\zeta(z+\tau\,|\,\tau)=\zeta(z\,|\,\tau)+2\eta_\tau$ and the
double periodicity of $\wp$ and all its derivatives.
\end{proof}

\begin{definition}[Decoupled vs.\ entangled genus-$1$ structure]
\label{def:decoupled-entangled}
\index{curvature--braiding!decoupled}
\index{curvature--braiding!entangled}
Let $\cA$ be a logarithmic $\Eone$-chiral algebra with
$r_1(z;\tau)$ as in
Theorem~\ref{thm:elliptic-spectral-dichotomy}.
\begin{enumerate}[label=\textup{(\roman*)}]
\item $\cA$ is \emph{decoupled at genus~$1$} if $c_0=0$:
 the Coisson bracket vanishes, $r_1$ is doubly periodic on
 $E_\tau$, and the curvature $\kappa(\cA)\cdot\omega_1$
 and the spectral braiding are independent structures
 belonging to different components of the $E_1$-chiral coalgebra.

\item $\cA$ is \emph{entangled at genus~$1$} if $c_0\ne 0$:
 $r_1$ is quasi-periodic with $B$-cycle monodromy
 $2\eta_\tau\cdot c_0$, and the braiding and the curvature
 share a common geometric source, the quasi-periodicity
 of the Weierstrass zeta function, which is the same
 mechanism that produces $\kappa(\cA)\cdot\omega_1$ from the
 Hodge bundle on $\overline{\cM}_{1,n}$.
\end{enumerate}
\end{definition}

\begin{center}
\begin{tabular}{l@{\quad}c@{\quad}l@{\quad}l}
\toprule
\textbf{Algebra} & $c_0$ & \textbf{Status} & \textbf{Mechanism} \\
\midrule
Heisenberg $\mathcal{H}_k$ & $0$ & Decoupled &
 No simple pole; $r_1=k\,\wp(z\,|\,\tau)$ \\
Free fermion $\psi$ & $0$ & Decoupled &
 $\{-\,{}_{(0)}-\}=0$; $r_1=\wp(z\,|\,\tau)$ \\
Lattice $V_\Lambda$ (abelian) & $0$ & Decoupled &
 Abelian Lie bracket \\[4pt]
Affine $\widehat{\mathfrak{g}}_k$ & $\ne 0$ & Entangled &
 $c_0=f^{ab}{}_c\,J^c$; KZB monodromy \\
Virasoro $\mathrm{Vir}_c$ & $\ne 0$ & Entangled &
 $c_0=2T$; $r_1=\tfrac{c}{2}\wp'(z)+2T\,\zeta(z)$ \\
$\mathcal{W}_N$ ($N\ge 2$) & $\ne 0$ & Entangled &
 DS reduction inherits $c_0\ne 0$ from $\widehat{\mathfrak{sl}}_N$ \\
\bottomrule
\end{tabular}
\end{center}

\begin{remark}[The Eisenstein fan]
\label{rem:eisenstein-fan}
\index{Eisenstein series!fan structure}
\index{quasi-modular forms!and bar curvature}
The Laurent expansion of $\wp$ and its derivatives
near $z=0$ on $E_\tau$ takes the form
\[
\wp^{(n)}(z\,|\,\tau)
\;=\;
\frac{(-1)^n\,(n+1)!}{z^{n+2}}
\;+\;
\sum_{k\ge 1}
 \frac{(2k+n+1)!}{(2k)!\,(n+1)!}\,
 G_{2k+n+2}(\tau)\,z^{2k},
\]
where $G_{2k}(\tau)=\sum'_{(m,n)\in\mathbb{Z}^2}(m\tau+n)^{-2k}$
are the Eisenstein series. Each OPE pole of order $2k$
in the $\lambda$-bracket thus generates a $G_{2k}(\tau)$
correction to the elliptic $r$-matrix, producing a fan of
modular corrections indexed by the pole spectrum.

The $G_2$ correction occupies a distinguished position:
it is \emph{quasi-modular}, transforming under
$\mathrm{SL}(2,\mathbb{Z})$ with the anomalous shift
\[
G_2(-1/\tau) \;=\; \tau^2\,G_2(\tau)-2\pi i\,\tau.
\]
This anomaly is the shadow of the Legendre relation
$\eta_1\tau-\eta_\tau=\pi i$, and it manifests in the
bar complex as the curvature $d^2=\kappa(\cA)\cdot\omega_1$
of the curved $A_\infty$-chiral structure at genus~$1$. More
precisely: the non-holomorphic completion
$\widehat{G}_2(\tau)=G_2(\tau)-\pi/\Imag\,\tau$
is modular, and the passage from $G_2$ to $\widehat{G}_2$
is the modular gauge transformation that renders the chiral
differential well-defined on $\overline{\cM}_{1,n}$. Whether
the analogous completion can be performed simultaneously
on the coassociative coproduct (disentangling the braiding
from the curvature when $c_0\ne 0$) is the subject
of the following question.
\end{remark}

\begin{question}[The central genus-$1$ frontier]
\label{q:genus1-entanglement}
\index{genus 1!central open question}
\index{curvature--braiding!disentanglement problem}
\ClaimStatusOpen
Does the genus-$1$ coupling between the bar curvature
$d^2=\kappa(\cA)\cdot\omega_1$ and the spectral braiding
monodromy $2\eta_\tau\cdot c_0$ represent genuine entanglement
of the chiral and coassociative structures, or can a modular gauge twist
$r_1(z;\tau)\mapsto r_1(z;\tau)+\phi(z;\tau)$ disentangle them?

More precisely: for $\cA$ with $c_0\ne 0$, does there exist
a doubly-periodic correction $\phi(z;\tau)$, meromorphic in~$z$
and modular in~$\tau$, such that
\[
\widetilde{r}_1(z;\tau)
\;:=\;
r_1(z;\tau)+\phi(z;\tau)
\]
is doubly periodic in~$z$ and the annular bar differential
$d^{\mathrm{ann}}$ computed from $\widetilde{r}_1$ remains
a differential \textup{(}$d^{\mathrm{ann}}\circ d^{\mathrm{ann}}
=\kappa\cdot\omega_1$, not something worse\textup{)}?

The obstruction lies in the Legendre relation. The
quasi-period $\eta_\tau$ satisfies
$\eta_1\tau-\eta_\tau=\pi i$, which constrains the
$A$- and $B$-cycle monodromies simultaneously. Any
correction $\phi$ that kills the $B$-cycle monodromy
$2\eta_\tau\cdot c_0$ must introduce an $A$-cycle defect of
magnitude $2\pi i\cdot c_0/\tau$, and this defect couples
to the modular parameter through the $\mathrm{SL}(2,\mathbb{Z})$
action. The question reduces to whether the resulting
cocycle in
$H^1(\mathrm{SL}(2,\mathbb{Z}),\,
\mathcal{O}^{\mathrm{mer}}(E_\tau)\otimes\End(\cA))$
is trivial.

For $c_0=0$ the answer is trivially yes: no correction
is needed. For $c_0\ne 0$, the Eisenstein fan suggests
the answer is no: the $G_2$ quasi-modularity anomaly
is intrinsic, and the entanglement between the chiral differential
and the coassociative coproduct
at genus~$1$ is a genuine feature of the theory, not a
coordinate artifact. But a proof (or a counterexample
at some specific algebra) remains open.
\end{question}



% ================================================================
% ACT V: THE LANDSCAPE
% ================================================================

\bigskip
The main line of computations treats the families that drive the
theory. The remaining standard families (free fields and lattice
algebras) provide complementary tests.

\section{The \texorpdfstring{$\beta\gamma$}{beta-gamma} and
 \texorpdfstring{$bc$}{bc} ordered bar complexes}
\label{subsec:betagamma-ordered-bar}

The $\beta\gamma$ system $\mathcal{F}_{\beta\gamma}$ has generators
$\beta, \gamma$ of conformal weights $(1,0)$ with
$\beta(z)\,\gamma(w) \sim 1/(z{-}w)$. The $bc$ ghost system
$\mathcal{F}_{bc}$ has generators $b,c$ of weights
$(\lambda,1{-}\lambda)$ with
$b(z)\,c(w) \sim 1/(z{-}w)$. Both are $\Einf$-chiral
(tier~(ii): vertex algebras with poles).

\begin{theorem}[Free-field ordered bar complexes;
\ClaimStatusProvedHere]
\label{thm:bg-ordered-bar}
\label{thm:bg-rmatrix}
\label{thm:bg-koszul-dual}
\label{thm:bc-ordered-bar}
Let $\mathcal{F}$ denote either $\mathcal{F}_{\beta\gamma}$ or
$\mathcal{F}_{bc}$, with generators $a,a'$ and OPE
$a(z)\,a'(w) \sim 1/(z{-}w)$.
\begin{enumerate}[label=\textup{(\alph*)}]
\item \emph{Collision residue.} The $d\log$ kernel absorbs
the simple OPE pole leaving the nilpotent kernel
\begin{equation}\label{eq:bg-theta}
\Theta \;:=\; a \otimes a' - a' \otimes a, \qquad
\Theta^2 = 0.
\end{equation}

\item \emph{Bar differential.} Every adjacent collision
inserts $a_{(0)}a' = \pm 1$, which the augmentation kills.
Hence $d = 0$ in the reduced complex.

\item \emph{R-matrix.} The connection
$\nabla = d - \Theta\,d(z_1{-}z_2)/(z_1{-}z_2)$ on
$\operatorname{Conf}_2^{\mathrm{ord}}(\bC)$ has monodromy
\begin{equation}\label{eq:bg-rmatrix}
R \;=\; \exp(2\pi i\,\Theta) \;=\; \id + 2\pi i\,\Theta,
\end{equation}
terminating at linear order because $\Theta^2 = 0$.

\item \emph{Koszul duals.} Duality exchanges statistics:
\begin{equation}\label{eq:bg-koszul-dual}
\mathcal{F}_{\beta\gamma}^{!,\mathrm{ord}}
\simeq \bigwedge\nolimits^{\bullet}(\susp\,\beta^*
 \oplus \susp\,\gamma^*),
\qquad
\mathcal{F}_{bc}^{!,\mathrm{ord}}
\simeq \mathrm{Sym}^{\bullet}(\susp\,b^*
 \oplus \susp\,c^*).
\end{equation}
Both are class~C with shadow depth $r_{\max} = 4$
\textup{(}Volume~I,
Theorem~\textup{\ref*{thm:betagamma-global-depth}}\textup{)},
and $K(bc_\lambda) = c_{bc}(\lambda)+c_{\beta\gamma}(\lambda) = 0$.
\end{enumerate}
\end{theorem}

\begin{proof}
Nilpotency: $\Theta^2 = 0$ because
$a_{(0)}a' = 1$ is scalar (central), so expanding
$(a\otimes a' - a'\otimes a)^2$ and using $a'_{(0)}a = -1$
gives cancellation. Bar differential: in the reduced complex
every collision residue $\pm 1$ lies in $\bC\cdot 1$,
which the augmentation kills. Monodromy: the parallel transport
$T(u,u_0) = \id + \Theta\log((u{+}x{-}y)/(u_0{+}x{-}y))$
terminates at linear order; the $d\log$ loop integral
contributes $2\pi i$. Koszul dual: the antisymmetric
$\Theta_{\beta\gamma}$ (bosonic generators) dualizes to
$\bigwedge$; $\Theta_{bc}$ (fermionic generators) dualizes
to $\mathrm{Sym}$.
\end{proof}

\begin{remark}[Contrast with Heisenberg]
\label{rem:bg-vs-heisenberg}
The Heisenberg R-matrix $R(z) = \exp(k\hbar/z)$
(Theorem~\ref{thm:heisenberg-rmatrix}) is an infinite power
series; the free-field R-matrix terminates at linear order.
This dichotomy reflects the pole-order classification:
double pole with scalar residue (class~G) vs.\ simple pole
with nilpotent residue (class~C).
\end{remark}

\begin{theorem}[Wakimoto bar complex descent;
\ClaimStatusProvedHere]
\label{thm:wakimoto-ordered-bar}
For $\fg$ simple of rank~$r$ and $k \neq -h^\vee$, the
Wakimoto module
$\mathcal{M}_{\mathrm{Wak}}
= \bigotimes_{\alpha\in\Delta_+}(\beta\gamma)_\alpha
\otimes \cH_{\mathfrak{h}}$
satisfies
\begin{equation}\label{eq:wakimoto-ordered-bar-descent}
H^\bullet\bigl(
 \Barch(\mathcal{M}_{\mathrm{Wak}}),\;
 d_{\Barch} + Q_{\mathrm{DS}}\bigr)
\;\simeq\; \Barch(\widehat{\fg}_k),
\end{equation}
and at the Koszul dual level,
$Y^{\mathrm{dg}}_\hbar(\fg)
\simeq H^\bullet\!\bigl(
 \bigotimes_\alpha
 \mathcal{F}_{\beta\gamma,\alpha}^{!,\mathrm{ord}}
 \otimes Y(\mathfrak{u}(1))^{\otimes r},\;
 Q_{\mathrm{DS}}^!\bigr)$.
\end{theorem}

\begin{proof}
$Q_{\mathrm{DS}}$ commutes with $d_{\Barch}$, so a spectral
sequence with $E_1 = H^\bullet(\Barch(\mathcal{M}_{\mathrm{Wak}}),
d_{\Barch})$ and $d_1 = Q_{\mathrm{DS}}$ converges to
$\Barch(\widehat{\fg}_k)$ (Volume~I,
Theorem~\ref*{thm:wakimoto-bar}).
Theorem~\ref{thm:ordered-associative-ds-principal}
ensures Koszul duality and DS commute.
\end{proof}

\begin{table}[ht]
\centering
\renewcommand{\arraystretch}{1.4}
\begin{tabular}{llccll}
\toprule
\textbf{System} & \textbf{$R$-matrix}
 & \textbf{Depth} & \textbf{Class}
 & \textbf{$A^{!,\mathrm{ord}}_{\mathrm{line}}$}
 & \textbf{$\Einf$}\\
\midrule
$\cH_k$ & $e^{k\hbar/z}$ \textup{(scalar)}
 & $2$ & $\mathbf{G}$ & $\bC[h_0,h_1,\ldots]$
 & $\checkmark$ \\[2pt]
$\beta\gamma$ & $\id + 2\pi i\,\Theta$
 \textup{(nilpotent)}
 & $4$ & $\mathbf{C}$
 & $\bigwedge^*(\beta^*,\gamma^*)$
 & $\checkmark$ \\[2pt]
$bc$ & $\id + 2\pi i\,\Theta_{bc}$
 \textup{(nilpotent)}
 & $4$ & $\mathbf{C}$
 & $\mathrm{Sym}^*(b^*,c^*)$
 & $\checkmark$ \\[2pt]
$\widehat{\fg}_k$ & $\exp(\hbar\,\Omega/z)$
 & $3$ & $\mathbf{L}$ & $Y_\hbar(\fg)$
 & $\checkmark$ \\[2pt]
$\mathrm{Vir}_c$ & \textup{[triple pole]}
 & $\infty$ & $\mathbf{M}$ & $Y^{\mathrm{dg}}(\mathrm{Vir})$
 & $\checkmark$ \\
\bottomrule
\end{tabular}
\caption{Ordered bar complex data for the standard families.
All entries are $\Einf$-chiral; the $R$-matrix is derived
from the local OPE in every case.}
\label{tab:free-field-contrast}
\end{table}

\begin{remark}[The Wakimoto escalation]
\label{rem:wakimoto-escalation}
Shadow depth escalates through the Wakimoto chain:
$\cH_k \otimes \beta\gamma^{\otimes |\Delta_+|}$
(depths $2$, $4$)
$\xrightarrow{Q_{\mathrm{DS}}} \widehat{\fg}_k$
(depth $3$)
$\xrightarrow{Q_{\mathrm{DS}}} \mathrm{Vir}_c$
(depth $\infty$).
The second reduction manufactures the quartic pole
via Sugawara, escalating from class~L to class~M.
\end{remark}

\section{Lattice vertex algebras}
\label{subsec:lattice-ordered-bar}

Let $\Lambda$ be an even lattice of rank $d$ with bilinear form
$\langle\cdot,\cdot\rangle$ and let $\varepsilon \colon \Lambda \times
\Lambda \to \bC^\times$ be a normalized bimultiplicative $2$-cocycle.
The lattice vertex algebra $\Vlat_\Lambda^\varepsilon$ has vertex operators
$Y(e^\alpha, z)$ for $\alpha \in \Lambda$ with OPE
\begin{equation}\label{eq:lattice-ope-ordered}
Y(e^\alpha, z)\,Y(e^\beta, w) \;\sim\;
 \varepsilon(\alpha,\beta)\,(z - w)^{\langle\alpha,\beta\rangle}\,
 Y(e^{\alpha+\beta}, w) + \cdots
\end{equation}
The $\Eone$/$\Einf$ character of $\Vlat_\Lambda^\varepsilon$ depends on
the symmetry of the cocycle, as established in
\S\ref{sec:lattice:e1-families} of the lattice foundations chapter.

\subsubsection{Symmetric cocycle: \texorpdfstring{$\Einf$}{E-infinity}-chiral}

\begin{theorem}[Ordered bar complex with symmetric cocycle;
\ClaimStatusProvedHere]
\label{thm:lattice-symmetric-ordered-bar}
Let $\varepsilon$ be a symmetric cocycle satisfying the Borcherds
symmetry $\varepsilon(\alpha,\beta) = (-1)^{\langle\alpha,\beta\rangle}
\varepsilon(\beta,\alpha)$. Then $\Vlat_\Lambda^\varepsilon$ is
$\Einf$-chiral, and the ordered bar complex
$\Barch(\Vlat_\Lambda^\varepsilon)$ has:

\begin{enumerate}[label=\textup{(\arabic*)}]
\item \textbf{Collision residue.}
The collision residue kernel on the diagonal $\alpha$-sector is:
\begin{equation}\label{eq:lattice-symmetric-residue}
r(z) \;=\; \frac{1}{z}\sum_{\alpha \in \Lambda}\,
 \langle\alpha,\alpha\rangle\,
 (e^\alpha \otimes e^{-\alpha}).
\end{equation}
This is the \emph{diagonal} residue: only the Heisenberg-type
contributions $\alpha \otimes (-\alpha)$ survive, weighted by the norm
$\langle\alpha,\alpha\rangle$.

\item \textbf{R-matrix.}
Since $\varepsilon$ is symmetric, the R-matrix is
determined by the local OPE \textup{(}tier~(ii) of the three-tier
picture\textup{)}.
The ordered-to-unordered descent
$\Barch \to \barBgeom$
is the R-matrix-twisted $\Sigma_n$-quotient, with twist
$R(z)$ derived from the OPE monodromy:
\begin{equation}\label{eq:lattice-symmetric-rmatrix}
R_{\alpha,\beta}(z) \;=\;
 (-1)^{\langle\alpha,\beta\rangle}\,
 e^{\pi i\langle\alpha,\beta\rangle},
\end{equation}
the standard lattice braiding phase. This is \emph{not} independent
input: it is the monodromy of the flat connection on
$\operatorname{Conf}_2(\bC)$ determined by the OPE poles,
confirming the $\Einf$ character.
\end{enumerate}
\end{theorem}

\begin{proof}
The OPE $Y(e^\alpha,z)\,Y(e^\beta,w) \sim
\varepsilon(\alpha,\beta)(z-w)^{\langle\alpha,\beta\rangle}
Y(e^{\alpha+\beta},w) + \cdots$ has a pole of order
$-\langle\alpha,\beta\rangle$ at $z = w$. After $d\log$ absorption,
the collision residue extracts the order-$1$ coefficient:
$\operatorname{Res}_{z=w}[(z-w)^{-m}\,d\log(z-w)] = \delta_{m,1}$.
For the diagonal sector $\beta = -\alpha$, the pole order is
$-\langle\alpha,-\alpha\rangle = \langle\alpha,\alpha\rangle$, which
is always even for an even lattice. The non-diagonal residue at
order~$1$ gives~\eqref{eq:lattice-symmetric-residue}.

For the R-matrix: the connection is
$\nabla = d - r(z)\,d(z_1 - z_2)$, and the monodromy around
$z_1 = z_2$ is $\exp(2\pi i \cdot r)$. On the lattice sector
$(e^\alpha, e^\beta)$, this is
$(-1)^{\langle\alpha,\beta\rangle}\,
e^{\pi i \langle\alpha,\beta\rangle}$, matching the standard
lattice braiding. The Borcherds symmetry of $\varepsilon$ ensures
this monodromy is \emph{derived from} the local OPE, confirming
$\Einf$.
\end{proof}

\subsubsection{Non-symmetric cocycle: genuinely \texorpdfstring{$\Eone$}{E1}-chiral}

\begin{theorem}[Ordered bar complex with non-symmetric cocycle;
\ClaimStatusProvedHere]
\label{thm:lattice-nonsymmetric-ordered-bar}
Let $\varepsilon$ be a non-symmetric cocycle: there exist
$\alpha, \beta \in \Lambda$ with
$\varepsilon(\alpha,\beta)/\varepsilon(\beta,\alpha) \neq
(-1)^{\langle\alpha,\beta\rangle}$. Then
$\Vlat_\Lambda^\varepsilon$ is genuinely $\Eone$-chiral
\textup{(}nonlocal\textup{)},
and the ordered bar complex $\Barch(\Vlat_\Lambda^\varepsilon)$ has:

\begin{enumerate}[label=\textup{(\arabic*)}]
\item \textbf{Differential.} In bar degree~$2$:
\begin{equation}\label{eq:lattice-nonsym-diff}
d([e^{\alpha_1} \mid e^{\alpha_2}])
\;=\; \varepsilon(\alpha_1,\alpha_2)\,
 \delta_{\langle\alpha_1,\alpha_2\rangle,\,-1}\,
 [e^{\alpha_1 + \alpha_2}].
\end{equation}
The cocycle $\varepsilon(\alpha_1,\alpha_2)$ appears with its value at
the \emph{specific ordered pair} $(\alpha_1, \alpha_2)$. For a
non-symmetric cocycle, the coefficients
$\varepsilon(\alpha,\beta) \neq
(-1)^{\langle\alpha,\beta\rangle}\varepsilon(\beta,\alpha)$, so
permuting the pair changes the differential. No $\Sigma_2$-quotient is
possible.

\item \textbf{R-matrix as independent input.}
The R-matrix $R_{\alpha,\beta}$ is part of the defining data of the
$\Eone$-chiral algebra, not derivable from a local OPE. The ratio
\begin{equation}\label{eq:lattice-nonsym-rmatrix-ratio}
\frac{\varepsilon(\alpha,\beta)}{\varepsilon(\beta,\alpha)}
\end{equation}
measures the asymmetry of the cocycle and enters the R-matrix directly.
When this ratio is not $(-1)^{\langle\alpha,\beta\rangle}$, the
braiding is genuinely nonlocal: it cannot be obtained as the monodromy
of any connection on $\operatorname{Conf}_2(\bC)$ determined by a local
OPE. This is tier~(iii) of the three-tier picture.
\end{enumerate}
\end{theorem}

\begin{proof}
The differential follows from the
lattice OPE~\eqref{eq:lattice-ope-ordered} and the ordered bar
construction: adjacent collisions in the ordered bar complex use the
$d\log$ kernel, selecting simple poles
($\langle\alpha_i,\alpha_{i+1}\rangle = -1$) with coefficient
$\varepsilon(\alpha_i, \alpha_{i+1})$ at the specific ordered pair.

For the R-matrix claim: the Borcherds symmetry
$\varepsilon(\alpha,\beta) =
(-1)^{\langle\alpha,\beta\rangle}\varepsilon(\beta,\alpha)$
is precisely the condition for the $\Sigma_2$-action on the ordered bar
complex to be compatible with the differential. When this fails, the
descent $\Barch \to \barBgeom$ requires independent braiding data, and
the algebra is genuinely $\Eone$.
\end{proof}

\subsubsection{Koszul dual of lattice algebras}

\begin{theorem}[Ordered Koszul dual of lattice algebras;
\ClaimStatusProvedHere]
\label{thm:lattice-ordered-koszul-dual}
The ordered Koszul dual of $\Vlat_\Lambda^\varepsilon$ depends on
the lattice:
\begin{enumerate}[label=\textup{(\roman*)}]
\item For a self-dual even lattice $\Lambda = \Lambda^*$
 with symmetric cocycle, the Koszul dual
 $(\Vlat_\Lambda^\varepsilon)^!$ is again a lattice vertex algebra
 $(\Vlat_\Lambda^{\varepsilon^{-1}})^c$ with inverse cocycle
 \textup{(}Theorem~\textup{\ref{thm:lattice:koszul-dual}}\textup{)}.
\item For a non-self-dual lattice or non-symmetric cocycle, the
 ordered Koszul dual carries strictly more structure: the
 ordered bar complex has independent generators in each ordering,
 and the dual algebra reflects this asymmetry.
\item In all cases, the lattice grading is preserved: the Koszul dual
 decomposes as $\bigoplus_{\gamma \in \Lambda}
 (\Vlat_\Lambda^\varepsilon)^!_\gamma$ with the differential
 preserving total lattice degree.
\end{enumerate}
\end{theorem}

\begin{proof}
Part~(i) follows from Theorem~\ref{thm:lattice:koszul-dual} and its
proof: the bar complex coproduct has structure constants
$\varepsilon(\alpha,\beta)^{-1}$, and linear duality recovers the
lattice algebra with inverse cocycle. Self-duality $\Lambda = \Lambda^*$
ensures the dual lives on the same lattice.

Part~(ii): for a non-symmetric cocycle, the ordered bar complex in
degree~$2$ has generators $[e^\alpha \mid e^\beta]$ and
$[e^\beta \mid e^\alpha]$ that are \emph{independent} (no
$\Sigma_2$-identification). Their duals generate the Koszul dual
algebra with product reflecting the asymmetric cocycle values.

Part~(iii): the bar differential preserves the total lattice degree
$\gamma = \sum_i \alpha_i \in \Lambda$
(Theorem~\ref{thm:lattice:bar-structure}(i)), so the Koszul dual
inherits the same decomposition.
\end{proof}

\begin{remark}[The three-tier picture for lattice algebras]
\label{rem:lattice-three-tier}
Lattice vertex algebras provide the cleanest illustration of the
three-tier picture:
\begin{enumerate}[label=(\roman*)]
\item \emph{Pole-free commutative} ($R = \tau$): degenerate case where
 $\langle\alpha,\beta\rangle = 0$ for all $\alpha, \beta$ (the lattice
 is trivial or the restriction to a sublattice with vanishing pairing).
\item \emph{$\Einf$ with poles} ($R \neq \tau$, derived from locality):
 symmetric cocycle, OPE with poles from
 $\langle\alpha,\beta\rangle < 0$. The R-matrix is the standard lattice
 braiding, determined by the OPE monodromy.
\item \emph{Genuinely $\Eone$} ($R \neq \tau$, independent input):
 non-symmetric cocycle. The R-matrix is part of the defining data.
 The descent from ordered to unordered configurations requires this
 independent braiding datum.
\end{enumerate}
The lattice family thus realizes all three tiers explicitly, with the
cocycle symmetry as the single discrete parameter controlling the
$\Einf$/$\Eone$ transition.
\end{remark}



\section{Evaluation modules and the line category}
\label{sec:evaluation-modules}

The ordered bar complex and the Yangian $Y_\hbar(\mathfrak{g})$ classify the
line operators of the corresponding holomorphic-topological theory. This
section develops the representation-theoretic side: evaluation modules, the
Drinfeld polynomial classification, tensor products and their braiding, and the
identification of the resulting module category with the line category
$\mathcal{C}_{\mathrm{line}}$.

Throughout, $\mathfrak{g}$ denotes a finite-dimensional simple Lie algebra
over~$\mathbb{C}$ of rank~$r$, with Cartan matrix $(a_{ij})$,
simple roots $\alpha_1,\dots,\alpha_r$, fundamental weights
$\omega_1,\dots,\omega_r$, and normalised Killing form
$(\cdot,\cdot)$ for which $(\alpha,\alpha)=2$ for long roots. We write
$\Omega\in\mathfrak{g}\otimes\mathfrak{g}$ for the quadratic Casimir element
and $P\in\operatorname{End}(V\otimes W)$ for the permutation operator
$v\otimes w\mapsto w\otimes v$.

\subsection{The Yangian $Y_\hbar(\mathfrak{g})$ and the evaluation homomorphism}

\begin{definition}[Drinfeld Yangian]
\label{def:yangian-drinfeld}
The \emph{Yangian} $Y_\hbar(\mathfrak{g})$ is the associative
$\mathbb{C}[\![\hbar]\!]$-algebra generated by
$\{J^a_0,\,J^a_1\mid a=1,\dots,\dim\mathfrak{g}\}$
subject to:
\begin{enumerate}[label=(\roman*)]
\item $J^a_0$ generate a copy of $U(\mathfrak{g})$;
\item $[J^a_0,\,J^b_1]=f^{ab}{}_c\,J^c_1$;
\item Drinfeld's ternary relation:
$[J^a_1,\,[J^b_0,\,J^c_1]]-[J^a_0,\,[J^b_1,\,J^c_1]]$
$=\hbar^2\,A^{abc}_{def}\{J^d_0,J^e_0,J^f_0\}$,
where $A^{abc}_{def}$ is the totally symmetric tensor
of \textup{\cite{Drinfeld-Yangians}} and
$\{x,y,z\}=\tfrac{1}{24}\sum_{\sigma\in S_3}x_{\sigma(1)}y_{\sigma(2)}z_{\sigma(3)}$.
\end{enumerate}
The coproduct is
$\Delta(J^a_0)=J^a_0\otimes 1+1\otimes J^a_0$,
$\Delta(J^a_1)=J^a_1\otimes 1+1\otimes J^a_1
+\tfrac{\hbar}{2}\,f^a{}_{bc}\,J^b_0\otimes J^c_0$.
\end{definition}

\begin{construction}[Evaluation homomorphism]
\ClaimStatusProvedHere
\label{constr:evaluation-map}
For each $a\in\mathbb{C}$, there is a surjective algebra
homomorphism
\[
\mathrm{ev}_a\colon
Y_\hbar(\mathfrak{g})\longrightarrow U(\mathfrak{g}),
\qquad
J^a_0\mapsto J^a,
\quad
J^a_1\mapsto a\,J^a.
\]
The evaluation map is an algebra homomorphism because
$[J^a,a\,J^b]=a\,f^{ab}{}_c\,J^c=a\,f^{ab}{}_c\,J^c_0$,
and the ternary relation reduces to the Jacobi identity in
$U(\mathfrak{g})$.
\end{construction}

\begin{definition}[Evaluation module]
\label{def:evaluation-module}
Let $(\rho,V)$ be a finite-dimensional
$\mathfrak{g}$-representation and $a\in\mathbb{C}$.
The \emph{evaluation module} $V_a$ is the
$Y_\hbar(\mathfrak{g})$-module obtained by pulling back
along $\mathrm{ev}_a$:
\[
V_a
:=
\mathrm{ev}_a^*(V).
\]
Concretely, $V_a=V$ as a vector space, with
\begin{equation}\label{eq:eval-transfer-matrix}
T^a(u)\cdot v
=
v+\frac{\hbar\,\rho(\Omega_{(1)})\otimes\Omega_{(2)}}{u-a}\cdot v
=
v+\frac{\hbar}{u-a}\sum_\alpha
\rho(J^\alpha)\,v\otimes J_\alpha,
\end{equation}
where $\Omega=\sum_\alpha J^\alpha\otimes J_\alpha$
is the Casimir tensor and the transfer matrix
$T(u)=1+\hbar\sum_{n\geq 0}J_n\, u^{-n-1}$.
\end{definition}

\subsection{Explicit evaluation modules for $\mathfrak{sl}_2$}

\begin{computation}[$\mathfrak{sl}_2$ evaluation module; \ClaimStatusProvedHere]
\label{comp:sl2-eval}
\index{evaluation module!sl2@$\mathfrak{sl}_2$}
Let $\mathfrak{g}=\mathfrak{sl}_2$ with standard basis
$e,f,h$ and $[h,e]=2e$, $[h,f]=-2f$, $[e,f]=h$.
The fundamental representation is $V=\mathbb{C}^2$ with
basis $v_+,v_-$ and
\[
\rho(e)=\begin{pmatrix}0&1\\0&0\end{pmatrix},
\quad
\rho(f)=\begin{pmatrix}0&0\\1&0\end{pmatrix},
\quad
\rho(h)=\begin{pmatrix}1&0\\0&-1\end{pmatrix}.
\]
The Casimir tensor is
$\Omega=\tfrac{1}{2}h\otimes h+e\otimes f+f\otimes e$,
so $\rho(\Omega)\in\operatorname{End}(\mathbb{C}^2)$
is $\rho(\Omega)=\tfrac{1}{2}C_2(\mathrm{fund})\cdot\id
=\tfrac{3}{4}\cdot\id$ (half the quadratic Casimir eigenvalue).

The evaluation module $V_a=\mathbb{C}^2$ has
\[
J^a_0\cdot v_\pm=\rho(J^a)\,v_\pm,
\qquad
J^a_1\cdot v_\pm=a\,\rho(J^a)\,v_\pm.
\]
The transfer matrix in the fundamental representation is
\begin{equation}\label{eq:sl2-transfer-fundamental}
T(u)
\;=\;
\id_2
\;+\;
\frac{\hbar}{u-a}
\begin{pmatrix}
\tfrac{1}{2}h & f \\[4pt]
e & -\tfrac{1}{2}h
\end{pmatrix},
\end{equation}
where the matrix entries are $\mathfrak{sl}_2$-generators
acting via $\rho$ in the quantum space and via the
fundamental in the auxiliary space.
\end{computation}

\begin{proof}
The evaluation homomorphism sends each Yangian generator to its
$\mathfrak{sl}_2$-action scaled by the spectral parameter, so the
displayed formulas for $J_0$ and $J_1$ are immediate from
$J_n^a\mapsto a^n\rho(J^a)$.
With the chosen matrices,
\[
\rho(\Omega)
=\tfrac12\rho(h)^2+\rho(e)\rho(f)+\rho(f)\rho(e)
=\tfrac34\,\id_{\mathbb C^2},
\]
which is the stated fundamental Casimir value.
Substituting this into the general transfer-matrix formula
\eqref{eq:eval-transfer-matrix} yields
\eqref{eq:sl2-transfer-fundamental}.
\end{proof}

\begin{theorem}[$R$-matrix on $V_a\otimes V_b$ for $\mathfrak{sl}_2$;
\ClaimStatusProvedHere]
\label{thm:sl2-R-matrix}
\index{R-matrix!sl2 evaluation@$\mathfrak{sl}_2$ evaluation}
The $R$-matrix governing the braiding of
$V_a\otimes V_b$ is
\begin{equation}\label{eq:sl2-R-eval}
R(a-b)
\;=\;
(a-b)\,\id\;+\;\hbar\,P
\;\in\;
\operatorname{End}(\mathbb{C}^2\otimes\mathbb{C}^2),
\end{equation}
where $P$ is the permutation operator $v\otimes w\mapsto w\otimes v$.
Equivalently, in the normalised form $\check{R}(u)=P\,R(u)$:
\[
\check{R}(u)
=
u\,P+\hbar\,\id.
\]
This is the Yang $R$-matrix.
\end{theorem}

\begin{proof}
The intertwining equation
$\Delta^{\mathrm{op}}(x)\,R=R\,\Delta(x)$ for all
$x\in Y_\hbar(\mathfrak{sl}_2)$, evaluated on
$V_a\otimes V_b$, uniquely determines $R$ up to scalar.
For $x=J^a_0$:
$[\rho^{(1)}(J^a)+\rho^{(2)}(J^a),R]=0$, so $R$
commutes with the diagonal $\mathfrak{sl}_2$-action. By
Schur's lemma applied to the decomposition
$\mathbb{C}^2\otimes\mathbb{C}^2=\mathrm{Sym}^2\oplus\bigwedge^2$,
$R=\alpha\,\Pi_{\mathrm{sym}}+\beta\,\Pi_{\mathrm{alt}}$
for scalars $\alpha,\beta$. Writing
$\Pi_{\mathrm{sym}}=\tfrac{1}{2}(\id+P)$ and
$\Pi_{\mathrm{alt}}=\tfrac{1}{2}(\id-P)$:
\[
R=\tfrac{\alpha+\beta}{2}\,\id
+\tfrac{\alpha-\beta}{2}\,P.
\]
Evaluating the intertwining equation for
$x=J^a_1$ with $\Delta(J^a_1)=J^a_1\otimes 1+1\otimes J^a_1
+\tfrac{\hbar}{2}f^a{}_{bc}J^b_0\otimes J^c_0$
yields $\alpha-\beta=\hbar$ and $\alpha+\beta=a-b$.
Hence $R=(a-b)\id+\hbar\,P$.
\end{proof}

\begin{corollary}[Clebsch--Gordan decomposition and non-semisimplicity;
\ClaimStatusProvedHere]
\label{cor:sl2-clebsch-gordan}
The tensor product $V_a\otimes V_b$ of fundamental
$\mathfrak{sl}_2$-evaluation modules is:
\begin{enumerate}[label=(\alph*)]
\item \textbf{Semisimple} when $a-b\notin\{0,\pm\hbar\}$:
the $R$-matrix~\eqref{eq:sl2-R-eval} is invertible, and
$V_a\otimes V_b\simeq W_{3,a+b}\oplus W_{1,a+b}$,
where $W_{d,c}$ denotes the evaluation module of dimension~$d$
at the averaged spectral parameter.

\item \textbf{Non-semisimple} when $a-b=\pm\hbar$:
$R(a-b)=\pm\hbar\,\id+\hbar\,P$ is not scalar but remains
invertible. The tensor product is reducible but
indecomposable: $V_a\otimes V_b$ has a two-step filtration
with Jordan--H\"older factors $W_3$ and $W_1$, but the
extension does not split.

\item \textbf{Singular} when $a=b$: $R(0)=\hbar\,P$
has eigenvalues $\pm\hbar$, and the tensor product $V_a\otimes V_a$
degenerates: the braiding becomes the graded
permutation, and the intertwiner is no longer generic.
The image of the bar complex at this point reflects the
coalescence singularity of the ordered configuration space.
\end{enumerate}
\end{corollary}

\begin{proof}
The eigenvalues of $R(u)=u\,\id+\hbar\, P$ on the
$\mathrm{Sym}^2$ and $\bigwedge^2$ summands are
$u+\hbar$ and $u-\hbar$, respectively. Part~(a):
when $u=a-b\neq 0,\pm\hbar$, both eigenvalues are nonzero
and the corresponding projectors split the tensor product.
Part~(b): when $u=\pm\hbar$, one eigenvalue is $2\hbar$
or $0$ and the other is $0$ or $-2\hbar$; the zero
eigenvalue signals a non-split extension. Part~(c) is
$u=0$.
\end{proof}

\subsection{Evaluation modules for $\mathfrak{sl}_3$}

\begin{computation}[$\mathfrak{sl}_3$ fundamental evaluation module;
\ClaimStatusProvedHere]
\label{comp:sl3-eval-fundamental}
\index{evaluation module!sl3@$\mathfrak{sl}_3$!fundamental}
Let $\mathfrak{g}=\mathfrak{sl}_3$ with standard Chevalley
generators $e_i,f_i,h_i$ ($i=1,2$). The fundamental
representation $V=\mathbb{C}^3$ has basis $v_1,v_2,v_3$
with weight vectors of weights $\omega_1,\,\omega_1-\alpha_1,\,
\omega_1-\alpha_1-\alpha_2$. The evaluation module
$V_a=\mathbb{C}^3$ has
\[
J^{e_i}_0=E_i,\quad J^{f_i}_0=F_i,\quad J^{h_i}_0=H_i,
\qquad
J^{e_i}_1=a\,E_i,\quad J^{f_i}_1=a\,F_i,\quad J^{h_i}_1=a\,H_i,
\]
where $E_i,F_i,H_i$ are the standard $3\times 3$ matrices
representing $e_i,f_i,h_i$.

The $R$-matrix on $V_a\otimes V_b$ in the fundamental takes
the universal Yang form:
\begin{equation}\label{eq:sl3-R-fund}
R(u)
=
u\,\id\;+\;\hbar\,P
\;\in\;
\operatorname{End}(\mathbb{C}^3\otimes\mathbb{C}^3),
\end{equation}
where $u=a-b$. Under $\mathfrak{sl}_3$,
$\mathbb{C}^3\otimes\mathbb{C}^3
\simeq\mathrm{Sym}^2(\mathbb{C}^3)\oplus\bigwedge^2(\mathbb{C}^3)$
\textup{(}dimensions $6+3$\textup{)}, and the eigenvalues of
$R(u)$ are $u+\hbar$ on $\mathrm{Sym}^2$ and $u-\hbar$
on $\bigwedge^2$, exactly as for $\mathfrak{sl}_2$.
\end{computation}

\begin{proof}
Applying the evaluation map to the defining representation of
$\mathfrak{sl}_3$ gives the displayed formulas for the level-$0$ and
level-$1$ generators.
The decomposition
$\mathbb C^3\otimes\mathbb C^3
=\mathrm{Sym}^2(\mathbb C^3)\oplus\bigwedge^2(\mathbb C^3)$
is the standard $\mathfrak{sl}_3$ tensor-product splitting.
Theorem~\ref{thm:yang-r-sl3} gives
$R(u)=u\,\id+\hbar P$ on the fundamental tensor square.
Since $P$ acts by $+1$ on the symmetric summand and by $-1$ on the
alternating summand, the two eigenvalues are $u+\hbar$ and $u-\hbar$.
\end{proof}

\begin{computation}[$\mathfrak{sl}_3$ adjoint evaluation module;
\ClaimStatusProvedHere]
\label{comp:sl3-eval-adjoint}
\index{evaluation module!sl3@$\mathfrak{sl}_3$!adjoint}
The adjoint representation $W=\mathfrak{sl}_3$ has
$\dim W=8$. The evaluation module $W_a$ is the pullback
of the adjoint $\mathfrak{sl}_3$-module along
$\mathrm{ev}_a$. Since
$\mathbb{C}^3\otimes(\mathbb{C}^3)^*\simeq
\mathfrak{sl}_3\oplus\mathbb{C}$, one has
\[
W_a\simeq (V_a\otimes V^*_a)/\mathbb{C},
\]
where $V^*_a$ is the contragredient evaluation module at the
same spectral parameter. The $R$-matrix on $V_a\otimes W_b$
\textup{(}fundamental $\otimes$ adjoint\textup{)} is determined by
the Yangian intertwining equation and decomposes into three
eigenvalues on the three $\mathfrak{sl}_3$-irreducible summands
of $\mathbb{C}^3\otimes\mathfrak{sl}_3$:
\[
\mathbb{C}^3\otimes\mathfrak{sl}_3
\;\simeq\;
V(2\omega_1+\omega_2)\;\oplus\;V(\omega_1)\;\oplus\;V(\omega_2),
\]
with eigenvalues $u+2\hbar$, $u-\hbar$, $u+\hbar$
respectively, where $u=a-b$.
\end{computation}

\begin{proof}
The identification
$W_a\simeq (V_a\otimes V_a^*)/\mathbb C$ is the usual realization of
the adjoint representation as the traceless part of
$\End(\mathbb C^3)$, transported through the evaluation map.
The tensor-product decomposition
$\mathbb C^3\otimes\mathfrak{sl}_3
\cong V(2\omega_1+\omega_2)\oplus V(\omega_1)\oplus V(\omega_2)$
is the standard Littlewood--Richardson decomposition
$\mathbf 3\otimes\mathbf 8=\mathbf{15}\oplus\mathbf 3\oplus\mathbf 6^*$.
Proposition~\ref{prop:r-matrix-eigenvalue} implies that the Yangian
$R$-matrix acts by scalars on these three irreducible summands.
Using the fusion realization of $W_a$ from the fundamental and dual
fundamental evaluation modules, those scalars are fixed by the
fundamental Yang $R$-matrix of
Computation~\ref{comp:sl3-eval-fundamental}; the three channels pick up
the shifts $+2\hbar$, $-\hbar$, and $+\hbar$, which gives the stated
eigenvalues.
\end{proof}

\subsection{Drinfeld polynomials and the classification of finite-dimensional modules}

\begin{theorem}[Drinfeld classification; \ClaimStatusProvedElsewhere]
\label{thm:drinfeld-classification}
\index{Drinfeld polynomial!classification of Yangian modules|textbf}
\index{evaluation module!Drinfeld polynomial}
For $\mathfrak{g}$ a finite-dimensional simple Lie algebra
of rank~$r$, the isomorphism classes of irreducible
finite-dimensional $Y_\hbar(\mathfrak{g})$-modules are in
bijection with $r$-tuples of monic polynomials
$(P_1(u),\dots,P_r(u))\in\mathbb{C}[u]^r$, called the
\emph{Drinfeld polynomials}.

The bijection sends an irreducible module~$L$ to the
$r$-tuple where $P_i(u)$ records the eigenvalues of the
level-$1$ generating series associated to the $i$-th simple
root on a highest-weight vector of~$L$:
\[
T_i(u)\cdot v_+
=
\frac{P_i(u-\hbar)}{P_i(u)}\,v_+.
\]
\end{theorem}

\begin{proposition}[Evaluation modules as single-root Drinfeld polynomials;
\ClaimStatusProvedHere]
\label{prop:eval-drinfeld}
\index{Drinfeld polynomial!of evaluation module}
Let $V(\omega_i)$ be the $i$-th fundamental
representation of~$\mathfrak{g}$ and $a\in\mathbb{C}$.
The evaluation module $V(\omega_i)_a$ has Drinfeld
polynomials
\[
P_j(u)=
\begin{cases}
u-a & \text{if }j=i,\\
1 & \text{if }j\neq i.
\end{cases}
\]
More generally, the evaluation of an irreducible
$\mathfrak{g}$-module $V(\lambda)$ with highest weight
$\lambda=\sum_i m_i\omega_i$ at points
$a_1,\dots,a_N\in\mathbb{C}$ produces the tensor product
evaluation module $V(\lambda_1)_{a_1}\otimes\cdots\otimes
V(\lambda_N)_{a_N}$, whose Drinfeld polynomials are
\[
P_j(u)=\prod_{k=1}^N (u-a_k)^{m_j^{(k)}},
\]
where $m_j^{(k)}$ is the multiplicity of $\omega_j$ in
$\lambda_k$.
\end{proposition}

\begin{proof}
By definition, $\mathrm{ev}_a$ sends $J^a_1\mapsto a\,J^a$,
so the generating series $T_i(u)$ acts on the highest-weight
vector $v_+$ of $V(\omega_i)_a$ as a rational function of~$u$.
In the normalisation where $P_i(u)=u-a$, one obtains
$T_i(u)\cdot v_+=(u-a-\hbar)/(u-a)\cdot v_+$ and
$T_j(u)\cdot v_+=v_+$ for $j\neq i$, matching the
classification. The general formula follows from the
tensor product rule for Drinfeld polynomials: the
$P_j$ of a tensor product is the product of the
$P_j$'s of the factors.
\end{proof}

\begin{example}[Drinfeld polynomials for evaluation chains]
\label{ex:drinfeld-eval-chain}
For $\mathfrak{g}=\mathfrak{sl}_2$, rank $r=1$, and an
evaluation chain at points $a_1,\dots,a_n\in\mathbb{C}$
in the fundamental representation:
\[
P(u)=\prod_{k=1}^n(u-a_k).
\]
The module $V_{a_1}\otimes\cdots\otimes V_{a_n}$ is
irreducible if and only if $a_i-a_j\notin\mathbb{Z}\hbar$
for all $i\neq j$ (the non-resonance condition).
When $a_i-a_j\in\mathbb{Z}\hbar$ for some pair, the tensor
product is reducible and the irreducible subquotients are
classified by combinatorial data (segments in the sense of
Zelevinsky).
\end{example}

\subsection{The line category}

The line operators of a holomorphic-topological theory on
$\mathbb{C}_z\times\mathbb{R}_t$ are operators supported on
lines $\{z=z_0\}\times\mathbb{R}_t$. These form a
category $\mathcal{C}_{\mathrm{line}}$, and the ordered bar
complex identifies it with the module category of the
$E_1$-chiral Koszul dual.

\begin{theorem}[Line category as Yangian modules;
\ClaimStatusProvedHere]
\label{thm:line-category}
\index{line category!as Yangian modules|textbf}
\index{Koszul dual!E1-chiral!line category}
Let $A$ be a strongly admissible $E_1$-chiral algebra
with ordered bar coalgebra $C=\Barchord(A)$ and
$E_1$-chiral Koszul dual $A^!_{\mathrm{line}}$. Then:
\begin{enumerate}[label=(\roman*)]
\item The line category is
\[
\mathcal{C}_{\mathrm{line}}
\;\simeq\;
A^!_{\mathrm{line}}\text{-}\mathrm{mod}^{\mathrm{fd}},
\]
the category of finite-dimensional modules over the
$E_1$-chiral Koszul dual.

\item On the chirally Koszul locus \textup{(}where the bar--cobar
equivalence holds and $A^!_{\mathrm{line}}$ is formal\textup{)},
for $A$ in the affine lineage
\textup{(}i.e., $A=\hat{\mathfrak{g}}_k$
for simple $\mathfrak{g}$ at non-critical level\textup{)},
\[
A^!_{\mathrm{line}}
\;\simeq\;
Y_\hbar(\mathfrak{g})
\]
and accordingly
\[
\mathcal{C}_{\mathrm{line}}
\;\simeq\;
Y_\hbar(\mathfrak{g})\text{-}\mathrm{mod}^{\mathrm{fd}}.
\]

\item The evaluation modules $V_a$ of
Definition~\textup{\ref{def:evaluation-module}} are the line
operators obtained by inserting a Wilson line in
representation~$V$ at position $z=a$. The spectral
parameter $a$ is the holomorphic position on $\mathbb{C}_z$.
\end{enumerate}
\end{theorem}

\begin{proof}
Part~(i) is the content of the one-sided module/comodule
duality (Hypothesis~\ref{hyp:inputs}(I)) combined with
the ordered bar--cobar equivalence: left $C$-comodules are
equivalently right $A^!_{\mathrm{line}}$-modules.
Part~(ii) follows from the identification of
$A^!_{\mathrm{line}}$ with the dg-shifted Yangian
(proved in Volume~II,
Theorem~\ref*{thm:Koszul_dual_Yangian}).
Part~(iii) is the physical interpretation: a Wilson line in
representation $\rho$ at position $z=a$ defines a line
operator whose endomorphism algebra is
$\mathrm{ev}_a^*\rho$, and the intertwining between
different Wilson lines is governed by the $R$-matrix.
\end{proof}

\subsection{Tensor products and the braided monoidal structure}

\begin{theorem}[Braiding from the $R$-matrix;
\ClaimStatusProvedHere]
\label{thm:eval-braiding}
\index{braiding!evaluation modules|textbf}
\index{R-matrix!braiding of evaluation modules}
The tensor product of evaluation modules $V_a\otimes W_b$,
with $V,W$ finite-dimensional $\mathfrak{g}$-representations,
is a $Y_\hbar(\mathfrak{g})$-module via the coproduct~$\Delta$.
The braiding isomorphism
\[
\sigma_{V_a,W_b}\colon
V_a\otimes W_b
\;\xrightarrow{\;\sim\;}
W_b\otimes V_a
\]
is given by $\sigma_{V_a,W_b}=P\circ R_{V,W}(a-b)$,
where $R_{V,W}(u)\in\operatorname{End}(V\otimes W)(\!(u^{-1})\!)$
is the universal $R$-matrix of $Y_\hbar(\mathfrak{g})$
evaluated in $V\otimes W$.

When $a=b$, the $R$-matrix $R_{V,W}(0)$ degenerates:
its kernel is nontrivial whenever $V\otimes W$ is not
multiplicity-free. This singularity reflects the
coalescence of two line operators at the same point,
corresponding to the boundary stratum
$z_1=z_2$ in $\overline{\mathrm{FM}}_2^{\mathrm{ord}}(\mathbb{C})$.
\end{theorem}

\begin{proof}
The coproduct $\Delta$ of $Y_\hbar(\mathfrak{g})$ defines
the tensor product module structure. The universal
$R$-matrix satisfies the intertwining equation
$\Delta^{\mathrm{op}}(x)\,R=R\,\Delta(x)$, which means
$R_{V,W}(a-b)$ intertwines the two orderings of the
tensor product. The braiding axioms (Yang--Baxter equation)
follow from the quantum group structure. The coalescence
statement follows from the fact that $R_{V,W}(u)$ has a
pole at $u=0$ of order determined by the Casimir eigenvalue
in the tensor product, and the bar differential at the
boundary stratum $\partial\overline{\mathrm{FM}}_2$
imposes exactly the residue of this pole.
\end{proof}

\subsection{The Grothendieck ring}

\begin{theorem}[Grothendieck ring of Yangian modules;
\ClaimStatusProvedElsewhere]
\label{thm:grothendieck-yangian}
\index{Grothendieck ring!Yangian modules|textbf}
The Grothendieck ring of finite-dimensional
$Y_\hbar(\mathfrak{g})$-modules is
\begin{equation}\label{eq:grothendieck-yangian}
K_0\!\bigl(Y_\hbar(\mathfrak{g})\text{-}\mathrm{mod}^{\mathrm{fd}}\bigr)
\;\simeq\;
K_0\!\bigl(\operatorname{Rep}(\mathfrak{g})\bigr)[u],
\end{equation}
the polynomial extension of the representation ring
of~$\mathfrak{g}$. The isomorphism sends the class of
$V(\lambda)_a$ to $[V(\lambda)]\cdot u^{\deg P_\lambda}$,
and the product is the convolution of Drinfeld polynomials.
\end{theorem}

\begin{remark}[Generators of the Grothendieck ring]
The Grothendieck ring is generated as a
$\mathbb{Z}[u]$-algebra by the classes
$[V(\omega_i)_a]$ ($i=1,\dots,r$, $a\in\mathbb{C}$).
The Drinfeld polynomial classification
(Theorem~\ref{thm:drinfeld-classification}) establishes
the bijection between irreducible classes and
$r$-tuples of monic polynomials. The ring structure
reflects the fact that the tensor product of evaluation
modules, while not generically irreducible as a
$Y_\hbar(\mathfrak{g})$-module, becomes irreducible in the
Grothendieck group.
\end{remark}

\begin{remark}[Relation to the open trace formalism]
\label{rem:eval-open-trace}
The ordered open trace formalism of
Theorem~\ref{thm:ordered-open} acquires concrete content
on the Yangian locus. The diagonal bicomodule $C_\Delta$
parametrises the identity line operator; evaluation modules
$V_a$ are line operators at position~$a$; the interval
composition $\square_C^{\mathbf{R}}$ is the sequential
fusion of line operators along $\mathbb{R}_t$; and the
ordered pair-of-pants $\mu_P$ is the OPE of line operators
at the same holomorphic position. The Grothendieck ring
$K_0$ of the line category is the natural receptacle for
the decategorified open trace.
\end{remark}



\section{Explicit $R$-matrices for the standard families}
\label{sec:explicit-r-matrices}

The $R$-matrix is the monodromy of the Kohno connection
$\nabla = d - \sum_{i<j} r_{ij}(z_{ij})\,dz_{ij}$.
For the Yangian $Y(\mathfrak{g})$ of a simple Lie algebra,
the rational $R$-matrix on $V = \mathbb{C}^N$ takes a closed form:
the Yang $R$-matrix for type~$A$, and the
Zamolodchikov--Zamolodchikov $R$-matrix for types $B$, $C$, $D$,
and the exceptionals. In every case, the classical limit
recovers $r(u) = \Omega/u$ with $\Omega = \sum_a I^a \otimes I_a$
the Casimir tensor. All $R$-matrices below are derived from the
local OPE via Construction~\ref{constr:r-matrix-monodromy-vol1};
they are \emph{not} independent input.

\begin{proposition}[Eigenvalue decomposition; \ClaimStatusProvedHere]
\label{prop:r-matrix-eigenvalue}
Let $V \otimes V = \bigoplus_\lambda V_\lambda$ be the
$\mathfrak{g}$-irreducible decomposition. Then
$R(u)|_{V_\lambda} = f_\lambda(u)\cdot\id_{V_\lambda}$
with $f_\lambda(u)$ rational, $f_\lambda(u) \sim u$ as
$u \to \infty$. For type~$A$, $f_\lambda$ is linear;
for types $B$/$C$/$D$ and the exceptionals, additional
primitive invariants introduce rational corrections.
\end{proposition}

\begin{proof}
The universal $R$-matrix intertwines the diagonal
$\mathfrak{g}$-action on $V\otimes V$, so each irreducible summand
$V_\lambda$ is $R(u)$-stable.
Schur's lemma therefore implies that the restriction of $R(u)$ to
$V_\lambda$ is scalar, say $f_\lambda(u)\cdot\id_{V_\lambda}$.
The rationality of $f_\lambda$ comes from the RTT presentation of the
Yangian, where $R(u)$ is a rational function of the spectral
parameter, and the normalization of the universal $R$-matrix gives
$R(u)=u\,\id+O(1)$ as $u\to\infty$, hence $f_\lambda(u)\sim u$.
For type~$A$, only the permutation operator appears, so the eigenvalue
functions are linear; the other types acquire extra invariant tensors,
which produce the stated rational corrections.
\end{proof}

\begin{theorem}[Yang $R$-matrix for $\mathfrak{sl}_N$;
\ClaimStatusProvedHere]
\label{thm:yang-r-sl2}
\label{thm:yang-r-sl3}
\label{thm:yang-r-sl4}
For all $\mathfrak{sl}_N$, the rational $R$-matrix on
$V = \mathbb{C}^N$ is
\begin{equation}\label{eq:yang-sl2}
R(u) = u\cdot\id_{V\otimes V} + P,
\end{equation}
where $P(v\otimes w) = w\otimes v$ is the permutation.
The eigenvalues are $u{+}1$ on
$\mathrm{Sym}^2(\mathbb{C}^N)$ and $u{-}1$ on
$\bigwedge^2(\mathbb{C}^N)$. In the standard basis of\/
$\mathbb{C}^2 \otimes \mathbb{C}^2$:
\begin{equation}\label{eq:yang-sl2-matrix}
R(u) =
\begin{pmatrix}
u{+}1 & 0 & 0 & 0 \\
0 & u & 1 & 0 \\
0 & 1 & u & 0 \\
0 & 0 & 0 & u{+}1
\end{pmatrix}.
\end{equation}
\end{theorem}

\begin{proof}
$P$ has eigenvalue $+1$ on symmetric and $-1$ on alternating
tensors. The YBE holds in
$\operatorname{End}(V^{\otimes 3})$ for any $V$, using only
$P_{ij}^2 = \id$ and
$P_{12}P_{13}P_{23} = P_{23}P_{13}P_{12}$.
The classical limit $r(u) = P/u = \Omega/u$ satisfies the
CYBE~\eqref{eq:cybe-vol1}.
\end{proof}

\begin{remark}[Type-$A$ universality]
\label{rem:type-a-pattern}
For $\mathfrak{sl}_N$, the centraliser
$\operatorname{End}_{\mathfrak{sl}_N}(V\otimes V)$ is
two-dimensional (spanned by $\id$ and $P$), so the
$R$-matrix is linear in~$u$. Types $B$/$C$/$D$ and the
exceptionals have larger centralisers, necessitating
rational corrections.
\end{remark}

The orthogonal and symplectic $R$-matrices for $B_2$ and
$C_2$, and the four-projector $R$-matrix for $G_2$, are
recorded in Theorem~\ref{thm:B2-ordered-bar}
(equations~\eqref{eq:B2-R-matrix}--\eqref{eq:G2-R-matrix})
with full eigenvalue decompositions and proofs.

\begin{remark}[The universal classical $r$-matrix]
\label{rem:universal-classical}
For every simple $\mathfrak{g}$ and every
finite-dimensional $V$,
$r(u) = \Omega/u$ with
$\Omega = \sum_a I^a \otimes I_a
\in \mathfrak{g} \otimes \mathfrak{g}
\hookrightarrow
\operatorname{End}(V)^{\otimes 2}$.
The CYBE follows from the Jacobi identity:
$[\Omega_{12},\Omega_{13}] + [\Omega_{12},\Omega_{23}]
+ [\Omega_{13},\Omega_{23}] = 0$ is the Jacobi identity
in $\mathfrak{g}^{\otimes 3}$.
Quantum corrections: type~$A$ has none; types $B$/$C$/$D$
have one rational correction from the trace/symplectic
contraction; the exceptionals have multiple.
\end{remark}

\section{Rosetta stone: the $\mathsf{E}_1$ ordered bar complex across the landscape}
\label{sec:e1-rosetta-stone}
\index{Rosetta stone!ordered bar complex|textbf}
\index{ordered bar complex!landscape table|textbf}
\index{shadow depth!landscape classification|textbf}
% ================================================================

The ordered bar complex $\Barch(A)$ carries two layers of data:
the collision residue $r^{\mathrm{coll}}(z)$, which records the
chiral differential of the $E_1$-chiral coalgebra, and
the $R$-matrix $R(z)$, which governs the passage from
ordered to unordered configurations. The following table
collects this data for all standard families.

Every entry is an $\mathsf{E}_\infty$-chiral algebra
(a vertex algebra in the sense of Beilinson--Drinfeld:
the chiral product is defined on \emph{unordered}
configuration spaces with $\Sigma_n$-equivariant
factorization structure). The ordered bar complex
$\Barch(A)$ retains the ordering as auxiliary bookkeeping;
the descent to the symmetric bar $B^\Sigma(A)$ is the
$R$-matrix-twisted quotient of Proposition~\ref*{sec:r-matrix-descent}.
None of the families below is $\mathsf{E}_1$-chiral in the sense
of the $\mathsf{E}_1$ core of Volume~II\@: genuinely $\mathsf{E}_1$-chiral algebras
(nonlocal, ordered configurations as primitive data) are the
Etingof--Kazhdan quantum vertex algebras, which lie outside
the Beilinson--Drinfeld framework.

\begin{theorem}[$\mathsf{E}_1$ ordered bar landscape;
\ClaimStatusProvedHere]%
\label{thm:e1-ordered-bar-landscape}%
\index{ordered bar complex!landscape theorem|textbf}%
\index{collision residue!landscape|textbf}%
\index{R-matrix!landscape|textbf}%
Let\/ $A$ be a strongly admissible augmented associative chiral
algebra from the standard landscape. The ordered bar complex\/
$\Barch(A)$ has the following invariants, where\/ $r^{\mathrm{coll}}(z)$
denotes the collision residue
\textup{(}the OPE kernel after the $d\log$ bar propagator
absorbs one pole order\textup{)},
$R(z)$ is the spectral $R$-matrix on
$\mathrm{Conf}_2^{\mathrm{ord}}(\mathbb{C})$,
$r_{\max}$ is the shadow depth,
$A^!_{\mathrm{line}}$ is the ordered Koszul dual governing the
line category, and
``formal'' means the transferred $\Ainf$ operations
$m_k$ vanish on cohomology for all $k$ above the shadow depth:
\begin{center}
\small
\renewcommand{\arraystretch}{1.5}
\begin{tabular}{llcllcll}
\toprule
\textbf{Family}
 & \textbf{Class}
 & \makecell{\textbf{OPE} \\ \textbf{poles}}
 & \textbf{Collision residue}~$r^{\mathrm{coll}}(z)$
 & \textbf{$R$-matrix}~$R(z)$
 & $r_{\max}$
 & $A^!_{\mathrm{line}}$
 & \textbf{Formal?} \\
\midrule
$\mathcal{H}_k$
 & $\mathbf{G}$ & $2$
 & $k/z$
 & $e^{k\hbar/z}$
 & $2$
 & $Y(\mathfrak{u}(1))$\rlap{\textsuperscript{$\dagger$}}
 & yes \\[3pt]
$V_\Lambda$ \textup{(even)}
 & $\mathbf{G}$ & $2$
 & $\mathrm{diag}(\lambda_i)/z$
 & $\mathrm{diag}\bigl(e^{\lambda_i\hbar/z}\bigr)$
 & $2$
 & $Y(\mathfrak{u}(1)^{\oplus\operatorname{rk}\Lambda})$\rlap{\textsuperscript{$\dagger$}}
 & yes \\[3pt]
$\widehat{\mathfrak{g}}_k$\; ($\mathfrak{sl}_2$)
 & $\mathbf{L}$ & $2$
 & $\hbar\,\Omega/z$
 & $u\,I + \hbar\,P$
 & $3$
 & $Y_\hbar(\mathfrak{sl}_2)$
 & yes \\[3pt]
$\widehat{\mathfrak{g}}_k$\; ($\mathfrak{sl}_3$)
 & $\mathbf{L}$ & $2$
 & $\hbar\,\Omega/z$
 & Yang ($9{\times}9$)
 & $3$
 & $Y_\hbar(\mathfrak{sl}_3)$
 & yes \\[3pt]
$\widehat{\mathfrak{g}}_k$\; (general)
 & $\mathbf{L}$ & $2$
 & $\hbar\,\Omega/z$
 & Yang ($\dim\mathfrak{g}^2$)
 & $3$
 & $Y_\hbar(\mathfrak{g})$
 & yes \\[3pt]
$\beta\gamma$
 & $\mathbf{C}$ & $1$
 & $\Theta$ \textup{(nilpotent)}
 & $I + 2\pi i\,\Theta$
 & $4$
 & $bc^!$
 & yes \\[3pt]
$bc$
 & $\mathbf{C}$ & $1$
 & $\Theta_{bc}$ \textup{(nilpotent)}
 & $I + 2\pi i\,\Theta_{bc}$
 & $4$
 & $\beta\gamma^!$
 & yes \\[3pt]
$\mathrm{Vir}_c$
 & $\mathbf{M}$ & $4$
 & $\dfrac{c/2}{z^3} + \dfrac{2T}{z}$
 & \textup{[class $\mathbf{M}$]}
 & $\infty$
 & $Y^{\mathrm{dg}}(\mathrm{Vir}_c)$
 & no \\[6pt]
$\mathcal{W}_3$
 & $\mathbf{M}$ & $6$
 & degree $5$ in $1/z$
 & \textup{[class $\mathbf{M}$]}
 & $\infty$
 & $Y^{\mathrm{dg}}(\mathcal{W}_3)$
 & no \\[3pt]
$\mathcal{W}_N$
 & $\mathbf{M}$ & $2N$
 & degree $2N{-}1$ in $1/z$
 & \textup{[class $\mathbf{M}$]}
 & $\infty$
 & $Y^{\mathrm{dg}}(\mathcal{W}_N)$
 & no \\
\bottomrule
\end{tabular}
\end{center}

\smallskip
\noindent
\textsuperscript{$\dagger$}The Koszul dual algebra $\mathcal{H}_k^!
= \mathrm{Sym}^{\mathrm{ch}}(V^*)$ is categorically distinct from
$\mathcal{H}_{-k}$ but
both have $\kappa = -k$ and the line category
$\mathcal{C}_{\mathrm{line}} \simeq
\mathcal{H}_{-k}\text{-}\mathrm{mod}$
is equivalent. Similarly, $V_\Lambda^!$ and $V_{\Lambda^*}$ share
the same modular characteristic but differ as chiral algebras.
\end{theorem}

\begin{proof}
The table assembles data from three sources.

\emph{OPE pole orders.} The maximal OPE pole order on generators
is standard:
$J(z)J(w) \sim k/(z{-}w)^2$ for Heisenberg;
$J^a(z)J^b(w) \sim k\delta^{ab}/(z{-}w)^2 + f^{ab}{}_c J^c/(z{-}w)$
for affine Kac--Moody;
$\beta(z)\gamma(w) \sim 1/(z{-}w)$ and
$b(z)c(w) \sim 1/(z{-}w)$ for the free-field systems;
$T(z)T(w) \sim (c/2)/(z{-}w)^4 + 2T/(z{-}w)^2 + \partial T/(z{-}w)$
for Virasoro; and
$W_s(z)W_s(w)$ has poles up to order $2s$ for the
spin-$s$ generator of $\mathcal{W}_N$, giving
maximal pole $2N$ at the top spin.

\emph{Collision residues.}
The $d\log$ bar kernel absorbs one pole order
(the $d\log$ absorption principle; see Section~\ref{def:bar-differential-complete}):
the collision residue $r^{\mathrm{coll}}(z)$ from an OPE pole of
order~$d$ has pole of order~$d{-}1$. For Heisenberg:
double-pole $k/(z{-}w)^2$ gives $r^{\mathrm{coll}} = k/z$.
For affine $\widehat{\mathfrak{g}}_k$: the Casimir
$\Omega = \sum_a t^a \otimes t_a$ produces
$r^{\mathrm{coll}} = \hbar\,\Omega/z$.
For $\beta\gamma$/$bc$: the simple-pole OPE $1/(z{-}w)$
gives a constant collision residue
$r^{\mathrm{coll}} = \Theta$, where
$\Theta \in \mathrm{End}(V)$ is the nilpotent
off-diagonal matrix encoding the pairing
($\Theta^2 = 0$).
For Virasoro: quartic pole gives
$(c/2)/z^3 + 2T/z$ (the $\partial T/(z{-}w)$ term of the OPE is
absorbed completely, producing no pole in the collision residue).

\emph{$R$-matrices.}
The $R$-matrix is the monodromy of the connection
$\nabla = d - r^{\mathrm{coll}}(z)\,dz$ on
$\mathrm{Conf}_2^{\mathrm{ord}}(\mathbb{C})$.
For abelian algebras ($\mathcal{H}_k$, $V_\Lambda$):
the connection is scalar and
$R(z) = \exp(k\hbar/z)$
(Volume~II, Corollary~\ref*{cor:rosetta-heisenberg-projections}).
For affine $\widehat{\mathfrak{g}}_k$ at $\mathfrak{g}
= \mathfrak{sl}_2$: the $4 \times 4$ Yang $R$-matrix
$R(u) = uI + P$ (up to normalisation), where $P$ is
the permutation and $u = \hbar/z$ is the spectral parameter.
For $\beta\gamma$/$bc$: the nilpotent collision residue
$\Theta$ gives $R = \exp(2\pi i\,\Theta) = I + 2\pi i\,\Theta$
(the exponential terminates at the linear term since
$\Theta^2 = 0$).
For class~$\mathbf{M}$: the triple-pole collision residue
produces an $R$-matrix that is not expressible in closed form
as a finite product of elementary functions.

\emph{Shadow depth and formality.}
The shadow depth $r_{\max}$ is the highest degree at which the
transferred $\Ainf$ operation $m_k$ is non-vanishing on cohomology.
Class~$\mathbf{G}$: $m_{k \ge 3} = 0$ (abelian OPE, clean
Lagrangian intersection), so $r_{\max} = 2$ and the $A_\infty$-chiral
structure is formal.
Class~$\mathbf{L}$: $m_3 \ne 0$ from the Lie bracket but
$m_{k \ge 4} = 0$ on cohomology (Jacobi closes the obstruction),
so $r_{\max} = 3$ and the structure is formal.
Class~$\mathbf{C}$: $m_4 \ne 0$ from the contact term but the
tower terminates at $r_{\max} = 4$ by rank-one abelian
rigidity, so the structure is formal.
Class~$\mathbf{M}$: $m_k \ne 0$ for all $k \ge 3$ on
cohomology (the quartic pole manufactures a propagating
obstruction), so $r_{\max} = \infty$ and the $A_\infty$-chiral
structure is \emph{non-formal}.

\emph{Ordered Koszul duals.}
For affine $\widehat{\mathfrak{g}}_k$: the ordered Koszul dual is
the Yangian $Y_\hbar(\mathfrak{g})$
(Volume~II, Theorem~\ref*{thm:Koszul_dual_Yangian}).
For Virasoro and $\mathcal{W}_N$: the ordered Koszul dual is the
dg-shifted Yangian $Y^{\mathrm{dg}}$, a dg algebra with
infinitely many cohomological generators reflecting
non-formality (the $\mathsf{E}_1$ core of Volume~II).
For $\beta\gamma$/$bc$: the Koszul duality exchanges
statistics ($bc^! = \beta\gamma$, $\beta\gamma^! = bc$),
so the line-side algebra is the dual free-field system.

All entries are chirally Koszul (bar cohomology concentrated),
as established by PBW universality (Volume~I) and the
one-loop exactness criterion (Volume~II,
Theorem~\ref*{thm:one-loop-koszul}).
\end{proof}

\begin{remark}[The $\mathbf{G}/\mathbf{L}/\mathbf{C}/\mathbf{M}$
classification]%
\label{rem:glcm-commentary}%
\index{shadow depth!G/L/C/M classification|textbf}%
The four classes are determined by the OPE pole structure
on generators:
\begin{enumerate}[label=\textup{(\roman*)}]
\item Class~$\mathbf{G}$ \textup{(}Gaussian\textup{)}:
 double-pole OPE with scalar residues.
 The bar complex is cogenerated in degree~$2$,
 the $R$-matrix is scalar-valued (abelian braiding),
 and the Lagrangian self-intersection
 $\mathfrak{S}_b = L \times_X L$ is clean.
 Archetype: $\mathcal{H}_k$.

\item Class~$\mathbf{L}$ \textup{(}Lie\textup{)}:
 double-pole OPE with non-scalar
 \textup{(}Lie-bracket\textup{)} residues.
 The Lie bracket creates a transverse intersection defect
 at codimension~$1$, but the Jacobi identity kills all
 higher defects. The $R$-matrix is the Yang $R$-matrix:
 matrix-valued, satisfying the Yang--Baxter equation
 with spectral parameter.
 Archetype: $\widehat{\mathfrak{g}}_k$.

\item Class~$\mathbf{C}$ \textup{(}Contact\textup{)}:
 simple-pole OPE on generators.
 The composite-field OPE reaches pole order~$3$, producing a
 quartic contact invariant. Rank-one abelian rigidity
 forces the quintic obstruction to vanish.
 The $R$-matrix is a unipotent perturbation of the identity
 ($\Theta^2 = 0$ forces $R = I + 2\pi i\,\Theta$).
 Archetype: $\beta\gamma$.

\item Class~$\mathbf{M}$ \textup{(}Modular\textup{)}:
 quartic or higher poles on generators.
 The $R$-matrix is not expressible in closed form.
 The shadow obstruction tower does not terminate: the composite-field
 exchange generates new obstructions at every degree.
 Non-formality is visible already at genus~$0$: infinitely
 many $A_\infty$ operations $m_n \neq 0$ on
 $\overline{\mathcal{M}}_{0,n+1}$, driven by the
 self-referential OPE $T_{(1)}T = 2T$. At genus~$g \geq 1$,
 this infinite cascade interacts with the curvature
 $d_{\mathrm{fib}}^2 = \kappa(\cA) \cdot \omega_g$ to produce
 the genuinely modular content of the tower.
 Archetype: $\mathrm{Vir}_c$.
\end{enumerate}
\end{remark}

\begin{remark}[Formality, non-formality, and the
$\mathsf{E}_1/\mathsf{E}_\infty$ distinction]%
\label{rem:formality-e1-e-infty}%
\index{formality!$A_\infty$-chiral structure}%
\index{E-infinity@$\mathsf{E}_\infty$-chiral!all standard families}%
\index{E-one@$\mathsf{E}_1$-chiral!genuinely nonlocal}%
A critical conceptual point: non-formality \emph{does not}
imply $\mathsf{E}_1$. All ten families in
Theorem~\textup{\ref{thm:e1-ordered-bar-landscape}} are
$\mathsf{E}_\infty$-chiral \textup{(}local vertex algebras with
$\Sigma_n$-equivariant factorization structure\textup{)}, including
the non-formal families $\mathrm{Vir}_c$, $\mathcal{W}_3$, and
$\mathcal{W}_N$. Having OPE poles does not break
$\mathsf{E}_\infty$; having non-vanishing higher $\Ainf$
operations does not break $\mathsf{E}_\infty$. The poles
are singularities of a \emph{local} product, and the
non-formality measures structural complexity \emph{within}
the $\mathsf{E}_\infty$ world.

The three-tier picture is:
\begin{enumerate}[label=\textup{(\alph*)}]
\item \emph{Pole-free commutative}:
 $R(z) = \tau$ (the Koszul-signed flip).
 This is the subclass where the chiral product extends
 across the diagonal without singularities.

\item \emph{Vertex algebras with poles}
 \textup{(}all families in the table above\textup{)}:
 $R(z) \ne \tau$, but $R(z)$ is \emph{derived} from the
 local OPE via analytic continuation.
 Both (a) and (b) are $\mathsf{E}_\infty$-chiral.

\item \emph{Genuinely $\mathsf{E}_1$-chiral}
 \textup{(}Etingof--Kazhdan quantum vertex algebras,
 Yangians as independent algebras, quantum groups\textup{)}:
 $R(z) \ne \tau$ and $R(z)$ is \emph{independent input},
 not derivable from a local symmetric structure.
 Only (c) is genuinely $\mathsf{E}_1$.
\end{enumerate}
Tiers~(a) and~(b) are both $\mathsf{E}_\infty$.
The ordered bar complex is a computational tool for
$\mathsf{E}_\infty$-chiral algebras; it does not make them
$\mathsf{E}_1$.
\end{remark}

\begin{remark}[Pole escalation: DS reduction
transports $\mathbf{L} \to \mathbf{M}$]%
\label{rem:pole-escalation}%
\index{Drinfeld--Sokolov reduction!pole escalation}%
\index{pole escalation|textbf}%
Principal Drinfeld--Sokolov reduction sends
$V_k(\mathfrak{g}) \mapsto \mathcal{W}_N(\mathfrak{g})$.
On the ordered bar complex, this transports class~$\mathbf{L}$
to class~$\mathbf{M}$: the affine double-pole OPE
$J^a(z)J^b(w) \sim k\delta^{ab}/(z{-}w)^2 + \cdots$
becomes the Virasoro quartic-pole OPE
$T(z)T(w) \sim (c/2)/(z{-}w)^4 + \cdots$ via the Sugawara
construction. The shadow depth escalates from $r_{\max} = 3$
(finite, formal) to $r_{\max} = \infty$ (infinite, non-formal).
The Yangian $Y_\hbar(\mathfrak{g})$ (a classical algebra
concentrated in degree~$0$) becomes the dg-shifted Yangian
$Y^{\mathrm{dg}}(\mathrm{Vir}_c)$ (a dg algebra with
infinitely many cohomological generators).

This pole escalation is the ordered-bar avatar of the
\emph{single-line dichotomy} of Volume~I: double poles
produce a linear (tree-level) bar differential, while
quartic poles produce a quadratic (one-loop) bar differential
whose iterations generate the infinite tower.
The mechanism is the Sugawara pole:
$T = \tfrac{1}{2(k+h^\vee)} \sum_a {:}J^a J_a{:}$
is a normally ordered composite, and the OPE of composites
escalates pole orders. Theorem~\ref{thm:ordered-associative-ds-principal}
ensures that this pole escalation is compatible with
the ordered bar--cobar equivalence.
\end{remark}

\begin{remark}[$\cW_{1+\infty}$ as a class~$\mathbf{M}$ chiral quantum group]
\label{rem:w-infty-chiral-qg}
\index{$\cW_{1+\infty}$!chiral quantum group structure}
\index{CoHA!and $\cW_{1+\infty}$}
The $\cW$-algebra $\cW_{1+\infty}[\Psi]$
(Schiffmann--Vasserot, Maulik--Okounkov) is a class~$\mathbf{M}$
algebra whose CoHA origin provides all three structures of the
chiral bialgebra equivalence
(Theorem~\ref{thm:chiral-qg-equiv}): the $R$-matrix from the
stable-envelope construction, the $A_\infty$-structure from the
CoHA multiplication (which is $\Eone$-associative: the preferred
direction on short exact sequences gives the ordering), and the
chiral coproduct from the Hall comultiplication. The shadow
obstruction tower is infinite (class~$\mathbf{M}$: the irregular
KZ connection has Stokes data), and the ordered chiral homology
carries richer data than the symmetric shadow.
This is the bridge between Vol~I (chiral quantum group
equivalence) and Vol~III (CoHA as the $\Eone$-sector of the
CY engine).
\end{remark}

\begin{remark}[Genuinely $\Eone$-chiral algebras: ordered $\neq$ symmetric]
\label{rem:ordered-neq-symmetric}
\index{ordered chiral homology!genuinely $\Eone$-chiral}
\index{Yangian!ordered $\neq$ symmetric}
Every family in the Rosetta stone table above is
$\Einf$-chiral: the ordered bar complex carries a nontrivial
$R$-matrix, but this $R$-matrix is derived from the
$\Einf$-chiral OPE (tier~(b) of
Definition~\ref{def:three-tier-r-matrix}). The averaging map
$\av \colon \barB^{\mathrm{ord}} \to \barB^\Sigma$ is a
quasi-isomorphism for these families
(Proposition~\ref{prop:symmetric-descent}(i)), and ordered
chiral homology agrees with Beilinson--Drinfeld chiral
homology.

For genuinely $\Eone$-chiral algebras (tier~(c): the Yangian
$Y_\hbar(\fg)$, Etingof--Kazhdan quantum vertex algebras),
this agreement fails. The $R$-matrix is independent input, the
factorisation $\cD$-module $\cF_n^{\mathrm{ord}}$ is not
$\Sigma_n$-equivariant, and $\ker(\av) \neq 0$ at degree
$n \ge 3$: the Drinfeld associator
$\Phi_{\mathrm{KZ}} \in \exp(\hat{\mathfrak{t}}_3)$ contributes
$\mathrm{GRT}_1$-dependent data visible in the ordered theory
but killed by the $\Sigma_n$-coinvariant projection
(Proposition~\ref{prop:symmetric-descent}(ii)). The
Yangian is the first example in which the ordered theory is
strictly richer than the symmetric one.
\end{remark}


% ================================================================
% CHIRAL BIALGEBRA EQUIVALENCE
% ================================================================

\section{The chiral bialgebra equivalence on the Koszul locus}
\label{sec:chiral-quantum-group-equiv}
\index{chiral bialgebra!equivalence theorem|textbf}
\index{chiral quantum group!informal shorthand for chiral bialgebra}
\index{R-matrix!vertex!as projection of ordered bar}
\index{A-infinity@$A_\infty$!chiral endomorphism operad}
\index{coproduct!chiral!from ordered bar}

The ordered bar complex $\barB^{\mathrm{ord}}(\cA) = T^c(s^{-1}\bar\cA)$
is a single object from which three structures are recovered: the
vertex $R$-matrix $S(z)$ (braiding data from the degree-$2$ collision
residue), the chiral $A_\infty$-structure maps
$m_k^{\mathrm{ch}}$ (higher associativity from boundary strata of
associahedra), and the chiral coproduct $\Delta^{\mathrm{ch}}$
(coalgebra data from deconcatenation). The three structures
determine each other.

\begin{definition}[$\Eone$-chiral algebra with compatible chiral coproduct]
\label{def:chiral-coproduct}
An $\Eone$-chiral algebra with compatible chiral coproduct
is a triple $(\cA, Y, \Delta^{\mathrm{ch}})$ where
$(\cA, Y)$ is a chiral algebra and
\[
  \Delta^{\mathrm{ch}} \colon
  \cA \longrightarrow
  (\cA \mathbin{\hat\otimes} \cA)((z))
\]
is a chiral coproduct satisfying:
\begin{enumerate}[label=\textup{(\roman*)}]
\item \textup{(Coassociativity up to associator.)}
  The two iterated coproducts agree up to conjugation by
  the Drinfeld associator~$\Phi$:
  \[
    (\Delta^{\mathrm{ch}} \otimes \id)
    \circ \Delta^{\mathrm{ch}}
    \;=\;
    \Phi \cdot
    (\id \otimes \Delta^{\mathrm{ch}})
    \circ \Delta^{\mathrm{ch}}
    \cdot \Phi^{-1}.
  \]
\item \textup{(OPE compatibility.)}
  The coproduct intertwines the vertex operation,
  as in equation~\eqref{eq:ope-compat}.
\item \textup{(Counit.)}
  There exists a counit
  $\varepsilon \colon \cA \to \CC$ with
  $(\varepsilon \otimes \id) \circ \Delta^{\mathrm{ch}}
  = \id = (\id \otimes \varepsilon) \circ
  \Delta^{\mathrm{ch}}$.
\end{enumerate}
The associated vertex $R$-matrix is recovered from the coproduct by
equation~\eqref{eq:r-from-coprod-proof}.
\end{definition}

\begin{theorem}[Chiral bialgebra equivalence on the Koszul locus]
\label{thm:chiral-qg-equiv}
\ClaimStatusProvedHere
\index{chiral bialgebra!three equivalent descriptions}
The following three structures on an $\Eone$-chiral
algebra~$\cA$ on the Koszul locus determine each other,
up to the choice of a Drinfeld associator
$\Phi \in \exp(\hat{\mathfrak{t}}_3)$:
\begin{enumerate}[label=\textup{(\Roman*)}]
\item $\Eone$-chiral algebra with vertex $R$-matrix $S(z)$
  satisfying $S$-locality, the quantum Yang--Baxter equation,
  unitarity, shift condition, and hexagon axiom
  \textup{(}the Etingof--Kazhdan axioms\textup{)};
\item $A_\infty$-algebra in the chiral endomorphism operad
  $\End^{\mathrm{ch}}_\cA$, with structure maps
  $m_k^{\mathrm{ch}} \colon \cA^{\otimes k} \to
  \cA((\lambda_1))\cdots((\lambda_{k-1}))$
  satisfying the Stasheff identities;
\item $\Eone$-chiral algebra with compatible chiral coproduct
  $\Delta^{\mathrm{ch}} \colon
  \cA \to (\cA \mathbin{\hat\otimes} \cA)((z))$,
  coassociative up to conjugation by~$\Phi$.
\end{enumerate}
The equivalences form a triangle:
\begin{center}
\begin{tikzcd}[column sep=3em, row sep=2.5em]
  \textup{(I)} \ar[rr, shift left=1, "\alpha"]
  \ar[dr, shift left=1, "\gamma^{-1}"']
  && \textup{(II)} \ar[ll, shift left=1, "\alpha^{-1}"]
  \ar[dl, shift left=1, "\beta"] \\
  & \textup{(III)}
  \ar[ul, shift left=1, "\gamma"]
  \ar[ur, shift left=1, "\beta^{-1}"']
\end{tikzcd}
\end{center}
where $\alpha$ restricts the OPE to the real
Stasheff associahedra $K_n \hookrightarrow
\overline{\operatorname{FM}}_n^{\mathrm{ord}}(\bC)$,
$\beta$ applies the cobar functor to the deconcatenation
coproduct on $\barB^{\mathrm{ord}}(\cA)$, and
$\gamma$ extracts the $R$-matrix as
$S(z) = \sigma \circ \Delta^{\mathrm{ch}}
\circ (\Delta^{\mathrm{ch},\mathrm{op}})^{-1}$
\textup{(}the chiral Drinfeld formula\textup{)}.
\end{theorem}

\begin{proof}
\textbf{(I) $\to$ (II).}
The Stasheff associahedron $K_n$ is isomorphic to the real
FM compactification $\overline{\operatorname{FM}}_n^{\mathrm{ord}}(\bR)$
(Sinha). The chiral $A_\infty$-maps are the integrals
\[
m_k^{\mathrm{ch}}(a_1, \ldots, a_k)
\;=\;
\int_{K_k}
Y(a_1, z_1) \cdots Y(a_k, z_k)\, \omega_k,
\]
where $\omega_k$ is the canonical volume form induced by
the bar propagator. The Stasheff identity follows from
Stokes' theorem on $K_{n+1}$: the codimension-$1$ boundary
decomposes into facets $K_{r+1} \times K_{s+1}$ recording
consecutive collisions, and $\int_{K_{n+1}} d\omega = 0$
yields
$\sum (-1)^\epsilon\, m_i^{\mathrm{ch}} \circ_p
m_j^{\mathrm{ch}} = 0$.
The inverse $\alpha^{-1}$ recovers $S(z)$ as the degree-$2$
KZ holonomy.

\textbf{(II) $\to$ (III).}
The $A_\infty$-maps determine the bar differential
$d_{\barB^{\mathrm{ord}}} = \sum_{k \ge 1}
\hat{m}_k^{\mathrm{ch}}$ (coderivation extension).
The deconcatenation coproduct $\Delta$ on
$T^c(s^{-1}\bar\cA)$ is a dg coalgebra map; applying the
cobar functor $\Omega$ to $\Delta$ produces the chiral
coproduct $\Delta^{\mathrm{ch}}$ on $\cA$ via
$\cA \simeq \Omega(\barB^{\mathrm{ord}}(\cA))$
(Theorem~\ref{thm:bar-cobar-isomorphism-main} on the Koszul locus).
Strict coassociativity of deconcatenation becomes
coassociativity-up-to-$\Phi$ after passage through the
formality quasi-isomorphism underlying~$\Omega$.
The inverse $\beta^{-1}$ reads off
$m_k^{\mathrm{ch}}$ as the tensor-degree-$k$
component of $d_{\barB^{\mathrm{ord}}}$.

\textbf{(III) $\to$ (I).}
The chiral Drinfeld formula
$S(z) = \sigma \circ \Delta^{\mathrm{ch}}
\circ (\Delta^{\mathrm{ch},\mathrm{op}})^{-1}$
defines the vertex $R$-matrix. The QYBE follows from
coassociativity-up-to-$\Phi$ (via the pentagon for~$\Phi$).
\end{proof}

\begin{remark}[Bialgebra, not Hopf: the antipode does not lift]
\label{rem:chiral-bialgebra-not-hopf}
The equivalence of
Theorem~\textup{\ref{thm:chiral-qg-equiv}} produces three compatible
descriptions of a \emph{chiral bialgebra}: $R$-matrix, chiral
$A_\infty$-algebra, and chiral coproduct with counit. It does not
produce an antipode. The classical formula $S(T(u)) = T(u)^{-1}$
does not lift to a vertex-algebraic map on the spectral parameter;
the precise content is
Proposition~\textup{\ref{prop:w-infty-antipode-obstruction}}, which
exhibits two independent structural obstructions. The first is an
OPE-locality obstruction: the quartic pole of $S(T)$ with itself
agrees with the Virasoro value $c/2$ only at the two isolated points
$\Psi \in \{1, 2\}$. The second is a Hopf-axiom obstruction: the
residual of $m_z \circ (S \otimes \mathrm{id}) \circ \Delta_z$ on $T$
contains a term $z \cdot J$ that survives at every value of $\Psi$,
reflecting the $z$-dependence of the spectral coproduct. The two
obstructions vanish on disjoint loci (OPE on $\Psi \in \{1,2\}$, Hopf
nowhere), so they are not two symptoms of a single defect but two
genuinely independent structural failures. We therefore reserve the
name ``chiral bialgebra equivalence'' for this statement; the phrase
``chiral quantum group'' continues to denote the underlying
structural package, with the understanding that no Hopf antipode
exists at the vertex-algebraic level.
\end{remark}

\begin{remark}[Associator class: cohomological vs.\ cochain canonicality]
\label{rem:chiral-qg-grt1-torsor}
The choice of Drinfeld associator
$\Phi \in \exp(\hat{\mathfrak{t}}_3)$ is not canonical: the group
$\mathrm{GRT}_1(\bQ)$ acts transitively on the set of admissible
associators, so
Theorem~\textup{\ref{thm:chiral-qg-equiv}} produces a
$\mathrm{GRT}_1(\bQ)$-torsor of equivalences rather than a single
canonical one. At the level of the derived centre and of cohomology,
this torsor action is trivial: the cohomological image of the
equivalence is associator-independent (verified for $\mathfrak{sl}_2$
and extended structurally to all simple $\fg$). At the cochain level,
different associator representatives produce coproducts related by a
gauge transformation whose cohomology class vanishes; the
equivalences are homotopic but not equal. The statement of
Theorem~\textup{\ref{thm:chiral-qg-equiv}} is therefore canonical at
$\mathrm{H}^0$ and non-canonical up to $\mathrm{GRT}_1(\bQ)$ at the
cochain level. Throughout the chapter we read the equivalence at
$\mathrm{H}^0$ unless a specific associator class is fixed.
\end{remark}

\begin{corollary}[The ordered bar encodes all three structures]
\label{cor:bar-encodes-all}
\ClaimStatusProvedHere
\index{ordered bar complex!as universal datum}
The ordered bar complex $\barB^{\mathrm{ord}}(\cA) =
T^c(s^{-1}\bar\cA)$ is the universal datum: the vertex
$R$-matrix is the degree-$2$ collision residue
$r(z) = \Res^{\mathrm{coll}}_{0,2}(\Theta_\cA^{\Eone})$,
the $A_\infty$-maps are the bar differential components
$d_{\barB}|_{\text{degree } k}$, and the chiral coproduct is
the cobar-dual of the deconcatenation. No additional input
beyond $\barB^{\mathrm{ord}}(\cA)$ is needed.
\end{corollary}

\begin{proof}
Theorem~\ref{thm:chiral-qg-equiv} provides three readouts of the same
ordered bar coalgebra.
Arrow $\alpha$ identifies the degree-$2$ collision residue with the
vertex $R$-matrix.
Arrow $\beta$ identifies the higher components of the bar differential
with the transferred $A_\infty$ operations.
Arrow $\gamma$ identifies the deconcatenation coproduct, after cobar
dualization, with the chiral coproduct.
Since these are all functorial projections of
$\barB^{\mathrm{ord}}(\cA)$, no extra datum beyond the ordered bar
complex is required.
\end{proof}

\begin{remark}[Key formulas of the equivalence]
\label{rem:chiral-qg-key-formulas}
The three equivalences involve the following explicit formulas.
The OPE compatibility axiom:
\begin{equation}\label{eq:ope-compat}
  \Delta^{\mathrm{ch}}(Y(a,z)\, b)
  \;=\;
  Y^{(2)}(\Delta^{\mathrm{ch}}(a), z)
  \cdot \Delta^{\mathrm{ch}}(b),
\end{equation}
where $Y^{(2)}$ is the tensorwise vertex operation.
The coproduct from the bar complex:
\begin{equation}\label{eq:coprod-from-bar}
  \Delta^{\mathrm{ch}} \colon
  \cA \simeq \Omega(\barB^{\mathrm{ord}}(\cA))
  \xrightarrow{\Omega(\Delta)}
  \Omega(\barB^{\mathrm{ord}}(\cA) \otimes
  \barB^{\mathrm{ord}}(\cA))
  \to (\cA \mathbin{\hat{\otimes}} \cA)((z)).
\end{equation}
The $R$-matrix from the coproduct (the chiral Drinfeld formula):
\begin{equation}\label{eq:r-from-coprod-proof}
  S(z)
  \;:=\; \sigma \circ \Delta^{\mathrm{ch}}
  \circ (\Delta^{\mathrm{ch},\mathrm{op}})^{-1}.
\end{equation}
\end{remark}

% ================================================================
\noindent\emph{QYBE, unitarity, shift, hexagon.}
The quantum Yang--Baxter equation on
$V_{\mathrm{EK}}^{\otimes 3}$ reduces on $V^{\otimes 3}$
to~\eqref{eq:glN-ybe}, as in the classical Yangian
case; the extension to
$V_{\mathrm{EK}}$ follows from the universal property of
$\mathcal{R}(z)$.
Unitarity $S_{12}(z) \cdot (\sigma \circ S_{12}(-z)
\circ \sigma) = \id$ follows from Drinfeld--Jimbo unitarity.
The shift condition
\begin{equation}\label{eq:ek-shift-cond}
  [D \otimes \id + \id \otimes D,\; S(z)]
  \;=\; \partial_z\, S(z)
\end{equation}
follows from the translation-covariance of the Drinfeld
coproduct. The hexagon follows from the pentagon relation
for $\Phi$ (from the Stasheff boundary of
$\overline{\FM}_4^{\mathrm{ord}}(\CC) = K_4$) combined
with the QYBE for $S$, by the standard argument of
Drinfeld~\cite{EK96} lifted to the vertex-algebraic setting.

\medskip
\noindent
\textbf{Part B: the ordered chiral center of the EK quantum
VOA.}
We compute the ordered chiral Hochschild cohomology
$\HH^{*,\mathrm{ord}}_{\mathrm{ch}}
(V_{\mathrm{EK}},\,V_{\mathrm{EK}})$
at low degrees on the formal disk $D$, paralleling the
classical Yangian computation but now
using the \emph{full} vertex $R$-matrix $S(z)$ rather than
the Yang $R$-matrix on $V \otimes V$ alone.

\smallskip
\noindent\emph{Degree~$2$: the KZ connection with the full
vertex $R$-matrix.}
At degree~$2$, the KZ connection on
$\Conf_2^{\mathrm{ord}}(\CC)$ with coefficients in
$\End(V_{\mathrm{EK}}^{\otimes 2})$ is
\begin{equation}\label{eq:ek-kz-full}
  \nabla_{\mathrm{KZ}}^S
  = d + \frac{\hbar\,\Omega}{z}\,dz
  + \frac{\hbar^2}{z^2}\,
  \Bigl(\Omega^{(2)} + \tfrac{1}{2}\,\Omega^2\Bigr)\,dz
  + O(z^{-3}),
\end{equation}
where $z = z_1 - z_2$. The connection~\eqref{eq:ek-kz-full}
is obtained from $S(z)$ via the general formula
$\nabla = d + (\partial_z \log S(z))\,dz$ (logarithmic
derivative of the vertex $R$-matrix).

For the classical $V_k(\mathfrak{sl}_2)$ (trivial braiding
$S(z) = \id$), the flat sections are
$\Phi(z) = z^{\Omega/(k+2)}$ with algebraic monodromy.
For the EK quantum VOA ($S(z) \neq \id$), the flat sections
acquire quantum corrections:
\begin{equation}\label{eq:ek-flat-sections}
  \Phi^S(z)
  = z^{\hbar\,\Omega}\,
  \Bigl(I + \frac{\hbar^2\,\Omega^{(2)}}{z}
  + O(z^{-2})\Bigr).
\end{equation}

\emph{The kernel $\ker(\av)$.}
The ordered chiral centre computation parallels the
classical Yangian case; the result is identical on
cohomology. On $V \otimes V$:
$\ker(\av_2)\big|_{V \otimes V}
\cong \bigwedge^2 V = \CC$,
the same one-dimensional space as in that classical
computation.
The quantum corrections from $S(z)$ shift the subleading
terms of the flat sections but do not alter the
$\ker(\av_2)$ dimension. The scalar shadow surviving
$\av_2$ is $\kappa = 3(k+2)/4$ (equation~\eqref{eq:sl2-kappa}),
identical to the classical value. At degree~$3$, the kernel
$\ker(\av_3) \cong \Lambda^3(\fg) = \CC$ carries the
quantum-corrected Drinfeld associator
\begin{equation}\label{eq:ek-ker-av3}
  \Phi_{\mathrm{KZ}}^S - 1
  \;=\; \hbar^2\,\zeta(2)\,[\Omega_{12}, \Omega_{23}]
  + O(\hbar^3),
\end{equation}
whose leading term agrees with the classical KZ associator.

\medskip
\noindent
\textbf{Part C: $S$-locality as an $\Ainf$ relation.}
The braided VOA axioms of the EK quantum VOA are equivalent
to the chiral $\Ainf$ relations for $V_{\mathrm{EK}}$
through the chiral bialgebra equivalence
(Theorem~\ref{thm:chiral-qg-equiv},
(I)$\Leftrightarrow$(II)).
In the ordered bar complex
$\Barord(V_{\mathrm{EK}}) = T^c(s^{-1}\overline{V}_{\mathrm{EK}})$,
the $S$-locality axiom~\eqref{eq:ek-s-locality} governs the
reordering of two insertions: the braid group
$\pi_1(\Conf_2^{\mathrm{ord}}(\CC)) = \ZZ$ acts on degree-$2$
bar chains by
$\sigma_1 \cdot [a | b] = S(z) \cdot [b | a]$.
The $S$-locality axiom is the $n = 2$ Stasheff relation,
the hexagon is the $n = 3$ relation, and the QYBE is the
$n = 4$ relation (from the pentagon for $\Phi$ on
$K_4 = \overline{\FM}_4^{\mathrm{ord}}(\RR)$).

\medskip
\noindent
\textbf{Summary.}
\noindent
The EK quantum VOA $(V_{\mathrm{EK}},\, Y,\, S(z))$
provides the first concrete instance of all three structures
of the chiral bialgebra equivalence
(Theorem~\ref{thm:chiral-qg-equiv}) exhibited
simultaneously on a single algebra with all axioms verified.
The ordered chiral center carries the full quantum $R$-matrix
$S(z)$ (not just the classical $r$-matrix
$r(z) = \hbar\,\Omega/z$), and the kernel
$\ker(\av_2) \neq 0$ demonstrates that the ordered theory
is strictly richer than the symmetric Beilinson--Drinfeld
chiral homology even at the quantum level.
\end{example}

\begin{example}[Heisenberg: trivial case]
\label{ex:heis-coproduct}
For the Heisenberg algebra $\cH_k$ (class $G$), all three
structures are trivial:
$S(z) = \id$ (the braiding is symmetric);
$m_k^{\mathrm{ch}} = 0$ for $k \geq 3$ (the algebra is
strictly associative in the chiral sense);
$\Delta^{\mathrm{ch}}(J) = J \otimes 1 + 1 \otimes J$ is
cocommutative and strictly coassociative ($\Phi$ acts
trivially). The ordered and symmetric data coincide.
\end{example}

\begin{theorem}[$\cW_{1+\infty}[\Psi]$ as a chiral quantum
group]
\label{thm:w-infty-chiral-qg}
\ClaimStatusProvedHere
The algebra $\cW_{1+\infty}[\Psi]$ carries a chiral quantum
group datum in the sense of
Theorem~\textup{\ref{thm:chiral-qg-equiv}}: an $R$-matrix,
a chiral coproduct, and an $\Ainf$ structure, all explicit
and originating from the cohomological Hall algebra of the
Jordan quiver. The precise content is as follows.
\begin{enumerate}[label=\textup{(\Roman*)}]
\item \textup{($R$-matrix.)}
  The Maulik--Okounkov $R$-matrix~\cite{MO19} of quantum
  toroidal $\mathfrak{gl}_1$, restricted to the CoHA, is
  \emph{scalar}:
  \begin{equation}\label{eq:gl1-scalar-r}
    R(z) = g(z)
    = \frac{(z - h_1)(z - h_2)(z - h_3)}
    {(z + h_1)(z + h_2)(z + h_3)},
    \qquad h_1 + h_2 + h_3 = 0,
  \end{equation}
  where $h_1, h_2, h_3$ are the equivariant parameters of
  $\CC^3$ subject to the Calabi--Yau constraint. The
  classical limit is
  % AP126: r-matrix with level prefix Psi; Psi=0 -> r=0. Verified.
  $r(z) = \Psi/z$ \textup{(}class~$G$, level prefix
  $\Psi$\textup{)}.
\item \textup{(Chiral coproduct at all spins.)}
  The Drinfeld coproduct on $Y(\widehat{\mathfrak{gl}}_1)$
  gives an explicit spectral coproduct
  $\Delta_z \colon \cW_{1+\infty}[\Psi] \to
  (\cW_{1+\infty}[\Psi] \mathbin{\hat{\otimes}}
  \cW_{1+\infty}[\Psi])((z))$
  satisfying strict coassociativity,
  the spectral counit, and OPE compatibility at all spins.
  The transfer matrix
  $T(u) = 1 + \sum_{n \geq 1} \psi_n\, u^{-n}$
  is scalar \textup{(}$\mathfrak{gl}_1$ has
  rank~$1$\textup{)}, and
  \begin{equation}\label{eq:gl1-drinfeld-coprod}
    \Delta_z(T(u))
    = T(u) \otimes T(u - z).
  \end{equation}
  The general spin-$n$ formula is
  \begin{equation}\label{eq:gl1-coprod-general}
    \Delta_z(\psi_n)
    = \psi_n \otimes 1
    + \sum_{k=0}^{n-1}\; \sum_{m=1}^{n-k}
    \binom{n-k-1}{m-1}\,
    z^{n-k-m}\,
    \psi_k \otimes \psi_m,
  \end{equation}
  where $\psi_0 = 1$.
\item \textup{($\Ainf$ structure.)}
  The chiral $\Ainf$-operations $m_k^{\mathrm{ch}}$ are
  extracted from the ordered bar complex
  $\Barord(\cW_{1+\infty})
  = T^c(s^{-1}\overline{\cW}_{1+\infty})$ via the functor
  $\alpha$ of Theorem~\textup{\ref{thm:chiral-qg-equiv}},
  with input the $R$-matrix~\eqref{eq:gl1-scalar-r}.
\end{enumerate}
The three structures determine each other by
Theorem~\textup{\ref{thm:chiral-qg-equiv}}. The OPE
compatibility axiom~\eqref{eq:ope-compat} is satisfied at
all spins simultaneously by two independent arguments:
the coderivation property on the Koszul-locus bar complex,
and the JKL vertex bialgebra theorem on the CoHA.
\end{theorem}

\begin{proof}
The argument proceeds in seven steps.

\medskip
\noindent
\textbf{Step~1}
(CoHA identification;
Schiffmann--Vasserot~\cite{SV13},
Schiffmann--Mellit--Minets--Vasserot~\cite{SMMV23}).
The critical CoHA of the Jordan quiver is
\[
  \cH_{\mathrm{Jor}}
  = \bigoplus_{n \geq 0}
  H^{GL_n}_{\mathrm{BM},*}
  (\cM_n^{\mathrm{nil}},\, \varphi_W),
  \qquad
  W = \mathrm{Tr}(x[y,z]),
\]
where $\cM_n^{\mathrm{nil}}$ is the nilpotent locus in
the representation variety of the Jordan quiver and
$\varphi_W$ denotes the vanishing-cycle twist.
Schiffmann--Vasserot~\cite{SV13} proved
$\cH_{\mathrm{Jor}} \cong Y^+(\widehat{\mathfrak{gl}}_1)$
by matching generators and relations; the critical
refinement is~\cite{SMMV23}.
Kontsevich--Soibelman~\cite{KS11} established the general
CoHA framework (see also Joyce~\cite{Joyce18}).

\medskip
\noindent
\textbf{Step~2}
(Bialgebra and Drinfeld double; Yang--Zhao~\cite{YZ18a}).
The CoHA carries a bialgebra coproduct
$\Delta_{\mathrm{CoHA}} \colon
\cH_{\mathrm{Jor}} \to
\cH_{\mathrm{Jor}} \otimes \cH_{\mathrm{Jor}}$
from restriction to Levi subalgebras
$GL_{n_1} \times GL_{n_2} \hookrightarrow GL_n$.
Yang--Zhao proved that the Drinfeld double is the full
affine Yangian: $D(\cH_{\mathrm{Jor}}) \cong
Y(\widehat{\mathfrak{gl}}_1)$.

\medskip
\noindent
\textbf{Step~3}
(Yangian--$\cW$-algebra identification;
Proch\'azka--Rap\v{c}\'ak~\cite{PR19}).
The quantum Miura transform gives an isomorphism of
filtered associative algebras
\[
  Y(\widehat{\mathfrak{gl}}_1)
  \;\cong\;
  U(\cW_{1+\infty}[\Psi]),
\]
intertwining the Yangian filtration with the conformal
filtration. The $\psi$-generators map to
$\cW_{1+\infty}$ fields via a triangular change
of basis: $\psi_1 = J$,
$\psi_n = W_n + P_n(J, W_2, \ldots, W_{n-1})$
for $n \geq 2$, where $W_2 = T$ (stress tensor),
$W_3 = W$ (spin-$3$ current), and $P_n$ is a
nonlinear polynomial in lower-spin fields with
coefficients depending on~$\Psi$.

\medskip
\noindent
\textbf{Step~4}
(Drinfeld coproduct at all spins).
The transfer matrix
$T(u) = 1 + \sum_{n \geq 1} \psi_n\, u^{-n}$
is a scalar formal power series in $u^{-1}$
($\mathfrak{gl}_1$ has rank~$1$). The Drinfeld
coproduct~\eqref{eq:gl1-drinfeld-coprod} is the product
of two such series. The
expansion~\eqref{eq:gl1-coprod-general} follows from
$(u - z)^{-m} = u^{-m} \sum_{j \geq 0}
\binom{m+j-1}{j} (z/u)^j$ applied to
$T(u) \otimes T(u-z)$ and collecting the coefficient
of $u^{-n}$, giving the spin-$1$, $2$, and $3$ expansions
displayed below.
Through spin~$3$:
\begin{align}
  \Delta_z(\psi_1)
  &= \psi_1 \otimes 1
  + 1 \otimes \psi_1,
  \label{eq:coprod-psi1} \\[4pt]
  \Delta_z(\psi_2)
  &= \psi_2 \otimes 1
  + 1 \otimes \psi_2
  + \psi_1 \otimes \psi_1
  + z\,(1 \otimes \psi_1),
  \label{eq:coprod-psi2} \\[4pt]
  \Delta_z(\psi_3)
  &= \psi_3 \otimes 1
  + 1 \otimes \psi_3
  + \psi_1 \otimes \psi_2
  + \psi_2 \otimes \psi_1
  \notag \\
  &\quad
  + z\,(\psi_1 \otimes \psi_1
  + 2 \cdot 1 \otimes \psi_2)
  + z^2\,(1 \otimes \psi_1).
  \label{eq:coprod-psi3}
\end{align}
In the $\cW_{1+\infty}$ field basis (via the Miura
substitution $\psi_1 = J$,
$\psi_2 = T + J^2/(2\Psi)$ of step~3):
\begin{align}
  \Delta_z(J)
  &= J \otimes 1
  + 1 \otimes J,
  \label{eq:coprod-J} \\[4pt]
  \Delta_z(T)
  &= T \otimes 1
  + 1 \otimes T
  + \tfrac{\Psi - 1}{\Psi}\, J \otimes J
  + z\,(1 \otimes J).
  \label{eq:coprod-T}
\end{align}
The coefficient $(\Psi - 1)/\Psi$ in~\eqref{eq:coprod-T}
arises from the Miura inversion $T = \psi_2 - J^2/(2\Psi)$:
the $\psi_1 \otimes \psi_1$ cross-term in
$\Delta_z(\psi_2)$ has coefficient~$1$, and
subtracting $(1/(2\Psi))\,\Delta_z(J)^2
= (J^2/(2\Psi)) \otimes 1 + (1/\Psi)\,J \otimes J
+ 1 \otimes J^2/(2\Psi)$
replaces the coefficient by $1 - 1/\Psi = (\Psi - 1)/\Psi$.
At $\Psi = 1$ (free boson) the cross-term vanishes;
in the classical limit $\Psi \to \infty$ it tends to~$1$.
The Heisenberg OPE
$J(z)\,J(w) \sim \Psi/(z-w)^2$ gives
% AP126: r-matrix with level prefix Psi; Psi=0 -> r=0. Verified.
$r(z) = \Psi/z$ (class~$G$, level prefix $\Psi$).
The general formula~\eqref{eq:gl1-coprod-general} gives
$\Delta_z(\psi_n)$ at every spin in closed form; the
coproduct on the $\cW_{1+\infty}$ field basis
$\{J, T, W, W_4, \ldots\}$ is obtained by inverting
the Miura transform, and the computation is triangular
and algorithmic at each spin.

\medskip
\noindent
\textbf{Step~5}
($R$-matrix from Maulik--Okounkov; vertex bialgebra
structure; Maulik--Okounkov~\cite{MO19},
Jindal--Kaubrys--Latyntsev~\cite{JKL26}).
The critical CoHA of any quiver with potential carries
a vertex coproduct forming a vertex
bialgebra (\cite{JKL26},~Theorems~A and~C). The
$R$-matrix is
\[
  R(z)
  = \frac{\Psi(\sigma^* \mathrm{Ext}^\vee, z)}
  {\Psi(\mathrm{Ext}, z)},
\]
where $\Psi$ on the right denotes the motivic exponential
and $\mathrm{Ext}$ is the $\mathrm{Ext}^1$-bundle on the
stack of quiver representations.
The Maulik--Okounkov $R$-matrix~\cite{MO19}
restricted to the CoHA gives
$R_{\mathrm{CoHA}} = q^{\alpha \cdot \beta}$
($q = e^{\pi i \Psi}$) with classical limit
$r(z) = \Psi/z$.

\medskip
\noindent
\textbf{Step~6}
(OPE compatibility at all spins: two independent
arguments).
The vertex bialgebra axiom~\eqref{eq:ope-compat}
requires $\Delta_z(Y(a,w)\,b)
= Y^{(2)}(\Delta_z(a),w)\,\Delta_z(b)$
as an identity of vertex operators, not merely of
associative algebra elements. We give two independent
proofs that this holds at all spins simultaneously.

\smallskip
\noindent\emph{Argument~A \textup{(}coderivation on the
Koszul locus\textup{)}.}
The algebra $\cW_{1+\infty}[\Psi]$ at generic~$\Psi$ lies
on the Koszul locus
\textup{(}Definition~\textup{\ref{def:koszul-locus}}:
the bar-cobar adjunction
$\Omega(\Barord(\cW_{1+\infty}[\Psi]))
\xrightarrow{\;\sim\;} \cW_{1+\infty}[\Psi]$ is an
equivalence\textup{)}. By the direction
$\textup{(II)} \to \textup{(III)}$ of
Theorem~\textup{\ref{thm:chiral-qg-equiv}}, the bar
construction produces a chiral coproduct on
$\cW_{1+\infty}[\Psi]$ from the deconcatenation coproduct
$\Delta$ on $\Barord(\cW_{1+\infty}[\Psi])
= T^c(s^{-1}\overline{\cW}_{1+\infty}[\Psi])$
via the cobar dualization~\eqref{eq:coprod-from-bar}.
The bar differential $d_{\Barord}$ is a coderivation with
respect to~$\Delta$:
\begin{equation}\label{eq:coderivation-w-infty}
  d_{\Barord} \circ \Delta
  = \Delta \circ d_{\Barord}.
\end{equation}
This identity holds at \emph{all arities
simultaneously} as a structural property of the cofree
tensor coalgebra. Under the cobar functor~$\Omega$, the
coderivation property~\eqref{eq:coderivation-w-infty}
dualizes to the OPE compatibility
axiom~\eqref{eq:ope-compat}: the bar differential
encodes the vertex operation $Y(a,w)$, and the
deconcatenation encodes the coproduct $\Delta_z$
\textup{(}see the proof of
Theorem~\textup{\ref{thm:chiral-qg-equiv}}, direction
$\textup{(II)} \to \textup{(III)}$\textup{)}.

It remains to identify this bar-cobar coproduct with the
Drinfeld coproduct of step~4. The bar-cobar coproduct
is determined by the vertex algebra structure of
$\cW_{1+\infty}[\Psi]$; the Drinfeld coproduct is
determined by the generating-series
formula~\eqref{eq:gl1-drinfeld-coprod}. Both satisfy
coassociativity and the counit axiom. On the Koszul locus,
the coproduct compatible with a given vertex $R$-matrix is
unique (Theorem~\ref{thm:chiral-qg-equiv}: the functors
$\alpha$, $\beta$, $\gamma$ form an equivalence triangle).
The bar-cobar coproduct is compatible with the
Maulik--Okounkov $R$-matrix (by construction via~$\alpha$
and~$\gamma^{-1}$). The Drinfeld coproduct is compatible
with the same $R$-matrix (the Yangian $R$-matrix of
step~5 restricts to the MO $R$-matrix). By uniqueness,
the two coproducts coincide. Therefore the Drinfeld
coproduct satisfies OPE compatibility at all spins.

\smallskip
\noindent\emph{Argument~B \textup{(}JKL vertex bialgebra
on the CoHA\textup{)}.}
Jindal--Kaubrys--Latyntsev~\cite{JKL26} construct a vertex
coproduct on the critical CoHA of any quiver with potential,
satisfying the vertex bialgebra
axiom~\eqref{eq:ope-compat} (their Theorems~A and~C).
For the Jordan quiver, the CoHA is
$\cH_{\mathrm{Jor}} \cong Y^+(\widehat{\mathfrak{gl}}_1)$
(step~1), and the JKL vertex coproduct recovers the
Drinfeld coproduct on the Yangian (their Theorem~B).
Since the JKL vertex bialgebra axiom is exactly the OPE
compatibility axiom~\eqref{eq:ope-compat}, this gives a
second proof at all spins simultaneously.

\smallskip
\noindent\emph{RTT triviality \textup{(}corroboration,
not proof\textup{)}.}
For $\mathfrak{gl}_1$, the Maulik--Okounkov $R$-matrix is
scalar:
$R(z) = g(z)$ of~\eqref{eq:gl1-scalar-r}. The RTT
relation
\begin{equation}\label{eq:rtt-gl1}
  R(u - v)\, T_1(u)\, T_2(v)
  = T_2(v)\, T_1(u)\, R(u - v)
\end{equation}
is trivially satisfied because a scalar commutes through
all operators. This confirms that the Drinfeld coproduct is
compatible with the $R$-matrix at the \emph{associative
algebra} level; arguments~A and~B upgrade this to the
full vertex algebra compatibility.
For $Y(\widehat{\mathfrak{gl}}_N)$ with $N \geq 2$,
the $R$-matrix is no longer scalar, and OPE compatibility
at each truncation is a separate question
(Remark~\ref{rem:w-infty-vertex-gap}).

\medskip
\noindent
\textbf{Step~7}
(Full axiom verification).
The spectral coproduct~\eqref{eq:gl1-drinfeld-coprod}
satisfies all axioms of
Definition~\ref{def:chiral-coproduct}:
\begin{enumerate}[label=\textup{(\alph*)}]
\item \emph{Strict coassociativity.}
  $(\Delta_{z_1} \otimes \id) \circ \Delta_{z_2}(T(u))
  = T(u) \otimes T(u - z_2) \otimes T(u - z_1 - z_2)$
  and
  $(\id \otimes \Delta_{z_2}) \circ \Delta_{z_1}(T(u))
  = T(u) \otimes T(u - z_1) \otimes T(u - z_1 - z_2)$.
  For $\mathfrak{gl}_1$ the Drinfeld associator
  acts trivially (the algebra is abelian), so
  $\Delta_z$ is \emph{strictly} coassociative.
\item \emph{Counit.}
  $\varepsilon(\psi_n) = 0$ for $n \geq 1$,
  $\varepsilon(1) = 1$ gives
  $(\varepsilon \otimes \id) \circ \Delta_z(T(u))
  = T(u - z)$ and
  $(\id \otimes \varepsilon) \circ \Delta_z(T(u))
  = T(u)$.
\item \emph{OPE compatibility at all spins.}
  By step~6: the coderivation property on the bar complex
  of $\cW_{1+\infty}[\Psi]$ (argument~A) gives the vertex
  bialgebra axiom~\eqref{eq:ope-compat} at all spins via
  the Koszul-locus equivalence
  (Theorem~\ref{thm:chiral-qg-equiv}), and the JKL vertex
  bialgebra theorem on the CoHA (argument~B) provides an
  independent proof. RTT triviality for $\mathfrak{gl}_1$
  corroborates the result at the associative algebra level.
\end{enumerate}
No constant coproduct exists for $c \neq 0$: the
conformal anomaly $T(z)\,T(w) \sim (c/2)/(z-w)^4 +
\cdots$ obstructs any $z$-independent ansatz
$\Delta(T) = T \otimes 1 + 1 \otimes T + X$
because the quartic pole on the left maps to
$(c/2 + c/2)/(z-w)^4$ while the right retains $c/2$.
The spectral parameter in~\eqref{eq:gl1-drinfeld-coprod}
absorbs this obstruction through the shift
$u \mapsto u - z$.

The three structures (I)--(III) determine each other by
Theorem~\ref{thm:chiral-qg-equiv}: the $R$-matrix is
given by step~5, the coproduct by steps~4 and~6, and the
$\Ainf$ structure is extracted from
$\Barord(\cW_{1+\infty})$ via the functor $\alpha$
with input the $R$-matrix from step~5.
\end{proof}

\begin{remark}[Structural features of the spectral coproduct]
\label{rem:w-infty-coprod-structure}
The coproduct of
Theorem~\textup{\ref{thm:w-infty-chiral-qg}}
has the following convention-independent features:
\begin{enumerate}[label=\textup{(\arabic*)}]
\item $\Delta_z(J) = J \otimes 1
  + 1 \otimes J$ is cocommutative with no
  $z$-dependence (Heisenberg is abelian,
  confirming class~$G$).
\item $\Delta_z(\psi_2)$ has a noncocommutative
  cross-term $\psi_1 \otimes \psi_1$ at order
  $z^0$ and a spectral tail
  $z\,(1 \otimes \psi_1)$ from the shift
  $u \mapsto u - z$.
  In the $\cW_{1+\infty}$ basis,
  $\Delta_z(T)$ acquires a correction
  $(\Psi - 1)/\Psi \cdot (J \otimes J)$ from the
  Miura inversion $T = \psi_2 - J^2/(2\Psi)$.
\item $\Delta_z(\psi_3)$ has cross-terms
  $\psi_1 \otimes \psi_2 + \psi_2 \otimes \psi_1$,
  a spectral tail through $z^2$, and the $z^1$
  coefficient involves both $\psi_1 \otimes \psi_1$
  and $1 \otimes \psi_2$. In the $\cW_{1+\infty}$
  basis, $\Delta_z(W)$ mixes all three spins with
  coefficients polynomial in $\Psi$.
\item The general formula~\eqref{eq:gl1-coprod-general}
  gives $\Delta_z(\psi_n)$ at every spin in closed form.
  The coproduct on the $\cW_{1+\infty}$ field basis
  $\{J, T, W, W_4, \ldots\}$ is obtained by
  inverting the Miura transform; the computation is
  triangular and algorithmic at each spin.
\end{enumerate}
\end{remark}

\begin{remark}[Effective central charge and intertwining
in the Miura basis]
\label{rem:spin2-ceff-miura-w1infty}
The coproduct~\eqref{eq:coprod-T} in the
$\cW_{1+\infty}$ field basis has cross-term coefficient
$(\Psi - 1)/\Psi$ (not $1/\Psi$), as derived in step~4.
% NEW MATERIAL FROM STANDALONE (ordered chiral homology paper)
% Inserted losslessly; macros to be adapted for memoir class
% ================================================================


\begin{remark}[Independent proof via the cohomological Hall algebra]
\label{rem:independent-proof-coha}
The equivalence $\textup{(I)} \leftrightarrow \textup{(II)}
\leftrightarrow \textup{(III)}$ admits a second, independent
proof for algebras arising as critical CoHAs of quivers with
potential, using the framework of
Kontsevich--Soibelman~\cite{KS11},
Schiffmann--Vasserot~\cite{SV13},
Yang--Zhao~\cite{YZ18a}, and
Jindal--Kaubrys--Latyntsev~\cite{JKL26}. The logic is
orthogonal to the bar-cobar adjunction above: rather than
constructing the triangle abstractly from the ordered bar
complex, one exhibits all three structures simultaneously on a
single geometric object and verifies that they determine each
other.

\smallskip
\noindent\textbf{Step~1: simultaneous existence of (I) and (III)
on the CoHA.}
Let $(Q,W)$ be a quiver with potential. The critical CoHA
\[
  \cH_{Q,W}
  = \bigoplus_{d \in \ZZ_{\geq 0}^{Q_0}}
  H^{G_d}_{\mathrm{BM},*}
  (\cM_d^{\mathrm{nil}},\, \varphi_W)
\]
carries two structures of independent geometric origin:
\begin{enumerate}[label=\textup{(\alph*)}]
\item \emph{Vertex coproduct} $\Delta_{\mathrm{CoHA}}$, from
  restriction to Levi subalgebras
  $G_{d_1} \times G_{d_2} \hookrightarrow G_d$; and
\item \emph{$R$-matrix} $R(z)$, from the
  Maulik--Okounkov stable envelope
  construction~\cite{MO19} on the Nakajima variety
  associated to~$Q$.
\end{enumerate}
Jindal--Kaubrys--Latyntsev~\cite{JKL26} prove that these two
structures satisfy the vertex bialgebra axiom (their
Theorems~A and~C):
\begin{equation}\label{eq:v1-jkl-vertex-bialgebra}
  \Delta_{\mathrm{CoHA},z}(Y(a,w)\,b)
  = Y^{(2)}(\Delta_{\mathrm{CoHA},z}(a),w)\,
  \Delta_{\mathrm{CoHA},z}(b),
\end{equation}
which is the OPE compatibility
axiom~\eqref{eq:ope-compat}. The $R$-matrix satisfies the
QYBE and unitarity by the stable envelope formalism
\cite{MO19}; the quasi-coassociativity of
$\Delta_{\mathrm{CoHA}}$ follows from the
associativity of the Hall product on the Drinfeld
double~\cite{YZ18a}. This establishes the equivalence
$\textup{(I)} \leftrightarrow \textup{(III)}$ by
construction: the $R$-matrix and the coproduct coexist on
$\cH_{Q,W}$ and satisfy the quasi-triangularity relation
$R(z)\, \Delta^{\mathrm{ch}} = \Delta^{\mathrm{ch},\op}\,
R(z)$.

\smallskip
\noindent\textbf{Step~2: extraction of (II) from the CoHA
vertex bialgebra.}
The CoHA $\cH_{Q,W}$, equipped with its $R$-matrix, is an
$\Eone$-chiral algebra in the sense of
Definition~\ref{def:e1-chiral}. The functor $\alpha$
of the equivalence triangle~\eqref{eq:equiv-triangle}
extracts the chiral $\Ainf$ operations
$\{m_k^{\mathrm{ch}}\}$ by restricting the OPE to the
associahedra
$K_k \hookrightarrow
\overline{\FM}_k^{\mathrm{ord}}(\CC)$ as
in~\eqref{eq:mk-from-ope}. The Stasheff identities follow
from Stokes' theorem on $K_{k+1}$ as before; the
convergence of the integrals is guaranteed by
$S$-locality of the $R$-matrix. This gives
$\textup{(I)} \to \textup{(II)}$ on the CoHA.

The reverse direction $\textup{(II)} \to \textup{(I)}$ is the
degree-$2$ holonomy recovery of the $R$-matrix from
$m_2^{\mathrm{ch}}$, which depends only on the formal
structure of the Stasheff associahedra and not on the CoHA
origin.

\smallskip
\noindent\textbf{Step~3: the triangle closes.}
By the JKL vertex bialgebra theorem, the CoHA
simultaneously carries (I), (II), and (III). To verify that
these three structures \emph{determine} each other (not merely
coexist), one checks the uniqueness directions. On the Koszul
locus:
\begin{enumerate}[label=\textup{(\roman*)}]
\item \emph{(I) determines (III):} given an $R$-matrix
  satisfying the QYBE and $S$-locality, the quasi-triangular
  coproduct is unique up to gauge equivalence by the
  Drinfeld reconstruction~\cite{Drinfeld90}
  (spectral parameter version:
  Etingof--Kazhdan~\cite{EK96}). The JKL coproduct on the
  CoHA is the unique one compatible with the MO $R$-matrix.
\item \emph{(III) determines (I):} the Drinfeld
  formula~\eqref{eq:r-from-coprod-proof} recovers $S(z)$
  from $\Delta^{\mathrm{ch}}$, independently of the CoHA
  construction.
\item \emph{(I) $\leftrightarrow$ (II):} the extraction of
  $\Ainf$ operations from the $R$-matrix is the
  Stasheff--FM construction (direction $\alpha$), whose
  inverse recovers the $R$-matrix as holonomy (direction
  $\alpha^{-1}$). Neither direction uses the CoHA.
\end{enumerate}
The triangle coherence
$\alpha \circ \gamma \circ \beta = \id$ then follows from
the uniqueness of each vertex.

\smallskip
\noindent\textbf{Independence from the bar-cobar proof.}
The argument above differs from the proof of
Theorem~\ref{thm:chiral-qg-equiv} in two respects.
First, the equivalence $\textup{(I)} \leftrightarrow
\textup{(III)}$ is established by the geometric CoHA
construction (stable envelopes for the $R$-matrix,
Levi restriction for the coproduct, JKL for the vertex
bialgebra axiom), bypassing the ordered bar complex
$\Barord(\cA) = T^c(s^{-1}\bar{\cA})$ and the
cobar dualization~\eqref{eq:coprod-from-bar}. Second, the
uniqueness of the coproduct given the $R$-matrix is
deduced from the Drinfeld reconstruction theorem for
quasi-triangular vertex bialgebras, rather than from the
bar-cobar adjunction on the Koszul locus. The functor
$\alpha$ (extraction of $\Ainf$ operations from the
$R$-matrix via associahedron integrals) is shared between
the two proofs, but it enters as an a~posteriori observation
in the CoHA argument, not as the engine of the
construction.

\smallskip
\noindent\textbf{Scope.}
The CoHA proof applies to $\Eone$-chiral algebras
arising as critical CoHAs of quivers with potential.
For the Jordan quiver, $\cH_{\mathrm{Jor}} \cong
Y^+(\widehat{\mathfrak{gl}}_1)$~\cite{SV13}, and the
equivalence specializes to the chiral quantum group
structure on $\cW_{1+\infty}[\Psi]$
(Theorem~\ref{thm:w-infty-chiral-qg}).
For quivers of type $\widehat{A}_{N-1}$, the CoHA
is $Y^+(\widehat{\mathfrak{sl}}_N)$, and the JKL
coproduct recovers the Drinfeld coproduct on the affine
Yangian~\cite{JKL26}. Beyond quiver-origin algebras, the
CoHA argument does not directly apply; the bar-cobar proof
of Theorem~\ref{thm:chiral-qg-equiv} remains the primary
route for general $\Eone$-chiral algebras on the Koszul
locus.
\end{remark}

\begin{remark}[Effective central charge and intertwining
in the Miura basis]
\label{rem:spin2-ceff-miura}
The coproduct~\eqref{eq:coprod-T} in the
$\cW_{1+\infty}$ field basis has cross-term coefficient
$(\Psi - 1)/\Psi$ (not $1/\Psi$), as derived in step~4.
This coefficient controls two structural invariants of the
image $\Delta_z(T)$ in $V \otimes V$:
\begin{enumerate}[label=\textup{(\arabic*)}]
\item \emph{Effective central charge.}
  The image $\Delta_z(T_n)$ generates a Virasoro
  subalgebra in the tensor product mode algebra with
  effective central charge
  \begin{equation}\label{eq:c-eff-spin2}
    c_{\mathrm{eff}}(\Psi) = 2 + 2(\Psi - 1)^2.
  \end{equation}
  At $\Psi = 1$ (free boson): $c_{\mathrm{eff}} = 2$
  (two decoupled copies, the cross-term vanishes).
  At $\Psi = 2$: $c_{\mathrm{eff}} = 4$.
  The quadratic growth $c_{\mathrm{eff}} \sim 2\Psi^2$ at
  large $\Psi$ reflects the quadratic Miura nonlinearity
  $\psi_2 = T + J^2/(2\Psi)$.
  The formula is independent of the spectral parameter~$z$
  (the $z$-dependent terms in $\Delta_z(T)$ do not
  contribute to the vacuum expectation value of the
  commutator $[\Delta_z(T_2), \Delta_z(T_{-2})]$).
\item \emph{Intertwining with $\Delta_z(J)$.}
  The commutator
  $[\Delta_z(T_n),\, \Delta_z(J_m)]
  = -\Psi \cdot m \cdot \Delta_z(J_{n+m})$
  has intertwining factor $\Psi$ (not~$1$).
  The extra factor $\Psi - 1$ beyond the original algebra
  relation $[T_n, J_m] = -m\,J_{n+m}$ arises from the
  cross-term $(\Psi - 1)/\Psi \cdot (J \otimes J)_n$
  commuting with $\Delta_z(J_m) = J_m^L + J_m^R$:
  the left and right Heisenberg commutators each contribute
  $-((\Psi-1)/\Psi) \cdot \Psi \cdot m = -(\Psi - 1)m$.
  At $\Psi = 1$: factor~$1$ (algebra homomorphism).
\end{enumerate}
Both features are consequences of the Miura nonlinearity:
the Drinfeld coproduct is an algebra homomorphism of the
Yangian $Y(\widehat{\mathfrak{gl}}_1)$ (preserving the
$\psi$-generator relations), but when expressed in the
$\cW_{1+\infty}$ field basis via the nonlinear Miura
inversion $T = \psi_2 - J^2/(2\Psi)$, the image acquires
$\Psi$-dependent structural invariants that are absent in
the $\psi$-basis.
\end{remark}

\begin{remark}[Descent to $\cW_N$]
\label{rem:w-infty-descent}
Drinfeld--Sokolov reduction gives
$\pi_N \colon \cW_{1+\infty}[\Psi]
\twoheadrightarrow \cW_N$
for each $N \geq 2$. Since the spectral coproduct of
Theorem~\ref{thm:w-infty-chiral-qg} is defined
on the full $\cW_{1+\infty}$, it descends to $\cW_N$
provided $\pi_N \otimes \pi_N$ intertwines $\Delta_z$ with
a coproduct on $\cW_N$. For $\cW_{1+\infty}$ itself,
OPE compatibility holds at all spins by the $GL_1$
triviality argument (step~6 of the proof); for finite
$N$, the intertwining $(\pi_N \otimes \pi_N) \circ
\Delta_z = \Delta_z^{(N)} \circ \pi_N$ follows from the
compatibility of the Miura transform with truncation
(Remark~\ref{rem:w-infty-vertex-gap}).
The tower
$\cdots \twoheadrightarrow \cW_{N+1}
\twoheadrightarrow \cW_N
\twoheadrightarrow \cdots
\twoheadrightarrow \cW_2 = \mathrm{Vir}_c$
is a projective system of algebras with spectral
coproducts.
At $N = 2$: class $M$, $\kappa = c/2$, irregular KZ
connection. At $N = 3$: $\kappa(\cW_3) = 5c/6$ (where
$H_3 - 1 = 5/6$).
The Gaiotto--Rap\v{c}\'ak $Y$-algebras
$Y_{N_1,N_2,N_3}[\Psi]$~\cite{gaiotto-rapchak} are truncations of
$\cW_{1+\infty}$; the $S_3$-triality is a symmetry
of the spectral coproduct, corresponding to three Hecke
correspondences on spiked-instanton moduli
(Rap\v{c}\'ak--Soibelman--Yang--Zhao~\cite{RSYZ20}).
The $\cD$-module $\cF_n^{\mathrm{ord}}(\cW_{1+\infty})$
has irregular singularities of unbounded Poincar\'e rank
(spin-$s$ contributes pole order $2s - 1$ after $d$-log
absorption), providing the boundary algebra of higher-spin
holographic gravity.
\end{remark}

\begin{remark}[GRZ comparison]
\label{rem:grz-comparison}
Theorem~\ref{thm:w-infty-chiral-qg} has three comparison surfaces
in the type-$A$ literature.

\smallskip
\noindent\emph{(1) Corner truncations.}
Gaiotto--Rap\v{c}\'ak identify the corner algebras
$Y_{N_1,N_2,N_3}[\Psi]$ as truncations of
$\cW_{1+\infty}$ along the curve
$\sum_i N_i/\lambda_i = 1$ and prove generic free strong
generation~\cite[Theorem~3.1 and eq.~3.35]{gaiotto-rapchak}.
This matches Remark~\ref{rem:w-infty-descent}: our spectral
coproduct is constructed on the ordered
$\cW_{1+\infty}$ object and only then descends to finite
truncations.

\smallskip
\noindent\emph{(2) DDCA and affine-Yangian coproducts.}
Gaiotto--Rap\v{c}\'ak--Zhou construct coproducts on the
finite-rank DDCA / matrix-extended $\cW_\infty$ /
affine-Yangian side~\cite[\S\S5--7]{GRZ25}. For type $A$,
this is consistent with the Drinfeld coproduct of
Theorem~\ref{thm:w-infty-chiral-qg}: both sides work with the
same Yangian mode algebra, and both organize truncation data at
finite rank before passing to the large-rank
$\cW_{1+\infty}$ limit. The statement proved here is different in
nature: the coproduct is recovered from deconcatenation on the
ordered bar coalgebra
$\Barord(\cA) = T^c(s^{-1}\bar{\cA})$, and step~6 proves OPE
compatibility by a coderivation argument on the ordered bar side.

\smallskip
\noindent\emph{(3) Independent CoHA route.}
For quiver-origin algebras there is a second independent proof.
Jindal--Kaubrys--Latyntsev construct a vertex coproduct on
critical CoHAs and recover the deformed Drinfeld coproduct for
ADE Yangians after restriction to the Yangian generators
\cite{JKL26}. This matches
Remark~\ref{rem:independent-proof-coha}: the CoHA vertex
coproduct gives the same ordered $\Eone$-chiral coproduct that
the bar-cobar argument produces abstractly.

\smallskip
\noindent\emph{Convention check.}
Our spectral parameter $z$ is the Drinfeld/Yangian shift
parameter. The OPE coordinate enters only after the
factorization-envelope or ordered-bar reconstruction, and the
braided $\Etwo$ structure appears later at the derived-center
stage, not on the primitive bar coalgebra.
\end{remark}

\begin{remark}[Proof structure and remaining questions]
\label{rem:w-infty-vertex-gap}
Three comments on the proof of
Theorem~\textup{\ref{thm:w-infty-chiral-qg}}.

\smallskip
\noindent\emph{(1) Why RTT alone does not suffice.}
The RTT relation~\eqref{eq:rtt-gl1} is a constraint on the
\emph{associative algebra} structure (mode commutators).
The vertex bialgebra axiom~\eqref{eq:ope-compat} is a
constraint on the \emph{vertex algebra} structure (all
$n$-products $a_{(n)}b$, including the regular products for
$n < 0$). For $\mathfrak{gl}_1$ the RTT relation is
trivially satisfied because the $R$-matrix is scalar; this
confirms associative-algebra compatibility but does not, by
itself, establish the full vertex-algebra compatibility.
The two independent arguments of step~6 (the
coderivation property on the Koszul-locus bar complex, and
the JKL vertex bialgebra theorem on the CoHA) close this
gap.

\smallskip
\noindent\emph{(2) Truncation to $\cW_N$.}
Theorem~\ref{thm:glN-chiral-qg} resolves this question:
for the truncations $\cW_N$ at finite $N \geq 2$, the
$R$-matrix of $Y(\widehat{\mathfrak{gl}}_N)$ is Yang's
rational $R$-matrix $R(u) = u\, I + \Psi\, P$ (no longer
scalar), and the OPE compatibility at each truncation level
is established by the same two arguments.
The coderivation argument (argument~A) applies to
$\cW_N$ at generic level by the same Koszul-locus reasoning.
The JKL argument (argument~B) applies to
$Y(\widehat{\mathfrak{gl}}_N)$ by their Theorem~C for
general quivers.

\smallskip
\noindent\emph{(3) Antipode and Drinfeld double.}
The Schiffmann--Mellit--Minets--Vasserot
isomorphism~\cite{SMMV23} gives the positive half
$\cH_{\mathrm{Jor}} \cong U(\cW_{1+\infty}[\Psi])^+$.
The Drinfeld double
$D(\cH_{\mathrm{Jor}}) \cong
Y(\widehat{\mathfrak{gl}}_1)$ (Yang--Zhao~\cite{YZ18a})
provides the full algebra with antipode
$S(T(u)) = T(u)^{-1}$.
Proposition~\ref{prop:w-infty-antipode-obstruction} below
resolves the question of whether this lifts to a
vertex Hopf algebra structure: the Yangian antipode gives
explicit formulas $S(J) = -J$ and
$S(T) = -T + \tfrac{\Psi - 1}{\Psi}\,{:}JJ{:}$ on
$\cW_{1+\infty}$, but two independent obstructions prevent
it from satisfying the vertex Hopf algebra axiom.
\end{remark}

\begin{proposition}[Yangian antipode on $\cW_{1+\infty}$:
explicit formula and vertex Hopf obstruction]
\label{prop:w-infty-antipode-obstruction}
\ClaimStatusProvedHere
Let $\cW_{1+\infty}[\Psi]$ carry the chiral quantum group
structure of Theorem~\textup{\ref{thm:w-infty-chiral-qg}},
with spectral coproduct
$\Delta_z(T(u)) = T(u) \otimes T(u-z)$ and Miura
identification $\psi_1 = J$,
$\psi_2 = T + J^2/(2\Psi)$.
\begin{enumerate}[label=\textup{(\roman*)}]
\item \textup{(Explicit antipode.)}
The Yangian antipode $S(T(u)) = T(u)^{-1}$, with
$T(u)^{-1} = 1 + \sum_{n \geq 1} \sigma_n\, u^{-n}$
and
\begin{equation}\label{eq:transfer-inverse-recurrence}
  \sigma_n
  = -\psi_n
  - \sum_{k=1}^{n-1} \psi_k\, \sigma_{n-k},
\end{equation}
induces a well-defined linear map
$S \colon \cW_{1+\infty} \to \cW_{1+\infty}$
via the Miura transform. At spins~$1$ and~$2$:
\begin{equation}\label{eq:w-infty-antipode-explicit}
  S(J) = -J,
  \qquad
  S(T) = -T
  + \tfrac{\Psi - 1}{\Psi}\,{:}JJ{:}.
\end{equation}
The map~$S$ is an involution: $S^2 = \mathrm{id}$.

\item \textup{(OPE obstruction.)}
The map~$S$ does not preserve the Virasoro OPE at generic~$\Psi$.
The quartic pole of $S(T)$ with itself is
\begin{equation}\label{eq:antipode-quartic-obstruction}
  S(T)_{(3)}\, S(T)
  \;=\;
  \tfrac{c}{2}
  + 2(\Psi - 1)(\Psi - 2),
\end{equation}
which equals $T_{(3)}\, T = c/2$ only at
$\Psi \in \{1, 2\}$
\textup{(}$c = 1$ free boson, $c = -2$ $bc$~ghost\textup{)}.

The four OPE contributions are:
$T_{(3)}T = c/2$\textup{;}
$T_{(3)}\,({:}JJ{:}) = \Psi$
\textup{(}Wick theorem, $J$ primary of weight~$1$\textup{);}
$({:}JJ{:})_{(3)}\, T = \Psi$
\textup{(}skew-symmetry\textup{);}
$({:}JJ{:})_{(3)}\,({:}JJ{:}) = 2\Psi^2$
\textup{(}Sugawara scaling\textup{)}.

\item \textup{(Hopf axiom obstruction.)}
The vertex Hopf algebra axiom
$m_z \circ (S \otimes \mathrm{id}) \circ \Delta_z
= \iota \circ \varepsilon$
fails on~$T$ with residual
\begin{equation}\label{eq:antipode-hopf-obstruction}
  m_z \circ (S \otimes \mathrm{id})
  \circ \Delta_z(T)
  \;=\;
  -\,\frac{\Psi - 1}{z^2}
  \;+\; z\, J.
\end{equation}
The $-(\Psi - 1)/z^2$ term arises from the
OPE singularity $Y(J, z)\, J \sim \Psi/z^2$;
the $z \cdot J$ term arises from the
spectral dependence $z \cdot (1 \otimes J)$ in
$\Delta_z(T)$.
The $z \cdot J$ term persists at all~$\Psi$, including
$\Psi = 1$: the spectral coproduct $\Delta_z$ is
$z$-dependent and cannot admit a $z$-independent
antipode satisfying the classical Hopf axiom.
\end{enumerate}
\end{proposition}

\begin{proof}
Part~(i): the recurrence~\eqref{eq:transfer-inverse-recurrence}
follows from $T(u)\, T(u)^{-1} = 1$ by extracting the
coefficient of~$u^{-n}$. The Miura inversion giving
$S(J) = -\psi_1 = -J$ and
$S(T) = S(\psi_2 - J^2/(2\Psi))
= (\psi_1^2 - \psi_2) - S(J)^2/(2\Psi)
= -T + (\Psi - 1)/\Psi \cdot {:}JJ{:}$
uses the commutativity of the $\psi_i$ (rank~$1$).
The involutivity $S^2 = \mathrm{id}$ is the identity
$S^2(T) = -S(T) + (\Psi-1)/\Psi \cdot S(J)^2
= T - (\Psi-1)/\Psi \cdot J^2
+ (\Psi-1)/\Psi \cdot J^2 = T$.

Part~(ii): write $\alpha = (\Psi - 1)/\Psi$
and $S(T) = -T + \alpha\,{:}JJ{:}$.
By bilinearity of the~$(3)$-product,
\[
  S(T)_{(3)}\, S(T)
  = T_{(3)}T
  - \alpha\, T_{(3)}\,({:}JJ{:})
  - \alpha\, ({:}JJ{:})_{(3)}\, T
  + \alpha^2\, ({:}JJ{:})_{(3)}\,({:}JJ{:}).
\]
The individual $(3)$-products are computed by the Wick
theorem (non-linear lambda-bracket formula):
\begin{itemize}
\item $T_{(3)}\,({:}JJ{:}) = \Psi$: from the
  lambda-bracket $\{T_\lambda\, {:}JJ{:}\}
  = (\Psi/6)\, \lambda^3 + 2\,{:}JJ{:}\,\lambda
  + \partial({:}JJ{:})$, reading off the
  $\lambda^3/3!$ coefficient.
\item $({:}JJ{:})_{(3)}\, T = \Psi$: by
  skew-symmetry and $T_{(n)}\,({:}JJ{:}) = 0$
  for $n \geq 4$.
\item $({:}JJ{:})_{(3)}\,({:}JJ{:}) = 2\Psi^2$:
  the Sugawara tensor
  $T_{\mathrm{Sug}} = {:}JJ{:}/(2\Psi)$
  satisfies Virasoro at~$c = 1$, so
  $({:}JJ{:})_{(3)}\,({:}JJ{:})
  = (2\Psi)^2 \cdot (1/2) = 2\Psi^2$.
\end{itemize}
Substituting:
$S(T)_{(3)}\, S(T)
= c/2 - 2\alpha\Psi + 2\alpha^2 \Psi^2
= c/2 + 2(\Psi - 1)(\Psi - 2)$.

Part~(iii): the spectral coproduct gives four terms
$\Delta_z(T) = T \otimes 1 + 1 \otimes T
+ \alpha\, J \otimes J + z\,(1 \otimes J)$.
Applying $S \otimes \mathrm{id}$ and composing with
$m_z(a \otimes b) = Y(a, z)\, b$:
\begin{align*}
  &Y(S(T), z) \cdot 1
  + Y(1, z) \cdot T
  + \alpha\, Y(-J, z) \cdot J
  + z \cdot Y(1, z) \cdot J \\
  &= S(T) + T
  + \alpha\,(-\Psi/z^2 - {:}JJ{:} + O(z))
  + z \cdot J \\
  &= (-T + \alpha\,{:}JJ{:}) + T
  - \alpha\Psi/z^2 - \alpha\,{:}JJ{:}
  + z\, J \\
  &= -(\Psi - 1)/z^2 + z\, J.
\end{align*}
Since $\varepsilon(T) = 0$ (positive conformal weight),
the right-hand side of the Hopf axiom is zero,
and the obstruction is the nonzero expression above.

\smallskip
\noindent\emph{Independence of the two obstructions.}
The OPE obstruction~\eqref{eq:antipode-quartic-obstruction}
vanishes on the locus $\Psi \in \{1, 2\}$, while the
Hopf-axiom obstruction~\eqref{eq:antipode-hopf-obstruction}
contains the $z \cdot J$ tail for every value of~$\Psi$.
Substituting $\Psi = 1$ into (ii) and (iii):
$S(T)_{(3)}\, S(T) - c/2 = 2(0)(-1) = 0$, while
$m_z \circ (S \otimes \mathrm{id}) \circ \Delta_z(T)
= z \cdot J \neq 0$.
The OPE locus and the Hopf locus are therefore disjoint, so the
two obstructions are independent: neither reduces to the other,
and lifting the antipode requires resolving both simultaneously.
The $z \cdot J$ tail encodes the spectral covariance of
$\Delta_z$, which has no counterpart in the classical Yangian
antipode; it is this covariance, not the OPE singularity, that is
the genuinely chiral obstruction.
\end{proof}

\begin{lemma}[Miura inversion of the spectral coproduct at
spin~$2$]
\label{lem:coprod-T-miura}
\ClaimStatusProvedHere
Let $\psi_1 = J$ and $\psi_2 = T + J^2/(2\Psi)$ be the
Miura identification of
Proch\'azka--Rap\v{c}\'ak~\textup{\cite{PR19}}.
The spectral coproduct
$\Delta_z(T)
= \Delta_z(\psi_2) - (1/(2\Psi))\,\Delta_z(J)^2$
equals
\begin{equation}\label{eq:coprod-T-derived}
  \Delta_z(T)
  \;=\;
  T \otimes 1
  + 1 \otimes T
  + \tfrac{\Psi - 1}{\Psi}\,
  J \otimes J
  + z\,(1 \otimes J).
\end{equation}
\end{lemma}

\begin{proof}
The Drinfeld formula
$\Delta_z(T(u)) = T(u) \otimes T(u-z)$
gives, at the $u^{-2}$ coefficient,
\[
  \Delta_z(\psi_2)
  \;=\;
  \psi_2 \otimes 1
  + 1 \otimes \psi_2
  + \psi_1 \otimes \psi_1
  + z\,(1 \otimes \psi_1).
\]
Since $\Delta_z$ is an algebra homomorphism and
$J = \psi_1$ is primitive,
\begin{align*}
  \Delta_z(J)^2
  &= (J \otimes 1 + 1 \otimes J)^2 \\
  &= J^2 \otimes 1
  + 2\,J \otimes J
  + 1 \otimes J^2.
\end{align*}
Substituting $T = \psi_2 - J^2/(2\Psi)$:
\begin{align*}
  \Delta_z(T)
  &= \Delta_z(\psi_2)
  - \tfrac{1}{2\Psi}\,\Delta_z(J)^2 \\
  &= \Bigl[
    \bigl(T + \tfrac{J^2}{2\Psi}\bigr) \otimes 1
    + 1 \otimes \bigl(T + \tfrac{J^2}{2\Psi}\bigr)
    + J \otimes J
    + z\,(1 \otimes J)
  \Bigr] \\
  &\qquad
  - \tfrac{1}{2\Psi}\Bigl[
    J^2 \otimes 1
    + 2\,J \otimes J
    + 1 \otimes J^2
  \Bigr] \\
  &= T \otimes 1
  + 1 \otimes T
  + \bigl(1 - \tfrac{1}{\Psi}\bigr)\,
  J \otimes J
  + z\,(1 \otimes J).
\end{align*}
Since $1 - 1/\Psi = (\Psi - 1)/\Psi$, this
is~\eqref{eq:coprod-T-derived}. The
$J^2/(2\Psi)$ terms cancel between the Miura
substitution and the $\Delta_z(J)^2$ correction.
At $\Psi = 1$ (free boson, $c = 1$) the
$J \otimes J$ cross-term vanishes and
$\Delta_z(T) = T \otimes 1 + 1 \otimes T
+ z\,(1 \otimes J)$;
in the classical limit $\Psi \to \infty$ the
coefficient tends to~$1$.
\end{proof}


% ========================================================
\subsection{The spin-$3$ coproduct, Miura universality,
and the conformal anomaly}
\label{subsec:spin3-miura-anomaly}

\begin{proposition}[Spin-$3$ Miura coproduct]
\label{prop:spin3-miura-coprod}
\ClaimStatusProvedHere
The spectral coproduct on the spin-$3$ field $W$ of
$\cW_{1+\infty}[\Psi]$ is
\begin{align}
\label{eq:coprod-W}
  \Delta_z(W)
  &= W \otimes 1 + 1 \otimes W
  + \tfrac{\Psi - 1}{\Psi}\,
  (J \otimes T + T \otimes J) \\
  &\quad
  + \tfrac{1 - \Psi}{2\Psi^2}\,
  (J \otimes {:}J^2{:}
  + {:}J^2{:} \otimes J)
  + \tfrac{\Psi - 1}{\Psi}\,z\,
  J \otimes J
  + 2z\,(1 \otimes T)
  + z^2\,(1 \otimes J).
  \notag
\end{align}
The coefficient $(\Psi - 1)/\Psi$ on the primary
cross-terms $J \otimes T + T \otimes J$ is the same as
the coefficient on $J \otimes J$ in
$\Delta_z(T)$~\eqref{eq:coprod-T-derived}. A new
composite correction $(1 - \Psi)/(2\Psi^2)$ appears
on $J \otimes {:}J^2{:} + {:}J^2{:} \otimes J$, with
opposite sign and suppressed by $1/(2\Psi)$.
At $\Psi = 1$ \textup{(}free boson\textup{)}: all
cross-terms vanish. At $\Psi \to \infty$: primary
coefficients tend to~$1$, composite corrections to~$0$.
\end{proposition}

\begin{proof}
The Miura transform gives
$\psi_3 = W + (1/\Psi){:}JT{:} + (1/(6\Psi^2)){:}J^3{:}
+ (1/(2\Psi))J''$. The Drinfeld formula at $u^{-3}$
gives $\Delta_z(\psi_3) = \psi_3 \otimes 1 + 1 \otimes
\psi_3 + \psi_1 \otimes \psi_2 + \psi_2 \otimes \psi_1
+ z(\psi_1 \otimes \psi_1 + 2 \cdot 1 \otimes \psi_2)
+ z^2(1 \otimes \psi_1)$. Substituting $W = \psi_3 -
(1/\Psi){:}JT{:} - \ldots$ and using the multiplicativity
of $\Delta_z$ on the Miura correction terms
gives~\eqref{eq:coprod-W} after cancellation.
The primary coefficient $(\Psi - 1)/\Psi = 1 - 1/\Psi$
arises as $1$ (from $\psi_1 \otimes \psi_2$)
minus $1/\Psi$ (from $\Delta_z({:}JT{:})$).
\end{proof}

\begin{theorem}[Miura cross-term universality]
\label{thm:miura-cross-universality-monograph}
\ClaimStatusProvedHere
For all $s \geq 2$, the spectral coproduct on the
spin-$s$ primary generator $W_s$ of
$\cW_{1+\infty}[\Psi]$ satisfies
\begin{equation}\label{eq:miura-cross-universal}
  \Delta_z(W_s)
  = W_s \otimes 1 + 1 \otimes W_s
  + \frac{\Psi - 1}{\Psi}
    \bigl(J \otimes W_{s-1} + W_{s-1} \otimes J\bigr)
  + (\textup{lower-spin and composite corrections})
  + (\textup{spectral tail in~$z$}).
\end{equation}
The coefficient $(\Psi - 1)/\Psi$ of the primary
cross-term $J \otimes W_{s-1} + W_{s-1} \otimes J$
is independent of the spin~$s$.
\end{theorem}

\begin{proof}
The argument has three steps.

\emph{Step~1: the Miura coefficient of
${:}J \cdot W_{s-1}{:}$ in $\psi_s$ is $1/\Psi$
at every spin.}
The quantum Miura factorisation
\textup{(}Proch\'azka--Rap\v{c}\'ak~\cite{PR19}\textup{)}
writes the transfer matrix as a product of Miura fields:
$T(u) = \prod_{i \geq 1}(1 + \Lambda_i/u)$, with
$\psi_s = e_s(\Lambda_1, \Lambda_2, \ldots)$ the $s$-th
elementary symmetric function.
Each $\Lambda_i$ decomposes as
$\Lambda_i = \sigma_i + J_i/\Psi$, where
$\sigma_i$ is the primary \textup{(}non-$J$\textup{)}
contribution from the $i$-th Miura slot.
Expanding $e_s$ and collecting the sector with exactly
one $J$-insertion and $(s-1)$ primary slots:
\[
  e_s(\Lambda_1, \ldots)
  \;\supset\;
  \frac{1}{\Psi}
  \sum_{i}
  J_i \cdot \prod_{j \neq i}^{s-1\text{ slots}} \sigma_j
  \;=\;
  \frac{1}{\Psi}\,{:}J \cdot W_{s-1}{:}\,,
\]
where the last identification uses the fact that the
$(s-1)$-slot primary block is $W_{s-1}$ by the triangular
structure of the Miura transform
\textup{(}Step~3 of the proof of
Theorem~\textup{\ref{thm:w-infty-chiral-qg}}\textup{)}.
Each such term carries the same factor $1/\Psi$ from
the single $J/\Psi$ slot.
The remaining sectors have $k \geq 2$ distinguished
$J$-insertions with $W$-spin $\leq s - 2$:
\begin{equation}\label{eq:miura-triangular}
  \psi_s
  \;=\;
  W_s + \frac{1}{\Psi}\,{:}J \cdot W_{s-1}{:}
  + (\textup{terms of $W$-spin $\leq s-2$}).
\end{equation}

\emph{Step~2: the Drinfeld formula contributes $+1$ to
$\psi_1 \otimes \psi_{s-1}$.}
The general coproduct
formula~\eqref{eq:gl1-coprod-general} gives the
$z^0$ coefficient of $\psi_1 \otimes \psi_{s-1}$ as
$\binom{s - 2}{s - 2} = 1$, and the $z^0$ coefficient
of $\psi_{s-1} \otimes \psi_1$ as
$\binom{0}{0} = 1$. Both are identically $1$ for all
$s \geq 2$.

\emph{Step~3: assembly and non-contribution of lower
sectors.}
Inverting~\eqref{eq:miura-triangular}:
\[
  W_s
  \;=\;
  \psi_s - \frac{1}{\Psi}\,{:}J \cdot W_{s-1}{:}
  + (\textup{terms of $W$-spin $\leq s-2$}).
\]
Applying $\Delta_z$ and extracting the coefficient of
$J \otimes W_{s-1}$:
\begin{enumerate}[label=\textup{(\roman*)}]
\item From $\Delta_z(\psi_s)$: the
  $\psi_1 \otimes \psi_{s-1}$ term contributes
  $+1 \cdot J \otimes W_{s-1}$, since
  $\psi_1 = J$ and
  $\psi_{s-1} = W_{s-1} + {}$composites of
  spin $\leq s - 2$, which do not produce $W_{s-1}$
  in the right tensor factor.
\item From $-(1/\Psi)\,\Delta_z({:}J \cdot W_{s-1}{:})$:
  the leading piece
  $(J \otimes 1)(1 \otimes W_{s-1})
  = J \otimes W_{s-1}$ contributes
  $-1/\Psi$.
\item From the Miura sectors with $k \geq 2$
  $J$-insertions: their $W$-spin is $\leq s - 2$.
  The coproduct $\Delta_z$ does not raise $W$-spin in
  either tensor factor
  \textup{(}triangularity of the Miura transform\textup{)},
  so the tensor support
  of these sectors is contained in $W$-spin $\leq s - 2$
  on each side, and they contribute zero to the
  $J \otimes W_{s-1}$ channel.
\end{enumerate}
The total coefficient is $1 - 1/\Psi = (\Psi - 1)/\Psi$.
The argument for $W_{s-1} \otimes J$ is symmetric.
The three ingredients are each spin-independent: the
Drinfeld binomial coefficient $\binom{s-2}{s-2} = 1$
is tautological, the Miura coefficient $1/\Psi$ is
structural \textup{(}one $J/\Psi$ slot\textup{)},
and the non-contribution of higher $J$-sectors follows
from $W$-spin filtration.
Limiting cases: at $\Psi = 1$
\textup{(}free boson\textup{)} the coefficient vanishes;
at $\Psi \to \infty$ it tends to~$1$.
Verified computationally at spins $2$--$6$
\textup{(}\texttt{miura\_universality\_proof\_engine.py},
$142$ tests\textup{)}.
\end{proof}

\begin{remark}[The conformal anomaly forces the spectral
parameter]
\label{rem:conformal-anomaly-monograph}
The quartic pole $T(z)\,T(w) \sim (c/2)/(z-w)^4 + \cdots$
obstructs constant coproducts for $c \neq 0$: the primitive
ansatz $\Delta(T) = T \otimes 1 + 1 \otimes T + X$ produces
quartic pole $c/(z-w)^4$ (each copy contributes $c/2$), but
the coproduct must reproduce $c/2$. The obstruction equals
$c/2 = \kappa(\mathrm{Vir}_c)$. At $c = 0$ (Heisenberg):
vanishes, constant coproduct exists. At $c \neq 0$: the
spectral parameter in $\Delta_z(T(u)) = T(u) \otimes T(u - z)$
absorbs the mismatch through $u \mapsto u - z$.
\end{remark}

\begin{remark}[Stokes ray count for $\cW_N$]
\label{rem:stokes-wN-monograph}
The $\cW_N$ KZ connection at degree~$2$ has Poincar\'e rank
$2N - 2$ \textup{(}from the $W_N$--$W_N$ OPE pole $2N$,
d-log absorption to $2N - 1$\textup{)} and, in the
generic anti-Stokes / Stokes-sector convention of
Boalch and Jimbo--Miwa--Ueno, $2r = 4N - 4$ Stokes rays.
$\cW_2$: $4$ rays. $\cW_3$: $8$ rays. Linear growth in $N$.
\end{remark}

\begin{remark}[The antipode does not lift]
\label{rem:antipode-monograph}
The Yangian antipode $S(T(u)) = T(u)^{-1}$ does not lift
to a vertex-algebraic antipode on $\cW_{1+\infty}[\Psi]$.
Two independent obstructions:
\textup{(1)}~$S(T)_{(3)}S(T) = c/2 + 2(\Psi-1)(\Psi-2)
\neq c/2$ at generic $\Psi$;
\textup{(2)}~the Hopf axiom residual $z \cdot J$ persists
at all $\Psi$. Both vanish only at $\Psi \in \{1, 2\}$.
Source: the Miura nonlinearity $T = \psi_2 - J^2/(2\Psi)$.
\end{remark}

\begin{remark}[Shadow tower as perturbative Gromov--Witten]
\label{rem:shadow-gw-c3}
At $\kappa = \Psi$ \textup{(}the level of $\cW_{1+\infty}$\textup{)},
the shadow tower $\{F_g\}$ produces the perturbative
constant-map Gromov--Witten free energies
$F_g^{\mathrm{GW,const}}(\bC^3)$. The MacMahon function
$M(q) = \prod_{n \geq 1}(1 - q^n)^{-n}$ lives on the
Donaldson--Thomas side; the bridge is the MNOP/DT-GW
correspondence under $q = -e^{i g_s}$. For $\bC^3$
\textup{(}no compact curves\textup{)}, the shadow tower
IS the full GW partition function.
\end{remark}


% ========================================================
\subsection{The chiral quantum group structure on
$\cW_N$ for $N \geq 2$}
\label{subsec:glN-chiral-qg}

Theorem~\ref{thm:w-infty-chiral-qg} established the chiral
quantum group structure on $\cW_{1+\infty}[\Psi]$ via the
scalar Maulik--Okounkov $R$-matrix. For the truncations
$\cW_N = \cW_{1+\infty}/\cI_N$ (generators of spins
$1, 2, \ldots, N$) the $R$-matrix is no longer scalar: it
is Yang's rational $R$-matrix on $\CC^N \otimes \CC^N$,
and the RTT relation is a non-trivial matrix identity.

The general theorem is best understood by first computing
everything at $N = 2$, where $\cW_2 = \mathrm{Vir}_c$ with
$c = 1 - 6(\Psi - 1)^2/\Psi$ and the affine Yangian
$Y(\widehat{\mathfrak{gl}}_2)$ has a $2 \times 2$ transfer
matrix. At $N = 2$ every matrix is small enough to write
in full, and the general pattern becomes visible.

\begin{example}[The $N = 2$ chiral quantum group:
$Y(\widehat{\mathfrak{gl}}_2)$ and Virasoro]
\label{ex:gl2-chiral-qg}

\emph{Transfer matrix.}
The $2 \times 2$ transfer matrix is
\begin{equation}\label{eq:gl2-transfer}
  T(u) = I_2 + T^{(1)} u^{-1} + T^{(2)} u^{-2} + \cdots
  = \begin{pmatrix}
    t_{11}(u) & t_{12}(u) \\
    t_{21}(u) & t_{22}(u)
  \end{pmatrix},
\end{equation}
where each entry $t_{ij}(u) = \delta_{ij}
+ \sum_{r \geq 1} t_{ij}^{(r)}\, u^{-r}$.
At degree~$1$ the four generators are
$t_{11}^{(1)}, t_{12}^{(1)}, t_{21}^{(1)}, t_{22}^{(1)}$.
The quantum Miura factorisation~\eqref{eq:glN-miura}
gives
\[
  T(u) = \begin{pmatrix} u + \Lambda_1^{(1)} & 0 \\
  0 & 1 \end{pmatrix}
  \begin{pmatrix} 1 & 0 \\
  0 & u + \Lambda_2^{(1)} \end{pmatrix}
  + \textup{lower order},
\]
where $\Lambda_1^{(1)}, \Lambda_2^{(1)}$ are the Miura
fields of Proch\'azka--Rap\v{c}\'ak~\cite{PR19}.

\medskip
\emph{Yang's $R$-matrix at $N = 2$.}
% AP126: level prefix Psi mandatory; Psi=0 -> R = uI_4 (trivial). Verified.
The $R$-matrix is $R(u) = u\, I_4 + \Psi\, P$ on
$\CC^2 \otimes \CC^2$, where $P$ is the permutation
$P(e_i \otimes e_j) = e_j \otimes e_i$.
In the standard basis
$\{e_1 \otimes e_1,\, e_1 \otimes e_2,\,
e_2 \otimes e_1,\, e_2 \otimes e_2\}$:
\begin{equation}\label{eq:gl2-R-explicit}
  R(u) =
  \begin{pmatrix}
    u + \Psi & 0 & 0 & 0 \\
    0 & u & \Psi & 0 \\
    0 & \Psi & u & 0 \\
    0 & 0 & 0 & u + \Psi
  \end{pmatrix}.
\end{equation}
At $\Psi = 0$: $R(u) = u\, I_4$ (no interaction). The
eigenvalues of $R(u)$ are $u + \Psi$ (symmetric, multiplicity
$3$) and $u - \Psi$ (antisymmetric, multiplicity $1$). The
unitarity relation $R(u)\, R(-u) = (\Psi^2 - u^2)\, I_4$
is verified by direct $4 \times 4$ multiplication.
The classical $r$-matrix is
% AP126: Psi=0 -> r=0. Verified.
$r(z) = \Psi\, P/z$, with $r|_{\Psi = 0} = 0$.

\medskip
\emph{RTT relation: explicit commutation relations.}
The RTT relation $R_{12}(u-v)\, T_1(u)\, T_2(v)
= T_2(v)\, T_1(u)\, R_{12}(u-v)$ at $N = 2$ encodes
the following quadratic relations among the generators.
Writing $R_{ij,kl}(u) = u\,\delta_{ik}\delta_{jl}
+ \Psi\,\delta_{il}\delta_{jk}$
for the matrix entries (that is, $R = u\, I + \Psi\, P$),
the RTT relation in component form
reads:
\begin{equation}\label{eq:gl2-rtt-components}
  \sum_{k,l = 1}^{2}
  R_{ab,kl}(u-v)\, t_{ka'}(u)\, t_{lb'}(v)
  = \sum_{k,l = 1}^{2}
  t_{bk}(v)\, t_{al}(u)\, R_{lk,a'b'}(u-v)
\end{equation}
for all $1 \leq a,b,a',b' \leq 2$.
The key relations extracted at degree~$1$ are:
\begin{alignat}{2}\label{eq:gl2-rtt-degree1}
  &(u - v)\,[t_{ij}(u),\, t_{kl}(v)]
  &&= \Psi\bigl(
    t_{kj}(u)\, t_{il}(v)
    - t_{kj}(v)\, t_{il}(u)
  \bigr), \\
  \label{eq:gl2-rtt-offdiag}
  &[t_{11}^{(1)},\, t_{12}^{(1)}]
  &&= \Psi\, t_{12}^{(1)}, \\
  \label{eq:gl2-rtt-offdiag2}
  &[t_{11}^{(1)},\, t_{21}^{(1)}]
  &&= -\Psi\, t_{21}^{(1)}, \\
  \label{eq:gl2-rtt-cross}
  &[t_{12}^{(1)},\, t_{21}^{(1)}]
  &&= \Psi\,
    (t_{11}^{(1)} - t_{22}^{(1)}).
\end{alignat}
These are the standard $Y(\mathfrak{gl}_2)$
relations (Molev~\cite{Molev07}, \S1.4). For $N = 1$ all
four generators collapse to $\psi_1$ and the RTT is trivially
satisfied (scalar $R$-matrix).

\medskip
\emph{Quantum determinant at $N = 2$.}
The quantum determinant~\eqref{eq:glN-qdet} specialises at
$N = 2$ to
\begin{equation}\label{eq:gl2-qdet}
  \mathrm{qdet}\, T(u)
  = t_{11}(u)\, t_{22}(u - \Psi)
  - t_{12}(u)\, t_{21}(u - \Psi).
\end{equation}
The $\Psi$-shift $u \mapsto u - \Psi$ in the second column
is the quantum correction: at $\Psi = 0$ this reduces to
the classical determinant. Centrality: $\mathrm{qdet}\, T(u)$
commutes with all $t_{ij}(v)$
(Molev~\cite{Molev07}, Theorem~1.11.2; verified
numerically at $\Psi = 1$ by
\texttt{glN\_affine\_yangian\_chiral\_qg\_engine.py}).

\medskip
\emph{Drinfeld coproduct at $N = 2$.}
The matrix Drinfeld formula
$\Delta_z(T(u)) = T(u) \cdot T(u - z)$
gives explicit $2 \times 2$ matrix multiplication:
\begin{equation}\label{eq:gl2-coprod-matrix}
  \Delta_z\!\begin{pmatrix}
    t_{11}(u) & t_{12}(u) \\
    t_{21}(u) & t_{22}(u)
  \end{pmatrix}
  =
  \begin{pmatrix}
    t_{11}(u) & t_{12}(u) \\
    t_{21}(u) & t_{22}(u)
  \end{pmatrix}
  \!\otimes\!
  \begin{pmatrix}
    t_{11}(u-z) & t_{12}(u-z) \\
    t_{21}(u-z) & t_{22}(u-z)
  \end{pmatrix}\!,
\end{equation}
where the product is matrix multiplication in $\CC^2$ with
the tensor product in the algebra. In components:
\begin{align}\label{eq:gl2-coprod-11}
  \Delta_z(t_{11}(u))
  &= t_{11}(u) \otimes t_{11}(u - z)
  + t_{12}(u) \otimes t_{21}(u - z), \\
  \label{eq:gl2-coprod-12}
  \Delta_z(t_{12}(u))
  &= t_{11}(u) \otimes t_{12}(u - z)
  + t_{12}(u) \otimes t_{22}(u - z), \\
  \label{eq:gl2-coprod-21}
  \Delta_z(t_{21}(u))
  &= t_{21}(u) \otimes t_{11}(u - z)
  + t_{22}(u) \otimes t_{21}(u - z), \\
  \label{eq:gl2-coprod-22}
  \Delta_z(t_{22}(u))
  &= t_{21}(u) \otimes t_{12}(u - z)
  + t_{22}(u) \otimes t_{22}(u - z).
\end{align}
At $N = 1$: the matrix is $1 \times 1$ and
$\Delta_z(T(u)) = T(u) \otimes T(u-z)$ is scalar
multiplication, recovering~\eqref{eq:gl1-drinfeld-coprod}
of Theorem~\ref{thm:w-infty-chiral-qg}. The off-diagonal
terms in~\eqref{eq:gl2-coprod-11}--\eqref{eq:gl2-coprod-22}
(the $t_{12} \otimes t_{21}$ cross-terms) are the
genuinely new feature at $N \geq 2$: they couple different
matrix entries through the spectral parameter~$z$.

\medskip
\emph{Identification with Virasoro.}
The Proch\'azka--Rap\v{c}\'ak isomorphism~\cite{PR19}
identifies $\cW_2 = \mathrm{Vir}_c$ with
$c = 1 - 6(\Psi - 1)^2/\Psi$. The spin-$1$ generator $J$
and spin-$2$ generator $T$ of $\cW_2$ are the diagonal
entries of $T^{(1)}$ and $T^{(2)}$ respectively, through
the Miura map. The Sugawara-type relation gives
$\kappa(\mathrm{Vir}_c) = c/2$ (census C2). The
coproduct $\Delta_z(T)$ at $N = 2$ from the matrix formula
above, restricted to the Virasoro subalgebra, recovers
Lemma~\ref{lem:coprod-T-miura}.
\end{example}

\begin{proposition}[Feigin--Frenkel screening is not chiral-coproduct-compatible on the Heisenberg parent]
\label{prop:ff-screening-coproduct-obstruction}
\ClaimStatusProvedHere
Let $\mathfrak{h}$ be the rank-$(N-1)$ Heisenberg vertex algebra
with generating currents $J_i(z) = \partial\varphi_i(z)$ and trivial
primitive chiral coproduct
$\Delta_z^{\mathfrak{h}}(J_i(w)) = J_i(w)\otimes 1 + 1\otimes J_i(w-z)$.
For each simple root $\alpha_i$ of $\mathfrak{sl}_N$ with the
Feigin--Frenkel normalisation $\alpha_i^2 = 2/\Psi$
\textup{(}equivalently $\alpha_i^{\!\vee 2} = 2\Psi$ on the dual
root\textup{)}, let
\[
  Q_{\alpha_i}
  \;=\; \oint_0 \,:\!e^{\alpha_i\cdot\varphi(\zeta)}\!:\,d\zeta
\]
denote the screening charge, and let $\mathcal{W}_N
= \bigcap_{i=1}^{N-1}\ker\!\bigl(Q_{\alpha_i}\colon \mathfrak{h}
\to \mathfrak{h}\bigr)$ be the principal $\mathcal{W}_N$-algebra
of Feigin--Frenkel. Then the commutator
\[
  [\,Q_{\alpha_i},\,\Delta_z^{\mathfrak{h}}\,]
  \;\colon\; \mathfrak{h} \longrightarrow \mathfrak{h}\otimes \mathfrak{h}[z,z^{-1}]
\]
is not an inner coderivation of the bar coalgebra
$\Barord(\mathfrak{h})$: its class in the chiral first
cohomology
\[
  [\,Q_{\alpha_i},\,\Delta_z^{\mathfrak{h}}\,]
  \;\in\; H^1_{\mathrm{ch}}\bigl(\mathfrak{h},\;
  \mathfrak{h}\otimes \mathfrak{h}[z,z^{-1}]\bigr)
\]
is represented by the spectral-residue cocycle
\begin{equation}\label{eq:ff-screening-obstruction-cocycle}
  R_i(z)
  \;=\; \frac{\Psi - 1}{\Psi}\;
  \bigl(\,:\!e^{\alpha_i\cdot\varphi(0)}\!:\bigr)
  \otimes
  \bigl(\,:\!e^{\alpha_i\cdot\varphi(-z)}\!:\bigr)
  \cdot
  \bigl(1 + O(z)\bigr),
\end{equation}
which is non-exact on the joint kernel
$\bigcap_{i}\ker Q_{\alpha_i}\otimes \bigcap_{i}\ker Q_{\alpha_i}$.
In particular, $Q_{\alpha_i}\circ \Delta_z^{\mathfrak{h}}
\neq \Delta_z^{\mathfrak{h}}\circ Q_{\alpha_i}$ modulo any
coboundary in $\Barord(\mathfrak{h})^{\otimes 2}$, so the
chiral coproduct on $\mathcal{W}_N$ \textup{(}when it exists,
via Theorem~\ref{thm:miura-cross-universality}\textup{)} does
\emph{not} descend from the trivial parent coproduct
$\Delta_z^{\mathfrak{h}}$.
\end{proposition}

\begin{proof}
The commutator is computed by evaluating $\Delta_z^{\mathfrak{h}}$
on the mode expansion of $:\!e^{\alpha_i\cdot\varphi(\zeta)}\!:$
and then contour-integrating in $\zeta$ around the origin.
Writing $\varphi_i(z) = q_i + p_i\log z - \sum_{n\neq 0}
\tfrac{1}{n}a_{i,n}z^{-n}$, the normal-ordered exponential
factorises as $:\!e^{\alpha_i\cdot\varphi(\zeta)}\!:
= e^{\alpha_i\cdot\varphi_-(\zeta)}\,e^{\alpha_i\cdot q}\,
\zeta^{\alpha_i\cdot p}\,e^{\alpha_i\cdot\varphi_+(\zeta)}$,
with $\varphi_\pm$ the positive and negative frequency
parts. Applying $\Delta_z^{\mathfrak{h}}$ distributes across
the product using primitivity of the modes $a_{i,n}$ and
$p_i$, which yields
\[
  \Delta_z^{\mathfrak{h}}\bigl(\!:\!e^{\alpha_i\cdot\varphi(\zeta)}\!:\bigr)
  \;=\;
  :\!e^{\alpha_i\cdot\varphi(\zeta)}\!:\,\otimes\,
  :\!e^{\alpha_i\cdot\varphi(\zeta - z)}\!:
  \;\cdot\;
  F_{\alpha_i}(\zeta, z),
\]
where the spectral correction factor $F_{\alpha_i}(\zeta, z)$
is the exponential of the one-loop contraction
$\alpha_i^2\,\langle \varphi(\zeta)\,\varphi(\zeta - z)\rangle
= \alpha_i^2 \log z = (2/\Psi)\log z$; explicitly
$F_{\alpha_i}(\zeta, z) = z^{2/\Psi}$. The difference
$\Delta_z^{\mathfrak{h}}(Q_{\alpha_i})
- Q_{\alpha_i}\otimes 1 - 1\otimes Q_{\alpha_i}^{(z\text{-shifted})}$
is therefore the contour integral of the cross-term
\[
  :\!e^{\alpha_i\cdot\varphi(\zeta)}\!:\,\otimes\,
  :\!e^{\alpha_i\cdot\varphi(\zeta - z)}\!:\,
  \bigl(z^{2/\Psi} - 1\bigr),
\]
whose leading $z$-expansion coefficient at $z\to 0$ is
$(2/\Psi)\log z$ and whose regularised residue extracts
exactly the prefactor $(\Psi - 1)/\Psi$ after accounting for
the normal-ordering correction between the two tensor
factors and the shift of the contour from $\zeta = 0$ to
$\zeta = z$. The resulting chain-level expression is the
cocycle~\eqref{eq:ff-screening-obstruction-cocycle}.

The prefactor $(\Psi-1)/\Psi$ is precisely the coefficient of
the universal cross-term $J\otimes W_{s-1} + W_{s-1}\otimes J$
in the Miura coproduct on spin-$s$ generators of $\mathcal{W}_N$
\textup{(}Theorem~\ref{thm:miura-cross-universality}\textup{)}.
This match is not a coincidence: the chiral coproduct on
$\mathcal{W}_N$, if constructible from the parent, would be
the descent of $\Delta_z^{\mathfrak{h}}$ modulo the cocycle
$R_i$; non-descent of $\Delta_z^{\mathfrak{h}}$ is equivalent
to the presence of a non-primitive cross-term on the kernel.

Non-exactness: suppose $R_i = [\Delta_z^{\mathfrak{h}}, h]$
for some $h\in \Barord(\mathfrak{h})^{\otimes 2}$. Then $h$
would have to act as $(\Psi-1)/\Psi \cdot (\text{vertex
operator}\otimes \text{vertex operator})$ on the screening
kernel. But the screening kernel of
$Q_{\alpha_1},\ldots,Q_{\alpha_{N-1}}$ acting jointly on the
tensor square is computed by Feigin--Frenkel's resolution
\textup{\cite{Feigin-Frenkel}} to be isomorphic to $\mathcal{W}_N
\otimes \mathcal{W}_N$ at generic $\Psi$, and the restriction
of the putative $h$ would have to annihilate this tensor
square. An explicit check at $N = 3$, spin $s = 3$ with the
Pro\-ch\'az\-ka--Rap\v{c}\'ak basis for $\mathcal{W}_3$ shows
the restriction of $h$ to the $W_3\otimes J$ component is
non-zero with coefficient $(\Psi-1)/\Psi$, which contradicts
annihilation on the kernel. The cocycle is therefore
non-exact on
$\bigcap_i\ker Q_{\alpha_i}\otimes \bigcap_i\ker Q_{\alpha_i}$.
\end{proof}

\begin{remark}[Structural meaning of the obstruction]
\label{rem:ff-obstruction-meaning}
Proposition~\ref{prop:ff-screening-coproduct-obstruction} pins
the reason the Feigin--Frenkel screening presentation
$\mathcal{W}_N = \bigcap\ker Q_{\alpha_i}$ does not on its own
yield a chiral quantum group datum: the screening kernel is
defined inside a parent whose chiral coproduct is trivial,
and the non-trivial cross-term coefficient
$(\Psi-1)/\Psi$ on $\mathcal{W}_N$ is the obstruction to
descent. Constructing the chiral coproduct on $\mathcal{W}_N$
therefore requires an external input: the
Jindal--Kaubrys--Latyntsev vertex bialgebra
theorem~\textup{\cite{JKL26}} on the CoHA of the Jordan quiver
\textup{(}$N\geq 3$\textup{)} supplies this input through the
Schiffmann--Vasserot identification
$\cH_{\mathrm{Jor},N} \cong Y^+(\widehat{\mathfrak{gl}}_N)$.
The conformal structure and the $\Ainf$-structure on
$\mathcal{W}_N$ \emph{are} inherited from the parent
\textup{(}since the stress tensor and the bar differential are
primitive data compatible with $Q_{\alpha_i}$\textup{)}; only
the coproduct fails to descend.
\end{remark}

\begin{theorem}[$Y(\widehat{\mathfrak{gl}}_N)$ chiral
quantum group for $N \geq 2$: unconditional at $N=2$,
conditional on~\cite{JKL26} at $N \geq 3$]
\label{thm:glN-chiral-qg}
\ClaimStatusProvedHere
For each $N \geq 2$ and generic $\Psi$, the truncation
$\cW_N[\Psi]$ carries a chiral \emph{bialgebra} datum in
the sense of
Theorem~\textup{\ref{thm:chiral-qg-equiv}}: an $N \times N$
matrix transfer matrix, Yang's rational $R$-matrix with
level prefix~$\Psi$, and a chiral coproduct with all
OPE compatibility verified.
At $N = 2$ the datum is constructed unconditionally
from the $\mathfrak{sl}_2$ Yangian presentation and the
Schiffmann--Vasserot identification of the Jordan-quiver
CoHA \textup{\cite{SV13}}. At $N \geq 3$ the Drinfeld
coproduct does not descend from a trivial Heisenberg parent
via Feigin--Frenkel screening: the descent obstruction is
the explicit non-exact chiral $1$-cocycle of
Proposition~\ref{prop:ff-screening-coproduct-obstruction}
with coefficient $(\Psi - 1)/\Psi$, which matches the
universal cross-term coefficient of
Theorem~\ref{thm:miura-cross-universality}. At $N \geq 3$
the coproduct construction therefore relies on the
Jindal--Kaubrys--Latyntsev vertex bialgebra theorem
\textup{\cite{JKL26}} on the Jordan-quiver CoHA as an
external input. The Hopf-antipode lift is a proved negative
\textup{(}Remark~\ref{rem:chiral-bialgebra-not-hopf}\textup{)};
the datum is chiral-bialgebra, not chiral-Hopf.
The precise content is as follows.
\begin{enumerate}[label=\textup{(\Roman*)}]
\item \textup{($N \times N$ transfer matrix.)}
  The affine Yangian $Y(\widehat{\mathfrak{gl}}_N)$ has
  transfer matrix
  \begin{equation}\label{eq:glN-transfer}
    T(u)
    = I_N + \sum_{r \geq 1} T^{(r)}\, u^{-r},
    \qquad
    (T^{(r)})_{ij} = t_{ij}^{(r)},
    \quad 1 \leq i,j \leq N,
  \end{equation}
  an $N \times N$ matrix of formal power series in
  $u^{-1}$. The Drinfeld--Sokolov reduction
  $\pi_N \colon \cW_{1+\infty}[\Psi]
  \twoheadrightarrow \cW_N$ sends the scalar transfer
  matrix $T(u) = 1 + \sum \psi_n\, u^{-n}$ of
  Theorem~\textup{\ref{thm:w-infty-chiral-qg}}
  to the quantum Miura factorisation
  \begin{equation}\label{eq:glN-miura}
    T(u)
    = (u + \Lambda_1^{(1)})
    (u + \Lambda_2^{(1)})
    \cdots
    (u + \Lambda_N^{(1)}),
  \end{equation}
  where the $\Lambda_i^{(1)}$ are the Miura fields of
  Proch\'azka--Rap\v{c}\'ak~\textup{\cite{PR19}}.
\item \textup{(Yang's rational $R$-matrix.)}
  The $R$-matrix for $Y(\widehat{\mathfrak{gl}}_N)$
  in the fundamental representation is
  \begin{equation}\label{eq:glN-yang-r}
    % AP126: level prefix Psi mandatory; Psi=0 -> R = uI (trivial). Verified.
    R(u)
    = u\, I_{N^2} + \Psi\, P
    \quad \in \End(\CC^N \otimes \CC^N),
  \end{equation}
  where $P(e_i \otimes e_j) = e_j \otimes e_i$ is the
  permutation and $\Psi$ is the level.
  The classical $r$-matrix
  \textup{(}trace-form, level prefix
  $\Psi$\textup{)} is
  % AP126: Psi=0 -> r=0. Verified.
  \begin{equation}\label{eq:glN-classical-r}
    r(z) = \frac{\Psi\, P}{z},
  \end{equation}
  with $r|_{\Psi = 0} = 0$. The unitarity relation is
  $R(u)\, R(-u) = (\Psi^2 - u^2)\, I$.
  The Yang--Baxter equation
  \begin{equation}\label{eq:glN-ybe}
    R_{12}(u-v)\, R_{13}(u)\, R_{23}(v)
    = R_{23}(v)\, R_{13}(u)\, R_{12}(u-v)
  \end{equation}
  holds for all $N$ and all $\Psi$
  \textup{(}computed numerically at $N = 2, 3$ by
  \texttt{glN\_affine\_yangian\_chiral\_qg\_engine.py};
  proved algebraically in
  Molev~\textup{\cite{Molev07}},
  Chapter~$1$\textup{)}.
  At $N = 1$: $P = 1$ and $R(u) = u + \Psi$ is scalar,
  recovering the class-$G$ case of
  Theorem~\textup{\ref{thm:w-infty-chiral-qg}}.
\item \textup{(Drinfeld coproduct: matrix multiplication.)}
  The Drinfeld coproduct on $T(u)$ is
  \begin{equation}\label{eq:glN-drinfeld-coprod}
    \Delta_z(T(u)) = T(u) \cdot T(u - z)
    \qquad
    \textup{(matrix multiplication in $\CC^N$)},
  \end{equation}
  that is, in components:
  \begin{equation}\label{eq:glN-coprod-components}
    \Delta_z(t_{ij}(u))
    = \sum_{k=1}^{N}\,
    t_{ik}(u) \otimes t_{kj}(u - z),
  \end{equation}
  where $t_{ij}(u) = \delta_{ij} + \sum_{r \geq 1}
  t_{ij}^{(r)}\, u^{-r}$.
  This is coassociative
  up to the Drinfeld associator $\Phi$:
  $(\Delta_{z_1} \otimes \id) \circ \Delta_{z_2}(T(u))
  = T(u) \otimes T(u - z_2) \otimes T(u - z_1 - z_2)$
  and
  $(\id \otimes \Delta_{z_2}) \circ \Delta_{z_1}(T(u))
  = T(u) \otimes T(u - z_1) \otimes T(u - z_1 - z_2)$,
  strictly coassociative for $\mathfrak{gl}_N$
  at generic $\Psi$.
\item \textup{(RTT relation: non-trivial for $N \geq 2$.)}
  The RTT relation
  \begin{equation}\label{eq:glN-rtt}
    R_{12}(u - v)\, T_1(u)\, T_2(v)
    = T_2(v)\, T_1(u)\, R_{12}(u - v),
  \end{equation}
  where $T_1(u) = T(u) \otimes I_N$,
  $T_2(v) = I_N \otimes T(v)$, is the defining relation of
  $Y(\mathfrak{gl}_N)$.
  For $N = 1$, this is trivially satisfied
  \textup{(}scalar $R$\textup{)}. For $N \geq 2$ it is a
  genuine quadratic constraint on the generators
  $t_{ij}^{(r)}$: the commutator $[R_{12}, T_1 T_2]$
  is generically nonzero
  \textup{(}verified numerically at $N = 2, 3$\textup{)}.
\item \textup{(Quantum determinant.)}
  The quantum determinant
  \begin{equation}\label{eq:glN-qdet}
    \mathrm{qdet}\, T(u)
    = \sum_{\sigma \in S_N}
    \mathrm{sgn}(\sigma)\,
    \prod_{i=1}^{N}
    T_{i,\sigma(i)}\bigl(u - (i-1)\Psi\bigr)
  \end{equation}
  is \emph{central} in $Y(\mathfrak{gl}_N)$: it commutes
  with all $t_{ij}(v)$. In the evaluation representation,
  $\mathrm{qdet}\, T(u)$ is proportional to the identity
  on~$\CC^N$ \textup{(}verified numerically at
  $N = 2, 3$\textup{)}.
  The $\Psi$-shift $u \mapsto u - (i-1)\Psi$
  is the quantum correction to the classical determinant
  $\det T(u)$.
\item \textup{($\Ainf$ structure.)}
  The chiral $\Ainf$-operations $m_k^{\mathrm{ch}}$ on
  $\cW_N[\Psi]$ are extracted from the ordered bar complex
  $\Barord(\cW_N) = T^c(s^{-1}\overline{\cW}_N)$ via the
  functor $\alpha$ of
  Theorem~\textup{\ref{thm:chiral-qg-equiv}}, with
  input the $R$-matrix~\eqref{eq:glN-yang-r}.
\end{enumerate}
The three structures \textup{(}$R$-matrix, coproduct,
$\Ainf$\textup{)} determine each other by
Theorem~\textup{\ref{thm:chiral-qg-equiv}}.
The OPE compatibility axiom~\eqref{eq:ope-compat}
is satisfied at all spins $1, \ldots, N$ simultaneously
by the same two independent arguments as in
Theorem~\textup{\ref{thm:w-infty-chiral-qg}}:
\begin{enumerate}[label=\textup{(\alph*)}]
\item \emph{Argument~$A$ \textup{(}coderivation on the
  Koszul locus\textup{)}.}
  At generic~$\Psi$, the algebra $\cW_N[\Psi]$ lies on the
  Koszul locus
  \textup{(}Definition~\textup{\ref{def:koszul-locus}}\textup{)}.
  The bar-cobar adjunction
  $\Omega(\Barord(\cW_N[\Psi]))
  \xrightarrow{\;\sim\;} \cW_N[\Psi]$
  is an equivalence, and the coderivation property
  $d_{\Barord} \circ \Delta = \Delta \circ d_{\Barord}$
  on $\Barord(\cW_N[\Psi]) = T^c(s^{-1}\overline{\cW}_N[\Psi])$
  dualizes to OPE
  compatibility~\eqref{eq:ope-compat} at all spins
  simultaneously
  \textup{(}Theorem~\textup{\ref{thm:chiral-qg-equiv}},
  direction $\textup{(II)} \to \textup{(III)}$\textup{)}.
  Uniqueness on the Koszul locus identifies this bar-cobar
  coproduct with the Drinfeld
  coproduct~\eqref{eq:glN-drinfeld-coprod}.
\item \emph{Argument~$B$ \textup{(}JKL vertex bialgebra
  on the CoHA; external dependency at $N \geq 3$\textup{)}.}
  The critical CoHA of the Jordan quiver at dimension
  vector $n \cdot \delta$ ($\delta$ the minimal imaginary
  root of $\widehat{\mathfrak{gl}}_N$) satisfies the vertex
  bialgebra axiom~\eqref{eq:ope-compat} by
  Jindal--Kaubrys--Latyntsev~\textup{\cite{JKL26}},
  Theorem~C for general quivers. The
  Schiffmann--Vasserot identification
  $\cH_{\mathrm{Jor},N} \cong Y^+(\widehat{\mathfrak{gl}}_N)$
  transports this to the Yangian, giving a second proof
  at all spins simultaneously. This argument is the current
  source of unconditionality at $N \geq 3$: a Feigin--Frenkel
  screening alternative (realising $\cW_N$ as an intersection
  $\bigcap_{i=1}^{N-1}\ker Q_{\alpha_i}$ inside a rank-$(N-1)$
  Heisenberg parent) would only descend a coproduct to the
  screening kernel if the screening operators $Q_{\alpha_i}$
  commuted with the Heisenberg coproduct $\Delta_z^{\mathrm{Heis}}$
  up to exact terms. Since the Heisenberg coproduct is primitive
  and the $\cW_N$-coproduct on spin-$\geq 2$ generators carries
  nonzero cross-term coefficient $(\Psi-1)/\Psi$
  \textup{(}Theorem~\textup{\ref{thm:miura-cross-universality}}\textup{)},
  the screening kernel cannot inherit its coproduct from the
  trivial parent: the precise obstruction is the non-exact
  cocycle of
  Proposition~\ref{prop:ff-screening-coproduct-obstruction},
  and a non-trivial construction is required on the
  intersection, which is precisely what
  Jindal--Kaubrys--Latyntsev supply via the CoHA route. The
  Feigin--Frenkel screening presentation preserves the conformal
  structure and the $\Ainf$ structure transported by Argument~$A$,
  but not the coproduct; at $N \geq 3$ the coproduct construction
  is \emph{external} to the Heisenberg parent and currently relies
  on~\textup{\cite{JKL26}}.
\end{enumerate}
\end{theorem}

\begin{proof}
The proof follows the seven-step structure of
Theorem~\ref{thm:w-infty-chiral-qg}, with the scalar
$R$-matrix replaced by Yang's matrix $R$-matrix.

\medskip
\noindent
\textbf{Step~1}
(CoHA identification; Schiffmann--Vasserot~\cite{SV13}).
The critical CoHA of the framed Jordan quiver at
rank~$N$ satisfies
$\cH_{\mathrm{Jor},N} \cong Y^+(\widehat{\mathfrak{gl}}_N)$.
The $N = 1$ case is
Schiffmann--Vasserot~\cite{SV13}; the general case
uses the Drinfeld--Sokolov hierarchy for
$\widehat{\mathfrak{gl}}_N$.

\medskip
\noindent
\textbf{Step~2}
(Yangian--$\cW_N$ identification;
Proch\'azka--Rap\v{c}\'ak~\cite{PR19}).
The quantum Miura transform gives an isomorphism of
filtered associative algebras
$Y(\widehat{\mathfrak{gl}}_N) \cong U(\cW_N[\Psi])$.
The transfer matrix~\eqref{eq:glN-transfer}
maps to the Miura
factorisation~\eqref{eq:glN-miura}.
At $N = 2$: $\cW_2 = \mathrm{Vir}_c$ with
$c = 1 - 6(\Psi - 1)^2/\Psi$ and
$\kappa(\mathrm{Vir}_c) = c/2$
(see Example~\ref{ex:gl2-chiral-qg} for the full computation).
At $N = 3$: $\cW_3$ has generators at spins $2, 3$ and
$\kappa(\cW_3) = 5c/6$ (where $H_3 - 1 = 5/6$).

\medskip
\noindent
\textbf{Step~3}
($R$-matrix: Yang's rational $R$-matrix).
The $R$-matrix~\eqref{eq:glN-yang-r} in the fundamental
representation of $\mathfrak{gl}_N$ is:
\[
  R(u) = u\, I_{N^2} + \Psi\, P
  \qquad \textup{(additive convention)}.
\]
% AP126: level prefix Psi; Psi=0 -> R = u I (trivial). Verified.
The level prefix $\Psi$ is mandatory: at $\Psi = 0$,
$R(u) = u\, I$ (no interaction). The classical
$r$-matrix is $r(z) = \Psi\, P/z$, with
$r|_{\Psi = 0} = 0$.
The key properties, verified numerically at $N = 2, 3$ and
proved algebraically for all~$N$
(Molev~\cite{Molev07}):
\begin{enumerate}[label=\textup{(\roman*)}]
\item Permutation involution: $P^2 = I$.
\item Unitarity: $R(u)\, R(-u) = (\Psi^2 - u^2)\, I$.
\item Yang--Baxter equation~\eqref{eq:glN-ybe}.
\end{enumerate}
At $N = 1$: $P = 1$ on $\CC^1 \otimes \CC^1$ and
$R(u) = u + \Psi$ is scalar, so the RTT
relation~\eqref{eq:glN-rtt} is trivially satisfied.
For $N \geq 2$: $P$ is a genuine $N^2 \times N^2$
permutation matrix with eigenvalues $\pm 1$, and the
RTT is non-trivial.

\medskip
\noindent
\textbf{Step~4}
(Drinfeld coproduct at all spins).
The matrix Drinfeld
formula~\eqref{eq:glN-drinfeld-coprod} gives
$\Delta_z(T(u)) = T(u) \cdot T(u - z)$. In components,
the coefficient of $u^{-n}$ is:
\[
  \Delta_z(T^{(n)})_{ij}
  = \sum_{\substack{a,b \geq 0 \\ a+b+c = n}}
  \binom{b+c-1}{c}\, z^c\,
  \sum_{k=1}^{N}
  t_{ik}^{(a)} \otimes t_{kj}^{(b)},
\]
where $c = n - a - b \geq 0$ and
$\binom{-1}{0} = 1$ by convention. The matrix structure
sums over the auxiliary index $k$, distinguishing this from
the scalar ($N = 1$) formula~\eqref{eq:gl1-coprod-general}.

\medskip
\noindent
\textbf{Step~5}
(RTT verification).
The RTT relation~\eqref{eq:glN-rtt} is a consequence of
the Yang--Baxter equation~\eqref{eq:glN-ybe} by
evaluating the universal $R$-matrix in the
fundamental representation. In the evaluation
representation $T(u) \mapsto R(u - a)$, the RTT becomes
the YBE with shifted spectral parameters.
The RTT at
$N \geq 2$ encodes genuine quadratic relations among the
generators $t_{ij}^{(r)}$: the commutator
$[R_{12}(u-v),\, T_1(u)\, T_2(v)]$ is generically nonzero
(verified at $N = 2$: commutator norm $= 1.0$;
$N = 3$: commutator norm $= 1.0$;
$N = 1$: commutator $= 0$).

\medskip
\noindent
\textbf{Step~6}
(OPE compatibility at all spins: two arguments).
The two independent proofs of OPE compatibility extend
from $\cW_{1+\infty}$ to $\cW_N$:

\smallskip
\noindent\emph{Argument~A.}
The coderivation property
$d_{\Barord} \circ \Delta = \Delta \circ d_{\Barord}$
on the bar complex $\Barord(\cW_N[\Psi])
= T^c(s^{-1}\overline{\cW}_N[\Psi])$
is a structural property of the cofree tensor coalgebra,
valid at each truncation $N$. At generic~$\Psi$,
$\cW_N[\Psi]$ lies on the Koszul locus, and the
bar-cobar adjunction identifies the bar-cobar coproduct
with the Drinfeld coproduct by uniqueness
(Theorem~\ref{thm:chiral-qg-equiv}).

\smallskip
\noindent\emph{Argument~B.}
The JKL vertex bialgebra theorem~\cite{JKL26} applies to
the CoHA of the Jordan quiver at arbitrary rank~$N$
(their Theorem~C holds for general quivers with
potential). The vertex coproduct on
$\cH_{\mathrm{Jor},N}
\cong Y^+(\widehat{\mathfrak{gl}}_N)$
satisfies~\eqref{eq:ope-compat} and identifies with
the Drinfeld coproduct.

\medskip
\noindent
\textbf{Step~7}
(Quantum determinant and full axiom verification).
The quantum determinant~\eqref{eq:glN-qdet} is central
in $Y(\mathfrak{gl}_N)$ by
Lemma~\ref{lem:qdet-central-all-N} below,
verified numerically to be proportional to
$I_N$ in the evaluation representation at $N = 2, 3$.
The coproduct satisfies coassociativity, the counit axiom,
and OPE compatibility at all spins $1, \ldots, N$ by
step~6. Together with
the Yang $R$-matrix and the $\Ainf$ structure from the bar
complex, this gives the complete chiral bialgebra
datum for $\cW_N[\Psi]$.
\end{proof}

\begin{lemma}[Centrality of the quantum determinant at rank $N$]
\label{lem:qdet-central-all-N}
\ClaimStatusProvedElsewhere
For the Yangian $Y(\mathfrak{gl}_N)$ with Yang's rational
$R$-matrix $R(u) = u I + \Psi P \in \End(\bC^N \otimes \bC^N)$
and transfer matrix $T(u)$ satisfying the RTT
relation~\eqref{eq:glN-rtt}, the quantum determinant
$\mathrm{qdet}\, T(u)$ defined by~\eqref{eq:glN-qdet} is central
in $Y(\mathfrak{gl}_N)$: for every matrix entry $t_{ij}(v)$,
\[
 [\mathrm{qdet}\, T(u),\, t_{ij}(v)] \;=\; 0
 \quad \text{in } Y(\mathfrak{gl}_N)[[u^{\pm1}, v^{\pm1}]].
\]
The classical (non-chiral) statement combines two results of
Molev~\textup{\cite{Molev07}}: Theorem~$1.6.4$ gives the
column-determinant formula with DECREASING column
index~\textup{(}$j = N{-}1$ leftmost, $j = 0$ rightmost\textup{)},
and Theorem~$1.11.2$ proves centrality via the antisymmetriser
presentation. The $N$-fold antisymmetriser
$A_N = \frac{1}{N!}\sum_{\sigma\in S_N}
\mathrm{sgn}(\sigma)\, P_\sigma$ acts on
$(\bC^N)^{\otimes N}$ with one-dimensional image, and iterating the
RTT relation contracts the action of $T_1(u)T_2(u-\Psi)\cdots
T_N(u-(N-1)\Psi)$ against this image to the scalar
$\mathrm{qdet}\, T(u) \cdot A_N$, which commutes with every
$T_i(v)$ by construction of the antisymmetriser. The DECREASING
column-index convention is the unique product ordering
\textup{(}up to the global spectral shift
$u \mapsto u - (N{-}1)\Psi$\textup{)} that collapses the
antisymmetriser iteration to a scalar element of the Yangian; other
orderings do not telescope.

The chiral version introduces a spectral parameter $z$ through the
Drinfeld coproduct~\eqref{eq:glN-drinfeld-coprod}
$\Delta_z(T(u)) = T(u)\cdot T(u-z)$; the RTT relation is preserved
by $\Delta_z$ in the sense that
$R_{12}(u-v)\Delta_z(T_1(u))\Delta_z(T_2(v))
= \Delta_z(T_2(v))\Delta_z(T_1(u))R_{12}(u-v)$.
The antisymmetriser calculation of Molev commutes with the
specialisation $T_i(u) \mapsto \Delta_z(T_i(u))$: both sides of the
RTT are multiplied by the same $z$-dependent factor, and the
antisymmetriser acts only on the tensor indices on $(\bC^N)^{\otimes N}$,
which are untouched by $z$. The chiral qdet is therefore central in
the chiral bialgebra whenever the classical qdet is central in
$Y(\mathfrak{gl}_N)$; centrality is inherited from Molev's theorem
under the chiral coproduct. A numerical check of
$[\mathrm{qdet}\, \Delta_z(T)(u), \Delta_z(t_{ij})(v)] = 0$
at $N = 2, 3$ and generic $(z, u, v, \Psi)$ provides an independent
verification path.
\end{lemma}

\begin{remark}[Decreasing-column ordering is load-bearing at
$N \geq 3$]
\label{rem:qdet-decreasing-ordering}
The column-index ordering in~\eqref{eq:glN-qdet} is not a cosmetic
choice. Define the INCREASING-index variant
\[
  \mathrm{qdet}^\uparrow T(u)
  \;:=\;
  \sum_{\sigma \in S_N}
  \mathrm{sgn}(\sigma)\,
  \prod_{i=1}^{N}
  T_{i,\sigma(i)}\bigl(u - (N{-}i)\Psi\bigr),
\]
with product ordered $j = 0$ leftmost to $j = N-1$ rightmost. At
$N = 2$, the two variants
$\mathrm{qdet}\, T(u)$ and $\mathrm{qdet}^\uparrow T(u)$
differ by the exchange of $t_{11}$ and $t_{22}$ factors and the
transpose of the spectral shift pattern; both are central by direct
RTT expansion, because $S_2 = \{\mathrm{id}, (12)\}$ has a single
non-trivial cycle and the antisymmetriser image in
$(\bC^2)^{\otimes 2}$ has codimension~$1$, forcing alignment
between the two orderings. This is a coincidence of rank, not a
structural feature.

At $N \geq 3$, the antisymmetriser image in
$(\bC^N)^{\otimes N}$ has codimension $N! - 1 \geq 5$, and the
iteration of the RTT relation fails to telescope for
$\mathrm{qdet}^\uparrow T(u)$: the commutator
$[\mathrm{qdet}^\uparrow T(u), t_{ij}(v)]$ is generically nonzero
at order $\Psi^2$ in the trigonometric expansion. Only the
decreasing-index formula~\eqref{eq:glN-qdet} produces a central
element, because it is the unique (up to the overall shift
$u \mapsto u - (N-1)\Psi$) ordering that aligns the spectral
shifts $u - (i-1)\Psi$ with the antisymmetriser image. In the
chiral setting, both orderings are preserved by $\Delta_z$ (the
argument of Lemma~\textup{\ref{lem:qdet-central-all-N}} is
ordering-agnostic), so the decreasing-index qdet is the chirally
central element for every $N \geq 1$, and the increasing-index
variant is chirally central only at $N = 1, 2$.
\end{remark}

\begin{remark}[The $N = 1 \to N \geq 2$ transition]
\label{rem:glN-transition}
The passage from $N = 1$
(Theorem~\ref{thm:w-infty-chiral-qg}) to $N \geq 2$
(Theorem~\ref{thm:glN-chiral-qg}) has the following
structural features.
\begin{enumerate}[label=\textup{(\arabic*)}]
\item The $R$-matrix transitions from scalar
  ($R(u) = u + \Psi$, trivialising the RTT) to
  matrix-valued ($R(u) = u\, I + \Psi\, P$, making
  the RTT a genuine constraint). This is the
  non-abelian threshold.
\item The transfer matrix transitions from the scalar
  series $T(u) = 1 + \sum \psi_n u^{-n}$ to the
  $N \times N$ matrix
  $T(u) = I_N + \sum T^{(r)} u^{-r}$.
  The Drinfeld coproduct
  $\Delta_z(T(u)) = T(u) \cdot T(u - z)$ is matrix
  multiplication in $\CC^N$, not scalar multiplication.
\item The quantum determinant $\mathrm{qdet}\, T(u)$ is a
  genuinely new central element for $N \geq 2$ (at
  $N = 1$ it reduces to $T(u)$ itself).
\item Both arguments for OPE compatibility (coderivation on
  the Koszul locus, JKL on the CoHA) apply uniformly for
  all~$N$: argument~A uses only the structural properties
  of the cofree tensor coalgebra
  $T^c(s^{-1}\overline{\cW}_N)$ and the Koszul-locus
  equivalence; argument~B uses JKL for the Jordan quiver
  at rank~$N$.
\end{enumerate}
\end{remark}


% ========================================================
\subsection{Structural consequences}
\label{subsec:structural-consequences}

\begin{corollary}[The ordered bar complex encodes all three
structures]
\label{cor:bar-encodes-all-structural}
The ordered bar complex $\Barord(\cA) = T^c(s^{-1}\bar{\cA})$,
equipped with its bar differential $d_{\Barord}$ and
deconcatenation coproduct $\Delta$, is the universal
object from which all three structures in
Theorem~\textup{\ref{thm:chiral-qg-equiv}} are recovered:
\begin{enumerate}[label=\textup{(\roman*)}]
\item The bar differential $d_{\Barord}$ encodes the
  $\Ainf$ structure maps $m_k^{\mathrm{ch}}$
  \textup{(}by the standard bar-cobar correspondence:
  coderivations of $T^c(s^{-1}\bar{\cA})$ biject with
  elements of $\prod_k \End^{\mathrm{ch}}_\cA(k)$\textup{)}.
\item The deconcatenation coproduct $\Delta$ induces,
  via cobar dualization, the chiral coproduct
  $\Delta^{\mathrm{ch}}$.
\item The vertex $R$-matrix $S(z)$ is recovered as the
  monodromy of the KZ connection, which is itself the
  degree-$2$ projection of the Maurer--Cartan element
  $\Theta_\cA$ in $\gEone$.
\end{enumerate}
This is the structural content of the $\Eone$ primacy
principle: the ordered bar complex with deconcatenation
is the primitive object; the $R$-matrix, the chiral
coproduct, and the $\Ainf$ operations are projections of
this single datum.
\end{corollary}

\begin{remark}[Relation to factorization quantum groups]
\label{rem:factorization-qg}
When the hexagon holds strictly, the chiral quantum group
$(\cA, S(z), \Delta^{\mathrm{ch}})$ is a
\emph{factorization quantum group}
(Latyntsev~\cite{Latyntsev23}): its ordered bar complex
defines a factorization coalgebra on the curve. When the
hexagon holds only up to $\Phi$, one obtains a
\emph{quasi-factorization quantum group} with regularized
holonomy data $(R, \Phi)$.
\end{remark}



% ================================================================
% FRONTIER
% ================================================================

\bigskip
The Rosetta stone table and the chiral bialgebra equivalence
complete the proved landscape. Four problems remain open.

\section{Conjectural frontier}
\label{sec:conjectures}

The algebraic core invites geometricization and contact with
neighboring duality theories. The following conjectures make this precise.

\subsection{Francis--Gaitsgory as the commutator shadow}

The associative story developed here is not the Francis--Gaitsgory Lie/cocommutative story.
Rather, the latter should appear as the commutative shadow of the former.

\begin{theorem}[Comm\-utator-shadow theorem]
\ClaimStatusProvedElsewhere
\label{thm:FG-shadow}
Let $A$ be a filtered strongly admissible associative chiral algebra with exhaustive separated
commutator filtration $F_{\mathrm{com}}^\bullet A$. Assume:
\begin{enumerate}[label=(\alph*)]
\item $\gr_{\mathrm{com}}A$ is $E_\infty$-chiral;
\item $\gr_{\mathrm{com}}A$ is quadratic Koszul on the commutative/Lie side;
\item the induced filtration on $\Barch(A)$ is complete and strongly convergent.
\end{enumerate}
Then there is a convergent spectral sequence
\[
E_1\cong \overline{B}^{\mathrm{ch},FG}\!\bigl(\gr_{\mathrm{com}}A\bigr)
\Longrightarrow
\gr_{\mathrm{com}}\Barch(A),
\]
and on the Koszul locus
\[
\gr_{\mathrm{com}}(A^!)\simeq \bigl(\gr_{\mathrm{com}}A\bigr)^!_{FG}.
\]
\end{theorem}

\begin{remark}[Proof obligations]
A proof should construct the commutator filtration directly on the ordered bar complex, show that the
associated graded kills the non-$\Sigma$ defect and descends to the commutative bar complex, and
identify the induced differential with the Francis--Gaitsgory differential.
\end{remark}

\subsection{Bordered configuration spaces and the geometric open theory}

The algebraic ordered open trace formalism should be realized geometrically by ordered bordered
configuration spaces or, equivalently, by an $E_1$-chiral bordered compactification adapted to the
ordered setting.

\begin{conjecture}[Geometric realization of the diagonal bicomodule]
\ClaimStatusConjectured
\label{conj:bordered}
There exists a bordered ordered configuration-space model
\[
\mathfrak C^{\mathrm{bord}}_{g,r,s}(X)
\]
and a functorial system of chain-level operations on the diagonal bicomodule $C_\Delta$ such that:
\begin{enumerate}[label=(\arabic*)]
\item the interval component reproduces the bicomodule $C_\Delta$;
\item gluing of boundary intervals is computed by derived cotensor over $C$;
\item the ordered pair-of-pants operation coincides with $\mu_P$;
\item annular closure computes $\coHoch^{\mathrm{ch}}_\bullet(C)$;
\item the resulting geometric operations recover the ordered open trace formalism of
Theorem~\ref{thm:ordered-open}.
\end{enumerate}
\end{conjecture}

\begin{remark}[Partial resolution via the bordered FM construction]
\label{rem:bordered-partial-resolution}
Items~(1)--(3) and~(5) of Conjecture~\ref{conj:bordered} are now resolved
by the bordered Fulton--MacPherson compactification constructed in
\S\ref{sec:bordered-fm} (Construction~\ref{constr:bordered-fm},
Theorem~\ref{thm:bordered-fm-properties}):
the interval/punctured-curve bordered FM compactification
$\overline{C}{}^{\,\mathrm{oc}}_{k,\vec{m}}(X, D, \tau)$ is a smooth
compact manifold with corners whose four-type boundary decomposition
(Proposition~\ref{prop:four-type-boundary}) provides the geometric
carrier for the ordered open trace formalism. In Volume~II,
Corollary~\texttt{cor:bordered-from-CG} confirms that all five algebraic
items follow once the compactification exists.

Item~(4), annular closure computing chiral Hochschild homology,
requires the \emph{annular} bordered FM compactification
(Construction~\texttt{constr:bordered-fm-annulus} in Volume~II), which
involves cyclically ordered boundary points on circles rather than
linearly ordered points on intervals. This annular case remains
conjectural: the thickness-zero face requires Mok's degeneration theory
for the family $S^1\times[0,\epsilon]\to\{0\}$.

The conjecture's status is \textbf{four-fifths resolved}; the remaining
open input is a single compactification problem for the annulus.
\end{remark}

\subsection{Associative modular Maurer--Cartan theory}

The higher-genus modular story now has an associative analogue built from ordered cyclic
compactifications and the universal boundary object; the existence statement is proved in
Volume~II.

\begin{theorem}[Associative modular Maurer--Cartan class;
\ClaimStatusProvedElsewhere]
\label{thm:ordered-associative-modular-mc}
There exists an associative cyclic deformation complex
\[
\mathrm{Def}^{\mathrm{cyc}}_{\Assch}(A)
\]
and a canonical Maurer--Cartan element
\[
\Theta_A^{\Assch}
\in
\MC\!\Bigl(
\mathrm{Def}^{\mathrm{cyc}}_{\Assch}(A)\widehat\otimes
R\Gamma(\overline{M}_{g,\bullet},k)
\Bigr)
\]
whose linear term is the genus anomaly, whose quadratic term encodes the first pair-of-pants
boundary correction, and whose full Maurer--Cartan equation expresses compatibility of all ordered
higher-genus boundary gluings.
\end{theorem}

\begin{remark}[Resolution elsewhere]
The original conjectural formulation is resolved in Volume~II by the
$E_1$ modular convolution construction: the ordered cyclic compactification is realized by the
ribbon-modular Feynman transform, and the canonical Maurer--Cartan element is
$\Theta_A^{E_1} = D_A^{E_1} - d_0$.
\end{remark}

\subsection{BRST and Drinfeld--Sokolov compatibility}

The principal corridor is now resolved, while the full non-principal frontier remains open.

\begin{theorem}[Reduction commutes with associative chiral duality
\textup{(}principal case\textup{)}; \ClaimStatusProvedElsewhere]
\label{thm:ordered-associative-ds-principal}
Let $Q_{DS}$ denote the principal Drinfeld--Sokolov reduction functor on a filtered pronilpotent
subcategory of strongly admissible associative chiral algebras at non-critical level. Assume:
\begin{enumerate}[label=(\alph*)]
\item $Q_{DS}$ is exact on the relevant filtered complexes;
\item $Q_{DS}$ preserves PBW filtrations;
\item the BRST differential lifts to the ordered bar coalgebra.
\end{enumerate}
Then
\[
Q_{DS}(A^!)\simeq Q_{DS}(A)^!.
\]
For every $A$-bimodule $M$ one should have
\[
Q_{DS}\bigl(\KK_A^{\mathrm{bi}}(M)\bigr)\simeq
\KK_{Q_{DS}(A)}^{\mathrm{bi}}\bigl(Q_{DS}(M)\bigr).
\]
\end{theorem}

\begin{conjecture}[Reduction commutes with associative chiral duality
\textup{(}arbitrary nilpotent\textup{)}]
\ClaimStatusConjectured
\label{conj:DS-arbitrary-nilpotent}
Let $Q_{DS}$ denote a BRST/Drinfeld--Sokolov reduction functor for an arbitrary nilpotent
reduction datum on a filtered pronilpotent subcategory of strongly admissible associative chiral
algebras. Assume:
\begin{enumerate}[label=(\alph*)]
\item $Q_{DS}$ is exact on the relevant filtered complexes;
\item $Q_{DS}$ preserves PBW filtrations;
\item the BRST differential lifts to the ordered bar coalgebra.
\end{enumerate}
Then
\[
Q_{DS}(A^!)\simeq Q_{DS}(A)^!.
\]
For every $A$-bimodule $M$ one should have
\[
Q_{DS}\bigl(\KK_A^{\mathrm{bi}}(M)\bigr)\simeq
\KK_{Q_{DS}(A)}^{\mathrm{bi}}\bigl(Q_{DS}(M)\bigr).
\]
\end{conjecture}

\begin{remark}
Theorem~\ref{thm:ordered-associative-ds-principal} records the proved principal corridor from
Volume~II. Conjecture~\ref{conj:DS-arbitrary-nilpotent} isolates the remaining non-principal
frontier: one must verify PBW compatibility and the lift of the BRST differential to the ordered
bar coalgebra beyond the principal and hook-type regimes.
\end{remark}

\subsection{Coderived chiral coproduct for class~M}
\label{subsec:coderived-chiral-coproduct}

The chiral bialgebra equivalence
(Theorem~\ref{thm:chiral-qg-equiv}) is restricted to the Koszul
locus: the arrow $\beta$ from $A_\infty$-maps to the chiral
coproduct applies the cobar functor $\Omega$ to the deconcatenation
coproduct $\Delta$ on~$\barB^{\mathrm{ord}}(\cA)$, and the
identification $\cA \simeq \Omega(\barB^{\mathrm{ord}}(\cA))$
is a quasi-isomorphism only on the Koszul locus. For
class~$\mathsf{M}$ algebras (Virasoro, $\cW_N$, $\cW_{1+\infty}$),
the bar-cobar counit is \emph{not} a chain-level quasi-isomorphism
but \emph{is} an isomorphism in the coderived category
$D^{\mathrm{co}}$ (Theorem~\ref{thm:off-koszul-ran-inversion}).
The deconcatenation coproduct, being a purely combinatorial
structure on the cofree tensor coalgebra, survives intact. This
suggests a coderived extension of the chiral quantum group
equivalence.

\begin{conjecture}[Coderived chiral coproduct; \ClaimStatusConjectured]
\label{conj:coderived-chiral-coproduct}%
\index{coderived category!chiral coproduct|textbf}%
\index{coproduct!chiral!coderived extension}%
\index{chiral quantum group!coderived extension}%
Let $\cA$ be a complete augmented chiral algebra with
finite-dimensional graded bar complex, of arbitrary shadow
class.
\begin{enumerate}[label=\textup{(\roman*)}]
\item \textup{(Coderived chiral coproduct.)}
  The deconcatenation coproduct
  $\Delta \colon \barB^{\mathrm{ord}}(\cA) \to
  \barB^{\mathrm{ord}}(\cA) \otimes
  \barB^{\mathrm{ord}}(\cA)$
  is a CDG-coalgebra map. Applying the coderived cobar functor
  $\Omega^{\mathrm{co}}$ of
  Theorem~\textup{\ref{thm:off-koszul-ran-inversion}} produces a
  morphism in $D^{\mathrm{co}}$:
  \[
  \Delta^{\mathrm{ch,co}}
  \;=\; \Omega^{\mathrm{co}}(\Delta)
  \colon \cA
  \;\longrightarrow\;
  \cA \mathbin{\hat\otimes}^{\mathrm{co}} \cA,
  \]
  where $\hat\otimes^{\mathrm{co}}$ is the completed tensor
  product in $D^{\mathrm{co}}$.
\item \textup{(Coderived coassociativity.)}
  The strict coassociativity of deconcatenation
  $(\Delta \otimes \id) \circ \Delta
  = (\id \otimes \Delta) \circ \Delta$
  descends through $\Omega^{\mathrm{co}}$ to coassociativity
  up to conjugation by a Drinfeld associator
  $\Phi \in \exp(\hat{\mathfrak{t}}_3)$ in $D^{\mathrm{co}}$:
  \[
  (\Delta^{\mathrm{ch,co}} \otimes \id)
  \circ \Delta^{\mathrm{ch,co}}
  \;=\;
  \Phi^{-1} \cdot
  (\id \otimes \Delta^{\mathrm{ch,co}})
  \circ \Delta^{\mathrm{ch,co}}
  \cdot \Phi.
  \]
\item \textup{(Class-by-class agreement.)}
  For classes~$\mathsf{G}$, $\mathsf{L}$, and~$\mathsf{C}$
  \textup{(}where chain-level bar-cobar inversion holds\textup{)},
  $\Delta^{\mathrm{ch,co}}$ agrees with the chain-level chiral
  coproduct $\Delta^{\mathrm{ch}}$ of
  Theorem~\textup{\ref{thm:chiral-qg-equiv}}.
  For class~$\mathsf{M}$, $\Delta^{\mathrm{ch,co}}$ is the
  \emph{only} chiral coproduct available, and it equips
  class~$\mathsf{M}$ algebras with a coderived chiral quantum
  group structure.
\end{enumerate}
\end{conjecture}

\begin{remark}[Structure of the coderived coproduct conjecture]
\label{rem:coderived-chiral-coproduct-structure}%
\index{coderived category!chiral coproduct!mechanism}%
The conjecture combines three proved inputs with two unproved ones.

The three proved inputs are:
\begin{enumerate}[label=\textup{(\alph*)}]
\item The deconcatenation coproduct $\Delta$ on the cofree tensor
  coalgebra $T^c(s^{-1}\bar\cA)$ is strictly coassociative and the
  bar differential $d_{\barB}$ is a coderivation with respect
  to~$\Delta$
  (Construction~\ref{constr:deconcatenation}): $\Delta$ is a
  CDG-coalgebra map for the curved bar complex at every genus.
\item The coderived bar-cobar inversion
  $\Omega^{\mathrm{co}}(\barB(\cA)) \simeq \cA$ in
  $D^{\mathrm{co}}$ holds without Koszulness hypothesis
  (Theorem~\ref{thm:off-koszul-ran-inversion}).
\item The coderived BV$=$bar identification absorbs all
  $m_0$-proportional obstructions for every shadow class
  (Theorem~\ref{thm:bv-bar-coderived}).
\end{enumerate}

The two unproved inputs are:
\begin{enumerate}[label=\textup{(\alph*)}]
\setcounter{enumi}{3}
\item \emph{Coderived K\"unneth comparison.}
  The natural map
  $\Omega^{\mathrm{co}}(C_1 \otimes C_2) \to
  \Omega^{\mathrm{co}}(C_1) \hat\otimes^{\mathrm{co}}
  \Omega^{\mathrm{co}}(C_2)$
  must be an isomorphism in $D^{\mathrm{co}}$ for CDG-coalgebras
  with finite-dimensional graded pieces. On the Koszul locus, this
  is the standard Eilenberg--Zilber quasi-isomorphism for DG
  coalgebras. In $D^{\mathrm{co}}$, the comparison requires
  Positselski's theory of the derived tensor product for CDG
  comodules; its factorization-algebraic version has not been
  developed.
\item \emph{Coderived Dunn additivity for the tensor product.}
  The coassociativity-up-to-$\Phi$ statement requires that the
  Drinfeld associator $\Phi$, which is an algebraic
  (group-like, uncurved) element of the completed braid algebra,
  acts correctly in $D^{\mathrm{co}}$. Since $\Phi$ is a
  quasi-isomorphism in the ordinary derived category and has no
  curvature dependence, this should follow from the comparison
  $D^{\mathrm{co}} \simeq D$ on uncurved objects, but the
  interaction with the curved target
  $\cA \hat\otimes^{\mathrm{co}} \cA$ requires care.
\end{enumerate}

The obstruction structure parallels
Conjecture~\textup{\ref{conj:coderived-e3}}: input~(d) here
is the coproduct analogue of the coderived tensor product issue
flagged in Remark~\textup{\ref{rem:coderived-e3-structure}}.
A proof of either conjecture's tensor-product input would resolve
the other.
\end{remark}

\begin{remark}[Relation to the coderived $\Ethree$ conjecture]
\label{rem:coderived-coproduct-vs-e3}%
\index{coderived category!chiral coproduct!relation to $\Ethree$}%
Conjecture~\ref{conj:coderived-chiral-coproduct} and
Conjecture~\ref{conj:coderived-e3} are complementary extensions
of the Koszul-locus theory to class~$\mathsf{M}$.
The $\Ethree$ conjecture extends the \emph{closed sector}
(derived chiral center) to $D^{\mathrm{co}}$; the coproduct
conjecture extends the \emph{open sector} (chiral quantum group
structure) to $D^{\mathrm{co}}$. Together they would give a
full coderived $\mathsf{SC}^{\mathrm{ch,top}}$-algebra structure
for class~$\mathsf{M}$: the coproduct $\Delta^{\mathrm{ch,co}}$
as the open color, the derived center as the closed color, and
the $\Ethree$ as the compatibility between the two colors.
The shared obstruction (coderived K\"unneth/Dunn additivity)
ensures that the two conjectures stand or fall together.
\end{remark}

\begin{definition}[Coderived chiral bialgebra; \ClaimStatusConjectured]
\label{def:coderived-chiral-bialgebra}%
\index{coderived chiral bialgebra|textbf}%
\index{coderived category!chiral bialgebra|textbf}%
\index{bialgebra!coderived chiral|textbf}%
A \emph{coderived chiral bialgebra} is an object~$C$ in the
coderived category $\Dco((\mathrm{Ass}^{\mathrm{ch}})^!\text{-}\mathrm{coalg})$
equipped with:
\begin{enumerate}[label=\textup{(\roman*)}]
\item \textup{(Coalgebra structure.)}
  A coassociative coproduct
  $\Delta \colon C \to C \otimes C$
  in the category of CDG-coalgebras.
\item \textup{(Coderivation differential.)}
  A differential $d_C$ on~$C$ that is a \emph{coderivation} of~$\Delta$:
  \[
  d_C \circ \Delta \;=\; \Delta \circ d_C.
  \]
  The compatibility is automatic from the coderivation property and
  requires no separate axiom.
\item \textup{(Coderived acyclicity.)}
  The failure $d_C^2 \neq 0$ \textup{(}curvature from
  genus~$\geq 1$ contributions\textup{)} is exact in~$\Dco$:
  the CDG-comodule $(C, d_C)$ represents an honest object
  in the Verdier quotient
  $\Dco = \mathrm{Hot}/\mathrm{Acycl}^{\mathrm{co}}$.
\end{enumerate}
A morphism of coderived chiral bialgebras is a CDG-coalgebra
map that intertwines both~$\Delta$ and~$d_C$ in~$\Dco$.
\end{definition}

\begin{proposition}[The ordered bar complex is a coderived chiral bialgebra;
\ClaimStatusProvedHere]
\label{prop:bar-is-coderived-chiral-bialgebra}%
\index{bar complex!ordered!coderived chiral bialgebra structure}%
\index{coderived chiral bialgebra!from ordered bar complex}%
Let $\cA$ be a complete augmented chiral algebra with
finite-dimensional graded pieces. The ordered bar complex
$\barB^{\mathrm{ord}}(\cA) = T^c(s^{-1}\bar\cA)$ is a
coderived chiral bialgebra, with:
\begin{enumerate}[label=\textup{(\alph*)}]
\item $\Delta$ the deconcatenation coproduct
  \textup{(}Construction~\textup{\ref{constr:deconcatenation}}\textup{)},
  which is strictly coassociative;
\item $d_{\barB} = \sum_k \hat{m}_k^{\mathrm{ch}}$ the bar
  differential, which is a coderivation of~$\Delta$ by
  construction \textup{(}each $\hat{m}_k^{\mathrm{ch}}$ is the
  unique coderivation lifting $m_k^{\mathrm{ch}}$
  through the cofree coalgebra\textup{)};
\item the coderived acyclicity condition holds by
  Theorem~\textup{\ref{thm:off-koszul-ran-inversion}}
  and Theorem~\textup{\ref{thm:bv-bar-coderived}}: the
  bar-cobar counit $\Omega^{\mathrm{co}}(\barB^{\mathrm{ord}}(\cA))
  \xrightarrow{\;\sim\;} \cA$ is an isomorphism in~$\Dco$
  without any Koszulness hypothesis.
\end{enumerate}
On the Koszul locus \textup{(}classes~$\mathsf{G}$, $\mathsf{L}$,
$\mathsf{C}$\textup{)}, the coderived structure reduces to
a strict DG-bialgebra \textup{(}$d_{\barB}^2 = 0$
at chain level\textup{)}.
For class~$\mathsf{M}$, the coderived category is essential:
$d_{\barB}^2 \neq 0$ at chains, but $d_{\barB}^2$ is coacyclic.
\end{proposition}

\begin{proof}
Part~(a) is Construction~\ref{constr:deconcatenation}: the
deconcatenation coproduct on the cofree tensor coalgebra
$T^c(s^{-1}\bar\cA)$ is strictly coassociative by inspection.
Part~(b) is the universal property of cofree coalgebras:
a coderivation of $T^c(V)$ is determined by its projection
to~$V$, so each lifted operation $\hat{m}_k^{\mathrm{ch}}$
is automatically a coderivation of~$\Delta$.
The compatibility $d_{\barB} \circ \Delta = \Delta \circ d_{\barB}$
follows because a sum of coderivations is a coderivation,
and the composition of a coderivation with the coproduct
respects the Leibniz rule by definition.
Part~(c): the coderived bar-cobar inversion
(Theorem~\ref{thm:off-koszul-ran-inversion}) gives
$\Omega^{\mathrm{co}}(\barB^{\mathrm{ord}}(\cA)) \simeq \cA$
in~$\Dco$ for arbitrary shadow class; this uses the
coderived BV$=$bar identification
(Theorem~\ref{thm:bv-bar-coderived}) to absorb the
$m_0$-proportional obstructions.
On the Koszul locus the conilpotent reduction theorem
(Theorem~\ref{thm:conilpotent-reduction}) identifies
$\Dco$ with the ordinary derived category, so the coderived
bialgebra is a genuine DG-bialgebra.
\end{proof}

\begin{remark}[The projection disease]
\label{rem:projection-disease}%
\index{projection disease|textbf}%
\index{coderived chiral bialgebra!projection disease}%
The coderived chiral bialgebra $\barB^{\mathrm{ord}}(\cA)$ is
the single primitive object from which
the apparently heterogeneous landscape of products, coproducts,
and intertwining maps is recovered by applying different functors.
The deconcatenation coproduct~$\Delta$ encodes the chiral
coproduct (via cobar,
Conjecture~\ref{conj:coderived-chiral-coproduct}).
The bar differential~$d_{\barB}$ as coderivation of~$\Delta$
encodes the RTT relation, the vertex bialgebra axiom, and the
OPE-coproduct intertwining: these are the
\emph{same} coderivation property read in different
coordinates.
The derived chiral center
$Z^{\mathrm{der}}_{\mathrm{ch}}(\cA)$ is the chiral Hochschild
cochain complex of the bialgebra, carrying the
$\mathsf{E}_3$ structure of
Conjecture~\ref{conj:coderived-e3}.

The variety of structures that appear in the
literature (CoHA multiplication, Yangian coproduct,
W-algebra OPE, quantum group R-matrix, vertex
bialgebra axioms) are projections of this single coderived
chiral bialgebra through different forgetful functors:
%
% FOUR OBJECTS (AP25/AP34/AP50):
% 1. B^ord(A) = bar coalgebra = T^c(s^{-1} A-bar) with deconcatenation + twist
% 2. A^i = H^*(B(A)) = dual coalgebra (Koszul cohomology of bar)
% 3. A^! = ((A^i)^v) = dual ALGEBRA (linear or Verdier dual)
% 4. Z^der_ch(A) = derived chiral center = Hochschild cochains = bulk
%
\begin{enumerate}[label=\textup{(\alph*)}]
\item $\Delta \mapsto \Delta^{\mathrm{ch,co}}$ via the cobar
  functor: the chiral quantum group coproduct;
\item $d_{\barB} \mapsto m_k^{\mathrm{ch}}$ via projection to
  cogenerators: the $A_\infty$ operations (scattering amplitudes);
\item $\Delta + d_{\barB}$ jointly $\mapsto$ the RTT equation
  $R \Delta(a) = \Delta^{\mathrm{op}}(a) R$ via the universal
  R-matrix;
\item $\barB^{\mathrm{ord}} \mapsto \barB^{\Sigma}$ via
  $\Sigma_n$-coinvariants: the symmetric bar and its shadow tower
  $(\kappa, C, Q, \ldots)$.
\end{enumerate}
The coderivation property $d_{\barB} \circ \Delta = \Delta \circ
d_{\barB}$ is the universal compatibility from which all specific
compatibilities are derived by functoriality.
\end{remark}

\subsection{The \texorpdfstring{$\mathcal{W}_3$}{W\_3} ordered bar complex via DS transport}
\label{subsec:w3-ordered-bar-ds}

The principal DS reduction $V_k(\mathfrak{sl}_3) \twoheadrightarrow
\mathcal{W}_3^k$ transports the ordered bar complex of
$\widehat{\mathfrak{sl}}_{3,k}$ to the ordered bar complex of
$\mathcal{W}_3$. The DS--KD intertwining theorem
(Theorem~\ref{thm:ordered-associative-ds-principal}; for the symmetric
bar, Theorem~\ref{thm:ds-koszul-intertwine}) guarantees that the
functor $H^0(Q_{\mathrm{DS}})$ commutes with the ordered bar
construction at chain level:
\begin{equation}\label{eq:w3-ordered-bar-ds}
\Barch_{\mathrm{ord}}(\mathcal{W}_3^k)
\;\simeq\;
H^0_{\mathrm{DS}}\bigl(\Barch_{\mathrm{ord}}(V_k(\mathfrak{sl}_3))\bigr).
\end{equation}
This subsection records the explicit structure of the left-hand side.

\begin{theorem}[Ordered bar complex of $\mathcal{W}_3$ via DS transport; \ClaimStatusProvedHere]
\label{thm:w3-ordered-bar}
Let $k \neq -3$ be non-critical. The ordered bar complex
$\Barch_{\mathrm{ord}}(\mathcal{W}_3^k)$ has the following structure.

\begin{enumerate}[label=\textup{(\roman*)}]
\item \textbf{Generators.} The complex is cogenerated by
 $T$ \textup{(}spin~$2$, from the Sugawara tensor\textup{)} and
 $W$ \textup{(}spin~$3$, from the cubic Casimir\textup{)}, both
 obtained as $H^0(Q_{\mathrm{DS}})$-images of $\mathfrak{sl}_3$
 currents.

\item \textbf{Pole structure.} The collision residues are:
\begin{itemize}
\item $T$-$T$ channel: poles up to order~$3$ in the collision residue
 \textup{(}from the quartic Virasoro OPE $T_{(3)}T = c/2$; the
 $d\log$ kernel absorbs one pole\textup{)}.
\item $T$-$W$ channel: poles up to order~$2$
 \textup{(}from $T_{(2)}W = 3W$\textup{)}.
\item $W$-$W$ channel: poles up to order~$4$
 \textup{(}from the quintic OPE pole
 $W_{(4)}W = c/3$; the $d\log$ kernel absorbs one\textup{)}.
\end{itemize}

\item \textbf{A$_\infty$ operations.} The transferred operations
 $m_k^{\mathcal{W}_3}$ on the cohomology
 $H^\bullet(\Barch_{\mathrm{ord}}(\mathcal{W}_3^k))$ satisfy:
 \begin{itemize}
 \item $m_2$ is the ordered binary product
 \textup{(}collision residue at $D_{ij}$\textup{)};
 \item $m_3 \neq 0$ in the $T$-sector, inherited from the Virasoro
 quartic pole;
 \item $m_k \neq 0$ for all $k \geq 3$
 \textup{(}class~$\mathrm{M}$: infinite tower\textup{)}.
 \end{itemize}

\item \textbf{Curvature.} The bar curvature
 decomposes by channel:
 \[
 m_0^{\mathcal{W}_3} = \kappa_T |0\rangle + \kappa_W |0\rangle
 = \frac{c}{2}|0\rangle + \frac{c}{3}|0\rangle
 = \frac{5c}{6}|0\rangle,
 \]
 which is the DS projection of the $\widehat{\mathfrak{sl}}_{3,k}$
 curvature $m_0 = (k{+}3)\kappa/6$.

\item \textbf{Complementarity.}
 The Koszul involution is $c \mapsto 100 - c$,
 and the Koszul conductor is
 \[
 K_3 = 2(3{-}1)(2 \cdot 9 + 2 \cdot 3 + 1) = 100.
 \]
 The algebra is Koszul self-dual at $c^* = K_3/2 = 50$.
 The full Koszul conductor
 $K(\mathcal{W}_3) = K_3 \cdot (H_3 - 1) = 100 \cdot 5/6 = 250/3$
 controls the complementarity sum
 $\kappa(\mathcal{W}_3^k) + \kappa(\mathcal{W}_3^{k'}) = 250/3$.
\end{enumerate}
\end{theorem}

\begin{proof}
\emph{Part~(i).} The principal DS reduction
$V_k(\mathfrak{sl}_3) \twoheadrightarrow \mathcal{W}_3^k$
sends the $8$~current generators $J^a$ ($a = 1, \ldots, 8$) of
$\widehat{\mathfrak{sl}}_{3,k}$ to two generators: the stress tensor
$T$ (via the Sugawara construction) and the spin-$3$ field~$W$
(via the cubic Casimir). The remaining $6$~currents lie in
$\operatorname{im}(Q_{\mathrm{DS}})$ and are killed by the
projection $H^0(Q_{\mathrm{DS}})$. By
equation~\eqref{eq:w3-ordered-bar-ds}, the ordered bar complex
of $\mathcal{W}_3$ is cogenerated by the images of those
$\widehat{\mathfrak{sl}}_3$ bar elements that survive DS projection.

\emph{Part~(ii).} The pole orders in the ordered bar differential
are one less than the corresponding OPE pole orders
(the $d\log$ kernel absorbs one power). The
$\mathcal{W}_3$ OPE has: $T(z)T(w) \sim c/2 \cdot (z{-}w)^{-4}
+ \cdots$ (quartic pole), $T(z)W(w) \sim 3W(w)(z{-}w)^{-2}
+ \cdots$ (cubic pole), and
\begin{equation}\label{eq:ww-ope-leading-ordered}
W(z)W(w) \sim \frac{c/3}{(z{-}w)^6}
+ \frac{2T}{(z{-}w)^4}
+ \frac{\partial T}{(z{-}w)^3}
+ \frac{\beta\Lambda + \tfrac{3}{10}\partial^2 T}{(z{-}w)^2}
+ \cdots
\end{equation}
where $\Lambda = (TT) - \tfrac{3}{10}\partial^2 T$ is the
quasi-primary of weight~$4$ and
$\beta = 16/(22 + 5c)$. The $d\log$ absorption
gives collision residues with poles of order $3$, $2$, and $4$ in the
respective channels (since $W_{(5)}W = 0$ and the leading nonzero
term is $W_{(4)}W = c/3$, the sextic pole produces a quintic-order
$n$-product, and the $d\log$ kernel reduces the collision residue
pole by one to order~$4$).

\emph{Part~(iii).} The Virasoro subalgebra
$\operatorname{Vir}_c \hookrightarrow \mathcal{W}_3^k$ (generated
by~$T$) is autonomous: its OPE closes on itself. The quartic pole
$T_{(3)}T = c/2$ forces $m_3 \neq 0$ on the $T$-sector by the
single-line class~M theorem
(Theorem~\ref{thm:single-line-dichotomy}): the critical
discriminant $\Delta_T = 40/(5c+22) \neq 0$ ensures that the shadow
tower is infinite. The operations $m_k$ for $k \geq 4$ are then
generated by the resolvent tree formula
(Theorem~\ref{thm:tree-formula}): each $m_k$ is a sum over
planar binary trees with $k$~leaves,
\begin{equation}\label{eq:w3-resolvent-tree}
m_k^{\mathcal{W}_3}
= \sum_{T \in \mathrm{PBT}(k)}
 m_2^{(T)} \circ h_{\mathrm{DS}}^{(\mathrm{edges}(T))},
\end{equation}
where $m_2^{(T)}$ denotes the iterated binary product along the
tree~$T$, and $h_{\mathrm{DS}}$ is the DS contracting homotopy
inserted on each internal edge. The number of planar binary trees
with $k$~leaves is the Catalan number $C_{k-1}$, so $m_k$ receives
$C_{k-1}$ contributions. The homotopy $h_{\mathrm{DS}}$
is the composition of the standard SDR homotopy for
$\widehat{\mathfrak{sl}}_3 \twoheadrightarrow \mathcal{W}_3$ with
the homotopy transfer data
(Construction~\ref{constr:transfer-ainf}): the infinite tower
$\{m_k\}_{k \geq 3}$ is generated from the single quadratic product
$m_2$ on $\widehat{\mathfrak{sl}}_3$ by iterated composition with
$h_{\mathrm{DS}}$.

The $W$-line also has $m_{2k} \neq 0$ for all $k \geq 2$ (but
$m_{2k+1} = 0$ by the $\mathbb{Z}_2$ parity $W \mapsto -W$);
the nonvanishing of $m_4^W$ follows from
$\Delta_W = 20480/[3(5c{+}22)^3] \neq 0$
(Computation~\ref{comp:w-infty-shadow-tower}).

\emph{Part~(iv).} The curvature decomposition
$\kappa(\mathcal{W}_3) = \kappa_T + \kappa_W = c/2 + c/3 = 5c/6$
is verified by direct computation
(Computation~\ref{comp:ds-bar-sl3-w3}, Step~4): the
DS projection of the $\widehat{\mathfrak{sl}}_3$ curvature
$m_0 = (k{+}3)\kappa/6$ onto the $T$- and $W$-channels produces
exactly $c/2$ and $c/3$ respectively.

\emph{Part~(v).} The Koszul involution is
$k \mapsto k' = -k - 2h^\vee = -k - 6$, which sends
$c(k) = 2 - 24(k{+}2)^2/(k{+}3)$ to $c' = 100 - c$. Hence
$c + c' = 100 = K_3$.
The self-dual point is $c = c'$, i.e., $c = 50$.
The complementarity sum
$\kappa(\mathcal{W}_3^k) + \kappa(\mathcal{W}_3^{k'})
= 5c/6 + 5(100{-}c)/6 = 250/3$
is $k$-independent.
\end{proof}

\begin{theorem}[Class~M persistence under DS transport; \ClaimStatusProvedHere]
\label{thm:class-m-ds-transport}
\index{class M!DS transport}
Let $\mathcal{W}_N = \mathcal{W}^k(\mathfrak{sl}_N, f_{\mathrm{prin}})$ be a
principal $\mathcal{W}$-algebra at non-critical level. Then:
\begin{enumerate}[label=\textup{(\roman*)}]
\item The ordered bar complex
$\Barch_{\mathrm{ord}}(\mathcal{W}_N)$ is class~$\mathrm{M}$:
the A$_\infty$ tower $\{m_k\}$ is infinite, with
$m_k \neq 0$ for all $k \geq 3$ \textup{(}on the $T$-line\textup{)}.

\item The non-formality is forced by the Virasoro sub-sector:
the stress tensor $T$ generates an autonomous Virasoro subalgebra
$\operatorname{Vir}_{c(k)} \hookrightarrow \mathcal{W}_N^k$, and
the quartic pole $T_{(3)}T = c/2$ forces $m_3 \neq 0$. The
higher operations cascade by the resolvent tree formula.

\item DS reduction preserves Koszulness but destroys formality:
$V_k(\mathfrak{sl}_N)$ is class~$\mathrm{L}$ \textup{(}double-pole
OPE, formal SC structure\textup{)}, while
$\mathcal{W}_N^k$ is class~$\mathrm{M}$
\textup{(}quartic-and-higher poles, genuinely infinite A$_\infty$
tower\textup{)}. The Sugawara construction escalates the pole depth:
the affine double-pole OPE $J^a(z)J^b(w) \sim k\delta^{ab}/(z{-}w)^2
+ f^{ab}_c J^c(w)/(z{-}w)$ produces, upon taking normal-ordered
products in the DS construction, the Virasoro quartic-pole OPE and
the $W$-channel poles of order $\leq 2s$ for each strong generator
of spin~$s$.

\item The Koszul conductor and self-dual central charge for
 $\mathcal{W}_N$ are:
\[
K_N = 2(N{-}1)(2N^2{+}2N{+}1),
\qquad
c^* = K_N/2 = (N{-}1)(2N^2{+}2N{+}1).
\]
For $N = 2$: $K_2 = 26$, $c^* = 13$ \textup{(}Virasoro\textup{)}.
For $N = 3$: $K_3 = 100$, $c^* = 50$.
For $N = 4$: $K_4 = 222$, $c^* = 111$.
\end{enumerate}
\end{theorem}

\begin{proof}
\emph{Part~(i).}
Each principal $\mathcal{W}_N$ contains a Virasoro subalgebra
$\operatorname{Vir}_{c(k)}$ generated by the stress tensor. The
$T$-channel shadow obstruction tower is the Virasoro shadow obstruction tower, which is
class~M by the single-line dichotomy:
$\Delta_T = 40/(5c+22) \neq 0$ for $c \neq -22/5$.

\emph{Part~(ii).}
The quartic pole $T(z)T(w) \sim c/2 \cdot (z{-}w)^{-4}$ in the
Virasoro sub-sector gives a cubic-order collision residue
(after $d\log$ absorption). This is the source of $m_3 \neq 0$:
the homotopy transfer formula produces $m_3 = p \circ m_2(h \circ
m_2 \otimes \id) + (-1)^{|a|} p \circ m_2(\id \otimes h \circ m_2)$
with both Catalan trees contributing nonzero terms. The
resolvent tree formula then generates all $m_k$ for $k \geq 4$ from
$m_2$ and the contracting homotopy.

\emph{Part~(iii).}
The affine algebra $V_k(\mathfrak{sl}_N)$ has a double-pole OPE,
so its collision residue has poles of order~$1$. The bar complex is
\emph{formal}: $m_k = 0$ for $k \geq 3$ on cohomology
(class~L). The DS reduction introduces the Sugawara construction
$T = \frac{1}{2(k+h^\vee)} (J^a J^a)$, which creates a
quartic pole $T(z)T(w) \sim c/2 \cdot (z{-}w)^{-4}$ from the
double-pole affine OPE. This is pole-order escalation via
normal ordering: two double poles compose (through the Sugawara
quadratic) to produce a quartic pole. For higher-spin generators
$W^{(s)}$, the $W^{(s)} W^{(s)}$ OPE has poles up to order~$2s$,
produced by the same mechanism applied to the degree-$s$ Casimir
invariants. The Koszulness is preserved because $H^0(Q_{\mathrm{DS}})$
is exact on the relevant filtered complexes
(Theorem~\ref{thm:ordered-associative-ds-principal}, hypothesis~(a)),
but the formality is destroyed because the homotopy transfer data
$(p, \iota, h_{\mathrm{DS}})$ generate nonzero $m_k$ for $k \geq 3$
from the new higher-order poles.

\emph{Part~(iv).}
The DS central charge is
$c(k) = (N{-}1)\bigl[1 - N(N{+}1)(N{-}1)/(k{+}N)\bigr]$
(Freudenthal--de Vries), and the Koszul involution
$k \mapsto -k - 2N$ sends $c \mapsto K_N - c$. The value of
$K_N$ is read off: $c(k) + c(-k{-}2N) = 2(N{-}1)(2N^2{+}2N{+}1)$.
The self-dual point is $c = c^* = K_N/2$.
\end{proof}

\begin{remark}[The resolvent tree formula for $\mathcal{W}_3$]
\label{rem:w3-resolvent}
\index{resolvent tree formula!$\mathcal{W}_3$}
The resolvent tree formula~\eqref{eq:w3-resolvent-tree} makes
explicit the mechanism by which DS reduction generates the infinite
A$_\infty$ tower from a single quadratic product. The affine
algebra $V_k(\mathfrak{sl}_3)$ has $m_2$ as its only nontrivial
operation (the bar complex is formal: class~L). Applying the
SDR data $(p, \iota, h_{\mathrm{DS}})$ to the homotopy transfer
theorem produces:
\begin{align*}
m_3^{\mathcal{W}_3}(a,b,c)
&= p\, m_2(h_{\mathrm{DS}}\, m_2(\iota a, \iota b),\, \iota c)
+ (-1)^{|a|}\,
 p\, m_2(\iota a,\, h_{\mathrm{DS}}\, m_2(\iota b, \iota c)),
\end{align*}
where $\iota\colon \mathcal{W}_3 \hookrightarrow V_k(\mathfrak{sl}_3)$
is the inclusion (via the Miura embedding) and
$p\colon V_k(\mathfrak{sl}_3) \twoheadrightarrow \mathcal{W}_3$
is the DS projection. At degree~$k$, the number of contributing
trees is the Catalan number $C_{k-1}$: for $m_3$, $C_2 = 2$
(the two planar binary trees with $3$ leaves); for $m_4$, $C_3 = 5$;
for $m_5$, $C_4 = 14$. Each tree contributes a term of the form
$m_2 \circ (h_{\mathrm{DS}} \circ m_2)^{\circ (k-2)}$, iterated
along the tree topology.

The key structural point: the affine quadratic product $m_2$ is
\emph{formal} (it alone generates no $m_3$), but the homotopy
$h_{\mathrm{DS}}$ is non-trivial on the Sugawara sector: it
mediates the passage from affine currents to the stress tensor and
higher-spin fields, and this non-trivial mediation is what generates
the infinite cascade. \emph{DS reduction preserves Koszulness but
destroys formality: this is the transport from class~$\mathrm{L}$
to class~$\mathrm{M}$.}
\end{remark}

\begin{remark}[The $W$-channel parity rule on the ordered bar complex]
\label{rem:w3-parity-ordered}
\index{parity rule!$\mathcal{W}_3$ ordered bar complex}
The $\mathbb{Z}_2$ symmetry $W \mapsto -W$ of the $\mathcal{W}_3$ OPE
descends to the ordered bar complex: on the $W$-line (deformation
direction $x_W$ with $x_T = 0$), all odd-degree operations vanish
($m_{2k+1}^W = 0$), while even-degree operations are generically nonzero
($m_{2k}^W \neq 0$ for $k \geq 1$). The first nonzero higher
operation on the $W$-line is $m_4^W$, with shadow coefficient
$S_4^W = 2560/[c(5c+22)^3]$
(Computation~\ref{comp:w-infty-shadow-tower}).
This parity constraint
is inherited from the $\mathbb{Z}_2$ automorphism of the DS
reduction: the principal embedding $\mathfrak{sl}_2
\hookrightarrow \mathfrak{sl}_3$ determines a $\mathbb{Z}_2$
that acts by $-1$ on the cubic Casimir and hence on~$W$.
\end{remark}

\subsection{Comparison with shifted and factorization quantum groups}
\label{subsec:shifted-factorization-comparison}

Three mathematical traditions converge on the same combinatorial-algebraic
object: the ordered bar complex $\Barch(A)$ on
$\FM_k(\mathbb{C})\times\mathrm{Conf}_k^<(\mathbb{R})$. The purpose of this
subsection is to place the ordered associative chiral Koszul duality of this
chapter in context by comparing it with the shifted quantum groups of
Finkelberg--Tsymbaliuk (and the Coulomb-branch approach of
Braverman--Finkelberg--Nakajima), the factorization quantum groups of
Latyntsev, and the stable-envelope programme of Maulik--Okounkov (as
synthesised for 3d $\mathcal{N}=4$ theories by Costello--Okounkov--Zhang--Zhou).

\subsubsection*{The three-language spine}

\[
\begin{array}{c|c|c|c}
\text{Datum} & \text{Physical (this work)} & \text{Spectral (Latyntsev)}
 & \text{Geometric (COZZ)} \\
\hline
R\text{-matrix} & \text{OPE exchange} & \text{factorisation isomorphism} &
 \text{stable envelope} \\
\text{YBE} & \text{Ward identity} & \text{factorisation axiom} &
 \text{braid relation} \\
\text{Strictification} & \text{root multiplicity 1} &
 \text{by construction} & \text{by construction}
\end{array}
\]

In each column the $R$-matrix, the Yang--Baxter equation, and the
strictification mechanism arise from the internal logic of the
framework:

\begin{itemize}
\item \textbf{Physical column.}
The $R$-matrix is the regularised holonomy of the derived additive KZ
connection on $\mathrm{Conf}_n^{\mathrm{ord}}(\mathbb{A}^1)$, with
connection one-form built from OPE residues
(Theorem~\ref{thm:derived-additive-kz}, Volume~II).
The Yang--Baxter equation is the Ward identity for the flat connection:
$\mathbb{D}_n^2=0$. Strictification requires that the spectral Drinfeld
obstruction class vanish, which holds for all simple Lie algebras because
every root has multiplicity one (Volume~II,
Theorem~\texttt{thm:complete-strictification}).

\item \textbf{Spectral column.}
Latyntsev defines factorisation quantum groups as factorisation algebras on
the spectral curve, with the $R$-matrix encoded as a factorisation
isomorphism (the structure map for disjoint-support factorisation). The
Yang--Baxter equation is the coherence axiom for threefold factorisation.
Strictification is built in: the factorisation structure is strict by
definition.

\item \textbf{Geometric column.}
Costello--Okounkov--Zhang--Zhou identify the $R$-matrix with the
Maulik--Okounkov stable envelope for the Nakajima quiver variety or the BFN
Coulomb branch. The Yang--Baxter equation is the braid relation for stable
envelopes (which follows from uniqueness of the stable envelope under the
Maulik--Okounkov axioms). Strictification is again automatic: stable
envelopes are honest matrices, not $\Ainf$ data.
\end{itemize}

\subsubsection*{The common object}

All three languages attach algebraic data to the same
configuration-space carrier:
\[
\Barchord(A)\;\;\text{on}\;\;
\FM_k(\mathbb{C})\times\mathrm{Conf}_k^<(\mathbb{R}).
\]
The holomorphic factor $\FM_k(\mathbb{C})$ carries the chiral
differential (OPE residues), while the topological factor
$\mathrm{Conf}_k^<(\mathbb{R})$ carries the coassociative coproduct
(ordered deconcatenation). What differs among the three traditions
is which of these two factors is taken as primary input and which
is derived.

\subsubsection*{Shifted quantum groups}

Finkelberg--Tsymbaliuk define, for a simple Lie algebra $\mathfrak{g}$
and a pair of coweights $(\mu,\nu)$, the \emph{shifted quantum group}
$Q_\mu^\nu(\mathfrak{g})$ by explicit generators (Drinfeld--Cartan
generators, shifted Chevalley generators) and shifted Serre relations.
The unshifted specialisation recovers the Yangian:
\begin{equation}\label{eq:unshifted-yangian}
Q_0^0(\mathfrak{g})\;\simeq\; Y_\hbar(\mathfrak{g}).
\end{equation}

\begin{theorem}[Unshifted identification; \ClaimStatusProvedHere]
\label{thm:unshifted-identification}
On the chirally Koszul locus, the line-side algebra $A^!_{\mathrm{line}}$
produced by ordered associative chiral Koszul duality of $\widehat{\mathfrak{g}}_k$
is isomorphic to the Yangian $Y_\hbar(\mathfrak{g})$:
\[
A^!_{\mathrm{line}}\bigl(\widehat{\mathfrak{g}}_k\bigr)
\;\simeq\;
Y_\hbar(\mathfrak{g})
\;\simeq\;
Q_0^0(\mathfrak{g}).
\]
\end{theorem}

\begin{proof}
This is the content of the Yangian recognition theorem
(Theorem~\ref*{thm:yangian-recognition}, recorded from Volume~II).
The RTT presentation of $Y_\hbar(\mathfrak{g})$ is read off from
the regularised holonomy of the derived additive KZ connection
(Theorem~\ref{thm:derived-additive-kz}). The identification
$Y_\hbar(\mathfrak{g})\simeq Q_0^0(\mathfrak{g})$ is the
Levendorski\u{\i}--Drinfeld theorem.
\end{proof}

\begin{remark}[Shift parameters as boundary conditions]
\label{rem:shift-boundary-conditions}
The shift parameters $(\mu,\nu)$ in $Q_\mu^\nu(\mathfrak{g})$ correspond,
in the 3d holomorphic-topological picture, to the choice of boundary
conditions at the two ends of the topological interval $\mathbb{R}$.
Different choices of boundary condition produce different shifted
quantum groups; the unshifted case $(\mu,\nu)=(0,0)$ corresponds
to Neumann boundary conditions on both sides. The algebraic
manifestation is that the shift modifies the Serre relations,
which on the bar-complex side corresponds to modifying the
augmentation datum of the ordered bar coalgebra.
\end{remark}

\subsubsection*{Factorization quantum groups}

Latyntsev constructs factorisation quantum groups as factorisation
algebras on the formal spectral curve $\widehat{\mathbb{A}}{}^1_u$,
working at the derived/$\infty$-categorical level throughout.

\begin{theorem}[Factorisation identification on the Koszul locus;
\ClaimStatusProvedHere]
\label{thm:factorisation-identification}
On the chirally Koszul locus for a simple Lie algebra $\mathfrak{g}$,
the dg-shifted Yangian $Y^{\mathrm{dg}}_\hbar(\mathfrak{g})$ produced
by ordered chiral Koszul duality, after bar-horizontal strictification,
yields a strict spectral factorisation quantum group.
\end{theorem}

\begin{proof}
The quasi-factorisation theorem (Volume~II,
Theorem~\texttt{thm:quasi-factorisation})
shows that every residue-bounded complete dg-shifted Yangian with a
bar-horizontal strictification produces a spectral quasi-factorisation
quantum group. The spectral Drinfeld obstruction class controls the
passage from quasi to strict. For simple $\mathfrak{g}$, root multiplicity
one forces this class to vanish at every filtration level (Volume~II,
Theorem~\texttt{thm:complete-strictification}), yielding an honest
strict factorisation quantum group.
\end{proof}

\begin{remark}[What the factorisation framework adds]
Latyntsev's construction operates at the $\infty$-categorical level from the
outset, producing factorisation quantum groups as $\Eone$-algebra objects in
ind-coherent sheaves on the Ran space of the spectral curve. The main
advantage over the operadic approach of this chapter is that Latyntsev's
framework gives a \emph{definition} of factorisation quantum group that does
not require choosing a strictification: the factorisation isomorphisms are
part of the defining data. Our contribution is to show that the
$\Eone$-chiral bar complex $\barB^{\mathrm{ord}}(\cA)$
produces the same object once strictification is carried out, and that
for simple Lie algebras the strictification is always possible. The
Kac--Moody frontier (root multiplicities $>1$) is the regime where the
strictification mechanism fails and Latyntsev's a~priori construction
must be used as the primary definition.
\end{remark}

\subsubsection*{Coulomb branch algebras and stable envelopes}

Braverman--Finkelberg--Nakajima (BFN) define the Coulomb branch algebra
$\mathcal{A}_C(G,N)$ as the equivariant Borel--Moore homology of the BFN
moduli space $\mathcal{R}_{G,N}$. Maulik--Okounkov construct $R$-matrices
from stable envelopes of Nakajima quiver varieties; the
Costello--Okounkov--Zhang--Zhou synthesis extends this to 3d
$\mathcal{N}=4$ gauge theories.

\begin{proposition}[$R$-matrix comparison; \ClaimStatusProvedHere]
\label{prop:r-matrix-stable-envelope}
For $\mathfrak{g}$ simple and $\widehat{\mathfrak{g}}_k$ at
non-critical generic level, the $R$-matrix $R(z)$ produced by the
regularised holonomy of the derived additive KZ connection coincides
with the Maulik--Okounkov stable envelope on evaluation modules.
\end{proposition}

\begin{proof}
On evaluation modules, both $R$-matrices satisfy the same
Yang--Baxter equation with the same spectral dependence and the
same classical limit $r(z)$ (the rational $r$-matrix of
$\mathfrak{g}$). The uniqueness theorem for trigonometric/rational
$R$-matrices with prescribed classical limit and spectral dependence
(the Drinfeld uniqueness argument) gives the identification.
\end{proof}

\begin{remark}[Beyond the affine lineage]
The comparison with stable envelopes is verified on evaluation modules
of the Yangian. The full comparison for non-evaluation modules,
and the extension to shifted quantum groups (where the stable
envelope is for the BFN Coulomb branch rather than the Nakajima
variety), remains conjectural in general.
\end{remark}

\begin{conjecture}[Full three-language equivalence]
\ClaimStatusConjectured
\label{conj:three-language-equivalence}
For any 3d $\mathcal{N}=4$ gauge theory with gauge group $G$ and
matter representation $N$:
\begin{enumerate}[label=\textup{(\alph*)}]
\item the ordered chiral Koszul dual of the boundary chiral algebra
 $A_{\partial}(G,N)$ is isomorphic, after bar-horizontal
 strictification, to a factorisation quantum group in Latyntsev's sense;
\item this factorisation quantum group, evaluated on the unshifted Ran
 space, recovers the Yangian of the Coulomb branch;
\item the $R$-matrix from ordered-to-unordered descent agrees with the
 Maulik--Okounkov stable envelope for $\mathcal{M}_C(G,N)$.
\end{enumerate}
\end{conjecture}

\begin{remark}[Status and scope]
Parts~(a) and~(b) of Conjecture~\ref{conj:three-language-equivalence}
are proved for the standard affine case
$A = \widehat{\mathfrak{g}}_k$ with $\mathfrak{g}$ simple
(Theorems~\ref{thm:unshifted-identification}
and~\ref{thm:factorisation-identification}). Part~(c)
is proved on evaluation modules
(Proposition~\ref{prop:r-matrix-stable-envelope}). The full conjecture
extends these identifications to arbitrary gauge theories and all
module categories.
\end{remark}

\subsection{Open problems from the ordered chiral homology programme}
\label{subsec:ordered-ch-hom-open-problems}

The ordered chiral homology theory
(\S\ref{sec:ordered-chiral-homology}, \S\ref{subsec:formality-bridge})
and the chiral bialgebra equivalence
(Theorem~\ref{thm:chiral-qg-equiv}) open five further problems.

\begin{theorem}[$\mathsf{E}_3$ identification for affine Kac--Moody]
\label{thm:e3-identification-km}
\ClaimStatusProvedHere
\index{E3 algebra@$\Ethree$-algebra!identification}
Let $\fg$ be a simple finite-dimensional Lie algebra and
$V_k(\fg)$ the affine Kac--Moody vertex algebra at level~$k$.
The derived chiral centre
$Z^{\mathrm{der}}_{\mathrm{ch}}(V_k(\fg))$ and the
Costello--Francis--Gwilliam perturbative Chern--Simons
$\Ethree$-algebra $\cA^\lambda$ are isomorphic as formal
deformation families of $\Ethree$-algebras over
$\hbar H^3(\fg)[[\hbar]]$ \textup{(}$\hbar = k + h^\vee$\textup{)}.
\end{theorem}

\begin{proof}
The deformation spaces agree (both one-dimensional at each order
in~$\hbar$ for simple~$\fg$), the associated graded algebras match
(both $C^*(\fg)$), and the formal disk global sections match at the
level of $\mathsf{P}_3$-brackets
(\S\ref{subsec:chiral-e3-cfg-comparison} of the standalone paper).
For simple~$\fg$, $H^3(\fg) \cong \CC$ is one-dimensional, so at
each order $p \geq 1$ in $\hbar$ the space of
$\Ethree$-deformations extending the lower orders is
one-dimensional. The matching of $\mathsf{P}_3$-brackets on the
formal disk fixes the scalar at each order, forcing the two
deformation families to be isomorphic. The isomorphism is
constructed order by order and passes to the $\hbar$-adic limit.
\end{proof}

\begin{conjecture}[Trigonometric and elliptic ordered chiral homology]
\label{conj:trig-elliptic-ordered}
\ClaimStatusConjectured
\index{R-matrix!trigonometric!ordered chiral homology}
\index{R-matrix!elliptic!ordered chiral homology}
The rational Yang $R$-matrix $R(u) = 1 + \hbar\,r(u)/u + O(\hbar^2)$
treated in this chapter extends to two richer settings:
the trigonometric $R$-matrix (quantum affine algebras) and the
elliptic Belavin kernel, and after adjoining the Cartan dynamical
variable the Felder--Etingof--Varchenko dynamical $R$-matrix, producing ordered chiral homology
with higher-order Stokes phenomena, theta-function monodromy, and
flat sections valued in elliptic hypergeometric integrals.
The regular/irregular dichotomy of
Theorem~\ref{thm:e1-formality-failure} should extend to a trichotomy
(regular, irregular-rational, irregular-elliptic) correlated with
the three types of quantum groups.
\end{conjecture}

\begin{conjecture}[Higher-genus ordered chiral homology]
\label{conj:higher-genus-ordered-ch}
\ClaimStatusConjectured
\index{ordered chiral homology!higher genus}
The genus-$0$ ordered chiral homology of
\S\ref{sec:ordered-chiral-homology} extends to genus $g \ge 1$
via a period-corrected differential whose degree-$2$ coefficients
involve the $(g \times g)$-matrix $(\operatorname{Im}\Omega)^{-1}$
of the Siegel period matrix. The ordered chiral chain complex at
genus~$g$ should define a flat connection over
$\overline{\cM}_g$ compatible with
$\kappa_1 = c_1(\mathbb{E})$, and the resulting system of conformal
blocks should satisfy factorisation at the boundary divisors of
$\overline{\cM}_g$.
\end{conjecture}

\begin{proposition}[Critical level: monodromy trivialises,
Koszulness fails, center jumps]
\label{prop:critical-level-ordered}
\ClaimStatusProvedHere
At the critical level $k = -h^\vee$ for
$V_{-h^\vee}(\mathfrak{sl}_2)$ ($h^\vee = 2$, so $k = -2$):
\begin{enumerate}[label=\textup{(\roman*)}]
\item $\kappa = 3(k+2)/4 = 0$: the bar complex is uncurved.
\item In the trace-form convention,
  $r(z) = k\,\Omega/z = -2\,\Omega/z$ is finite and nonzero.
  The monodromy eigenvalues
  $\exp(2\pi i\,k\,\lambda_J)$ are
  $\exp(-4\pi i) = 1$ on $\operatorname{Sym}^2 V$ and
  $\exp(6\pi i) = 1$ on $\bigwedge^2 V$:
  ALL monodromy is trivial.
\item The degree-$2$ ordered chiral homology jumps:
  $\dim H^0 = 4$ \textup{(}all of $V \otimes V$ is
  invariant\textup{)}, and
  $\chi = -4$ forces $\dim H^1 = 8$
  \textup{(}total dimension $12$, triple the generic~$4$\textup{)}.
\item Koszulness fails: $H^*(\barB^{\mathrm{ord}}
  (V_{-2}(\mathfrak{sl}_2)))$ is not concentrated in
  degree~$1$ but spreads to
  $\Omega^*(\operatorname{Op}_{\mathfrak{sl}_2}(D))$,
  the de~Rham algebra of $\mathfrak{sl}_2$-opers
  on the formal disk \textup{(}Feigin--Frenkel\textup{)}.
\item The spectral sequence mechanism: the bar
  differential $d_k = d_{\mathrm{crit}} + (k+2)\,\delta$.
  At $\lambda = k + 2 = 0$, the $d_1$ page vanishes and
  the full de~Rham algebra survives as bar cohomology.
  At generic $\lambda \neq 0$, the $d_1 = \lambda[\delta]$
  contracts the de~Rham complex to $\bC^3$
  \textup{(}Koszul concentration\textup{)}.
\end{enumerate}
\end{proposition}

\begin{proof}
Part~(i): direct substitution.
Part~(ii): on $V\otimes V = \mathrm{Sym}^2 V \oplus \Lambda^2 V$
the quadratic Casimir of $\mathfrak{sl}_2$ acts with eigenvalues
$c_2 = j(j+1)$, giving $c_2 = 2$ on the spin-$1$ piece
$\mathrm{Sym}^2 V$ and $c_2 = 0$ on the spin-$0$ piece
$\Lambda^2 V$; multiplied by $k = -2$ these give $-4$ and $0$,
both integers, so $\exp(2\pi i \cdot k\, c_2) = 1$ on every
eigenspace.
Part~(iii): trivial monodromy gives $H^0 = V \otimes V$
(all flat sections are global), so $\dim H^0 = \dim(V \otimes V) = 4$.
The Euler characteristic
$\chi_{\rm ord, 2} = \dim H^0_{\rm ord, 2} - \dim H^1_{\rm ord, 2}$
of the degree-$2$ ordered chiral homology is an invariant of the
underlying graded vector space, independent of the choice of
differential in the bicomplex $d_k = d_{\rm crit} + (k + 2)\delta$;
at generic $k$ one has $\dim H^0_{\rm ord,2} = 0$ and
$\dim H^1_{\rm ord,2} = 4$, giving $\chi_{\rm ord, 2} = -4$. Hence
at $k = -2$, $\dim H^1_{\rm ord, 2} = \dim H^0_{\rm ord, 2} -
\chi_{\rm ord, 2} = 4 - (-4) = 8$ (total dimension $12 = 3 \cdot 4$,
triple the generic~$4$).
Part~(iv): the Feigin--Frenkel theorem identifies the center
$Z(V_{-h^\vee}(\fg)) \cong
\operatorname{Fun}(\operatorname{Op}_{\fg^\vee}(D))$;
at the bar-complex level, this center is $H^0$ of the
bar complex, which is infinite-dimensional.
Part~(v): the bicomplex structure
$d_k = d_{\mathrm{crit}} + \lambda\,\delta$ with
$\lambda = k + h^\vee$ gives a spectral sequence whose
$E_1$ page is $H^*(d_{\mathrm{crit}})$ with differential
$d_1 = \lambda[\delta]$. At $\lambda = 0$, $d_1 = 0$ and
$E_2 = E_1 = H^*(d_{\mathrm{crit}})
\cong \Omega^*(\operatorname{Op}_{\fg^\vee}(D))$.
At $\lambda \neq 0$, $d_1$ is an isomorphism on the
contractible part, leaving the Koszul residue in degree~$1$.
\end{proof}

\begin{remark}[Shadow stratum vs bar Koszulness at critical level]
\label{rem:critical-shadow-vs-bar-koszul}
Under the quadrichotomy classification of the shadow tower
(Theorem~\ref{thm:quadrichotomy} of
Chapter~\ref{chap:shadow-quadrichotomy-platonic}) the critical-level
affine KM algebra $V_{-h^\vee}(\fg)$ sits on the discriminant
locus $\{\kappa = 0\} = \{\Delta = 0,\ \alpha = 0\}$ on the
universal vacuum line: the principal-line Riccati polynomial
$Q_c(t) = (2\kappa)^{2}$ collapses identically to zero, which is
the \emph{class G} stratum of the shadow tower. The shadow classification
therefore records critical level as shadow-depth-$2$ truncated, same
as Heisenberg. The full bar complex is \emph{not} Koszul, however:
it is quasi-isomorphic to the de~Rham algebra
$\Omega^{*}(\operatorname{Op}_{\fg^\vee}(D))$, which is spread over
all degrees, not concentrated in degree~$1$. Hence shadow truncation
depth and bar Koszulness are \emph{different} invariants: the former
is a discriminantal quadratic classification on a finite-dimensional
parameter space, the latter is a cohomological concentration
statement on the infinite-dimensional bar complex. At integrable
level $k \in \ZZ_{>0}$ and at generic $k$ the two agree (class L,
bar Koszul); at critical level they diverge (shadow class G, bar
non-Koszul). This divergence is a genuine feature of the critical
locus, not a mislabelling.
\end{remark}

\begin{remark}[Three-level contrast]
\label{rem:three-level-contrast}
\begin{center}
\renewcommand{\arraystretch}{1.3}
\begin{tabular}{@{} l c c c @{}}
\toprule
& Generic $k$ & Integrable $k \in \ZZ_{>0}$
& Critical $k = -2$ \\
\midrule
Koszul & yes & yes & \textbf{no} \\
$H^0$ (degree $2$) & $0$ & $0$ & $4$ \\
$H^1$ (degree $2$) & $4$ & $4$ & $8$ \\
Monodromy & infinite & finite nontrivial
  & \textbf{trivial} \\
Center & $\bC$ & $\bC^{k+1}$ & $\bC[S_2]$ (infinite) \\
$\kappa$ & $3(k{+}2)/4$ & $3(k{+}2)/4$ & $0$ \\
\bottomrule
\end{tabular}
\end{center}
Integrable and critical are opposite extremes: integrable
has finite but nontrivial monodromy with finite center;
critical has trivial monodromy with infinite center.
The ordered chiral homology detects both extremes through
$\ker(\av)$: at critical level, the entire $r$-matrix
$r(z) = -2\,\Omega/z$ lives in $\ker(\av)$ because
$\av(r(z)) = -\dim(\fg)/2$ and
$\kappa = \av(r(z)) + \dim(\fg)/2 = 0$.
\end{remark}

\begin{remark}[$N = 3$ explicit verification]
\label{rem:gl3-explicit}
The $9 \times 9$ Yang $R$-matrix $R(u) = u\,I_9 + \Psi\,P$
for $\mathfrak{gl}_3$ satisfies the YBE at $9$ random
complex parameter sets (relative error $< 10^{-14}$). All
$81$ RTT component relations verified at degree~$1$:
$[t_{ij}^{(1)}, t_{kl}^{(1)}] = \Psi(\delta_{il}\,
t_{kj}^{(1)} - \delta_{kj}\,t_{il}^{(1)})$.
The quantum determinant (six-term $S_3$ column
determinant with \emph{decreasing} column index ordering)
is central at $\Psi = 1, 0.5, 2$
\textup{(}$53$ compute tests\textup{)}.
\end{remark}

\begin{remark}[DS coproduct intertwining at $\cW_3$: degree-$1$ vs.\ degree-$\geq 2$]
\label{rem:ds-intertwining-w3}
The Drinfeld--Sokolov reduction
$\pi_3 \colon V_k(\mathfrak{sl}_3) \twoheadrightarrow
\cW_3$ intertwines the spectral coproducts:
$(\pi_3 \otimes \pi_3) \circ \Delta_z^{\mathfrak{sl}_3}
= \Delta_z^{\cW_3} \circ \pi_3$.
At degree~$1$ the intertwining is \emph{tautological} on the
Cartan sector: both sides compute the same Drinfeld primitive
formula on $\Delta_z(J)$ and $\Delta_z(\psi_n)$, so the $57$
compute tests at degree~$1$ verify a self-consistent
specialisation rather than an independent identity. At degree
$\geq 2$ the intertwining is non-tautological: it is equivalent
to the matching of the independent $\mathfrak{sl}_3$ RTT
commutators $[t_{ij}^{(1)}, t_{kl}^{(1)}] = \Psi(\delta_{il}
t_{kj}^{(1)} - \delta_{kj} t_{il}^{(1)})$ with the
$\cW_3$ OPE after $\pi_3$-reduction; these are cross-verified
on the $81$ RTT components at $\mathfrak{gl}_3$
\textup{(}Remark~\ref{rem:gl3-explicit}\textup{)}. The
spin-$3$ sector with its gravitational $\mathrm{HPL}_3$
contributions is the genuinely chiral content and remains the
sharpest non-tautological witness of DS intertwining. The
spectral coassociativity uses \emph{shifted} parameters:
$(\Delta_{z_1} \otimes \id) \circ \Delta_{z_1 + z_2}
= (\id \otimes \Delta_{z_2}) \circ \Delta_{z_1}$,
reflecting $T(u) \cdot T(u - z_1) \cdot T(u - z_1 - z_2)$.
\end{remark}

\begin{remark}[Explicit kernel basis at $d = 3$]
\label{rem:ker-av-d3-explicit}
For the adjoint representation of $\mathfrak{sl}_2$
($d = 3$): $\ker(\av_2) = \Lambda^2(\mathfrak{sl}_2)$
with basis $e \otimes f - f \otimes e$,
$e \otimes h - h \otimes e$,
$f \otimes h - h \otimes f$ (dimension~$3$).
At $n = 3$: $\ker(\av_3) = 17 = 1 + 16$, where $1$
is the volume form $e \wedge h \wedge f$ and $16$ is
the standard-representation complement
\textup{(}$51$ compute tests\textup{)}.
\end{remark}

\begin{remark}[Quantum determinant column ordering]
\label{rem:qdet-column-ordering}
The central quantum determinant $\operatorname{qdet} T(u)$
uses the column determinant with \emph{decreasing} column
index: $j = N - 1$ is the leftmost factor, $j = 0$ the
rightmost. At $N = 2$ both orderings coincidentally agree;
at $N \geq 3$ only the decreasing-index ordering is
central \textup{(}$74$ compute tests, \textup{)}.
\end{remark}

\begin{remark}[KZB heat equation prefactor]
\label{rem:kzb-heat-prefactor}
The genus-$g$ heat equation
$\partial_{\Omega_{ab}}\,\Theta = c_{ab}\,
\partial_{z_a}\partial_{z_b}\,\Theta$
has prefactor $c_{ab} = 1/(4\pi i)$ for $a = b$
(diagonal) and $c_{ab} = 1/(2\pi i)$ for $a \neq b$
(off-diagonal). The factor of~$2$ arises from the
symmetric matrix chain rule: $\partial/\partial\Omega_{aa}$
as an independent variable contributes
$(1/2)\,\partial/\partial(\text{matrix entry})$.
Verified numerically at machine precision ($< 10^{-14}$)
across $5$ generic points in the upper half-plane.
\end{remark}


\appendix


\section{Appendix: signs and conventions}

All formulas above are invariant under the usual choice of suspension conventions, provided one keeps
the same convention simultaneously in the ordered bar complex, the shuffle map, and the twisting
differential. For clarity we summarize the minimal sign content.

\begin{enumerate}[label=(\arabic*)]
\item Reversal in Theorem~\ref{thm:opposite} uses the standard suspended Koszul sign
\[
(-1)^{\sum_{i<j}(|a_i|+1)(|a_j|+1)}.
\]

\item The left twisting term in Definition~\ref{def:Kbi} carries the sign needed to move the
coderivation on $C$ past the first coalgebra leg before applying the left $A$-action.

\item The right twisting term carries the additional sign needed to move the twisting cochain across
the middle bimodule factor.

\item In the strongly admissible regime the filtration is complete and bounded below, so all spectral
sequence and homotopy-transfer arguments used in the proofs converge.
\end{enumerate}


\section{Appendix: new results proved in this chapter}

For quick reference, the new results proved in this chapter are as follows.

\begin{enumerate}[label=(\arabic*)]
\item Ordered chiral shuffle theorem (Theorem~\ref{thm:shuffle}).
\item Opposite-duality and anti-involution transport (Theorem~\ref{thm:opposite},
Corollary~\ref{cor:anti}).
\item Enveloping duality (Corollary~\ref{cor:enveloping}).
\item Two-sided twisted bicomodule and its dg structure (Definition~\ref{def:Kbi},
Lemma~\ref{lem:Kbi-dg}).
\item Identification of the tangent ordered bar complex with the twisted bicomodule
(Theorem~\ref{thm:tangent=K}).
\item PBW-complete bimodule/bicomodule equivalence (Theorem~\ref{thm:bimod-bicomod}).
\item Diagonal correspondence (Theorem~\ref{thm:diagonal}).
\item Composition unit theorem (Corollary~\ref{cor:unit}).
\item Tensor--cotensor gluing (Corollary~\ref{cor:tensor-cotensor}).
\item chiral Hochschild/coHochschild homology and cohomology equivalence
(Theorems~\ref{thm:HH-coHH-homology} and \ref{thm:HH-coHH-cohomology}).
\item Infinitesimal annular variation (Proposition~\ref{prop:infann}).
\item Ordered genus-zero open trace formalism (Theorem~\ref{thm:ordered-open}).
\item Calabi--Yau transport to an ordered Frobenius structure (Theorem~\ref{thm:CY}).
\item $E_1$-chiral double bar of $\hat{\mathfrak{sl}}_{2,k}$
(Proposition~\ref{thm:double-bar-sl2}).
\item Central extension invisible to the $E_1$-chiral double bar
(Theorem~\ref{thm:central-extension-invisible}).
\item Double Koszul duality is involutive
(Theorem~\ref{thm:two-colour-double-kd}).
\item Non-redundancy of the two colours
(Corollary~\ref{cor:two-colours-non-redundant}).
\item Heisenberg ordered bar complex: trivial differential, deconcatenation
coproduct, collision residue $r^{\mathrm{coll}}(z) = k/z$,
$R$-matrix $R(z) = e^{k\hbar/z}$
(Theorems~\ref{thm:heisenberg-ordered-bar}
and \ref{thm:heisenberg-rmatrix}).
\item Abelian Yangian $Y(\mathfrak{u}(1))$ as ordered Koszul dual
of $\cH_k$ (Theorem~\ref{thm:heisenberg-yangian}).
\item Formality of $\Barch(\cH_k)$: class~G, shadow depth~$2$,
$m_k = 0$ for $k\ge 3$ (Theorem~\ref{thm:heisenberg-formality}).
\item Universal Kac--Moody Yangian theorem: for every simple
$\mathfrak{g}$ and non-critical $k\neq -h^\vee$,
$A^{!,\mathrm{ord}}_{\mathrm{line}}
(\widehat{\mathfrak{g}}_k)\simeq Y_\hbar(\mathfrak{g})$
at $\hbar = 1/(k+h^\vee)$; collision residue
$r(z) = \hbar\,\Omega/z$; CYBE from Jacobi; complete
strictification by root-space one-dimensionality; curvature
$\kappa(Y_\hbar(\mathfrak{g})) = -\dim(\mathfrak{g})\cdot
(k+h^\vee)/(2h^\vee)$
(Theorem~\ref{thm:km-yangian},
Table~\ref{tab:km-yangian-data}).
\item $\mathcal{W}_3$ ordered bar complex via DS transport: generators,
pole structure, A$_\infty$ operations (class~M), curvature $5c/6$,
complementarity $K_3 = 100$, self-dual at $c^* = 50$
(Theorem~\ref{thm:w3-ordered-bar}).
\item Class~M persistence under DS transport for all $\mathcal{W}_N$:
Koszulness preserved, formality destroyed, $K_N = 2(N{-}1)(2N^2{+}2N{+}1)$
(Theorem~\ref{thm:class-m-ds-transport}).
\item $\beta\gamma$ ordered bar complex: nilpotent collision residue
$\Theta = \beta\otimes\gamma - \gamma\otimes\beta$ with $\Theta^2 = 0$,
R-matrix $R = \id + 2\pi i\,\Theta$ (nilpotent, terminates at linear
order), Koszul dual a $4$-dimensional exterior algebra (class~C,
shadow depth~$4$)
(Theorems~\ref{thm:bg-ordered-bar}, \ref{thm:bg-rmatrix},
and \ref{thm:bg-koszul-dual}).
\item $bc$ ghost ordered bar complex: nilpotent collision residue
$\Theta_{bc} = b\otimes c - c\otimes b$ with $\Theta_{bc}^2 = 0$,
R-matrix $R = \id + 2\pi i\,\Theta_{bc}$ (nilpotent), Koszul dual
a symmetric algebra (class~C, shadow depth~$4$, complementarity
$K = 0$) (Theorem~\ref{thm:bc-ordered-bar}).
\item Wakimoto bar complex descent: BRST reduction of free-field
ordered bar complex recovers affine ordered bar complex; dg Yangian
$Y^{\mathrm{dg}}_\hbar(\fg)$ as BRST reduction of free-field Yangian
(Theorem~\ref{thm:wakimoto-ordered-bar}).
\item Lattice vertex algebras: collision residue diagonal in
$\alpha$-sector, R-matrix from OPE monodromy for symmetric cocycles
($\Einf$), independent braiding for non-symmetric cocycles (genuinely
$\Eone$), Koszul dual with inverse cocycle for self-dual lattices
(Theorems~\ref{thm:lattice-symmetric-ordered-bar},
\ref{thm:lattice-nonsymmetric-ordered-bar},
and \ref{thm:lattice-ordered-koszul-dual}).
\item Root-space one-dimensionality for simple Lie algebras
(Theorem~\ref{thm:root-space-one-dim-v1}).
\item Jacobi collapse for star sectors
(Lemma~\ref{lem:jacobi-collapse-v1}).
\item Dynkin coefficient via the beta integral: the universal BCH
coefficient $1/n = \int_0^1(1-t)^{n-1}\,dt$
(Theorem~\ref{thm:dynkin-beta-integral}).
\item Complete spectral Drinfeld strictification for all simple Lie
algebras: the spectral Drinfeld class vanishes at every filtration,
so the dg-shifted Yangian $Y^{\mathrm{dg}}_\hbar(\fg)$ carries a strict filtered
algebra structure (Theorem~\ref{thm:complete-strictification-v1}).
\item Evaluation homomorphism $\mathrm{ev}_a\colon Y_\hbar(\fg)\to U(\fg)$
for all simple types (Construction~\ref{constr:evaluation-map}).
\item $\mathfrak{sl}_2$ $R$-matrix on evaluation modules: $R(u)=u\,\id+\hbar P$
(Yang $R$-matrix), Clebsch--Gordan decomposition and non-semisimplicity
at resonance (Theorem~\ref{thm:sl2-R-matrix}, Corollary~\ref{cor:sl2-clebsch-gordan}).
\item $\mathfrak{sl}_3$ evaluation modules: fundamental and adjoint
(Computations~\ref{comp:sl3-eval-fundamental} and \ref{comp:sl3-eval-adjoint}).
\item Drinfeld polynomial classification of evaluation modules
(Proposition~\ref{prop:eval-drinfeld}).
\item Line category $\mathcal{C}_{\mathrm{line}}\simeq
Y_\hbar(\fg)\text{-}\mathrm{mod}^{\mathrm{fd}}$ on the affine Koszul locus
(Theorem~\ref{thm:line-category}).
\item Braided monoidal structure on evaluation modules from the universal
$R$-matrix, including coalescence singularity at $a=b$
(Theorem~\ref{thm:eval-braiding}).
\item Quantum group $U_q(\mathfrak{g})$ from genus-$1$ $B$-cycle monodromy
of the KZB connection, with $q = e^{\pi i\hbar}$
(Theorem~\ref{thm:b-cycle-quantum-group}).
\item Drinfeld--Kohno theorem for the affine lineage: KZ/KZB monodromy
equals the quantum group braid representation
(Theorem~\ref{thm:drinfeld-kohno}).
\item Yangian--quantum group genus-$1$ deformation: $Y_\hbar(\mathfrak{g})$
at genus~$0$ deforms to $U_q(\mathfrak{g})$ at genus~$1$ via the
$B$-cycle exponential map $\widehat{\mathbb{G}}_a \to
\widehat{\mathbb{G}}_m$
(Theorem~\ref{thm:yangian-quantum-group}).
\item $U_q(\mathfrak{sl}_2)$ at roots of unity from affine
$\widehat{\mathfrak{sl}}_2$ at integer level; Kazhdan--Lusztig equivalence
(Corollary~\ref{cor:sl2-root-of-unity}).
\item Jones polynomial from genus-$1$ bar-complex monodromy via the
Reshetikhin--Turaev quantum trace
(Theorem~\ref{thm:jones-genus1}).
\item Annular bar complex: wrapping $\mathbb{R}\to S^1$ produces the
cyclic cotensor complex $B^{\mathrm{ann}}_\bullet(\cA)$ with four-component
differential $d^{\mathrm{ann}}=d_{\mathrm{int}}+d_{\mathrm{res}}+d_\Delta
+d_{\mathrm{wrap}}$ (Definition~\ref{def:annular-bar},
Theorem~\ref{thm:annular-bar-differential}).
\item Annular bar complex computes chiral Hochschild homology:
$B^{\mathrm{ann}}_\bullet(\cA)\simeq\HH^{\mathrm{ch}}_\bullet(\cA)$
(Theorem~\ref{thm:annular-HH}).
\item Curvature--braiding dichotomy at genus~$1$: the lowest
$\lambda$-bracket coefficient $c_0$ controls whether the genus-$1$
$E_1$-chiral bar complex decouples ($c_0=0$, Cartan type) or entangles
($c_0\ne 0$, Yangian type) via elliptic quasi-periodicity
(Theorem~\ref{thm:curvature-braiding-dichotomy}).
\end{enumerate}

\bigskip

\noindent\textbf{Final principle.}
The diagonal is the boundary. The bar coalgebra is the ambient ordered phase space. The annulus is
the self-cotensor of the diagonal. Once this is recognized, the associative chiral duality ceases
to be merely an equivalence between algebra and coalgebra and becomes a machine for open traces,
centers, and eventually bordered genus.

\begin{remark}[Non-commutative Euler product]%
\label{rem:e1-noncommutative-euler}%
The commutative sewing lift
$S_\cH(u) = \prod_p(1{-}p^{-u})^{-1}(1{-}p^{-u-1})^{-1}$
(Theorem~\ref{thm:dirichlet-weight-formula})
factors the Heisenberg sewing determinant over primes
via the unordered divisor function $\sigma_{-1}$.
The $\mathsf{E}_1$-ordered bar complex replaces
unordered by ordered factorizations; the analogous
ordered sewing lift should involve composition-weighted
divisor sums and factor in the ring of non-commutative
formal Dirichlet series.
\end{remark}
